%
\begin{isabellebody}%
\setisabellecontext{Paper}%
%
\isadelimtheory
%
\endisadelimtheory
%
\isatagtheory
%
\endisatagtheory
{\isafoldtheory}%
%
\isadelimtheory
%
\endisadelimtheory
%
\isadelimdocument
%
\endisadelimdocument
%
\isatagdocument
%
\isamarkupchapter{Introduction%
}
\isamarkuptrue%
%
\endisatagdocument
{\isafolddocument}%
%
\isadelimdocument
%
\endisadelimdocument
%
\begin{isamarkuptext}%
In the early stages of the compilation process, compilers perform \emph{syntactic analysis},
otherwise known as \emph{parsing}, which in essence consists of verifying that the input program
adheres to the syntax rules of the source language. Wilhelm, Seidl, and Hack (from this point onward
``Wilhelm et al.'') present the \isa{LR{\isacharparenleft}{\kern0pt}k{\isacharparenright}{\kern0pt}} parsing algorithms and their underlying theory in their book 
``\emph{Compiler Design: Syntactic and Semantic Analysis}''~\cite{Wilhelm}. This
thesis focuses on one variant of \isa{LR{\isacharparenleft}{\kern0pt}k{\isacharparenright}{\kern0pt}} parsing, namely \isa{LR{\isacharparenleft}{\kern0pt}{\isadigit{0}}{\isacharparenright}{\kern0pt}}, and presents two major results:
a formalization~\footnote{Some notes on terminology: the word ``formal'' when used throughout this
thesis carries exactly this meaning, i.e., machine-checked. Similarly, a ``formalization'' denotes
the translation of definitions, theorems, and proofs into code that is then machine-checked, in the
case of this thesis by Isabelle.} of this section of the book using the Isabelle Proof Assistant,
and a formal proof of the correctness of the parser constructed from the canonical \isa{LR{\isacharparenleft}{\kern0pt}{\isadigit{0}}{\isacharparenright}{\kern0pt}} automaton,
a property which Wilhelm et al.\ do not discuss in the original text. The official Isabelle
documentation~\footnote{\url{https://isabelle.in.tum.de/documentation.html}} provides detailed
information on the Isabelle Proof Assistant, and Nipkow and Klein offer an introduction to the
system in their book \emph{``Concrete Semantics with Isabelle/HOL''}~\cite{Nipkow-Book}.

The motivation for formalizing the work of Wilhelm et al.\ is one that is quite possibly shared among
all formal verification projects. In few words, informal proofs often contain errors.
Authors naturally rely on intuition to some extent to write
such proofs, and some arguments may seem to be correct on the surface, but turn out to be false if
one looks at them more closely. Wilhelm et al., naturally, are no exception to this problem, and one
can therefore use a proof assistant like Isabelle to discover errors and fill missing gaps in their
work more easily. Proof assistants prevent intuition from betraying the author: in machine-checked
proofs, every statement and every logical step, however simple, is verified by the proof assistant,
thereby confirming the correctness of the statements being proved. Thus, formal verification
guarantees that the conclusions one draws are indeed a consequence of the given premises by demanding 
logically sound reasoning. This leads to one of the two main questions I intend to answer: is \isa{LR{\isacharparenleft}{\kern0pt}{\isadigit{0}}{\isacharparenright}{\kern0pt}}
parsing theory as presented by Wilhelm et al.\ correct, and if it is not, how must it be modified to
achieve correctness?

Formal verification, naturally, is not infallible. In particular, the question that remains afterward
is whether the axioms, definitions and assumptions used are faithful to their informal counterparts.
This is precisely one of the areas where a formalization has the most room for debate, and one of
the main challenges in such a project. There are other questions which machine-checked proofs raise,
such as how trustworthy the proof assistants themselves are, but these are not in the scope of this
thesis.

The motivation for the second result, namely proving the parser's correctness, boils down to the
fact that in essence, the goal of building this parser is defining a deterministic automaton
that can decide whether a word is in a particular context-free grammar 
(CFG)~\footnote{In practice, parsers perform other tasks
such as building an internal representation of the syntactic structure of the input, but this is out
of the scope of the \isa{LR{\isacharparenleft}{\kern0pt}{\isadigit{0}}{\isacharparenright}{\kern0pt}} parsing section of the book.}. In the \isa{LR{\isacharparenleft}{\kern0pt}{\isadigit{0}}{\isacharparenright}{\kern0pt}} section of the book, the
main objective of Wilhelm et al.\ is to show one particular property: the parser's determinism. The
other core property of the parser, namely the fact that it recognizes the grammar's language, is not
discussed by them. However, this question of correctness is just as fundamental as the one about
determinism, since a parser that is deterministic but does not recognize the language of its grammar
is not useful. It is therefore of great interest to answer a second question on top of the formalization
of the authors' existing work: what is the language accepted by the \isa{LR{\isacharparenleft}{\kern0pt}{\isadigit{0}}{\isacharparenright}{\kern0pt}} parser as defined by
Wilhelm et al., and in particular, how does this language relate to that of the grammar from which it
is constructed?

For the first main question, on the correctness of the book, I formalized two main theorems. 
The first one, which highlights the relation between the characteristic finite automaton, rightmost
derivations, and the item pushdown automaton, is not fully correct, partly due to unclear wording
and variable naming, but the consequences the authors derive from it are correct. The proof
of the theorem is circular and consists of three implications. The first one is erroneous, and I
corrected it, while the remaining two are correct.

The second theorem, which states that the canonical \isa{LR{\isacharparenleft}{\kern0pt}{\isadigit{0}}{\isacharparenright}{\kern0pt}} automaton has no \isa{LR{\isacharparenleft}{\kern0pt}{\isadigit{0}}{\isacharparenright}{\kern0pt}}-inadequate states
if and only if the grammar is \isa{LR{\isacharparenleft}{\kern0pt}{\isadigit{0}}{\isacharparenright}{\kern0pt}}, is correct, but the proof has major errors in one direction,
and I present a corrected variant. 

Besides the two main theorems, there are other issues that I address throughout the thesis. The one
worth highlighting is definition of the \isa{LR{\isacharparenleft}{\kern0pt}k{\isacharparenright}{\kern0pt}} condition. Wilhelm et al.\ define this condition for
grammars extended by a new start symbol, but I present a proof that the satisfaction of the \isa{LR{\isacharparenleft}{\kern0pt}k{\isacharparenright}{\kern0pt}}
condition is not preserved in general when extending a grammar in this way. Since the \isa{LR{\isacharparenleft}{\kern0pt}{\isadigit{0}}{\isacharparenright}{\kern0pt}} parser
is defined only for extended grammars, this result shows that it is possible for an unextended grammar
to satisfy the \isa{LR{\isacharparenleft}{\kern0pt}{\isadigit{0}}{\isacharparenright}{\kern0pt}} condition while having a nondeterministic \isa{LR{\isacharparenleft}{\kern0pt}{\isadigit{0}}{\isacharparenright}{\kern0pt}} parser, since its extension
violates the condition, making the parser applicable only for a subset of \isa{LR{\isacharparenleft}{\kern0pt}{\isadigit{0}}{\isacharparenright}{\kern0pt}} grammars.

To answer the question regarding the parser's correctness, I have formally proved that it
is indeed correct, meaning that it accepts exactly the language of its grammar. For this proof, I
first introduce a concept I have named \emph{stack words}. This concept is based on an intuitive
explanation by Wilhelm et al.\ of how symbols on the parser's stack relate to each other, and with it,
I show the parser's soundness. Afterward, I formalize the \emph{shift-reduce pushdown
automaton} following the lecture slides of Petter~\cite{Petter}. I then prove the parser's
completeness by showing a relation between this automaton and the aforementioned stack words.%
\end{isamarkuptext}\isamarkuptrue%
%
\isadelimdocument
%
\endisadelimdocument
%
\isatagdocument
%
\isamarkupsection{Previous Work%
}
\isamarkuptrue%
%
\endisatagdocument
{\isafolddocument}%
%
\isadelimdocument
%
\endisadelimdocument
%
\begin{isamarkuptext}%
This thesis builds on the formalization of CFG theory by Nipkow et 
al.~\cite{Nipkow-AFP}, which is part of the Archive of Formal Proofs
(AFP)~\footnote{\url{https://www.isa-afp.org/}}, an archive of libraries and formal proof developments for
Isabelle. Many elements of this formalization, such as derivations (particularly rightmost 
derivations) and the elimination of unreachable and unproductive nonterminals, are essential for the
formal development of \isa{LR{\isacharparenleft}{\kern0pt}{\isadigit{0}}{\isacharparenright}{\kern0pt}} theory I present.

On the topic of \isa{LR} parsing in particular, Barthwal et al.~\cite{Barthwal} verified an executable
\isa{SLR} parser in the HOL4 proof assistant. Additionally, Jourdan et al.~\cite{Jourdan} formalized an
\isa{LR{\isacharparenleft}{\kern0pt}{\isadigit{1}}{\isacharparenright}{\kern0pt}} validator, which determines whether a given \isa{LR{\isacharparenleft}{\kern0pt}{\isadigit{1}}{\isacharparenright}{\kern0pt}} parser is correct w.r.t.\ a given CFG,
using the Coq Proof Assistant. 

In other areas of general parsing theory, Lasser et al.~\cite{Lasser} formalized an \isa{LL{\isacharparenleft}{\kern0pt}{\isadigit{1}}{\isacharparenright}{\kern0pt}} parser
generator, also using Coq; Blaudeau and Shankar~\cite{Blaudeau} formalized parsing expression grammars
(PEGs) and developed a recursive parser interpreter for these grammars in the PVS system, and lastly,
Rau and Nipkow~\cite{Rau} formalized an executable Early parser in Isabelle, which is able
to parse arbitrary CFGs.

This thesis presents the first formalization of \isa{LR} parsing theory in Isabelle, and is the first
to address \isa{LR{\isacharparenleft}{\kern0pt}{\isadigit{0}}{\isacharparenright}{\kern0pt}} parsing theory in particular.%
\end{isamarkuptext}\isamarkuptrue%
%
\isadelimdocument
%
\endisadelimdocument
%
\isatagdocument
%
\isamarkupsection{Isabelle Notation%
}
\isamarkuptrue%
%
\isamarkupsubsection{General Notation%
}
\isamarkuptrue%
%
\endisatagdocument
{\isafolddocument}%
%
\isadelimdocument
%
\endisadelimdocument
%
\begin{isamarkuptext}%
A term \isa{t} of type \isa{{\isasymtau}} is notated as \isa{t\ {\isacharcolon}{\kern0pt}{\isacharcolon}{\kern0pt}\ {\isasymtau}}, with type variables \isa{\isatfree{{\isacharprime}{\kern0pt}a}}, \isa{\isatfree{{\isacharprime}{\kern0pt}b}},
\isa{\isatfree{{\isacharprime}{\kern0pt}c}}, etc. Tuple types are notated using \isa{{\isasymtimes}}: for \isa{x\isactrlsub {\isadigit{0}}\ {\isacharcolon}{\kern0pt}{\isacharcolon}{\kern0pt}\ {\isacharprime}{\kern0pt}a\isactrlsub {\isadigit{0}}{\isacharcomma}{\kern0pt}\ x\isactrlsub {\isadigit{1}}\ {\isacharcolon}{\kern0pt}{\isacharcolon}{\kern0pt}\ {\isacharprime}{\kern0pt}a\isactrlsub {\isadigit{1}}{\isacharcomma}{\kern0pt}\ {\isasymdots}{\isacharcomma}{\kern0pt}\ x\isactrlsub n\ {\isacharcolon}{\kern0pt}{\isacharcolon}{\kern0pt}\ {\isacharprime}{\kern0pt}a\isactrlsub n}, we write
\isa{{\isacharparenleft}{\kern0pt}x\isactrlsub {\isadigit{0}}{\isacharcomma}{\kern0pt}\ x\isactrlsub {\isadigit{1}}{\isacharcomma}{\kern0pt}\ {\isasymdots}{\isacharcomma}{\kern0pt}\ x\isactrlsub n{\isacharparenright}{\kern0pt}\ {\isacharcolon}{\kern0pt}{\isacharcolon}{\kern0pt}\ {\isacharprime}{\kern0pt}a\isactrlsub {\isadigit{0}}\ {\isasymtimes}\ {\isacharprime}{\kern0pt}a\isactrlsub {\isadigit{1}}\ {\isasymtimes}\ {\isasymdots}\ {\isasymtimes}\ {\isacharprime}{\kern0pt}a\isactrlsub n}. We denote functions with the arrow \isa{{\isasymRightarrow}}, and we notate the
image of a set \isa{A} under function \isa{f} as \isa{\isafree{f}\ {\isacharbackquote}{\kern0pt}\ \isafree{A}}.

Type constructors are usually written postfix, as one can see in types such as
\isa{\isatfree{{\isacharprime}{\kern0pt}a}\ \isatconst{set}}, and in cases of multiple types \isa{{\isacharprime}{\kern0pt}a\isactrlsub {\isadigit{0}}{\isacharcomma}{\kern0pt}\ {\isacharprime}{\kern0pt}a\isactrlsub {\isadigit{1}}{\isacharcomma}{\kern0pt}\ {\isasymdots}{\isacharcomma}{\kern0pt}{\isacharprime}{\kern0pt}a\isactrlsub n}, they are written as 
\mbox{\isa{{\isacharparenleft}{\kern0pt}{\isacharprime}{\kern0pt}a\isactrlsub {\isadigit{0}}{\isacharcomma}{\kern0pt}\ {\isacharprime}{\kern0pt}a\isactrlsub {\isadigit{1}}{\isacharcomma}{\kern0pt}\ {\isasymdots}{\isacharcomma}{\kern0pt}\ {\isacharprime}{\kern0pt}a\isactrlsub n{\isacharparenright}{\kern0pt}}}. One of the most common types in this formalization, \isa{{\isacharparenleft}{\kern0pt}\isatfree{{\isacharprime}{\kern0pt}n}{\isacharcomma}{\kern0pt}\ \isatfree{{\isacharprime}{\kern0pt}t}{\isacharparenright}{\kern0pt}\ \isatconst{sym}}, 
is an example for this.

The keyword \isakeyword{datatype} is used to declare algebraic data types. This can be seen in the 
commonly used type of natural numbers \isa{\isatconst{nat}}, which we define recursively:
\begin{quote}
\isa{\isakeywordONE{datatype}\ \isatconst{nat}\ {\isacharequal}{\kern0pt}\ {\isadigit{0}}\ {\isacharbar}{\kern0pt}\ \isaconst{Suc}\ \isatconst{nat}}
\end{quote}
Another example is the type for lists, \isa{\isatfree{{\isacharprime}{\kern0pt}a}\ \isatconst{list}}, which we will also be using frequently:
\begin{quote}
\isa{\isakeywordONE{datatype}\ \isatfree{{\isacharprime}{\kern0pt}a}\ \isatconst{list}\ {\isacharequal}{\kern0pt}\ {\isacharbrackleft}{\kern0pt}{\isacharbrackright}{\kern0pt}\ {\isacharbar}{\kern0pt}\ {\isacharparenleft}{\kern0pt}{\isacharhash}{\kern0pt}{\isacharparenright}{\kern0pt}\ \isatfree{{\isacharprime}{\kern0pt}a}\ {\isacharparenleft}{\kern0pt}\isatfree{{\isacharprime}{\kern0pt}a}\ \isatconst{list}{\isacharparenright}{\kern0pt}}
\end{quote}

The list constructor \isa{{\isacharhash}{\kern0pt}} is usually written infix. An explicit list can be either
written as \isa{x\isactrlsub {\isadigit{0}}\ {\isacharhash}{\kern0pt}\ x\isactrlsub {\isadigit{1}}\ {\isacharhash}{\kern0pt}\ {\isasymdots}\ {\isacharhash}{\kern0pt}\ x\isactrlsub n\ {\isacharhash}{\kern0pt}\ {\isacharbrackleft}{\kern0pt}{\isacharbrackright}{\kern0pt}} or as \isa{{\isacharbrackleft}{\kern0pt}x\isactrlsub {\isadigit{0}}{\isacharcomma}{\kern0pt}\ x\isactrlsub {\isadigit{1}}{\isacharcomma}{\kern0pt}\ {\isasymdots}{\isacharcomma}{\kern0pt}\ x\isactrlsub n{\isacharbrackright}{\kern0pt}}. If 
\isa{xs\ {\isacharequal}{\kern0pt}\ y\ {\isacharhash}{\kern0pt}\ ys}, \isa{hd\ xs\ {\isacharequal}{\kern0pt}\ y} and \isa{\mbox{\isaconst{tl}\ \isafree{xs}}\ {\isacharequal}{\kern0pt}\ \isafree{ys}}. Furthermore, lists are concatenated
with the operator \isa{{\isacharat}{\kern0pt}}, \isa{\isaconst{rev}} reverses a list, \isa{\isaconst{set}} converts a list to a set, \isa{\isafree{xs}\ensuremath{_{[\mathit{\isafree{n}}]}}} returns the 
\isa{n}-th element of the list \isa{xs} (with 0-indexing), \isa{\isafree{xs}{\isacharbar}{\kern0pt}\isactrlbsub \isafree{n}\isactrlesub } is the prefix of length \isa{n} 
of \isa{xs}, and \isa{\isactrlbsub \isafree{n}\isactrlesub {\isacharbar}{\kern0pt}\isafree{xs}} is the suffix of \isa{xs} starting at index \isa{n}.
\isa{\isaconst{map}\ \isafree{f}\ \isafree{xs}} maps every element in \isa{xs\ {\isacharcolon}{\kern0pt}{\isacharcolon}{\kern0pt}\ {\isacharprime}{\kern0pt}a\ list} with the function \isa{f\ {\isacharcolon}{\kern0pt}{\isacharcolon}{\kern0pt}\ {\isacharprime}{\kern0pt}a\ {\isasymRightarrow}\ {\isacharprime}{\kern0pt}b}, i.e., 
\isa{map\ f\ {\isacharbrackleft}{\kern0pt}x\isactrlsub {\isadigit{0}}{\isacharcomma}{\kern0pt}\ {\isasymdots}{\isacharcomma}{\kern0pt}\ x\isactrlsub n{\isacharbrackright}{\kern0pt}} = \mbox{\isa{{\isacharbrackleft}{\kern0pt}f\ x\isactrlsub {\isadigit{0}}{\isacharcomma}{\kern0pt}\ {\isasymdots}{\isacharcomma}{\kern0pt}\ f\ x\isactrlsub n{\isacharbrackright}{\kern0pt}}}, and lastly, \isa{\isaconst{concat}} concatenates a list of 
lists: \isa{concat\ {\isacharbrackleft}{\kern0pt}xs\isactrlsub {\isadigit{0}}{\isacharcomma}{\kern0pt}\ {\isasymdots}{\isacharcomma}{\kern0pt}\ xs\isactrlsub n{\isacharbrackright}{\kern0pt}\ {\isacharequal}{\kern0pt}\ xs\isactrlsub {\isadigit{0}}\ {\isacharat}{\kern0pt}\ {\isasymdots}\ {\isacharat}{\kern0pt}\ xs\isactrlsub n}.

Pattern matching on a term \isa{{\isasymtau}} is done with the keywords \isakeyword{case} and \isakeyword{of}. For 
patterns \isa{{\isasympi}\isactrlsub {\isadigit{1}}{\isacharcomma}{\kern0pt}\ {\isasympi}\isactrlsub {\isadigit{2}}{\isacharcomma}{\kern0pt}\ {\isasymdots}{\isacharcomma}{\kern0pt}\ {\isasympi}\isactrlsub n} and expressions \isa{e\isactrlsub {\isadigit{1}}{\isacharcomma}{\kern0pt}\ e\isactrlsub {\isadigit{2}}{\isacharcomma}{\kern0pt}\ {\isasymdots}{\isacharcomma}{\kern0pt}\ e\isactrlsub n}, the expression
\begin{quote}
\isa{\textsf{case}\ \isafree{\isasymtau}\ \textsf{of}} \isa{{\isasympi}\isactrlsub {\isadigit{1}}\ {\isasymRightarrow}\ e\isactrlsub {\isadigit{1}}\ {\isacharbar}{\kern0pt}\ {\isasympi}\isactrlsub {\isadigit{2}}\ {\isasymRightarrow}\ e\isactrlsub {\isadigit{2}}\ {\isacharbar}{\kern0pt}\ {\isasymdots}\ {\isacharbar}{\kern0pt}\ {\isasympi}\isactrlsub n\ {\isasymRightarrow}\ e\isactrlsub n}
\end{quote}
returns \isa{e\isactrlsub i} if \isa{{\isasymtau}} matches \isa{{\isasympi}\isactrlsub i} and it does not match any \isa{{\isasympi}\isactrlsub j} with \isa{{\isadigit{0}}\ {\isacharless}{\kern0pt}\ j\ {\isacharless}{\kern0pt}\ i}. If \isa{{\isasymtau}} matches no
\isa{{\isasympi}\isactrlsub i}, \isa{\isaconst{undefined}} is returned. It is worth noting that variables introduced in \isa{{\isasympi}\isactrlsub i} are bound
on \isa{e\isactrlsub i}. 

A simple example on lists returning \isa{{\isacharbrackleft}{\kern0pt}{\isacharbrackright}{\kern0pt}} if the input is \isa{{\isacharbrackleft}{\kern0pt}{\isacharbrackright}{\kern0pt}}, and a
singleton list containing the first element otherwise:
\begin{quote}
\isa{\textsf{case}\ \isafree{xs}\ \textsf{of}\ {\isacharbrackleft}{\kern0pt}{\isacharbrackright}{\kern0pt}\ {\isasymRightarrow}\ {\isacharbrackleft}{\kern0pt}{\isacharbrackright}{\kern0pt}\ {\isacharbar}{\kern0pt}\ \isabound{y}\ {\isacharhash}{\kern0pt}\ \isabound{ys}\ {\isasymRightarrow}\ {\isacharbrackleft}{\kern0pt}\isabound{y}{\isacharbrackright}{\kern0pt}}
\end{quote}

If premises \isa{A\isactrlsub {\isadigit{1}}{\isacharcomma}{\kern0pt}\ A\isactrlsub {\isadigit{2}}{\isacharcomma}{\kern0pt}\ {\isasymdots}{\isacharcomma}{\kern0pt}\ A\isactrlsub n} imply \isa{B}, we write \mbox{\isa{{\isasymlbrakk}A\isactrlsub {\isadigit{1}}{\isacharsemicolon}{\kern0pt}\ A\isactrlsub {\isadigit{2}}{\isacharsemicolon}{\kern0pt}\ {\isasymdots}{\isacharsemicolon}{\kern0pt}\ A\isactrlsub n{\isasymrbrakk}\ {\isasymLongrightarrow}\ B}.}

When inducting on a reflexive transitive closure, we distinguish two types of induction:
with \emph{forward} induction (or simply ``induction'' without a qualifier), we induct by growing the
chain on the right, i.e., if \isa{r{\isacharasterisk}{\kern0pt}} is the reflexive transitive closure of \isa{r} and we assume \isa{r{\isacharasterisk}{\kern0pt}\ a\ c},
we show \isa{P\ a} for the base case. For the transitive case, we assume \isa{r{\isacharasterisk}{\kern0pt}\ a\ b}, \isa{r\ b\ c}, and \isa{P\ b}, 
and show that this implies \isa{P\ c}.

On the other hand, \emph{backwards} induction grows the chain on the left, meaning that for 
\isa{\mbox{\isafree{r}\ {\isacharasterisk}{\kern0pt}\ \isafree{a}\ \isafree{c}}},
we show \isa{P\ c} in the base case, and in the transitive case we assume \isa{r\ a\ b}, \isa{P\ b}, and \isa{r{\isacharasterisk}{\kern0pt}\ b\ c},
and we prove that \isa{P\ a} holds.

Lastly, the function \isa{fresh{\isadigit{0}}\ {\isacharcolon}{\kern0pt}{\isacharcolon}{\kern0pt}\ {\isacharprime}{\kern0pt}a\ set\ {\isasymRightarrow}\ {\isacharprime}{\kern0pt}a} takes a set, and returns a fresh item for this set,
i.e., some \isa{x\ {\isacharcolon}{\kern0pt}{\isacharcolon}{\kern0pt}\ {\isacharprime}{\kern0pt}a} that is not in the set passed to \isa{fresh{\isadigit{0}}}.%
\end{isamarkuptext}\isamarkuptrue%
%
\isadelimdocument
%
\endisadelimdocument
%
\isatagdocument
%
\isamarkupsubsection{Records and Finite Automata%
}
\isamarkuptrue%
%
\endisatagdocument
{\isafolddocument}%
%
\isadelimdocument
%
\endisadelimdocument
%
\begin{isamarkuptext}%
A record R whose record type has field names \isa{{\isasymphi}\isactrlsub {\isadigit{1}}{\isacharcomma}{\kern0pt}\ {\isasymphi}\isactrlsub {\isadigit{2}}{\isacharcomma}{\kern0pt}\ {\isasymdots}{\isacharcomma}{\kern0pt}\ {\isasymphi}\isactrlsub n} 
is defined through the notation
\[ \isa{R\ {\isacharequal}{\kern0pt}\ {\isasymlparr}{\isasymphi}\isactrlsub {\isadigit{1}}\ {\isacharequal}{\kern0pt}\ v\isactrlsub {\isadigit{1}}{\isacharcomma}{\kern0pt}\ {\isasymphi}\isactrlsub {\isadigit{2}}\ {\isacharequal}{\kern0pt}\ v\isactrlsub {\isadigit{2}}{\isacharcomma}{\kern0pt}\ {\isasymphi}\isactrlsub n\ {\isacharequal}{\kern0pt}\ v\isactrlsub n{\isasymrparr}} \]
if R has values \isa{v\isactrlsub {\isadigit{1}}{\isacharcomma}{\kern0pt}\ v\isactrlsub {\isadigit{2}}{\isacharcomma}{\kern0pt}\ {\isasymdots}{\isacharcomma}{\kern0pt}\ v\isactrlsub n} such that \isa{{\isasymphi}\isactrlsub i\ {\isacharequal}{\kern0pt}\ v\isactrlsub i} for every \isa{{\isadigit{1}}\ {\isasymle}\ i\ {\isasymle}\ n}. Furthermore, notation 
\isa{{\isasymphi}\isactrlsub i\ R} refers to the value of field \isa{{\isasymphi}\isactrlsub i} of \isa{R} (in this example, \isa{v\isactrlsub i}). Note that this 
causes a record field consisting of a function \isa{f\ {\isacharcolon}{\kern0pt}{\isacharcolon}{\kern0pt}\ {\isacharprime}{\kern0pt}a\ {\isasymRightarrow}\ {\isacharprime}{\kern0pt}b} to have an additional parameter for 
the record field, i.e., \isa{f\ {\isacharcolon}{\kern0pt}{\isacharcolon}{\kern0pt}\ {\isasymrho}\ {\isasymRightarrow}\ {\isacharprime}{\kern0pt}a\ {\isasymRightarrow}\ {\isacharprime}{\kern0pt}b}, where \isa{{\isasymrho}} is the record type of which \isa{f} is a field.

In this thesis, we work with multiple finite automata, which we define and implement on top of 
Paulson's formalization of both nondeterministic and deterministic finite automata
(NFAs and DFAs respectively)~\cite{Paulson}.

NFAs are defined by record \isa{{\isacharparenleft}{\kern0pt}\isatfree{{\isacharprime}{\kern0pt}a}{\isacharcomma}{\kern0pt}\ \isatfree{{\isacharprime}{\kern0pt}q}{\isacharparenright}{\kern0pt}\ \isatconst{nfa}} with alphabet type \isa{\isatfree{{\isacharprime}{\kern0pt}a}} and state type 
\isa{\isatfree{{\isacharprime}{\kern0pt}q}}, with the following fields:
\begin{itemize}
\item \isa{states\ {\isacharcolon}{\kern0pt}{\isacharcolon}{\kern0pt}\ {\isacharprime}{\kern0pt}q\ set}: a finite set of states.
\item \isa{init\ {\isacharcolon}{\kern0pt}{\isacharcolon}{\kern0pt}\ {\isacharprime}{\kern0pt}q\ set}: the set of initial states, with \isa{init\ {\isasymsubseteq}\ states}.
\item \isa{final\ {\isacharcolon}{\kern0pt}{\isacharcolon}{\kern0pt}\ {\isacharprime}{\kern0pt}q\ set}: the set of final states, also with \isa{final\ {\isasymsubseteq}\ states}.
\item \isa{nxt\ {\isacharcolon}{\kern0pt}{\isacharcolon}{\kern0pt}\ {\isacharprime}{\kern0pt}q\ {\isasymRightarrow}\ {\isacharprime}{\kern0pt}a\ {\isasymRightarrow}\ {\isacharprime}{\kern0pt}q\ set}: the transition function for reading transitions.
\item \isa{eps\ {\isacharcolon}{\kern0pt}{\isacharcolon}{\kern0pt}\ {\isacharparenleft}{\kern0pt}{\isacharprime}{\kern0pt}q\ {\isasymtimes}\ {\isacharprime}{\kern0pt}q{\isacharparenright}{\kern0pt}\ set}: the \isa{{\isasymepsilon}}-transition relation.
\end{itemize}

Furthermore, function \isa{nextl\ {\isacharcolon}{\kern0pt}{\isacharcolon}{\kern0pt}\ {\isacharprime}{\kern0pt}q\ set\ {\isasymRightarrow}\ {\isacharprime}{\kern0pt}a\ list\ {\isasymRightarrow}\ {\isacharprime}{\kern0pt}q\ set} for a NFA denotes the extension of 
\isa{nxt} to sets of states and words. 

DFAs are defined by a record \isa{{\isacharparenleft}{\kern0pt}\isatfree{{\isacharprime}{\kern0pt}a}{\isacharcomma}{\kern0pt}\ \isatfree{{\isacharprime}{\kern0pt}q}{\isacharparenright}{\kern0pt}\ \isatconst{dfa}} where \isa{\isatfree{{\isacharprime}{\kern0pt}a}} and \isa{\isatfree{{\isacharprime}{\kern0pt}q}} are again the 
alphabet and state types respectively. The \isa{dfa} record type has the same fields as 
\isa{nfa} except that it naturally lacks an \isa{{\isasymepsilon}}-transition relation. Moreover, in a DFA,
\isa{init\ {\isacharcolon}{\kern0pt}{\isacharcolon}{\kern0pt}\ {\isacharprime}{\kern0pt}q} is a unique initial state, and \isa{nxt\ {\isacharcolon}{\kern0pt}{\isacharcolon}{\kern0pt}\ {\isacharprime}{\kern0pt}q\ {\isasymRightarrow}\ {\isacharprime}{\kern0pt}a\ {\isasymRightarrow}\ {\isacharprime}{\kern0pt}q} returns a unique state as well.

We also define a \isa{nextl} function for DFAs, lifting the \isa{nxt} function of a DFA from symbols to 
words with \isa{nextl\ {\isacharcolon}{\kern0pt}{\isacharcolon}{\kern0pt}\ {\isacharprime}{\kern0pt}q\ {\isasymRightarrow}\ {\isacharprime}{\kern0pt}a\ list\ {\isasymRightarrow}\ {\isacharprime}{\kern0pt}q}.%
\end{isamarkuptext}\isamarkuptrue%
%
\isadelimdocument
%
\endisadelimdocument
%
\isatagdocument
%
\isamarkupsubsection{Context-Free Grammars%
}
\isamarkuptrue%
%
\endisatagdocument
{\isafolddocument}%
%
\isadelimdocument
%
\endisadelimdocument
%
\begin{isamarkuptext}%
In their formalization of context-free grammars, Nipkow et al.\ introduce type
\isa{{\isacharparenleft}{\kern0pt}\isatfree{{\isacharprime}{\kern0pt}n}{\isacharcomma}{\kern0pt}\ \isatfree{{\isacharprime}{\kern0pt}t}{\isacharparenright}{\kern0pt}\ \isatconst{sym}} for context-free grammar \concept{symbols} as a tagged union consisting of
nonterminals (\isa{\isaconst{Nt}}) and terminals (\isa{\isaconst{Tm}}) of type \isa{\isatfree{{\isacharprime}{\kern0pt}n}} and \isa{\isatfree{{\isacharprime}{\kern0pt}t}} respectively:
\begin{quote}
\isa{\isakeywordONE{datatype}\ {\isacharparenleft}{\kern0pt}\isatfree{{\isacharprime}{\kern0pt}n}{\isacharcomma}{\kern0pt}\ \isatfree{{\isacharprime}{\kern0pt}t}{\isacharparenright}{\kern0pt}\ \isatconst{sym}\ {\isacharequal}{\kern0pt}\ \isaconst{Nt}\ \isatfree{{\isacharprime}{\kern0pt}n}\ {\isacharbar}{\kern0pt}\ \isaconst{Tm}\ \isatfree{{\isacharprime}{\kern0pt}t}}
\end{quote}

Besides defining this type for symbols, they also define the following abbreviations:
\begin{quote}
\begin{tabular}{ll}
Lists of symbols & \isa{{\isacharparenleft}{\kern0pt}{\isacharprime}{\kern0pt}n{\isacharcomma}{\kern0pt}{\isacharprime}{\kern0pt}t{\isacharparenright}{\kern0pt}\ syms\ {\isacharequal}{\kern0pt}\ {\isacharparenleft}{\kern0pt}{\isacharprime}{\kern0pt}n{\isacharcomma}{\kern0pt}\ {\isacharprime}{\kern0pt}t{\isacharparenright}{\kern0pt}\ sym\ list}\\
Productions & \isa{{\isacharparenleft}{\kern0pt}{\isacharprime}{\kern0pt}n{\isacharcomma}{\kern0pt}{\isacharprime}{\kern0pt}t{\isacharparenright}{\kern0pt}\ prod\ {\isacharequal}{\kern0pt}\ {\isacharprime}{\kern0pt}n\ {\isasymtimes}\ {\isacharparenleft}{\kern0pt}{\isacharprime}{\kern0pt}n{\isacharcomma}{\kern0pt}{\isacharprime}{\kern0pt}t{\isacharparenright}{\kern0pt}\ syms}\\
Sets of productions & \isa{{\isacharparenleft}{\kern0pt}{\isacharprime}{\kern0pt}n{\isacharcomma}{\kern0pt}{\isacharprime}{\kern0pt}t{\isacharparenright}{\kern0pt}\ Prods\ {\isacharequal}{\kern0pt}\ {\isacharparenleft}{\kern0pt}{\isacharprime}{\kern0pt}n{\isacharcomma}{\kern0pt}{\isacharprime}{\kern0pt}t{\isacharparenright}{\kern0pt}\ prod\ set}
\end{tabular}
\end{quote}
where we informally write \isa{{\isacharparenleft}{\kern0pt}A{\isacharcomma}{\kern0pt}\ {\isasymalpha}{\isacharparenright}{\kern0pt}\ {\isacharcolon}{\kern0pt}{\isacharcolon}{\kern0pt}\ {\isacharparenleft}{\kern0pt}{\isacharprime}{\kern0pt}n{\isacharcomma}{\kern0pt}\ {\isacharprime}{\kern0pt}t{\isacharparenright}{\kern0pt}\ prod} as \isa{A\ {\isasymrightarrow}\ {\isasymalpha}}. For \isa{{\isasymalpha}\ {\isacharcolon}{\kern0pt}{\isacharcolon}{\kern0pt}\ {\isacharparenleft}{\kern0pt}{\isacharprime}{\kern0pt}n{\isacharcomma}{\kern0pt}\ {\isacharprime}{\kern0pt}t{\isacharparenright}{\kern0pt}\ syms}, 
\isa{\mbox{\isaconst{Nts{\isacharunderscore}{\kern0pt}syms}\ \isafree{{\isasymalpha}}}} returns the set of all \isa{X\ {\isacharcolon}{\kern0pt}{\isacharcolon}{\kern0pt}\ {\isacharprime}{\kern0pt}n} such that \isa{\isaconst{Nt}\ \isafree{X}\ {\isasymin}\ \isaconst{set}\ \isafree{{\isasymalpha}}}. Similarly 
for \isa{P\ {\isacharcolon}{\kern0pt}{\isacharcolon}{\kern0pt}\ {\isacharparenleft}{\kern0pt}{\isacharprime}{\kern0pt}n{\isacharcomma}{\kern0pt}\ {\isacharprime}{\kern0pt}t{\isacharparenright}{\kern0pt}\ Prods}, Nipkow et al.\ define \isa{\isaconst{Nts}\ \isafree{P}}:
\begin{equation*}
\isa{\isaconst{Nts}\ \isafree{P}\ {\isacharequal}{\kern0pt}\ {\isacharparenleft}{\kern0pt}{\isasymUnion}\isactrlbsub {\isacharparenleft}{\kern0pt}\isabound{A}{\isacharcomma}{\kern0pt}\ \isabound{{\isasymalpha}}{\isacharparenright}{\kern0pt}{\isasymin}\isafree{P}\isactrlesub \ {\isacharbraceleft}{\kern0pt}\isabound{A}{\isacharbraceright}{\kern0pt}\ {\isasymunion}\ \isaconst{Nts{\isacharunderscore}{\kern0pt}syms}\ \isabound{{\isasymalpha}}{\isacharparenright}{\kern0pt}}.
\end{equation*}

They further define a datatype for context-free grammars:
\begin{quote}
\isakeyword{datatype}\ \isa{{\isacharparenleft}{\kern0pt}{\isacharprime}{\kern0pt}n{\isacharcomma}{\kern0pt}\ {\isacharprime}{\kern0pt}t{\isacharparenright}{\kern0pt}\ Cfg\ {\isacharequal}{\kern0pt}\ Cfg\ {\isacharparenleft}{\kern0pt}{\isacharparenleft}{\kern0pt}{\isacharprime}{\kern0pt}n{\isacharcomma}{\kern0pt}{\isacharprime}{\kern0pt}t{\isacharparenright}{\kern0pt}\ Prods{\isacharparenright}{\kern0pt}\ {\isacharprime}{\kern0pt}n}
\end{quote}
where \isa{\isaconst{Cfg}\ \isafree{P}\ \isafree{S}} denotes a context-free grammar with production set \isa{P} and start symbol \isa{S}.
If \isa{\isafree{G}\ {\isacharequal}{\kern0pt}\ \isaconst{Cfg}\ \isafree{P}\ \isafree{S}}, \isa{\isaconst{Prods}\ \isafree{G}} refers to \isa{P}, and analogously, \isa{\isaconst{Start}\ \isafree{G}} refers to 
\isa{S}.

A derivation step from \isa{{\isasymphi}} to \isa{{\isasympsi}} under production set \isa{P} is notated as \isa{\isafree{P}\ {\isasymturnstile}\ \isafree{{\isasymphi}}\ {\isasymRightarrow}\ \isafree{{\isasympsi}}}.
More formally, for \isa{A\ {\isacharcolon}{\kern0pt}{\isacharcolon}{\kern0pt}\ {\isacharprime}{\kern0pt}n} and \isa{{\isasymalpha}{\isacharcomma}{\kern0pt}\ {\isasymbeta}{\isacharcomma}{\kern0pt}\ {\isasymgamma}\ {\isacharcolon}{\kern0pt}{\isacharcolon}{\kern0pt}\ {\isacharparenleft}{\kern0pt}{\isacharprime}{\kern0pt}n{\isacharcomma}{\kern0pt}\ {\isacharprime}{\kern0pt}t{\isacharparenright}{\kern0pt}\ syms}, Nipkow et al.\ define:
\begin{equation*} 
\isa{{\isacharparenleft}{\kern0pt}\isafree{A}{\isacharcomma}{\kern0pt}\ \isafree{{\isasymbeta}}{\isacharparenright}{\kern0pt}\ {\isasymin}\ \isafree{P}\ {\isasymLongrightarrow}\ \isafree{P}\ {\isasymturnstile}\ \isafree{{\isasymalpha}}\ {\isacharat}{\kern0pt}\ {\isacharbrackleft}{\kern0pt}\isaconst{Nt}\ \isafree{A}{\isacharbrackright}{\kern0pt}\ {\isacharat}{\kern0pt}\ \isafree{{\isasymgamma}}\ {\isasymRightarrow}\ \isafree{{\isasymalpha}}\ {\isacharat}{\kern0pt}\ \isafree{{\isasymbeta}}\ {\isacharat}{\kern0pt}\ \isafree{{\isasymgamma}}}.
\end{equation*}
Similarly, they denote the reflexive transitive closure of derivations by \isa{\isafree{P}\ {\isasymturnstile}\ \isafree{{\isasymphi}}\ {\isasymRightarrow}{\isacharasterisk}{\kern0pt}\ \isafree{{\isasympsi}}},
and derivations of length \isa{n} by \isa{\isafree{P}\ {\isasymturnstile}\ \isafree{{\isasymphi}}\ {\isasymRightarrow}{\isacharparenleft}{\kern0pt}\isafree{n}{\isacharparenright}{\kern0pt}\ \isafree{{\isasympsi}}}. Rightmost derivations are notated analogously, 
with \isa{{\isasymRightarrow}r}, \isa{{\isasymRightarrow}r{\isacharasterisk}{\kern0pt}} and \isa{{\isasymRightarrow}r{\isacharparenleft}{\kern0pt}n{\isacharparenright}{\kern0pt}} respectively.

Lastly, they define the language of a nonterminal w.r.t.\ a set of productions
\begin{equation*}
\isa{\isaconst{Lang}\ \isavar{P}\ \isavar{A}\ {\isacharequal}{\kern0pt}\ {\isacharbraceleft}{\kern0pt}\isabound{w}\ {\isacharbar}{\kern0pt}\ \isavar{P}\ {\isasymturnstile}\ {\isacharbrackleft}{\kern0pt}\isaconst{Nt}\ \isavar{A}{\isacharbrackright}{\kern0pt}\ {\isasymRightarrow}{\isacharasterisk}{\kern0pt}\ \isaconst{map}\ \isaconst{Tm}\ \isabound{w}{\isacharbraceright}{\kern0pt}},
\end{equation*}
and based on this, the language of a CFG \isa{G}
\begin{equation*}
\isa{L{\isacharparenleft}{\kern0pt}\isafree{G}{\isacharparenright}{\kern0pt}} = \isa{\isaconst{Lang}\ {\isacharparenleft}{\kern0pt}\isaconst{Prods}\ \isafree{G}{\isacharparenright}{\kern0pt}\ {\isacharparenleft}{\kern0pt}\isaconst{Start}\ \isafree{G}{\isacharparenright}{\kern0pt}}.
\end{equation*}
 
Besides type variables \isa{\isatfree{{\isacharprime}{\kern0pt}n}} for nonterminals and \isa{\isatfree{{\isacharprime}{\kern0pt}t}} for terminals, we use the following 
variable conventions: for brevity, we refer to \isa{{\isacharparenleft}{\kern0pt}{\isacharprime}{\kern0pt}n{\isacharcomma}{\kern0pt}\ {\isacharprime}{\kern0pt}t{\isacharparenright}{\kern0pt}\ sym} and \isa{{\isacharparenleft}{\kern0pt}{\isacharprime}{\kern0pt}n{\isacharcomma}{\kern0pt}\ {\isacharprime}{\kern0pt}t{\isacharparenright}{\kern0pt}\ syms} simply as \isa{sym}
and \isa{syms} respectively; \isa{A{\isacharcomma}{\kern0pt}\ B{\isacharcomma}{\kern0pt}\ C\ {\isacharcolon}{\kern0pt}{\isacharcolon}{\kern0pt}\ {\isacharprime}{\kern0pt}n}; \isa{a{\isacharcomma}{\kern0pt}\ b{\isacharcomma}{\kern0pt}\ c\ {\isacharcolon}{\kern0pt}{\isacharcolon}{\kern0pt}\ {\isacharprime}{\kern0pt}t}; \isa{u{\isacharcomma}{\kern0pt}\ v{\isacharcomma}{\kern0pt}\ w\ {\isacharcolon}{\kern0pt}{\isacharcolon}{\kern0pt}\ {\isacharprime}{\kern0pt}t\ list}; and finally
\isa{{\isasymalpha}{\isacharcomma}{\kern0pt}\ {\isasymbeta}{\isacharcomma}{\kern0pt}\ {\isasymgamma}\ {\isacharcolon}{\kern0pt}{\isacharcolon}{\kern0pt}\ {\isacharparenleft}{\kern0pt}{\isacharprime}{\kern0pt}n{\isacharcomma}{\kern0pt}\ {\isacharprime}{\kern0pt}t{\isacharparenright}{\kern0pt}\ syms}. Further definitions and detailed explanations on the formalization can be 
found on the AFP entry by Nipkow et al.~\cite{Nipkow-AFP}, as well as their paper on this
formalization~\cite{Nipkow-Paper}.%
\end{isamarkuptext}\isamarkuptrue%
%
\isadelimdocument
%
\endisadelimdocument
%
\isatagdocument
%
\isamarkupsection{Overview%
}
\isamarkuptrue%
%
\endisatagdocument
{\isafolddocument}%
%
\isadelimdocument
%
\endisadelimdocument
%
\begin{isamarkuptext}%
The thesis follows a similar structure to Wilhelm et al.\ First, we introduce some foundational
definitions and assumptions regarding the CFGs we will be working with. We then define the item
pushdown automaton, followed by a new section not present in the book: a formalized decomposition of
rightmost derivations. Afterward, we return to the book's structure with the characteristic finite
automaton, the canonical \isa{LR{\isacharparenleft}{\kern0pt}{\isadigit{0}}{\isacharparenright}{\kern0pt}} automaton, and lastly, the \isa{LR{\isacharparenleft}{\kern0pt}{\isadigit{0}}{\isacharparenright}{\kern0pt}} parser. In the parser's section,
we first finish the formal verification of Wilhelm et al., and we conclude with the
proof of the parser's correctness.%
\end{isamarkuptext}\isamarkuptrue%
%
\isadelimdocument
%
\endisadelimdocument
%
\isatagdocument
%
\isamarkupchapter{Basic Definitions%
}
\isamarkuptrue%
%
\isamarkupsection{Extended and Reduced Grammars%
}
\isamarkuptrue%
%
\isamarkupsubsection{Extending Grammars by a New Start Symbol%
}
\isamarkuptrue%
%
\endisatagdocument
{\isafolddocument}%
%
\isadelimdocument
%
\endisadelimdocument
%
\isadelimproof
%
\endisadelimproof
%
\isatagproof
%
\endisatagproof
{\isafoldproof}%
%
\isadelimproof
%
\endisadelimproof
%
\begin{isamarkuptext}%
Throughout this thesis, we will define several automata which will aid us in our goal of 
constructing a deterministic \isa{LR{\isacharparenleft}{\kern0pt}{\isadigit{0}}{\isacharparenright}{\kern0pt}} parser for a context-free grammar \isa{G}. Naturally, our 
definitions will be defined based on the productions of the grammar, and the starting state of all 
these automata relate to the start symbol \isa{S} of \isa{G}, specifically the productions where 
\isa{\mbox{\isaconst{S}}} is
on the left-hand side. In general, however, a context-free grammar's start symbol can appear on the right-hand 
side of a production, and in many cases we would like to be able to identify that a word has been 
derived when a production of the form \isa{{\isacharparenleft}{\kern0pt}\isaconst{S}{\isacharcomma}{\kern0pt}\ \isafree{{\isasymalpha}}{\isacharparenright}{\kern0pt}} is encountered. If \isa{\isaconst{S}} is
on the RHS of a production, this is not possible, since such a production could be encountered because
of an instance of \isa{S} that was produced at some point.

We therefore extend a context-free grammar \isa{G} with finite set of productions \isa{P} by
a fresh start symbol \isa{S{\isacharprime}{\kern0pt}} and a single production \isa{{\isacharparenleft}{\kern0pt}S{\isacharprime}{\kern0pt}{\isacharcomma}{\kern0pt}\ {\isacharbrackleft}{\kern0pt}Nt\ S{\isacharbrackright}{\kern0pt}{\isacharparenright}{\kern0pt}}. With this, we are guaranteed that
our new start symbol is never produced, and we can identify a complete derivation by this unique
production. The resulting grammar, which we define to be \isa{G{\isacharprime}{\kern0pt}}, is the \concept{extended grammar}, or the \concept{extension}, of \isa{G}. We 
analogously refer to \isa{Prods\ G{\isacharprime}{\kern0pt}}, which we abbreviate as \isa{P{\isacharprime}{\kern0pt}}, as the extension of \isa{P} or the 
\concept{extended set of productions} of \isa{G}. Formally:
\begin{gather*}
\isa{\isaconst{S{\isacharprime}{\kern0pt}}\ {\isasymequiv}\ \isaconst{fresh{\isadigit{0}}}\ {\isacharparenleft}{\kern0pt}\isaconst{Nts}\ \isaconst{P}\ {\isasymunion}\ {\isacharbraceleft}{\kern0pt}\isaconst{S}{\isacharbraceright}{\kern0pt}{\isacharparenright}{\kern0pt}}\\
\isa{\isaconst{G{\isacharprime}{\kern0pt}}\ {\isasymequiv}\ \isaconst{Cfg}\ {\isacharparenleft}{\kern0pt}\isaconst{P}\ {\isasymunion}\ {\isacharbraceleft}{\kern0pt}{\isacharparenleft}{\kern0pt}\isaconst{S{\isacharprime}{\kern0pt}}{\isacharcomma}{\kern0pt}\ {\isacharbrackleft}{\kern0pt}\isaconst{Nt}\ \isaconst{S}{\isacharbrackright}{\kern0pt}{\isacharparenright}{\kern0pt}{\isacharbraceright}{\kern0pt}{\isacharparenright}{\kern0pt}\ \isaconst{S{\isacharprime}{\kern0pt}}}
\end{gather*}

We now prove that extending the grammar preserves the language. 

\begin{lemma}\label{G'_deriven_Suc_imp_G_deriven}
If there exists a derivation in \isa{G{\isacharprime}{\kern0pt}}
\[ \isa{\isaconst{P{\isacharprime}{\kern0pt}}\ {\isasymturnstile}\ {\isacharbrackleft}{\kern0pt}\isaconst{Nt}\ \isaconst{S{\isacharprime}{\kern0pt}}{\isacharbrackright}{\kern0pt}\ {\isasymRightarrow}{\isacharparenleft}{\kern0pt}\isaconst{Suc}\ \isafree{n}{\isacharparenright}{\kern0pt}\ \isafree{{\isasymbeta}}}, \]
then there also exists a derivation in \isa{G}
\[ \isa{\isaconst{P}\ {\isasymturnstile}\ {\isacharbrackleft}{\kern0pt}\isaconst{Nt}\ \isaconst{S}{\isacharbrackright}{\kern0pt}\ {\isasymRightarrow}{\isacharparenleft}{\kern0pt}\isafree{n}{\isacharparenright}{\kern0pt}\ \isafree{{\isasymbeta}}}. \]
\begin{proof} 
We induct on \isa{n}. If \isa{n\ {\isacharequal}{\kern0pt}\ {\isadigit{0}}}, \isa{{\isasymbeta}\ {\isacharequal}{\kern0pt}\ {\isacharbrackleft}{\kern0pt}Nt\ S{\isacharbrackright}{\kern0pt}}, and the implication holds.

If \isa{n\ {\isacharequal}{\kern0pt}\ Suc\ m} for some \isa{m}, the induction hypothesis tells us that sentential form \isa{{\isasymalpha}} that 
derives \isa{{\isasymbeta}} in the final derivation step can be derived by \isa{S} in \isa{G} in \isa{m} steps. This means that 
all nonterminals in \isa{{\isasymalpha}} are in \isa{G}. Since the only production in \isa{P{\isacharprime}{\kern0pt}\ {\isacharbackslash}{\kern0pt}\ P} has 
\isa{S{\isacharprime}{\kern0pt}\ {\isasymnotin}\ Nts\ P} in the LHS, the final derivation step \isa{\isaconst{P{\isacharprime}{\kern0pt}}\ {\isasymturnstile}\ \isafree{{\isasymalpha}}\ {\isasymRightarrow}\ \isafree{{\isasymbeta}}} also 
exists in \isa{G}, completing the proof.
\end{proof}
\end{lemma}

\begin{lemma}\label{G'_derivern_Suc_imp_G_derivern}
If there exists a derivation in \isa{G{\isacharprime}{\kern0pt}}
\[ \isa{\isaconst{P{\isacharprime}{\kern0pt}}\ {\isasymturnstile}\ {\isacharbrackleft}{\kern0pt}\isaconst{Nt}\ \isaconst{S{\isacharprime}{\kern0pt}}{\isacharbrackright}{\kern0pt}\ {\isasymRightarrow}r{\isacharparenleft}{\kern0pt}\isaconst{Suc}\ \isafree{n}{\isacharparenright}{\kern0pt}\ \isafree{{\isasymbeta}}}, \]
then there also exists a derivation in \isa{G}
\[ \isa{\isaconst{P}\ {\isasymturnstile}\ {\isacharbrackleft}{\kern0pt}\isaconst{Nt}\ \isaconst{S}{\isacharbrackright}{\kern0pt}\ {\isasymRightarrow}r{\isacharparenleft}{\kern0pt}\isafree{n}{\isacharparenright}{\kern0pt}\ \isafree{{\isasymbeta}}}. \]
\begin{proof}
Analogous to the proof of Lemma~\ref{G'_deriven_Suc_imp_G_deriven}.
\end{proof}
\end{lemma}


\begin{theorem}[Preservation of language]\label{Lang_preserved}
\isa{L{\isacharparenleft}{\kern0pt}\isaconst{G{\isacharprime}{\kern0pt}}{\isacharparenright}{\kern0pt}\ {\isacharequal}{\kern0pt}\ L{\isacharparenleft}{\kern0pt}\isafree{G}{\isacharparenright}{\kern0pt}}
\begin{proof} 
Let \isa{\isafree{w}\ {\isasymin}\ L{\isacharparenleft}{\kern0pt}\isaconst{G{\isacharprime}{\kern0pt}}{\isacharparenright}{\kern0pt}}. Then there exists a derivation of \isa{\isaconst{map}\ \isaconst{Tm}\ \isafree{w}} of nonzero length,
i.e., of the form
\[\isa{\isaconst{P{\isacharprime}{\kern0pt}}\ {\isasymturnstile}\ {\isacharbrackleft}{\kern0pt}\isaconst{Nt}\ \isaconst{S{\isacharprime}{\kern0pt}}{\isacharbrackright}{\kern0pt}\ {\isasymRightarrow}{\isacharparenleft}{\kern0pt}\isaconst{Suc}\ \isafree{n}{\isacharparenright}{\kern0pt}\ \isaconst{map}\ \isaconst{Tm}\ \isafree{w}} \]
for some \isa{n}. Lemma~\ref{G'_deriven_Suc_imp_G_deriven} then implies \isa{\isafree{w}\ {\isasymin}\ L{\isacharparenleft}{\kern0pt}\isafree{G}{\isacharparenright}{\kern0pt}}.

Conversely, let \isa{\isafree{w}\ {\isasymin}\ L{\isacharparenleft}{\kern0pt}\isafree{G}{\isacharparenright}{\kern0pt}}. Then there exists a derivation
\[ \isa{\isaconst{P}\ {\isasymturnstile}\ {\isacharbrackleft}{\kern0pt}\isaconst{Nt}\ \isaconst{S}{\isacharbrackright}{\kern0pt}\ {\isasymRightarrow}{\isacharasterisk}{\kern0pt}\ \isaconst{map}\ \isaconst{Tm}\ \isafree{w}}. \] 
Since \isa{\isaconst{P}\ {\isasymsubseteq}\ \isaconst{P{\isacharprime}{\kern0pt}}}, this derivation also exists in \isa{G{\isacharprime}{\kern0pt}}. This, along with the fact
that \isa{S{\isacharprime}{\kern0pt}} derives \isa{S} in \isa{G{\isacharprime}{\kern0pt}}, implies that \isa{\isafree{w}\ {\isasymin}\ L{\isacharparenleft}{\kern0pt}\isaconst{G{\isacharprime}{\kern0pt}}{\isacharparenright}{\kern0pt}}.
\end{proof}
\end{theorem}

We can therefore use \isa{G{\isacharprime}{\kern0pt}} construct our parser in order to recognize the language of \isa{G}.%
\end{isamarkuptext}\isamarkuptrue%
%
\isadelimdocument
%
\endisadelimdocument
%
\isatagdocument
%
\isamarkupsubsection{Reduced Grammars%
}
\isamarkuptrue%
%
\endisatagdocument
{\isafolddocument}%
%
\isadelimdocument
%
\endisadelimdocument
%
\begin{isamarkuptext}%
In general, CFGs can contain problematic nonterminals which can be removed from the grammar 
without altering the language. Working with grammars that lack such nonterminals is ideal, since 
having them increases computational complexity and makes the grammar less well-behaved.

\begin{example}\label{ex:useless symbols}
Let \isa{G} be a CFG with \isa{S\ {\isacharequal}{\kern0pt}\ Start\ G} and productions:
\begin{center}
\begin{tabular}{cc}
\isa{S\ {\isasymrightarrow}\ A\ {\isacharbar}{\kern0pt}\ AB} & \isa{A\ {\isasymrightarrow}\ aA\ {\isacharbar}{\kern0pt}\ a}\\
\isa{C\ {\isasymrightarrow}\ ac\ {\isacharbar}{\kern0pt}\ BCD} & \isa{D\ {\isasymrightarrow}\ BC\ {\isacharbar}{\kern0pt}\ D}
\end{tabular}
\end{center}
Each nonterminal except for \isa{S} and \isa{A} carries problems with it:
\begin{itemize}
\item There are no productions where \isa{B} is on the left-hand side. This means that if \isa{S} 
reaches a sentential form \isa{{\isasymalpha}} such that \isa{Nt\ B\ {\isasymin}\ set\ {\isasymalpha}}, no word will be derived from \isa{{\isasymalpha}}. 
\item \isa{S} cannot reach \isa{C}, meaning no productions containing \isa{C}, or reachable only 
through \isa{C} (e.g. \isa{D}, which can only be reached using production \mbox{\isa{C\ {\isasymrightarrow}\ BCD}}), can be used to
derive words in \isa{L{\isacharparenleft}{\kern0pt}\isafree{G}{\isacharparenright}{\kern0pt}}.
\item \isa{D}, as opposed to \isa{B}, does show up on the LHS of certain productions, but none of these productions
can lead to a word: \isa{D\ {\isasymrightarrow}\ BC} contains \isa{B}, which cannot derive a \isa{{\isacharprime}{\kern0pt}t\ list}, and \isa{D\ {\isasymrightarrow}\ D} has no effect.
Furthermore, similarly to \isa{C}, \isa{D} cannot be reached by \isa{S}.
\end{itemize}
\end{example}

Nipkow et al.\ define \concept{useful} nonterminals w.r.t.\ a set of productions and a start symbol:
\begin{gather*}
\isa{\isaconst{productives}\ {\isacharparenleft}{\kern0pt}\isaconst{Prods}\ \isafree{G}{\isacharparenright}{\kern0pt}\ \isafree{{\isasymalpha}}\ {\isasymequiv}\ {\isasymexists}\isabound{w}{\isachardot}{\kern0pt}\ \isaconst{Prods}\ \isafree{G}\ {\isasymturnstile}\ \isafree{{\isasymalpha}}\ {\isasymRightarrow}{\isacharasterisk}{\kern0pt}\ \isaconst{map}\ \isaconst{Tm}\ \isabound{w}}\\
\begin{multlined}
\isa{useful\ P\ S\ A\ {\isacharequal}{\kern0pt}}\ (\isa{{\isasymexists}\isabound{{\isasymbeta}}{\isachardot}{\kern0pt}\ \isaconst{P}\ {\isasymturnstile}\ {\isacharbrackleft}{\kern0pt}\isaconst{Nt}\ \isaconst{S}{\isacharbrackright}{\kern0pt}\ {\isasymRightarrow}{\isacharasterisk}{\kern0pt}\ \isabound{{\isasymbeta}}}\\
  {} \wedge \isa{\isaconst{Nt}\ \isafree{A}\ {\isasymin}\ \isaconst{set}\ \isafree{{\isasymbeta}}\ {\isasymand}\ \isaconst{productives}\ \isaconst{P}\ \isafree{{\isasymbeta}}})
\end{multlined}
\end{gather*}
For a CFG \isa{G} with productions \isa{P}, \isa{A\ {\isacharcolon}{\kern0pt}{\isacharcolon}{\kern0pt}\ {\isacharprime}{\kern0pt}n} is \concept{reachable} if there exists a \isa{{\isasymbeta}\ {\isacharcolon}{\kern0pt}{\isacharcolon}{\kern0pt}\ syms}
such that \mbox{\isa{A\ {\isasymin}\ set\ {\isasymbeta}}} and \isa{\isaconst{P}\ {\isasymturnstile}\ {\isacharbrackleft}{\kern0pt}\isaconst{Nt}\ \isaconst{S}{\isacharbrackright}{\kern0pt}\ {\isasymRightarrow}{\isacharasterisk}{\kern0pt}\ \isafree{{\isasymbeta}}}. Otherwise, it is
\concept{unreachable}. Similarly, it is \concept{productive} if \isa{productives\ {\isacharparenleft}{\kern0pt}Prods\ G{\isacharparenright}{\kern0pt}\ {\isacharbrackleft}{\kern0pt}Nt\ A{\isacharbrackright}{\kern0pt}} holds, as 
defined by Nipkow et al.\ as well:
\begin{equation*}
\isa{\isaconst{productive}\ {\isacharparenleft}{\kern0pt}\isaconst{Prods}\ \isafree{G}{\isacharparenright}{\kern0pt}\ \isafree{A}\ {\isasymequiv}\ \isaconst{productives}\ {\isacharparenleft}{\kern0pt}\isaconst{Prods}\ \isafree{G}{\isacharparenright}{\kern0pt}\ {\isacharbrackleft}{\kern0pt}\isaconst{Nt}\ \isafree{A}{\isacharbrackright}{\kern0pt}}.
\end{equation*}
Similarly to reachable terminals, a nonterminal that is not productive is \concept{unproductive}. A 
useful nonterminal is therefore a nonterminal that is both reachable and productive.

Nipkow et al.\ have also proved that removing all nonterminals that are unreachable or unproductive, 
i.e. all non-useful nonterminals, preserves the language~\cite[Lemma Lang\_restrict\_useful]{Nipkow-AFP}:
\begin{lemma}
Let 
\[\isa{restrict{\isacharunderscore}{\kern0pt}Nts\ R\ P\ {\isacharequal}{\kern0pt}\ {\isacharbraceleft}{\kern0pt}{\isacharparenleft}{\kern0pt}A{\isacharcomma}{\kern0pt}{\isasymalpha}{\isacharparenright}{\kern0pt}\ {\isasymin}\ P{\isachardot}{\kern0pt}\ {\isasymforall}B\ {\isasymin}\ {\isacharbraceleft}{\kern0pt}A{\isacharbraceright}{\kern0pt}\ {\isasymunion}\ Nts{\isacharunderscore}{\kern0pt}syms\ {\isasymalpha}{\isachardot}{\kern0pt}\ R\ B{\isacharbraceright}{\kern0pt}}. \]
Then \isa{\isaconst{Lang}\ {\isacharparenleft}{\kern0pt}\isaconst{restrict{\isacharunderscore}{\kern0pt}Nts}\ {\isacharparenleft}{\kern0pt}\isaconst{useful}\ \isavar{P}\ \isavar{S}{\isacharparenright}{\kern0pt}\ \isavar{P}{\isacharparenright}{\kern0pt}\ \isavar{S}\ {\isacharequal}{\kern0pt}\ \isaconst{Lang}\ \isavar{P}\ \isavar{S}}.\qed
\end{lemma}

Consider Example~\ref{ex:useless symbols} once again. None of the nonterminals \isa{B{\isacharcomma}{\kern0pt}\ C{\isacharcomma}{\kern0pt}} and \isa{D} is 
useful:
\begin{itemize}
\item \isa{B} is reachable but unproductive
\item \isa{C} is productive but unreachable
\item \isa{D} is both unreachable and unproductive.
\end{itemize}
If we remove all productions containing non-useful nonterminals from this grammar, the productions 
that remain are:
\begin{center}
\begin{tabular}{cc}
\isa{S\ {\isasymrightarrow}\ A} & \isa{A\ {\isasymrightarrow}\ aA\ {\isacharbar}{\kern0pt}\ a}
\end{tabular}
\end{center}

As we can see, by applying this restriction to arbitrary grammars, the resulting set of 
productions can potentially become much smaller than the original one. This also guarantees that all
nonterminals are more well-behaved; for example, since every nonterminal is productive, we know 
that any part sentential form that can be derived from \isa{S} contains only productive nonterminals and can
therefore derive a word in the language.

With this, we now define the notion of \concept{reduced grammars}.
\begin{definition}[Reduced grammar]
\begin{equation*}
\isa{\isaconst{reduced}\ \isavar{G}\ {\isasymequiv}\ {\isasymforall}\isabound{A}{\isasymin}\isaconst{Nts}\ {\isacharparenleft}{\kern0pt}\isaconst{Prods}\ \isavar{G}{\isacharparenright}{\kern0pt}{\isachardot}{\kern0pt}\ \isaconst{useful}\ {\isacharparenleft}{\kern0pt}\isaconst{Prods}\ \isavar{G}{\isacharparenright}{\kern0pt}\ {\isacharparenleft}{\kern0pt}\isaconst{Start}\ \isavar{G}{\isacharparenright}{\kern0pt}\ \isabound{A}}.
\end{equation*}
\end{definition}

\begin{lemma}[Preservation of reduction]\label{G'_reduced}
If \isa{G} is reduced and \isa{L{\isacharparenleft}{\kern0pt}\isafree{G}{\isacharparenright}{\kern0pt}} is nonempty, \isa{G{\isacharprime}{\kern0pt}} is reduced.
\begin{proof}
Since \isa{G} is reduced, all nonterminals in \isa{P} are useful, and the fact that \isa{L{\isacharparenleft}{\kern0pt}\isafree{G}{\isacharparenright}{\kern0pt}\ {\isasymnoteq}\ {\isasymemptyset}}
implies that there exist \isa{{\isasymalpha}\ {\isacharcolon}{\kern0pt}{\isacharcolon}{\kern0pt}\ syms} and \isa{w\ {\isasymin}\ LangS\ G} such that
\begin{equation*}
\isa{P\ {\isasymturnstile}\ {\isacharbrackleft}{\kern0pt}Nt\ S{\isacharbrackright}{\kern0pt}\ {\isasymRightarrow}\ {\isasymalpha}\ {\isasymRightarrow}{\isacharasterisk}{\kern0pt}\ map\ Tm\ w}.
\end{equation*}
This implies that \isa{S\ {\isasymin}\ Nts\ P}. Since \isa{P\ {\isasymsubseteq}\ P{\isacharprime}{\kern0pt}}, this
means that all nonterminals in \isa{Nts\ P} are useful in \isa{P{\isacharprime}{\kern0pt}}. Therefore, to show that 
\isa{G{\isacharprime}{\kern0pt}} is reduced, it suffices to show that \isa{S{\isacharprime}{\kern0pt}} is useful in \isa{P{\isacharprime}{\kern0pt}}, i.e., reachable and productive. 
Reachability is trivial by reflexivity. To show that it is productive, we need to show 
that there exists a \isa{w\ {\isacharcolon}{\kern0pt}{\isacharcolon}{\kern0pt}\ {\isacharprime}{\kern0pt}t\ list} such that \isa{\isaconst{P{\isacharprime}{\kern0pt}}\ {\isasymturnstile}\ {\isacharbrackleft}{\kern0pt}\isaconst{Nt}\ \isaconst{S{\isacharprime}{\kern0pt}}{\isacharbrackright}{\kern0pt}\ {\isasymRightarrow}{\isacharasterisk}{\kern0pt}\ \isaconst{map}\ \isaconst{Tm}\ \isafree{w}}, which is 
equivalent to showing there exists some \isa{w\ {\isasymin}\ LangS\ G{\isacharprime}{\kern0pt}}. This follows directly from our assumption 
that \isa{L{\isacharparenleft}{\kern0pt}\isafree{G}{\isacharparenright}{\kern0pt}\ {\isasymnoteq}\ {\isasymemptyset}} and Theorem~\ref{Lang_preserved}.
\end{proof}
\end{lemma}

With this, we have now defined a way to extend any CFG by a new start symbol while preserving both 
language and the usefulness of all nonterminals.%
\end{isamarkuptext}\isamarkuptrue%
%
\isadelimdocument
%
\endisadelimdocument
%
\isatagdocument
%
\isamarkupsection{Context-Free Items%
}
\isamarkuptrue%
%
\endisatagdocument
{\isafolddocument}%
%
\isadelimdocument
%
\endisadelimdocument
%
\begin{isamarkuptext}%
\begin{definition}[Context-free item]
A \concept{context-free item} for a CFG \isa{G}, denoted by type \isa{{\isacharparenleft}{\kern0pt}\isatfree{{\isacharprime}{\kern0pt}n}{\isacharcomma}{\kern0pt}\ \isatfree{{\isacharprime}{\kern0pt}t}{\isacharparenright}{\kern0pt}\ \isatconst{item}}, is a triple 
\mbox{\isa{{\isacharparenleft}{\kern0pt}A{\isacharcomma}{\kern0pt}\ {\isasymalpha}{\isacharcomma}{\kern0pt}\ {\isasymbeta}{\isacharparenright}{\kern0pt}\ {\isacharcolon}{\kern0pt}{\isacharcolon}{\kern0pt}\ {\isacharprime}{\kern0pt}n\ {\isasymtimes}\ {\isacharparenleft}{\kern0pt}{\isacharprime}{\kern0pt}n{\isacharcomma}{\kern0pt}\ {\isacharprime}{\kern0pt}t{\isacharparenright}{\kern0pt}\ syms\ {\isasymtimes}\ {\isacharparenleft}{\kern0pt}{\isacharprime}{\kern0pt}n{\isacharcomma}{\kern0pt}\ {\isacharprime}{\kern0pt}t{\isacharparenright}{\kern0pt}\ syms}} such that 
\isa{{\isacharparenleft}{\kern0pt}\isafree{A}{\isacharcomma}{\kern0pt}\ \isafree{{\isasymalpha}}\ {\isacharat}{\kern0pt}\ \isafree{{\isasymbeta}}{\isacharparenright}{\kern0pt}\ {\isasymin}\ \isaconst{Prods}\ \isafree{G}}. We write the item \isa{{\isacharparenleft}{\kern0pt}A{\isacharcomma}{\kern0pt}\ {\isasymalpha}{\isacharcomma}{\kern0pt}\ {\isasymbeta}{\isacharparenright}{\kern0pt}} as \isa{{\isacharbrackleft}{\kern0pt}\isafree{A}\ {\isasymrightarrow}\ \isafree{{\isasymalpha}}\ {\isasymcdot}\ \isafree{{\isasymbeta}}{\isacharbrackright}{\kern0pt}}, and akin to
\isa{sym} and \isa{syms}, we often abbreviate the type \isa{{\isacharparenleft}{\kern0pt}\isatfree{{\isacharprime}{\kern0pt}n}{\isacharcomma}{\kern0pt}\ \isatfree{{\isacharprime}{\kern0pt}t}{\isacharparenright}{\kern0pt}\ \isatconst{item}} to simply
``\isa{item}'' for brevity.
\end{definition}

Context-free items allow tracking the current state of the parsing process. Generally, as we
work towards deriving a string, the symbols to the right of the bullet (e.g. \isa{{\isasymbeta}} in 
\isa{{\isacharbrackleft}{\kern0pt}\isafree{A}\ {\isasymrightarrow}\ \isafree{{\isasymalpha}}\ {\isasymcdot}\ \isafree{{\isasymbeta}}{\isacharbrackright}{\kern0pt}}) are shifted towards the left. If \isa{{\isacharparenleft}{\kern0pt}A{\isacharcomma}{\kern0pt}\ {\isasymalpha}\ {\isacharat}{\kern0pt}\ {\isasymbeta}{\isacharparenright}{\kern0pt}\ {\isasymin}\ Prods\ G}, the item
\isa{{\isacharbrackleft}{\kern0pt}\isafree{A}\ {\isasymrightarrow}\ \isafree{{\isasymalpha}}\ {\isasymcdot}\ \isafree{{\isasymbeta}}{\isacharbrackright}{\kern0pt}} denotes the situation where a word has already been derived from the substring 
\isa{{\isasymalpha}}, with a suffix still left to be derived from \isa{{\isasymbeta}}. We call the symbols that have already been 
shifted the \concept{history} of the item:
\[ \isa{\isaconst{history}\ \isavar{i}\ {\isasymequiv}\ \textsf{case}\ \isavar{i}\ \textsf{of}\ {\isacharbrackleft}{\kern0pt}\isabound{A}\ {\isasymrightarrow}\ \isabound{{\isasymalpha}}\ {\isasymcdot}\ \isabound{{\isasymbeta}}{\isacharbrackright}{\kern0pt}\ {\isasymRightarrow}\ \isabound{{\isasymalpha}}} \]

For \isa{{\isacharbrackleft}{\kern0pt}\isafree{A}\ {\isasymrightarrow}\ \isafree{{\isasymalpha}}\ {\isasymcdot}\ \isafree{{\isasymbeta}}{\isacharbrackright}{\kern0pt}}, \isa{{\isasymalpha}\ {\isacharequal}{\kern0pt}\ {\isasymepsilon}} denotes the situation where nothing has been 
derived from \isa{A} yet. Analogously, \isa{{\isasymbeta}\ {\isacharequal}{\kern0pt}\ {\isasymepsilon}} denotes the situation where a substring of the 
input has been completely derived from \isa{A}. These items are therefore called \concept{initial} and 
\concept{complete} items respectively. 
We often write the empty list implicitly on either side of the bullet, e.g., instead of
\mbox{\isa{{\isacharbrackleft}{\kern0pt}A\ {\isasymrightarrow}\ {\isasymalpha}\ {\isasymcdot}\ {\isacharbrackleft}{\kern0pt}{\isacharbrackright}{\kern0pt}{\isacharbrackright}{\kern0pt}}}, we write 
\isa{{\isacharbrackleft}{\kern0pt}\isafree{A}\ {\isasymrightarrow}\ \isafree{{\isasymalpha}}\ {\isasymcdot}\ {\isacharbrackright}{\kern0pt}}. We also use this convention if both sides of the bullet are empty.

Additionally, we denote the set of all complete items in a set of items \isa{I} by \isa{\isaconst{completes}\ \isafree{I}}:
\begin{equation*}
\isa{\isaconst{completes}\ \isavar{I}\ {\isasymequiv}\ {\isacharbraceleft}{\kern0pt}\isabound{i}\ {\isasymin}\ \isavar{I}\ {\isacharbar}{\kern0pt}\ \textsf{case}\ \isabound{i}\ \textsf{of}\ {\isacharbrackleft}{\kern0pt}\isabound{X}\ {\isasymrightarrow}\ \isabound{{\isasymalpha}}\ {\isasymcdot}\ \isabound{{\isasymbeta}}{\isacharbrackright}{\kern0pt}\ {\isasymRightarrow}\ \isabound{{\isasymbeta}}\ {\isacharequal}{\kern0pt}\ {\isacharbrackleft}{\kern0pt}{\isacharbrackright}{\kern0pt}{\isacharbraceright}{\kern0pt}}.
\end{equation*}
An item that is not complete is referred to as \concept{incomplete}, and we correspondingly define
\isa{\isaconst{incompletes}} as the complement of \isa{\isaconst{completes}}:
\begin{equation*}
\isa{\isaconst{incompletes}\ \isafree{I}\ {\isasymequiv}\ \isafree{I}\ {\isacharminus}{\kern0pt}\ \isaconst{completes}\ \isafree{I}}.
\end{equation*}
We also lift the definition of \isa{\isaconst{history}} from items to lists of items:
\begin{equation*}
\isa{\isaconst{hist}\ \isavar{{\isasymrho}}\ {\isasymequiv}\ \isaconst{concat}\ {\isacharparenleft}{\kern0pt}\isaconst{map}\ \isaconst{history}\ \isavar{{\isasymrho}}{\isacharparenright}{\kern0pt}},
\end{equation*}
and lastly, we refer to the set of all items of a CFG \isa{G} as \isa{\ensuremath{\mathrm{It}_{\isafree{G}}}}:
\begin{equation*}
\isa{\ensuremath{\mathrm{It}_{\isavar{G}}}\ {\isacharequal}{\kern0pt}\ {\isacharbraceleft}{\kern0pt}{\isacharbrackleft}{\kern0pt}\isabound{A}\ {\isasymrightarrow}\ \isabound{{\isasymalpha}}\ {\isasymcdot}\ \isabound{{\isasymbeta}}{\isacharbrackright}{\kern0pt}\ {\isacharbar}{\kern0pt}\ {\isacharparenleft}{\kern0pt}\isabound{A}{\isacharcomma}{\kern0pt}\ \isabound{{\isasymalpha}}\ {\isacharat}{\kern0pt}\ \isabound{{\isasymbeta}}{\isacharparenright}{\kern0pt}\ {\isasymin}\ \isaconst{Prods}\ \isavar{G}{\isacharbraceright}{\kern0pt}}.
\end{equation*}

\begin{lemma}\label{finite_It}
If \isa{G} is a CFG such that \isa{Prods\ G} is finite, \isa{\ensuremath{\mathrm{It}_{\isafree{G}}}} is finite.
\begin{proof}
The definition of \isa{\ensuremath{\mathrm{It}_{\isafree{G}}}} is clearly equivalent to the union of the sets of items for 
each production in \isa{Prods\ G}. Formally:
\[ \isa{It\isactrlsub G\ {\isacharequal}{\kern0pt}\ {\isacharparenleft}{\kern0pt}{\isasymUnion}{\isacharparenleft}{\kern0pt}A{\isacharcomma}{\kern0pt}w{\isacharparenright}{\kern0pt}{\isasymin}Prods\ G{\isachardot}{\kern0pt}\ {\isacharbraceleft}{\kern0pt}{\isacharbrackleft}{\kern0pt}A\ {\isasymrightarrow}\ {\isasymalpha}\ {\isasymcdot}\ {\isasymbeta}{\isacharbrackright}{\kern0pt}\ {\isacharbar}{\kern0pt}\ {\isasymalpha}\ {\isasymbeta}{\isachardot}{\kern0pt}\ {\isasymalpha}{\isacharat}{\kern0pt}{\isasymbeta}\ {\isacharequal}{\kern0pt}\ w{\isacharbraceright}{\kern0pt}{\isacharparenright}{\kern0pt}}. \]
We can show this union is finite by showing that the individual sets are finite. In order to do 
this, we prove that for arbitrary \isa{A} and \isa{w}, the set \isa{{\isacharbraceleft}{\kern0pt}{\isacharbrackleft}{\kern0pt}\isafree{A}\ {\isasymrightarrow}\ \isabound{{\isasymalpha}}\ {\isasymcdot}\ \isabound{{\isasymbeta}}{\isacharbrackright}{\kern0pt}\ {\isacharbar}{\kern0pt}\ \isabound{{\isasymalpha}}\ {\isacharat}{\kern0pt}\ \isabound{{\isasymbeta}}\ {\isacharequal}{\kern0pt}\ \isafree{w}{\isacharbraceright}{\kern0pt}}
is finite. 

We prove this property by showing there exists a bijection between this set and the first
\isa{{\isacharbar}{\kern0pt}\isafree{w}{\isacharbar}{\kern0pt}} natural numbers using the mapping \isa{f\ {\isacharcolon}{\kern0pt}{\isacharcolon}{\kern0pt}\ nat\ {\isasymRightarrow}\ {\isacharparenleft}{\kern0pt}{\isacharprime}{\kern0pt}n{\isacharcomma}{\kern0pt}\ {\isacharprime}{\kern0pt}t{\isacharparenright}{\kern0pt}\ item}
\[ \isa{\isafree{f}\ \isafree{n}\ {\isacharequal}{\kern0pt}\ {\isacharbrackleft}{\kern0pt}\isafree{A}\ {\isasymrightarrow}\ \isafree{w}{\isacharbar}{\kern0pt}\isactrlbsub \isafree{n}\isactrlesub \ {\isasymcdot}\ \isactrlbsub \isafree{n}\isactrlesub {\isacharbar}{\kern0pt}\isafree{w}{\isacharbrackright}{\kern0pt}}. \]
\end{proof}
\end{lemma}%
\end{isamarkuptext}\isamarkuptrue%
%
\isadelimdocument
%
\endisadelimdocument
%
\isatagdocument
%
\isamarkupsection{Generalized Pushdown Automata\label{sec:gpdas}%
}
\isamarkuptrue%
%
\endisatagdocument
{\isafolddocument}%
%
\isadelimdocument
%
\endisadelimdocument
%
\begin{isamarkuptext}%
In later sections, we will define several automata to lay the foundations for the \isa{LR{\isacharparenleft}{\kern0pt}{\isadigit{0}}{\isacharparenright}{\kern0pt}}
parser. Most of these automata, including the parser itself, require a stack to operate, but 
unlike conventional pushdown automata, it is sometimes necessary for them to read multiple stack 
symbols in a single transition step.
\begin{definition}[Generalized pushdown automata]
A generalized pushdown automaton (GPDA) is a record of type \isa{{\isacharparenleft}{\kern0pt}\isatfree{{\isacharprime}{\kern0pt}q}{\isacharcomma}{\kern0pt}\ \isatfree{{\isacharprime}{\kern0pt}a}{\isacharparenright}{\kern0pt}\ \isatconst{gpda}} where 
\isa{\isatfree{{\isacharprime}{\kern0pt}q}} is the type of stack symbols, \isa{\isatfree{{\isacharprime}{\kern0pt}a}} the type of alphabet symbols, and the fields
\begin{itemize}
\item \isa{states\ {\isacharcolon}{\kern0pt}{\isacharcolon}{\kern0pt}\ {\isacharprime}{\kern0pt}q\ set}: a finite set of \concept{states}.
\item \isa{init\ {\isacharcolon}{\kern0pt}{\isacharcolon}{\kern0pt}\ {\isacharprime}{\kern0pt}q}: the \concept{initial state} with \isa{init\ {\isasymin}\ states}.
\item \isa{final\ {\isacharcolon}{\kern0pt}{\isacharcolon}{\kern0pt}\ {\isacharprime}{\kern0pt}q\ set}: the set of \concept{final states} with \isa{final\ {\isasymsubseteq}\ states}.
\item \isa{nxt\ {\isacharcolon}{\kern0pt}{\isacharcolon}{\kern0pt}\ {\isacharparenleft}{\kern0pt}{\isacharprime}{\kern0pt}q\ list\ {\isasymtimes}\ {\isacharprime}{\kern0pt}a\ {\isasymtimes}\ {\isacharprime}{\kern0pt}q\ list{\isacharparenright}{\kern0pt}\ set}: the \concept{reading transition relation}, i.e., 
the relation of transitions that read and consume the leftmost symbol of the remaining input.
\item \isa{eps\ {\isacharcolon}{\kern0pt}{\isacharcolon}{\kern0pt}\ {\isacharparenleft}{\kern0pt}{\isacharprime}{\kern0pt}q\ list\ {\isasymtimes}\ {\isacharprime}{\kern0pt}q\ list{\isacharparenright}{\kern0pt}\ set}: the transition relation for \concept{\isa{{\isasymepsilon}}-transitions}, 
i.e., transitions that do not read the input.
\end{itemize}
Lastly, if \isa{M} is a GPDA we assume for all \isa{{\isacharparenleft}{\kern0pt}\isafree{ps}{\isacharcomma}{\kern0pt}\ \isafree{a}{\isacharcomma}{\kern0pt}\ \isafree{qs}{\isacharparenright}{\kern0pt}\ {\isasymin}\ \isaconst{nxt}\ \isafree{M}} that \isa{ps} is nonempty, 
and both \isa{ps} and \isa{qs} are subsets of the states of \isa{M}. We make the same assumption for all 
\isa{{\isacharparenleft}{\kern0pt}\isafree{ps}{\isacharcomma}{\kern0pt}\ \isafree{qs}{\isacharparenright}{\kern0pt}\ {\isasymin}\ \isaconst{eps}\ \isafree{M}}.
\end{definition}

It is worth noting that, differently from traditional PDAs, GPDAs do not have a dedicated state. 
Instead, a variable amount of topmost stack symbols are used to determine the transition. 
Therefore, if \isa{M} is a GPDA, talking about ``the state'' of \isa{M} is simply a
shorthand to refer to the topmost state on \isa{M}'s stack.

Another important aspect is that as opposed to our definition of GPDAs, Wilhelm et al.\ define them
as having finite transition relations. This assumption is of importance in particular if one wishes
to show the equivalence between GPDAs and PDAs, which is out of the scope of our formalization.
Besides the parser, our formalization defines the item pushdown automaton as a GPDA. Proving the
finiteness of this automaton's transition relation is a burden that yields a fact we will not use,
so we omit the finiteness assumption in the generic GPDA definition for the sake of simplicity. For
the \isa{LR{\isacharparenleft}{\kern0pt}{\isadigit{0}}{\isacharparenright}{\kern0pt}} parser itself, however, the finiteness of the transition relations is of great interest
due to the question of executability. We will address this in a later section.

For \isa{M\ {\isacharcolon}{\kern0pt}{\isacharcolon}{\kern0pt}\ {\isacharparenleft}{\kern0pt}{\isacharprime}{\kern0pt}q{\isacharcomma}{\kern0pt}\ {\isacharprime}{\kern0pt}a{\isacharparenright}{\kern0pt}\ gpda} we define a \concept{configuration} as a tuple 
\isa{{\isacharparenleft}{\kern0pt}qs{\isacharcomma}{\kern0pt}\ w{\isacharparenright}{\kern0pt}\ {\isacharcolon}{\kern0pt}{\isacharcolon}{\kern0pt}\ {\isacharprime}{\kern0pt}q\ list\ {\isasymtimes}\ {\isacharprime}{\kern0pt}a\ list} where \isa{qs} denotes the current stack, and \isa{w} the remaining input to 
be read. In accordance with the Isabelle/HOL list datatype, we define the topmost stack symbol as 
the leftmost list element, deviating from Wilhelm et al.\ in this regard.

A configuration of \isa{M} is \concept{initial} if the stack consists of a singleton list containing 
the initial state \isa{\isaconst{init}\ \isafree{M}}, while a \concept{final} configuration for \isa{M} consists of a 
singleton list with some final state on the stack after completely consuming the input, 
i.e., a configuration of the form \isa{{\isacharparenleft}{\kern0pt}{\isacharbrackleft}{\kern0pt}\isafree{f}{\isacharbrackright}{\kern0pt}{\isacharcomma}{\kern0pt}\ {\isacharbrackleft}{\kern0pt}{\isacharbrackright}{\kern0pt}{\isacharparenright}{\kern0pt}} for some \isa{f\ {\isasymin}\ final\ M}.

We now define the step relation \isa{{\isasymturnstile}} for GPDAs:
\begin{gather*}
\isa{{\isacharparenleft}{\kern0pt}\isavar{ps}{\isacharcomma}{\kern0pt}\ \isavar{a}{\isacharcomma}{\kern0pt}\ \isavar{qs}{\isacharparenright}{\kern0pt}\ {\isasymin}\ \isaconst{nxt}\ \isafree{M}\ {\isasymLongrightarrow}\ {\isacharparenleft}{\kern0pt}\isavar{ps}\ {\isacharat}{\kern0pt}\ \isavar{rs}{\isacharcomma}{\kern0pt}\ \isavar{a}\ {\isacharhash}{\kern0pt}\ \isavar{w}{\isacharparenright}{\kern0pt}\ {\isasymturnstile}\ {\isacharparenleft}{\kern0pt}\isavar{qs}\ {\isacharat}{\kern0pt}\ \isavar{rs}{\isacharcomma}{\kern0pt}\ \isavar{w}{\isacharparenright}{\kern0pt}}\\
\isa{{\isacharparenleft}{\kern0pt}\isavar{ps}{\isacharcomma}{\kern0pt}\ \isavar{qs}{\isacharparenright}{\kern0pt}\ {\isasymin}\ \isaconst{eps}\ \isafree{M}\ {\isasymLongrightarrow}\ {\isacharparenleft}{\kern0pt}\isavar{ps}\ {\isacharat}{\kern0pt}\ \isavar{rs}{\isacharcomma}{\kern0pt}\ \isavar{w}{\isacharparenright}{\kern0pt}\ {\isasymturnstile}\ {\isacharparenleft}{\kern0pt}\isavar{qs}\ {\isacharat}{\kern0pt}\ \isavar{rs}{\isacharcomma}{\kern0pt}\ \isavar{w}{\isacharparenright}{\kern0pt}}
\end{gather*}
We refer to sequences of steps as \concept{computations}. Computations can naturally be described
with the reflexive transitive closure of steps, which we denote with the infix operator \isa{{\isasymturnstile}{\isacharasterisk}{\kern0pt}}, and
similarly, we notate \isa{n}-step computations with \isa{{\isasymturnstile}{\isacharparenleft}{\kern0pt}n{\isacharparenright}{\kern0pt}}.

Finally, we define the \concept{language} \isa{\isafree{L}} for \isa{M} as the set of words for which \isa{M} 
can reach a final configuration from the corresponding initial configuration:
\begin{equation*}
\isa{L{\isacharparenleft}{\kern0pt}M{\isacharparenright}{\kern0pt}\ {\isasymequiv}\ {\isacharbraceleft}{\kern0pt}\isabound{w}\ {\isacharbar}{\kern0pt}\ {\isasymexists}\isabound{f}{\isasymin}\isaconst{final}\ \isafree{M}{\isachardot}{\kern0pt}\ {\isacharparenleft}{\kern0pt}{\isacharbrackleft}{\kern0pt}\isaconst{init}\ \isafree{M}{\isacharbrackright}{\kern0pt}{\isacharcomma}{\kern0pt}\ \isabound{w}{\isacharparenright}{\kern0pt}\ {\isasymturnstile}{\isacharasterisk}{\kern0pt}\ {\isacharparenleft}{\kern0pt}{\isacharbrackleft}{\kern0pt}\isabound{f}{\isacharbrackright}{\kern0pt}{\isacharcomma}{\kern0pt}\ {\isacharbrackleft}{\kern0pt}{\isacharbrackright}{\kern0pt}{\isacharparenright}{\kern0pt}{\isacharbraceright}{\kern0pt}}.
\end{equation*}%
\end{isamarkuptext}\isamarkuptrue%
%
\isadelimdocument
%
\endisadelimdocument
%
\isatagdocument
%
\isamarkupchapter{The Item Pushdown Automaton%
}
\isamarkuptrue%
%
\endisatagdocument
{\isafolddocument}%
%
\isadelimdocument
%
\endisadelimdocument
%
\isadelimproof
%
\endisadelimproof
%
\isatagproof
%
\endisatagproof
{\isafoldproof}%
%
\isadelimproof
%
\endisadelimproof
%
\isadelimdocument
%
\endisadelimdocument
%
\isatagdocument
%
\isamarkupsection{Definition%
}
\isamarkuptrue%
%
\endisatagdocument
{\isafolddocument}%
%
\isadelimdocument
%
\endisadelimdocument
%
\begin{isamarkuptext}%
One of the main objectives in the construction of our parser is determinism. Despite the ability of
PDAs of recognizing CFLs, they are nondeterministic in general, which means they are not easily
implemented in practice. In this section, we define the item pushdown automaton to a 
context-free grammar, from which we will later derive a deterministic parser. From this point onwards, 
unless stated otherwise, let \isa{G} be a CFG with production set \isa{P} and start symbol \isa{S} such that

\begin{itemize}
\item \isa{\isaconst{P}} is finite
\item \isa{L{\isacharparenleft}{\kern0pt}\isafree{G}{\isacharparenright}{\kern0pt}\ {\isasymnoteq}\ {\isasymemptyset}}
\item \isa{G} is reduced
\item \isa{G{\isacharprime}{\kern0pt}} is the extension of \isa{G} with start symbol \isa{S{\isacharprime}{\kern0pt}}. We again abbreviate \isa{Prods\ G{\isacharprime}{\kern0pt}} as \isa{P{\isacharprime}{\kern0pt}}.
\end{itemize}

\begin{definition}[Item pushdown automaton]
The \concept{item pushdown automaton} (IPDA) to \isa{G} is the \isa{{\isacharparenleft}{\kern0pt}{\isacharparenleft}{\kern0pt}\isatfree{{\isacharprime}{\kern0pt}n}{\isacharcomma}{\kern0pt}\ \isatfree{{\isacharprime}{\kern0pt}t}{\isacharparenright}{\kern0pt}\ \isatconst{item}{\isacharcomma}{\kern0pt}\ \isatfree{{\isacharprime}{\kern0pt}t}{\isacharparenright}{\kern0pt}\ \isatconst{gpda}}:
\begin{multline*}
  \isa{I\isactrlsub G\ {\isacharequal}{\kern0pt}\ {\isasymlparr}states} = \isa{\ensuremath{\mathrm{It}_{\isaconst{G{\isacharprime}{\kern0pt}}}}},\ \isa{init\ {\isacharequal}{\kern0pt}\ {\isacharbrackleft}{\kern0pt}S{\isacharprime}{\kern0pt}\ {\isasymrightarrow}\ {\isasymcdot}\ {\isacharbrackleft}{\kern0pt}Nt\ S{\isacharbrackright}{\kern0pt}{\isacharbrackright}{\kern0pt}{\isacharcomma}{\kern0pt}\ final\ {\isacharequal}{\kern0pt}\ {\isacharbraceleft}{\kern0pt}{\isacharbrackleft}{\kern0pt}S{\isacharprime}{\kern0pt}\ {\isasymrightarrow}\ {\isacharbrackleft}{\kern0pt}Nt\ S{\isacharbrackright}{\kern0pt}\ {\isasymcdot}\ {\isacharbrackright}{\kern0pt}{\isacharbraceright}{\kern0pt}{\isacharcomma}{\kern0pt}\ nxt\ {\isacharequal}{\kern0pt}\ {\isasymDelta}\isactrlsub I{\isacharcomma}{\kern0pt}}\\
    \isa{eps\ {\isacharequal}{\kern0pt}\ {\isasymE}\isactrlsub I{\isasymrparr}}
\end{multline*}
where 
\begin{multline*}
\isa{{\isasymDelta}\isactrlsub I\ {\isacharequal}{\kern0pt}\ {\isacharbraceleft}{\kern0pt}{\isacharparenleft}{\kern0pt}{\isacharbrackleft}{\kern0pt}{\isacharbrackleft}{\kern0pt}X\ {\isasymrightarrow}\ {\isasymbeta}\ {\isasymcdot}\ Tm\ a\ {\isacharhash}{\kern0pt}\ {\isasymgamma}{\isacharbrackright}{\kern0pt}{\isacharbrackright}{\kern0pt}{\isacharcomma}{\kern0pt}\ a{\isacharcomma}{\kern0pt}\ {\isacharbrackleft}{\kern0pt}{\isacharbrackleft}{\kern0pt}X\ {\isasymrightarrow}\ {\isasymbeta}\ {\isacharat}{\kern0pt}\ {\isacharbrackleft}{\kern0pt}Tm\ a{\isacharbrackright}{\kern0pt}\ {\isasymcdot}\ {\isasymgamma}{\isacharbrackright}{\kern0pt}{\isacharbrackright}{\kern0pt}{\isacharparenright}{\kern0pt}}\\
\isa{{\isacharbar}{\kern0pt}\ X\ {\isasymbeta}\ a\ {\isasymgamma}{\isachardot}{\kern0pt}\ {\isacharparenleft}{\kern0pt}X{\isacharcomma}{\kern0pt}\ {\isasymbeta}\ {\isacharat}{\kern0pt}\ Tm\ a\ {\isacharhash}{\kern0pt}\ {\isasymgamma}{\isacharparenright}{\kern0pt}\ {\isasymin}\ P{\isacharprime}{\kern0pt}{\isacharbraceright}{\kern0pt}}
\end{multline*}
and \isa{{\isasymE}\isactrlsub I\ {\isacharequal}{\kern0pt}\ E\ {\isasymunion}\ R} for
\begin{gather*}
\begin{multlined}
  \isa{E\ {\isacharequal}{\kern0pt}\ {\isacharbraceleft}{\kern0pt}{\isacharparenleft}{\kern0pt}{\isacharbrackleft}{\kern0pt}{\isacharbrackleft}{\kern0pt}X\ {\isasymrightarrow}\ {\isasymbeta}\ {\isasymcdot}\ Nt\ Y\ {\isacharhash}{\kern0pt}\ {\isasymgamma}{\isacharbrackright}{\kern0pt}{\isacharbrackright}{\kern0pt}{\isacharcomma}{\kern0pt}\ {\isacharbrackleft}{\kern0pt}{\isacharbrackleft}{\kern0pt}Y\ {\isasymrightarrow}\ {\isasymcdot}\ {\isasymalpha}{\isacharbrackright}{\kern0pt}{\isacharcomma}{\kern0pt}\ {\isacharbrackleft}{\kern0pt}X\ {\isasymrightarrow}\ {\isasymbeta}\ {\isasymcdot}\ Nt\ Y\ {\isacharhash}{\kern0pt}\ {\isasymgamma}{\isacharbrackright}{\kern0pt}{\isacharbrackright}{\kern0pt}{\isacharparenright}{\kern0pt}}\\ 
  \isa{{\isacharbar}{\kern0pt}\ X\ {\isasymbeta}\ Y\ {\isasymgamma}\ {\isasymalpha}{\isachardot}{\kern0pt}\ {\isacharparenleft}{\kern0pt}X{\isacharcomma}{\kern0pt}\ {\isasymbeta}\ {\isacharat}{\kern0pt}\ Nt\ Y\ {\isacharhash}{\kern0pt}\ {\isasymgamma}{\isacharparenright}{\kern0pt}\ {\isasymin}\ P{\isacharprime}{\kern0pt}\ {\isasymand}\ {\isacharparenleft}{\kern0pt}Y{\isacharcomma}{\kern0pt}\ {\isasymalpha}{\isacharparenright}{\kern0pt}\ {\isasymin}\ P{\isacharprime}{\kern0pt}{\isacharbraceright}{\kern0pt}}
\end{multlined}\\
\begin{multlined}
  \isa{R\ {\isacharequal}{\kern0pt}\ {\isacharbraceleft}{\kern0pt}{\isacharparenleft}{\kern0pt}{\isacharbrackleft}{\kern0pt}{\isacharbrackleft}{\kern0pt}Y\ {\isasymrightarrow}\ {\isasymalpha}\ {\isasymcdot}\ {\isacharbrackright}{\kern0pt}{\isacharcomma}{\kern0pt}\ {\isacharbrackleft}{\kern0pt}X\ {\isasymrightarrow}\ {\isasymbeta}\ {\isasymcdot}\ Nt\ Y\ {\isacharhash}{\kern0pt}\ {\isasymgamma}{\isacharbrackright}{\kern0pt}{\isacharbrackright}{\kern0pt}{\isacharcomma}{\kern0pt}\ {\isacharbrackleft}{\kern0pt}{\isacharbrackleft}{\kern0pt}X\ {\isasymrightarrow}\ {\isasymbeta}\ {\isacharat}{\kern0pt}\ {\isacharbrackleft}{\kern0pt}Nt\ Y{\isacharbrackright}{\kern0pt}\ {\isasymcdot}\ {\isasymgamma}{\isacharbrackright}{\kern0pt}{\isacharbrackright}{\kern0pt}{\isacharparenright}{\kern0pt}}\\
  \isa{{\isacharbar}{\kern0pt}\ Y\ {\isasymalpha}\ X\ {\isasymbeta}\ {\isasymgamma}{\isachardot}{\kern0pt}\ {\isacharparenleft}{\kern0pt}X{\isacharcomma}{\kern0pt}\ {\isasymbeta}\ {\isacharat}{\kern0pt}\ Nt\ Y\ {\isacharhash}{\kern0pt}\ {\isasymgamma}{\isacharparenright}{\kern0pt}\ {\isasymin}\ P{\isacharprime}{\kern0pt}\ {\isasymand}\ {\isacharparenleft}{\kern0pt}Y{\isacharcomma}{\kern0pt}\ {\isasymalpha}{\isacharparenright}{\kern0pt}\ {\isasymin}\ P{\isacharprime}{\kern0pt}{\isacharbraceright}{\kern0pt}}
\end{multlined}
\end{gather*}
\end{definition}

Overall, the IPDA has three types of transitions. We call transitions in \isa{\isaconst{nxt}} \concept{shifting} 
transitions, transitions in \isa{\isafree{E}\ {\isasymsubseteq}\ \isafree{{\isasymE}\isactrlsub I}} \concept{expanding} transitions, and transitions in 
\isa{\isafree{R}\ {\isasymsubseteq}\ \isafree{{\isasymE}\isactrlsub I}} \concept{reducing} transitions. We denote \isa{\isaconst{IPDA}} steps by \isa{{\isasymturnstile}I}, \isa{{\isasymturnstile}I{\isacharasterisk}{\kern0pt}} and
\isa{{\isasymturnstile}I{\isacharparenleft}{\kern0pt}n{\isacharparenright}{\kern0pt}} analogously to the general GPDA steps. 

Our definitions differ slightly from those by Wilhelm et al.: in all
transition sets, we explicitly restrict the elements to items that correspond to productions of \isa{G{\isacharprime}{\kern0pt}}.
In their book, Wilhelm et al.\ define the transition relation of a GPDA with state set \isa{Q} and input
alphabet \isa{V\isactrlsub T} to be a subset of \mbox{\isa{Q\isactrlsup {\isacharplus}{\kern0pt}\ {\isasymtimes}\ V\isactrlsub T\ {\isasymtimes}\ Q\isactrlsup {\isacharasterisk}{\kern0pt}}}. We approximate this in the record type of GPDAs, 
as we stated before, by defining \mbox{\isa{nxt\ {\isacharcolon}{\kern0pt}{\isacharcolon}{\kern0pt}\ {\isacharparenleft}{\kern0pt}{\isacharprime}{\kern0pt}q\ list\ {\isasymtimes}\ {\isacharprime}{\kern0pt}a\ {\isasymtimes}\ {\isacharprime}{\kern0pt}q\ list{\isacharparenright}{\kern0pt}\ set}} and 
\mbox{\isa{eps\ {\isacharcolon}{\kern0pt}{\isacharcolon}{\kern0pt}\ {\isacharparenleft}{\kern0pt}{\isacharprime}{\kern0pt}q\ list\ {\isasymtimes}\ {\isacharprime}{\kern0pt}q\ list{\isacharparenright}{\kern0pt}\ set}} for a \isa{{\isacharparenleft}{\kern0pt}{\isacharprime}{\kern0pt}q{\isacharcomma}{\kern0pt}\ {\isacharprime}{\kern0pt}a{\isacharparenright}{\kern0pt}\ gpda}. However, since we define the states
of the GPDA as a set instead of as the type itself, these definitions for \isa{nxt} and \isa{eps} allow 
lists containing elements of type \isa{{\isacharprime}{\kern0pt}q} but that are not in the set of states. Our definitions of 
\isa{nxt} and \isa{eps} for \isa{I\isactrlsub G} circumvent this problem by explicitly restricting the sets to lists with
items whose corresponding productions are in \isa{P{\isacharprime}{\kern0pt}}, which is equivalent to the items themselves being
in \isa{\ensuremath{\mathrm{It}_{\isaconst{G{\isacharprime}{\kern0pt}}}}}.

We will now look at how the automaton we just defined operates. Intuitively, \isa{I\isactrlsub G} accepts a word \isa{w}
by finding a rightmost derivation \isa{\isaconst{P{\isacharprime}{\kern0pt}}\ {\isasymturnstile}\ {\isacharbrackleft}{\kern0pt}\isaconst{Nt}\ \isaconst{S{\isacharprime}{\kern0pt}}{\isacharbrackright}{\kern0pt}\ {\isasymRightarrow}r{\isacharasterisk}{\kern0pt}\ \isaconst{map}\ \isaconst{Tm}\ \isafree{w}}.
If the current state of \isa{I\isactrlsub G} is \isa{{\isacharbrackleft}{\kern0pt}\isafree{A}\ {\isasymrightarrow}\ \isafree{{\isasymalpha}}\ {\isasymcdot}\ \isaconst{Tm}\ \isafree{a}\ {\isacharhash}{\kern0pt}\ \isafree{{\isasymbeta}}{\isacharbrackright}{\kern0pt}} for any \mbox{\isa{a\ {\isacharcolon}{\kern0pt}{\isacharcolon}{\kern0pt}\ {\isacharprime}{\kern0pt}t}}, \isa{I\isactrlsub G} will 
invariably shift \isa{Tm\ a} by replacing this topmost item by \isa{\mbox{{\isacharbrackleft}{\kern0pt}\isafree{A}\ {\isasymrightarrow}\ \isafree{{\isasymalpha}}\ {\isacharat}{\kern0pt}\ {\isacharbrackleft}{\kern0pt}\isaconst{Tm}\ \isafree{a}{\isacharbrackright}{\kern0pt}\ {\isasymcdot}\ \isafree{{\isasymbeta}}{\isacharbrackright}{\kern0pt}}}. 
Similarly, if the state is some complete item \isa{{\isacharbrackleft}{\kern0pt}\isafree{Y}\ {\isasymrightarrow}\ \isafree{{\isasymalpha}}\ {\isasymcdot}\ {\isacharbrackright}{\kern0pt}}, it will reduce the item 
and shift \isa{Nt\ Y} on the second-topmost item if possible. If the stack is 
\isa{{\isacharbrackleft}{\kern0pt}\isafree{Y}\ {\isasymrightarrow}\ \isafree{{\isasymalpha}}\ {\isasymcdot}\ {\isacharbrackright}{\kern0pt}\ {\isacharhash}{\kern0pt}\ {\isacharbrackleft}{\kern0pt}\isafree{X}\ {\isasymrightarrow}\ \isafree{{\isasymbeta}}\ {\isasymcdot}\ \isaconst{Nt}\ \isafree{Y}\ {\isacharhash}{\kern0pt}\ \isafree{{\isasymgamma}}{\isacharbrackright}{\kern0pt}\ {\isacharhash}{\kern0pt}\ \isafree{{\isasymrho}}}, the act of replacing the first two items by 
\isa{{\isacharbrackleft}{\kern0pt}\isafree{X}\ {\isasymrightarrow}\ \isafree{{\isasymbeta}}\ {\isacharat}{\kern0pt}\ {\isacharbrackleft}{\kern0pt}\isaconst{Nt}\ \isafree{Y}{\isacharbrackright}{\kern0pt}\ {\isasymcdot}\ \isafree{{\isasymgamma}}{\isacharbrackright}{\kern0pt}} is equivalent to the backwards (i.e. right-to-left)
application of a rightmost derivation step for some \isa{u\ {\isacharcolon}{\kern0pt}{\isacharcolon}{\kern0pt}\ {\isacharprime}{\kern0pt}t\ list} of the form
\begin{equation*}
\isa{\isaconst{P{\isacharprime}{\kern0pt}}\ {\isasymturnstile}\ \isafree{{\isasymbeta}}\ {\isacharat}{\kern0pt}\ \isaconst{Nt}\ \isafree{Y}\ {\isacharhash}{\kern0pt}\ \isaconst{map}\ \isaconst{Tm}\ \isafree{u}\ {\isasymRightarrow}r\ \isafree{{\isasymbeta}}\ {\isacharat}{\kern0pt}\ \isafree{{\isasymalpha}}\ {\isacharat}{\kern0pt}\ \isaconst{map}\ \isaconst{Tm}\ \isafree{u}}.
\end{equation*}

Lastly, the expanding case is the only transition type where nondeterministic behavior actually 
presents itself. While reducing transitions correspond to the IPDA applying a backward step in a 
rightmost derivation, the expanding step is essentially the IPDA nondeterministically choosing 
\emph{which} production it will attempt to reduce: as we will later prove, if \isa{I\isactrlsub G} performs the
expansion
\[ \isa{{\isacharparenleft}{\kern0pt}{\isacharbrackleft}{\kern0pt}\isafree{X}\ {\isasymrightarrow}\ \isafree{{\isasymbeta}}\ {\isasymcdot}\ \isaconst{Nt}\ \isafree{Y}\ {\isacharhash}{\kern0pt}\ \isafree{{\isasymgamma}}{\isacharbrackright}{\kern0pt}\ {\isacharhash}{\kern0pt}\ \isafree{{\isasymrho}}{\isacharcomma}{\kern0pt}\ \isafree{w}{\isacharparenright}{\kern0pt}\ {\isasymturnstile}I\ {\isacharparenleft}{\kern0pt}{\isacharbrackleft}{\kern0pt}\isafree{Y}\ {\isasymrightarrow}\ \ {\isasymcdot}\ \isafree{{\isasymalpha}}\ {\isacharbrackright}{\kern0pt}\ {\isacharhash}{\kern0pt}\ {\isacharbrackleft}{\kern0pt}\isafree{X}\ {\isasymrightarrow}\ \isafree{{\isasymbeta}}\ {\isasymcdot}\ \isaconst{Nt}\ \isafree{Y}\ {\isacharhash}{\kern0pt}\ \isafree{{\isasymgamma}}{\isacharbrackright}{\kern0pt}\ {\isacharhash}{\kern0pt}\ \isafree{{\isasymrho}}{\isacharcomma}{\kern0pt}\ \isafree{w}{\isacharparenright}{\kern0pt}} \]
it will only complete \isa{{\isasymalpha}}, i.e., reach a configuration with stack
\isa{{\isacharbrackleft}{\kern0pt}\isafree{Y}\ {\isasymrightarrow}\ \isafree{{\isasymalpha}}\ {\isasymcdot}\ {\isacharbrackright}{\kern0pt}\ {\isacharhash}{\kern0pt}\ \mbox{{\isacharbrackleft}{\kern0pt}\isafree{X}\ {\isasymrightarrow}\ \isafree{{\isasymbeta}}\ {\isasymcdot}\ \isaconst{Nt}\ \isafree{Y}\ {\isacharhash}{\kern0pt}\ \isafree{{\isasymgamma}}{\isacharbrackright}{\kern0pt}}\ {\isacharhash}{\kern0pt}\ \isafree{{\isasymrho}}}
if for some prefix \isa{u} of the remaining input word \isa{w\ {\isacharequal}{\kern0pt}\ uv} holds \isa{\isaconst{P{\isacharprime}{\kern0pt}}\ {\isasymturnstile}\ \isafree{{\isasymalpha}}\ {\isasymRightarrow}{\isacharasterisk}{\kern0pt}\ \isaconst{map}\ \isaconst{Tm}\ \isafree{u}}.
We will now work towards proving that \isa{I\isactrlsub G} accepts exactly \isa{L{\isacharparenleft}{\kern0pt}\isafree{G}{\isacharparenright}{\kern0pt}}.%
\end{isamarkuptext}\isamarkuptrue%
%
\isadelimdocument
%
\endisadelimdocument
%
\isatagdocument
%
\isamarkupsection{Language Equivalence%
}
\isamarkuptrue%
%
\endisatagdocument
{\isafolddocument}%
%
\isadelimdocument
%
\endisadelimdocument
%
\begin{isamarkuptext}%
We will first show that the IPDA \isa{I\isactrlsub G} is sound, that is, its language is a subset of
\isa{L{\isacharparenleft}{\kern0pt}\isafree{G}{\isacharparenright}{\kern0pt}}. We first prove an invariant following the original proof by Wilhelm et al.

\begin{lemma}[IPDA invariant]\label{ipda.invariant}
\isa{{\isacharparenleft}{\kern0pt}{\isacharbrackleft}{\kern0pt}\isaconst{init}\ \isafree{M}{\isacharbrackright}{\kern0pt}{\isacharcomma}{\kern0pt}\ \isafree{u}\ {\isacharat}{\kern0pt}\ \isafree{v}{\isacharparenright}{\kern0pt}\ {\isasymturnstile}I{\isacharasterisk}{\kern0pt}\ {\isacharparenleft}{\kern0pt}\isaconst{rev}\ \isafree{{\isasymrho}}{\isacharcomma}{\kern0pt}\ \isafree{v}{\isacharparenright}{\kern0pt}} implies \isa{\isaconst{P}\ {\isasymturnstile}\ \isaconst{hist}\ \isafree{{\isasymrho}}\ {\isasymRightarrow}{\isacharasterisk}{\kern0pt}\ \mbox{\isaconst{map}\ \isaconst{Tm}\ \isafree{u}}}.
\begin{proof}
We proceed by induction on the length \isa{n} of the computation for arbitrary \isa{u{\isacharcomma}{\kern0pt}\ v{\isacharcomma}{\kern0pt}} and \isa{{\isasymrho}}.

If \isa{\isafree{n}\ {\isacharequal}{\kern0pt}\ {\isadigit{0}}}, then 
$\isa{{\isacharbrackleft}{\kern0pt}init\ M{\isacharbrackright}{\kern0pt}} = \isa{\mbox{{\isacharbrackleft}{\kern0pt}{\isacharbrackleft}{\kern0pt}\isaconst{S{\isacharprime}{\kern0pt}}\ {\isasymrightarrow}\ \ {\isasymcdot}\ {\isacharbrackleft}{\kern0pt}\isaconst{Nt}\ \isaconst{S}{\isacharbrackright}{\kern0pt}\ {\isacharbrackright}{\kern0pt}{\isacharbrackright}{\kern0pt}}} = \isa{rev\ {\isasymrho}}$ and \isa{\isafree{u}\ {\isacharat}{\kern0pt}\ \isafree{v}\ {\isacharequal}{\kern0pt}\ \isafree{v}} hold. This in
turn implies \isa{\isaconst{hist}\ \isafree{{\isasymrho}}\ {\isacharequal}{\kern0pt}\ {\isacharbrackleft}{\kern0pt}{\isacharbrackright}{\kern0pt}} and \isa{\isafree{u}\ {\isacharequal}{\kern0pt}\ {\isacharbrackleft}{\kern0pt}{\isacharbrackright}{\kern0pt}}, fulfilling the invariant.

For the inductive step we distinguish cases on the final step of the computation.

If the last step was a shifting transition there exist \isa{A{\isacharcomma}{\kern0pt}\ {\isasymalpha}{\isacharcomma}{\kern0pt}\ a{\isacharcomma}{\kern0pt}\ {\isasymbeta}{\isacharcomma}{\kern0pt}\ {\isasymtau}{\isacharcomma}{\kern0pt}\ a{\isacharcomma}{\kern0pt}} and \isa{x} such that
the second to last configuration was of the form
\begin{gather}
\isa{{\isacharparenleft}{\kern0pt}{\isacharbrackleft}{\kern0pt}\isafree{A}\ {\isasymrightarrow}\ \isafree{{\isasymalpha}}\ {\isasymcdot}\ \isaconst{Tm}\ \isafree{a}\ {\isacharhash}{\kern0pt}\ \isafree{{\isasymbeta}}{\isacharbrackright}{\kern0pt}\ {\isacharhash}{\kern0pt}\ \isafree{{\isasymtau}}{\isacharcomma}{\kern0pt}\ \isafree{a}\ {\isacharhash}{\kern0pt}\ \isafree{v}{\isacharparenright}{\kern0pt}}\label{eq:ipda.invariant.shift}
\intertext{and}
\isa{\isaconst{rev}\ \isafree{{\isasymrho}}\ {\isacharequal}{\kern0pt}\ {\isacharbrackleft}{\kern0pt}\isafree{A}\ {\isasymrightarrow}\ \isafree{{\isasymalpha}}\ {\isacharat}{\kern0pt}\ {\isacharbrackleft}{\kern0pt}\isaconst{Tm}\ \isafree{a}{\isacharbrackright}{\kern0pt}\ {\isasymcdot}\ \isafree{{\isasymbeta}}{\isacharbrackright}{\kern0pt}\ {\isacharhash}{\kern0pt}\ \isafree{{\isasymtau}}}\label{eq:ipda.invariant.rho_shift}.
\end{gather}
This implies the existence of some \isa{y\ {\isacharcolon}{\kern0pt}{\isacharcolon}{\kern0pt}\ {\isacharprime}{\kern0pt}t\ list} such that the initial input was of the form
\isa{uv\ {\isacharequal}{\kern0pt}\ yav}. This, together with \eqref{eq:ipda.invariant.shift}, and the induction hypothesis implies 
\[ \isa{P\ {\isasymturnstile}\ hist\ {\isacharparenleft}{\kern0pt}rev\ {\isasymtau}\ {\isacharat}{\kern0pt}\ {\isacharbrackleft}{\kern0pt}{\isacharbrackleft}{\kern0pt}A\ {\isasymrightarrow}\ {\isasymalpha}\ {\isasymcdot}\ Tm\ a\ {\isacharhash}{\kern0pt}\ {\isasymbeta}{\isacharbrackright}{\kern0pt}{\isacharbrackright}{\kern0pt}{\isacharparenright}{\kern0pt}} = \isa{hist\ {\isacharparenleft}{\kern0pt}rev\ {\isasymtau}{\isacharparenright}{\kern0pt}\ {\isacharat}{\kern0pt}\ {\isasymalpha}\ {\isasymRightarrow}{\isacharasterisk}{\kern0pt}\ map\ Tm\ y}. \] 
With \isa{uv\ {\isacharequal}{\kern0pt}\ yav} and by substituting \eqref{eq:ipda.invariant.rho_shift}, it follows that
\isa{P\ {\isasymturnstile}\ hist\ {\isasymrho}\ {\isasymRightarrow}{\isacharasterisk}{\kern0pt}\ map\ Tm\ y\ {\isacharat}{\kern0pt}}\ \isa{\mbox{{\isacharbrackleft}{\kern0pt}\isaconst{Tm}\ \isafree{a}{\isacharbrackright}{\kern0pt}}\ {\isacharequal}{\kern0pt}\ \isafree{u}} holds, fulfilling the invariant.

For the reducing case, we have a second-to-last configuration 
\isa{{\isacharparenleft}{\kern0pt}{\isacharbrackleft}{\kern0pt}\isafree{Y}\ {\isasymrightarrow}\ \isafree{{\isasymalpha}}\ {\isasymcdot}\ {\isacharbrackright}{\kern0pt}\ {\isacharhash}{\kern0pt}\ \mbox{{\isacharbrackleft}{\kern0pt}\isafree{X}\ {\isasymrightarrow}\ \isafree{{\isasymbeta}}\ {\isasymcdot}\ \isaconst{Nt}\ \isafree{Y}\ {\isacharhash}{\kern0pt}\ \isafree{{\isasymgamma}}{\isacharbrackright}{\kern0pt}}\ {\isacharhash}{\kern0pt}\ \isafree{{\isasymtau}}{\isacharcomma}{\kern0pt}\ \isafree{v}{\isacharparenright}{\kern0pt}}, and 
\isa{\isaconst{rev}\ \isafree{{\isasymrho}}} = \isa{\mbox{{\isacharbrackleft}{\kern0pt}\isafree{X}\ {\isasymrightarrow}\ \isafree{{\isasymbeta}}\ {\isacharat}{\kern0pt}\ {\isacharbrackleft}{\kern0pt}\isaconst{Nt}\ \isafree{Y}{\isacharbrackright}{\kern0pt}\ {\isasymcdot}\ \isafree{{\isasymgamma}}{\isacharbrackright}{\kern0pt}}\ {\isacharhash}{\kern0pt}\ \isafree{{\isasymtau}}} holds for some \isa{Y{\isacharcomma}{\kern0pt}\ {\isasymalpha}{\isacharcomma}{\kern0pt}\ X{\isacharcomma}{\kern0pt}\ {\isasymbeta}{\isacharcomma}{\kern0pt}\ {\isasymgamma}} and \isa{{\isasymtau}}.
Moreover, \isa{{\isacharparenleft}{\kern0pt}\isafree{Y}{\isacharcomma}{\kern0pt}\ \isafree{{\isasymalpha}}{\isacharparenright}{\kern0pt}\ {\isasymin}\ \isaconst{P}} must hold; otherwise, \isa{{\isacharbrackleft}{\kern0pt}\isafree{X}\ {\isasymrightarrow}\ \isafree{{\isasymbeta}}\ {\isasymcdot}\ \isaconst{Nt}\ \isafree{Y}\ {\isacharhash}{\kern0pt}\ \isafree{{\isasymgamma}}{\isacharbrackright}{\kern0pt}\ {\isasymin}\ \ensuremath{\mathrm{It}_{\isaconst{G{\isacharprime}{\kern0pt}}}}} 
would imply that \isa{S{\isacharprime}{\kern0pt}} is on the RHS of a production in \isa{G{\isacharprime}{\kern0pt}}, which cannot be the case by definition.

With all this and the induction hypothesis follows 
\[ \isa{\isaconst{P}\ {\isasymturnstile}\ \isaconst{hist}\ \isafree{{\isasymrho}}\ {\isasymRightarrow}r\ \isaconst{hist}\ {\isacharparenleft}{\kern0pt}\isaconst{rev}\ \isafree{{\isasymtau}}{\isacharparenright}{\kern0pt}\ {\isacharat}{\kern0pt}\ \isafree{{\isasymbeta}}\ {\isacharat}{\kern0pt}\ \isafree{{\isasymalpha}}} 
  \overset{(IH)}{\Rightarrow\mkern-5mu r} \isa{map\ Tm\ u}. \] 
Therefore, the invariant holds once again.

Finally, in the expanding case we have \isa{{\isacharparenleft}{\kern0pt}{\isacharbrackleft}{\kern0pt}\isafree{X}\ {\isasymrightarrow}\ \isafree{{\isasymbeta}}\ {\isasymcdot}\ \isaconst{Nt}\ \isafree{Y}\ {\isacharhash}{\kern0pt}\ \isafree{{\isasymgamma}}{\isacharbrackright}{\kern0pt}\ {\isacharhash}{\kern0pt}\ \isafree{{\isasymtau}}{\isacharcomma}{\kern0pt}\ \isafree{v}{\isacharparenright}{\kern0pt}} and 
\isa{\isaconst{rev}\ \isafree{{\isasymrho}}\ {\isacharequal}{\kern0pt}\ {\isacharbrackleft}{\kern0pt}\isafree{Y}\ {\isasymrightarrow}\ \ {\isasymcdot}\ \isafree{{\isasymalpha}}\ {\isacharbrackright}{\kern0pt}\ {\isacharhash}{\kern0pt}\ {\isacharbrackleft}{\kern0pt}\isafree{X}\ {\isasymrightarrow}\ \isafree{{\isasymbeta}}\ {\isasymcdot}\ \isaconst{Nt}\ \isafree{Y}\ {\isacharhash}{\kern0pt}\ \isafree{{\isasymgamma}}{\isacharbrackright}{\kern0pt}\ {\isacharhash}{\kern0pt}\ \isafree{{\isasymtau}}}

By the induction hypothesis we have once more \isa{P\ {\isasymturnstile}\ hist\ {\isacharparenleft}{\kern0pt}rev\ {\isasymtau}{\isacharparenright}{\kern0pt}\ {\isacharat}{\kern0pt}\ {\isasymbeta}\ {\isasymRightarrow}{\isacharasterisk}{\kern0pt}\ map\ Tm\ u}. We then 
have
\begin{align*}
\isa{P\ {\isasymturnstile}\ hist\ {\isasymrho}} &\ \isa{{\isacharequal}{\kern0pt}\ hist\ {\isacharparenleft}{\kern0pt}{\isacharparenleft}{\kern0pt}rev\ {\isasymtau}{\isacharparenright}{\kern0pt}\ {\isacharat}{\kern0pt}\ {\isacharbrackleft}{\kern0pt}X\ {\isasymrightarrow}\ {\isasymbeta}\ {\isasymcdot}\ Nt\ Y\ {\isacharhash}{\kern0pt}\ {\isasymgamma}{\isacharbrackright}{\kern0pt}\ {\isacharat}{\kern0pt}\ {\isacharbrackleft}{\kern0pt}Y\ {\isasymrightarrow}\ {\isacharbrackleft}{\kern0pt}{\isacharbrackright}{\kern0pt}\ {\isasymcdot}\ {\isasymalpha}{\isacharbrackright}{\kern0pt}{\isacharparenright}{\kern0pt}}\\
&\ \isa{{\isacharequal}{\kern0pt}\ hist\ {\isacharparenleft}{\kern0pt}rev\ {\isasymtau}{\isacharparenright}{\kern0pt}\ {\isacharat}{\kern0pt}\ {\isasymbeta}\ {\isasymRightarrow}{\isacharasterisk}{\kern0pt}\ map\ Tm\ u}. 
\end{align*}
The invariant is therefore satisfied for all cases.
\end{proof}
\end{lemma}

\begin{lemma}\label{ipda.Lang_subst_Lang_G}
\isa{\ensuremath{L(I_G)}\ {\isasymsubseteq}\ L{\isacharparenleft}{\kern0pt}\isafree{G}{\isacharparenright}{\kern0pt}}
\begin{proof}
Assume \isa{\isafree{w}\ {\isasymin}\ \ensuremath{L(I_G)}}. Then, 
\[ \isa{{\isacharparenleft}{\kern0pt}{\isacharbrackleft}{\kern0pt}init\ I\isactrlsub G{\isacharbrackright}{\kern0pt}{\isacharcomma}{\kern0pt}\ w{\isacharparenright}{\kern0pt}\ {\isacharequal}{\kern0pt}}\ \isa{{\isacharparenleft}{\kern0pt}{\isacharbrackleft}{\kern0pt}\isaconst{init}\ \isaconst{I\isactrlsub G}{\isacharbrackright}{\kern0pt}{\isacharcomma}{\kern0pt}\ \isafree{w}\ {\isacharat}{\kern0pt}\ {\isacharbrackleft}{\kern0pt}{\isacharbrackright}{\kern0pt}{\isacharparenright}{\kern0pt}\ {\isasymturnstile}I{\isacharasterisk}{\kern0pt}\ {\isacharparenleft}{\kern0pt}{\isacharbrackleft}{\kern0pt}{\isacharbrackleft}{\kern0pt}\isaconst{S{\isacharprime}{\kern0pt}}\ {\isasymrightarrow}\ {\isacharbrackleft}{\kern0pt}\isaconst{Nt}\ \isaconst{S}{\isacharbrackright}{\kern0pt}\ {\isasymcdot}\ {\isacharbrackright}{\kern0pt}{\isacharbrackright}{\kern0pt}{\isacharcomma}{\kern0pt}\ {\isacharbrackleft}{\kern0pt}{\isacharbrackright}{\kern0pt}{\isacharparenright}{\kern0pt}}. \] 
By Lemma~\ref{ipda.invariant}, this implies \isa{\isaconst{P}\ {\isasymturnstile}\ \isaconst{hist}\ {\isacharbrackleft}{\kern0pt}{\isacharbrackleft}{\kern0pt}\isaconst{S{\isacharprime}{\kern0pt}}\ {\isasymrightarrow}\ {\isacharbrackleft}{\kern0pt}\isaconst{Nt}\ \isaconst{S}{\isacharbrackright}{\kern0pt}\ {\isasymcdot}\ {\isacharbrackright}{\kern0pt}{\isacharbrackright}{\kern0pt}\ {\isasymRightarrow}{\isacharasterisk}{\kern0pt}\ \isaconst{map}\ \isaconst{Tm}\ \isafree{w}}, and thus
\isa{\isafree{w}\ {\isasymin}\ L{\isacharparenleft}{\kern0pt}\isafree{G}{\isacharparenright}{\kern0pt}}.  
\end{proof}
\end{lemma}

And now, we prove the other direction:

\begin{lemma}\label{completes_Tms}
If \isa{{\isacharparenleft}{\kern0pt}\isafree{A}{\isacharcomma}{\kern0pt}\ \isafree{{\isasymalpha}}\ {\isacharat}{\kern0pt}\ \isaconst{map}\ \isaconst{Tm}\ \isafree{u}\ {\isacharat}{\kern0pt}\ \isafree{{\isasymbeta}}{\isacharparenright}{\kern0pt}\ {\isasymin}\ \isaconst{P{\isacharprime}{\kern0pt}}}, then 
\[ \isa{{\isacharparenleft}{\kern0pt}{\isacharbrackleft}{\kern0pt}\isafree{A}\ {\isasymrightarrow}\ \isafree{{\isasymalpha}}\ {\isasymcdot}\ \isaconst{map}\ \isaconst{Tm}\ \isafree{u}\ {\isacharat}{\kern0pt}\ \isafree{{\isasymbeta}}{\isacharbrackright}{\kern0pt}\ {\isacharhash}{\kern0pt}\ \isafree{{\isasymrho}}{\isacharcomma}{\kern0pt}\ \isafree{u}\ {\isacharat}{\kern0pt}\ \isafree{v}{\isacharparenright}{\kern0pt}\ {\isasymturnstile}I{\isacharasterisk}{\kern0pt}\ {\isacharparenleft}{\kern0pt}{\isacharbrackleft}{\kern0pt}\isafree{A}\ {\isasymrightarrow}\ \isafree{{\isasymalpha}}\ {\isacharat}{\kern0pt}\ \isaconst{map}\ \isaconst{Tm}\ \isafree{u}\ {\isasymcdot}\ \isafree{{\isasymbeta}}{\isacharbrackright}{\kern0pt}\ {\isacharhash}{\kern0pt}\ \isafree{{\isasymrho}}{\isacharcomma}{\kern0pt}\ \isafree{v}{\isacharparenright}{\kern0pt}}. \]
\begin{proof}
Trivial by induction on the length of \isa{u}.
\end{proof}
\end{lemma}

\begin{lemma}[Derivation implies IPDA completion]\label{derives_imp_completes}
If \isa{\isaconst{P{\isacharprime}{\kern0pt}}\ {\isasymturnstile}\ \isafree{{\isasymbeta}}\ {\isasymRightarrow}{\isacharasterisk}{\kern0pt}\ \isaconst{map}\ \isaconst{Tm}\ \isafree{w}} and \isa{\mbox{{\isacharparenleft}{\kern0pt}\isafree{A}{\isacharcomma}{\kern0pt}\ \isafree{{\isasymalpha}}\ {\isacharat}{\kern0pt}\ \isafree{{\isasymbeta}}\ {\isacharat}{\kern0pt}\ \isafree{{\isasymgamma}}{\isacharparenright}{\kern0pt}}\ {\isasymin}\ \isaconst{P{\isacharprime}{\kern0pt}}}, then for any \isa{{\isasymrho}} and \isa{x} holds: 
\[ \isa{{\isacharparenleft}{\kern0pt}{\isacharbrackleft}{\kern0pt}\isafree{A}\ {\isasymrightarrow}\ \isafree{{\isasymalpha}}\ {\isasymcdot}\ \isafree{{\isasymbeta}}\ {\isacharat}{\kern0pt}\ \isafree{{\isasymgamma}}{\isacharbrackright}{\kern0pt}\ {\isacharhash}{\kern0pt}\ \isafree{{\isasymrho}}{\isacharcomma}{\kern0pt}\ \isafree{w}\ {\isacharat}{\kern0pt}\ \isafree{x}{\isacharparenright}{\kern0pt}\ {\isasymturnstile}I{\isacharasterisk}{\kern0pt}\ {\isacharparenleft}{\kern0pt}{\isacharbrackleft}{\kern0pt}\isafree{A}\ {\isasymrightarrow}\ \isafree{{\isasymalpha}}\ {\isacharat}{\kern0pt}\ \isafree{{\isasymbeta}}\ {\isasymcdot}\ \isafree{{\isasymgamma}}{\isacharbrackright}{\kern0pt}\ {\isacharhash}{\kern0pt}\ \isafree{{\isasymrho}}{\isacharcomma}{\kern0pt}\ \isafree{x}{\isacharparenright}{\kern0pt}}. \]
\begin{proof}
We perform strong induction on the length of the derivation \isa{n} for arbitrary \isa{{\isasymbeta}{\isacharcomma}{\kern0pt}\ w{\isacharcomma}{\kern0pt}\ A{\isacharcomma}{\kern0pt}\ {\isasymalpha}{\isacharcomma}{\kern0pt}\ {\isasymgamma}{\isacharcomma}{\kern0pt}\ {\isasymrho}{\isacharcomma}{\kern0pt}}
and \isa{x}.
If \isa{n\ {\isacharequal}{\kern0pt}\ {\isadigit{0}}}, \isa{{\isasymbeta}\ {\isacharequal}{\kern0pt}\ map\ Tm\ w}, and the implication holds by Lemma~\ref{completes_Tms}.

If \isa{n\ {\isacharequal}{\kern0pt}\ Suc\ m} for some \isa{m\ {\isacharcolon}{\kern0pt}{\isacharcolon}{\kern0pt}\ nat}, \isa{{\isasymbeta}} must be of the form 
\begin{equation}\label{d_imp_c.b_decomp(1)}
\isa{\isafree{{\isasymbeta}}\ {\isacharequal}{\kern0pt}\ \isafree{{\isasymdelta}\isactrlsub {\isadigit{1}}}\ {\isacharat}{\kern0pt}\ \isaconst{Nt}\ \isafree{X}\ {\isacharhash}{\kern0pt}\ \isafree{{\isasymdelta}\isactrlsub {\isadigit{2}}}}
\end{equation}
for \isa{X\ {\isacharcolon}{\kern0pt}{\isacharcolon}{\kern0pt}\ {\isacharprime}{\kern0pt}n} and \isa{{\isasymdelta}\isactrlsub {\isadigit{1}}{\isacharcomma}{\kern0pt}\ {\isasymdelta}\isactrlsub {\isadigit{2}}\ {\isacharcolon}{\kern0pt}{\isacharcolon}{\kern0pt}\ syms}. Furthermore, Nipkow et al.\ have proved
\begin{multline*}
 \isa{P\ {\isasymturnstile}\ u\ {\isacharat}{\kern0pt}\ v\ {\isasymRightarrow}{\isacharparenleft}{\kern0pt}n{\isacharparenright}{\kern0pt}\ w}\\
  \isa{{\isasymlongleftrightarrow}\ {\isacharparenleft}{\kern0pt}{\isasymexists}n{\isadigit{1}}\ n{\isadigit{2}}\ w{\isadigit{1}}\ w{\isadigit{2}}{\isachardot}{\kern0pt}\ n\ {\isacharequal}{\kern0pt}\ n{\isadigit{1}}\ {\isacharplus}{\kern0pt}\ n{\isadigit{2}}\ {\isasymand}\ w\ {\isacharequal}{\kern0pt}\ w{\isadigit{1}}\ {\isacharat}{\kern0pt}\ w{\isadigit{2}}}\\ 
  \isa{{\isasymand}\ P\ {\isasymturnstile}\ u\ {\isasymRightarrow}{\isacharparenleft}{\kern0pt}n{\isadigit{1}}{\isacharparenright}{\kern0pt}\ w{\isadigit{1}}\ {\isasymand}\ P\ {\isasymturnstile}\ v\ {\isasymRightarrow}{\isacharparenleft}{\kern0pt}n{\isadigit{2}}{\isacharparenright}{\kern0pt}\ w{\isadigit{2}}{\isacharparenright}{\kern0pt}}.
\end{multline*}

By applying this lemma twice, the derivation of \isa{w} can be decomposed such that:
\begin{subequations}
\begin{gather}
\isa{\isafree{w}\ {\isacharequal}{\kern0pt}\ \isafree{u}\ {\isacharat}{\kern0pt}\ \isafree{v}\ {\isacharat}{\kern0pt}\ \isafree{y}}\label{d_imp_c.b_decomp(2)}\\
\isa{\isaconst{P{\isacharprime}{\kern0pt}}\ {\isasymturnstile}\ \isafree{{\isasymdelta}\isactrlsub {\isadigit{1}}}\ {\isasymRightarrow}{\isacharparenleft}{\kern0pt}\isafree{i}{\isacharparenright}{\kern0pt}\ \isaconst{map}\ \isaconst{Tm}\ \isafree{u}}\label{d_imp_c.d1}\\
\isa{\isaconst{P{\isacharprime}{\kern0pt}}\ {\isasymturnstile}\ {\isacharbrackleft}{\kern0pt}\isaconst{Nt}\ \isafree{X}{\isacharbrackright}{\kern0pt}\ {\isasymRightarrow}{\isacharparenleft}{\kern0pt}\isafree{j}{\isacharparenright}{\kern0pt}\ \isaconst{map}\ \isaconst{Tm}\ \isafree{v}}\label{d_imp_c.X}\\
\isa{\isaconst{P{\isacharprime}{\kern0pt}}\ {\isasymturnstile}\ \isafree{{\isasymdelta}\isactrlsub {\isadigit{2}}}\ {\isasymRightarrow}{\isacharparenleft}{\kern0pt}\isafree{k}{\isacharparenright}{\kern0pt}\ \isaconst{map}\ \isaconst{Tm}\ \isafree{y}}\label{d_imp_c.d2}\\
\isa{\isafree{n}\ {\isacharequal}{\kern0pt}\ \isafree{i}\ {\isacharplus}{\kern0pt}\ \isafree{j}\ {\isacharplus}{\kern0pt}\ \isafree{k}}.\label{d_imp_c.b_decomp(6)}
\end{gather}
\end{subequations}
For some \isa{u{\isacharcomma}{\kern0pt}\ v{\isacharcomma}{\kern0pt}\ y\ {\isacharcolon}{\kern0pt}{\isacharcolon}{\kern0pt}\ {\isacharprime}{\kern0pt}t\ list} and \isa{i{\isacharcomma}{\kern0pt}\ j{\isacharcomma}{\kern0pt}\ k\ {\isacharcolon}{\kern0pt}{\isacharcolon}{\kern0pt}\ nat}.

Since \isa{\isaconst{P{\isacharprime}{\kern0pt}}\ {\isasymturnstile}\ {\isacharbrackleft}{\kern0pt}\isaconst{Nt}\ \isafree{X}{\isacharbrackright}{\kern0pt}\ {\isasymRightarrow}{\isacharparenleft}{\kern0pt}\isafree{j}{\isacharparenright}{\kern0pt}\ \isaconst{map}\ \isaconst{Tm}\ \isafree{v}}, \isa{{\isadigit{0}}\ {\isacharless}{\kern0pt}\ \isafree{j}} must hold. With 
\eqref{d_imp_c.b_decomp(6)}, we know that there are only two cases: either \isa{\isafree{j}\ {\isacharequal}{\kern0pt}\ \isafree{n}} and 
\isa{i\ {\isacharequal}{\kern0pt}\ k\ {\isacharequal}{\kern0pt}\ {\isadigit{0}}}, or \isa{i}, \isa{j}, and \isa{k} are all strictly less than \isa{n}. We now distinguish these cases.

If \isa{j\ {\isacharequal}{\kern0pt}\ n} and \isa{i\ {\isacharequal}{\kern0pt}\ k\ {\isacharequal}{\kern0pt}\ {\isadigit{0}}}, 
\begin{gather}\label{d_imp_c.d1u_d2y}
\isa{{\isasymdelta}\isactrlsub {\isadigit{1}}\ {\isacharequal}{\kern0pt}\ map\ Tm\ u} \text{ and } \isa{{\isasymdelta}\isactrlsub {\isadigit{2}}\ {\isacharequal}{\kern0pt}\ map\ Tm\ y}
\end{gather} 
hold by \eqref{d_imp_c.d1} and \eqref{d_imp_c.d2}. \isa{j\ {\isacharequal}{\kern0pt}\ n} also implies the existence of some 
\isa{{\isasymbeta}{\isacharprime}{\kern0pt}\ {\isacharcolon}{\kern0pt}{\isacharcolon}{\kern0pt}\ syms} such that 
\begin{equation}\label{eq:d_imp_c.stepm}
\isa{P{\isacharprime}{\kern0pt}\ {\isasymturnstile}\ {\isacharbrackleft}{\kern0pt}Nt\ X{\isacharbrackright}{\kern0pt}\ {\isasymRightarrow}\ {\isasymbeta}{\isacharprime}{\kern0pt}\ {\isasymRightarrow}{\isacharparenleft}{\kern0pt}m{\isacharparenright}{\kern0pt}\ map\ Tm\ v}.
\end{equation}

By Lemma~\ref{completes_Tms}, and substituting \eqref{d_imp_c.b_decomp(1)}, 
\eqref{d_imp_c.b_decomp(2)}, and \eqref{d_imp_c.d1u_d2y}, \isa{{\isacharparenleft}{\kern0pt}{\isacharbrackleft}{\kern0pt}\isafree{A}\ {\isasymrightarrow}\ \isafree{{\isasymalpha}}\ {\isasymcdot}\ \isafree{{\isasymbeta}}\ {\isacharat}{\kern0pt}\ \isafree{{\isasymgamma}}{\isacharbrackright}{\kern0pt}\ {\isacharhash}{\kern0pt}\ \isafree{{\isasymrho}}{\isacharcomma}{\kern0pt}\ \isafree{w}\ {\isacharat}{\kern0pt}\ \isafree{x}{\isacharparenright}{\kern0pt}} 
reaches
\[ \isa{{\isacharparenleft}{\kern0pt}{\isacharbrackleft}{\kern0pt}\isafree{A}\ {\isasymrightarrow}\ \isafree{{\isasymalpha}}\ {\isacharat}{\kern0pt}\ \isaconst{map}\ \isaconst{Tm}\ \isafree{u}\ {\isasymcdot}\ \isaconst{Nt}\ \isafree{X}\ {\isacharhash}{\kern0pt}\ \isaconst{map}\ \isaconst{Tm}\ \isafree{y}\ {\isacharat}{\kern0pt}\ \isafree{{\isasymgamma}}{\isacharbrackright}{\kern0pt}\ {\isacharhash}{\kern0pt}\ \isafree{{\isasymrho}}{\isacharcomma}{\kern0pt}\ \isafree{v}\ {\isacharat}{\kern0pt}\ \isafree{y}\ {\isacharat}{\kern0pt}\ \isafree{x}{\isacharparenright}{\kern0pt}}, \]
and since \isa{X} derives \isa{{\isasymbeta}{\isacharprime}{\kern0pt}}, the item \isa{{\isacharbrackleft}{\kern0pt}X\ {\isasymrightarrow}\ {\isasymcdot}\ {\isasymbeta}{\isacharprime}{\kern0pt}{\isacharbrackright}{\kern0pt}} can then be pushed onto the stack through an 
expanding transition. \eqref{eq:d_imp_c.stepm} and the induction hypothesis then imply that 
\isa{I\isactrlsub G} reaches
\[ \isa{{\isacharparenleft}{\kern0pt}{\isacharbrackleft}{\kern0pt}\isafree{X}\ {\isasymrightarrow}\ \isafree{{\isasymbeta}{\isacharprime}{\kern0pt}}\ {\isasymcdot}\ {\isacharbrackright}{\kern0pt}\ {\isacharhash}{\kern0pt}\ {\isacharbrackleft}{\kern0pt}\isafree{A}\ {\isasymrightarrow}\ \isafree{{\isasymalpha}}\ {\isacharat}{\kern0pt}\ \isaconst{map}\ \isaconst{Tm}\ \isafree{u}\ {\isasymcdot}\ \isaconst{Nt}\ \isafree{X}\ {\isacharhash}{\kern0pt}\ \isaconst{map}\ \isaconst{Tm}\ \isafree{y}\ {\isacharat}{\kern0pt}\ \isafree{{\isasymgamma}}{\isacharbrackright}{\kern0pt}\ {\isacharhash}{\kern0pt}\ \isafree{{\isasymrho}}{\isacharcomma}{\kern0pt}\ \isafree{y}\ {\isacharat}{\kern0pt}\ \isafree{x}{\isacharparenright}{\kern0pt}}, \]
since \isa{m\ {\isacharless}{\kern0pt}\ Suc\ m\ {\isacharequal}{\kern0pt}\ n}. \isa{I\isactrlsub G} then trivially reaches \isa{{\isacharparenleft}{\kern0pt}{\isacharbrackleft}{\kern0pt}\isafree{A}\ {\isasymrightarrow}\ \isafree{{\isasymalpha}}\ {\isacharat}{\kern0pt}\ \isafree{{\isasymbeta}}\ {\isasymcdot}\ \isafree{{\isasymgamma}}{\isacharbrackright}{\kern0pt}\ {\isacharhash}{\kern0pt}\ \isafree{{\isasymrho}}{\isacharcomma}{\kern0pt}\ \isafree{x}{\isacharparenright}{\kern0pt}} after a 
reducing transition and applying Lemma~\ref{completes_Tms} once again.

Finally, we now consider the case where \isa{\isafree{j}\ {\isasymnoteq}\ \isafree{n}}:

With \isa{\isafree{n}\ {\isacharequal}{\kern0pt}\ \isafree{i}\ {\isacharplus}{\kern0pt}\ \isafree{j}\ {\isacharplus}{\kern0pt}\ \isafree{k}} and \isa{{\isadigit{0}}\ {\isacharless}{\kern0pt}\ \isafree{j}}, we know that \isa{\isafree{i}\ {\isacharless}{\kern0pt}\ \isafree{n}}, \isa{\isafree{j}\ {\isacharless}{\kern0pt}\ \isafree{n}}, and
\isa{\isafree{k}\ {\isacharless}{\kern0pt}\ \isafree{n}} all hold. Therefore, we can use the IH for each of the derivations \eqref{d_imp_c.d1}, 
\eqref{d_imp_c.X} and \eqref{d_imp_c.d2} after decomposing \isa{{\isasymbeta}} and \isa{w} as we did in the previous
cases with \eqref{d_imp_c.b_decomp(1)} and \eqref{d_imp_c.b_decomp(2)}. This completes the proof.
\end{proof}
\end{lemma}

With this lemma we can finally prove the second direction in the language equivalence proof:

\begin{lemma}\label{ipda.Lang_G_subst_Lang}
\isa{L{\isacharparenleft}{\kern0pt}\isafree{G}{\isacharparenright}{\kern0pt}\ {\isasymsubseteq}\ \ensuremath{L(I_G)}}
\begin{proof}
Let \isa{w} be a word in \isa{L{\isacharparenleft}{\kern0pt}\isafree{G}{\isacharparenright}{\kern0pt}}. Then \isa{\isaconst{P{\isacharprime}{\kern0pt}}\ {\isasymturnstile}\ {\isacharbrackleft}{\kern0pt}\isaconst{Nt}\ \isaconst{S}{\isacharbrackright}{\kern0pt}\ {\isasymRightarrow}{\isacharasterisk}{\kern0pt}\ \isaconst{map}\ \isaconst{Tm}\ \isafree{w}}
since \isa{\isaconst{P}\ {\isasymsubseteq}\ \isaconst{P{\isacharprime}{\kern0pt}}}.
With Lemma~\ref{derives_imp_completes} follows
\[ \isa{{\isacharparenleft}{\kern0pt}{\isacharbrackleft}{\kern0pt}{\isacharbrackleft}{\kern0pt}\isaconst{S{\isacharprime}{\kern0pt}}\ {\isasymrightarrow}\ \ {\isasymcdot}\ {\isacharbrackleft}{\kern0pt}\isaconst{Nt}\ \isaconst{S}{\isacharbrackright}{\kern0pt}\ {\isacharbrackright}{\kern0pt}{\isacharbrackright}{\kern0pt}{\isacharcomma}{\kern0pt}\ \isafree{w}{\isacharparenright}{\kern0pt}\ {\isasymturnstile}I{\isacharasterisk}{\kern0pt}\ {\isacharparenleft}{\kern0pt}{\isacharbrackleft}{\kern0pt}{\isacharbrackleft}{\kern0pt}\isaconst{S{\isacharprime}{\kern0pt}}\ {\isasymrightarrow}\ {\isacharbrackleft}{\kern0pt}\isaconst{Nt}\ \isaconst{S}{\isacharbrackright}{\kern0pt}\ {\isasymcdot}\ {\isacharbrackright}{\kern0pt}{\isacharbrackright}{\kern0pt}{\isacharcomma}{\kern0pt}\ {\isacharbrackleft}{\kern0pt}{\isacharbrackright}{\kern0pt}{\isacharparenright}{\kern0pt}}. \] 

This is equivalent to \isa{\isafree{w}\ {\isasymin}\ \ensuremath{L(I_G)}}.
\end{proof}
\end{lemma}

And thus with Lemma~\ref{ipda.Lang_subst_Lang_G}:

\begin{theorem}[Correctness of \isa{I\isactrlsub G}]\label{ipda.Lang_eq_Lang_G}
\isa{\ensuremath{L(I_G)}\ {\isacharequal}{\kern0pt}\ L{\isacharparenleft}{\kern0pt}\isafree{G}{\isacharparenright}{\kern0pt}}
\qed
\end{theorem}

Our completeness proof differs somewhat from that of Wilhelm et al. Rather than
Lemma~\ref{derives_imp_completes}, they prove a special case 
thereof~\cite[p. 61]{Wilhelm}, namely the statement
\begin{quote}
\textit{For each derivation \isa{P\ {\isasymturnstile}\ {\isacharbrackleft}{\kern0pt}Nt\ A{\isacharbrackright}{\kern0pt}\ {\isasymRightarrow}\ {\isasymalpha}\ {\isasymRightarrow}{\isacharasterisk}{\kern0pt}\ map\ Tm\ w} with \isa{A\ {\isacharcolon}{\kern0pt}{\isacharcolon}{\kern0pt}\ {\isacharprime}{\kern0pt}n}, 
\[ \isa{{\isacharparenleft}{\kern0pt}{\isacharbrackleft}{\kern0pt}\isafree{A}\ {\isasymrightarrow}\ \ {\isasymcdot}\ \isafree{{\isasymalpha}}\ {\isacharbrackright}{\kern0pt}\ {\isacharhash}{\kern0pt}\ \isafree{{\isasymrho}}{\isacharcomma}{\kern0pt}\ \isafree{w}\ {\isacharat}{\kern0pt}\ \isafree{v}{\isacharparenright}{\kern0pt}\ {\isasymturnstile}I{\isacharasterisk}{\kern0pt}\ {\isacharparenleft}{\kern0pt}{\isacharbrackleft}{\kern0pt}\isafree{A}\ {\isasymrightarrow}\ \isafree{{\isasymalpha}}\ {\isasymcdot}\ {\isacharbrackright}{\kern0pt}\ {\isacharhash}{\kern0pt}\ \isafree{{\isasymrho}}{\isacharcomma}{\kern0pt}\ \isafree{v}{\isacharparenright}{\kern0pt}} \] 
for arbitrary \isa{{\isasymrho}\ {\isacharcolon}{\kern0pt}{\isacharcolon}{\kern0pt}\ item\ list} and \isa{v\ {\isacharcolon}{\kern0pt}{\isacharcolon}{\kern0pt}\ {\isacharprime}{\kern0pt}t\ list}.}
\end{quote}
Despite this weaker statement being sufficient for the IPDA's correctness proof, we will need the
stronger lemma in order to prove the first of the two major theorems from the book, which we will
formalize in a later section.

We have now defined a nondeterministic GPDA that works with the items of a CFG \isa{G} and accepts
exactly the language of \isa{G}. Note that the IPDA is the first of multiple automata where we can see the benefit of
extending the grammar: if we were to define the set of final states to be the set of all complete
items of the form \isa{{\isacharbrackleft}{\kern0pt}\isaconst{S}\ {\isasymrightarrow}\ \isafree{{\isasymalpha}}\ {\isasymcdot}\ {\isacharbrackright}{\kern0pt}}, the automaton could encounter a final state before the end
of a computation. Since \isa{S{\isacharprime}{\kern0pt}} is in not on the RHS of any production, this scenario can not occur
in extended grammars.%
\end{isamarkuptext}\isamarkuptrue%
%
\isadelimdocument
%
\endisadelimdocument
%
\isatagdocument
%
\isamarkupchapter{Rightmost Chains: Decomposing Rightmost Derivations%
}
\isamarkuptrue%
%
\endisatagdocument
{\isafolddocument}%
%
\isadelimdocument
%
\endisadelimdocument
%
\begin{isamarkuptext}%
Wilhelm et al.\ informally assert that for a rightmost derivation 
\isa{P{\isacharprime}{\kern0pt}\ {\isasymturnstile}\ {\isacharbrackleft}{\kern0pt}Nt\ S{\isacharprime}{\kern0pt}{\isacharbrackright}{\kern0pt}\ {\isasymRightarrow}r{\isacharasterisk}{\kern0pt}\ {\isasymgamma}{\isacharprime}{\kern0pt}\ {\isacharat}{\kern0pt}\ Nt\ A\ {\isacharhash}{\kern0pt}\ map\ Tm\ w\ {\isasymRightarrow}r\ {\isasymgamma}{\isacharprime}{\kern0pt}\ {\isacharat}{\kern0pt}\ {\isasymalpha}\ {\isacharat}{\kern0pt}\ {\isasymbeta}\ {\isacharat}{\kern0pt}\ map\ Tm\ w}, there exists a decomposition
of the form
\begin{equation}\label{WSH rm chain}
\begin{split}
\isa{P{\isacharprime}{\kern0pt}}\ & \isa{{\isasymturnstile}\ {\isacharbrackleft}{\kern0pt}Nt\ S{\isacharprime}{\kern0pt}{\isacharbrackright}{\kern0pt}\ {\isasymRightarrow}r\ {\isasymalpha}\isactrlsub {\isadigit{1}}\ {\isacharat}{\kern0pt}\ Nt\ X\isactrlsub {\isadigit{1}}\ {\isacharhash}{\kern0pt}\ {\isasymbeta}\isactrlsub {\isadigit{1}}\ {\isasymRightarrow}r{\isacharasterisk}{\kern0pt}\ {\isasymalpha}\isactrlsub {\isadigit{1}}\ {\isacharat}{\kern0pt}\ Nt\ X\isactrlsub {\isadigit{1}}\ {\isacharhash}{\kern0pt}\ map\ Tm\ v\isactrlsub {\isadigit{1}}}\\
        & \isa{{\isasymRightarrow}r\ {\isasymalpha}\isactrlsub {\isadigit{1}}{\isasymalpha}\isactrlsub {\isadigit{2}}\ {\isacharat}{\kern0pt}\ Nt\ X\isactrlsub {\isadigit{2}}\ {\isacharhash}{\kern0pt}\ {\isasymbeta}\isactrlsub {\isadigit{2}}\ {\isacharat}{\kern0pt}\ map\ Tm\ v\isactrlsub {\isadigit{1}}}\\ 
        & \isa{{\isasymRightarrow}r{\isacharasterisk}{\kern0pt}\ {\isasymdots}\ {\isasymRightarrow}r{\isacharasterisk}{\kern0pt}\ {\isasymalpha}\isactrlsub {\isadigit{1}}\ {\isasymdots}\ {\isasymalpha}\isactrlsub n\ {\isacharat}{\kern0pt}\ Nt\ X\isactrlsub n\ {\isacharhash}{\kern0pt}\ map\ Tm\ {\isacharparenleft}{\kern0pt}v\isactrlsub n\ {\isasymdots}\ v\isactrlsub {\isadigit{1}}{\isacharparenright}{\kern0pt}}\\
        & \isa{{\isasymRightarrow}r\ {\isacharparenleft}{\kern0pt}{\isasymalpha}\isactrlsub {\isadigit{1}}\ {\isasymdots}\ {\isasymalpha}\isactrlsub n{\isacharparenright}{\kern0pt}\ {\isasymalpha}{\isasymbeta}\ {\isacharat}{\kern0pt}\ map\ Tm\ {\isacharparenleft}{\kern0pt}v\isactrlsub n\ {\isasymdots}\ v\isactrlsub {\isadigit{1}}{\isacharparenright}{\kern0pt}}.
\end{split}
\end{equation}
where \isa{X\isactrlsub n\ {\isacharequal}{\kern0pt}\ A}~\footnote{This expression denotes concatenation mostly by juxtaposition instead of 
the usual concatenation operator \isa{{\isacharparenleft}{\kern0pt}{\isacharat}{\kern0pt}{\isacharparenright}{\kern0pt}} for compactness.}. We now formalize this concept by defining
\concept{rightmost chains} inductively. If sentential form \isa{{\isasymalpha}} reaches sentential form \isa{{\isasymbeta}} with
rightmost chain \isa{{\isasymrho}} under production set \isa{P}, we write \isa{\isaconst{P}\ {\isasymturnstile}\ \isafree{{\isasymalpha}}\ {\isasymmidarrow}\isafree{{\isasymrho}}{\isasymrightarrow}r{\isacharasterisk}{\kern0pt}\ \isafree{{\isasymbeta}}}. For a fixed \isa{P},
we define a \concept{reflexive} rule:
\begin{gather*}
\isa{\isavar{P}\ {\isasymturnstile}\ \isavar{{\isasymalpha}}\ {\isasymmidarrow}{\isacharbrackleft}{\kern0pt}{\isacharbrackright}{\kern0pt}{\isasymrightarrow}r{\isacharasterisk}{\kern0pt}\ \isavar{{\isasymalpha}}}\\
\intertext{and a \concept{step} rule:}
\isa{\mbox{}\inferrule{\mbox{\isavar{P}\ {\isasymturnstile}\ \isavar{{\isasymalpha}\isactrlsub {\isadigit{0}}}\ {\isasymmidarrow}\isavar{{\isasymrho}}{\isasymrightarrow}r{\isacharasterisk}{\kern0pt}\ \isavar{{\isasymalpha}}\ {\isacharat}{\kern0pt}\ \isaconst{Nt}\ \isavar{X}\ {\isacharhash}{\kern0pt}\ \isaconst{map}\ \isaconst{Tm}\ \isavar{v}}\\\ \mbox{\isavar{P}\ {\isasymturnstile}\ \isavar{{\isasymalpha}}\ {\isacharat}{\kern0pt}\ \isaconst{Nt}\ \isavar{X}\ {\isacharhash}{\kern0pt}\ \isaconst{map}\ \isaconst{Tm}\ \isavar{v}\ {\isasymRightarrow}r\ \isavar{{\isasymalpha}}\ {\isacharat}{\kern0pt}\ \isavar{{\isasymalpha}{\isacharprime}{\kern0pt}}\ {\isacharat}{\kern0pt}\ \isaconst{Nt}\ \isavar{Y}\ {\isacharhash}{\kern0pt}\ \isavar{{\isasymbeta}}\ {\isacharat}{\kern0pt}\ \isaconst{map}\ \isaconst{Tm}\ \isavar{v}}\\\ \mbox{\isavar{P}\ {\isasymturnstile}\ \isavar{{\isasymbeta}}\ {\isasymRightarrow}r{\isacharasterisk}{\kern0pt}\ \isaconst{map}\ \isaconst{Tm}\ \isavar{u}}}{\mbox{\isavar{P}\ {\isasymturnstile}\ \isavar{{\isasymalpha}\isactrlsub {\isadigit{0}}}\ {\isasymmidarrow}{\isacharbrackleft}{\kern0pt}\isavar{X}\ {\isasymrightarrow}\ \isavar{{\isasymalpha}{\isacharprime}{\kern0pt}}\ {\isasymcdot}\ \isaconst{Nt}\ \isavar{Y}\ {\isacharhash}{\kern0pt}\ \isavar{{\isasymbeta}}{\isacharbrackright}{\kern0pt}\ {\isacharhash}{\kern0pt}\ \isavar{{\isasymrho}}{\isasymrightarrow}r{\isacharasterisk}{\kern0pt}\ \isavar{{\isasymalpha}}\ {\isacharat}{\kern0pt}\ \isavar{{\isasymalpha}{\isacharprime}{\kern0pt}}\ {\isacharat}{\kern0pt}\ \isaconst{Nt}\ \isavar{Y}\ {\isacharhash}{\kern0pt}\ \isaconst{map}\ \isaconst{Tm}\ \isavar{u}\ {\isacharat}{\kern0pt}\ \isaconst{map}\ \isaconst{Tm}\ \isavar{v}}}}
\end{gather*}

Essentially, rightmost chains allow us to store a list where each nonterminal produces the 
next, and the \isa{{\isasymbeta}\isactrlsub i} strings derive words in-between the \isa{X\isactrlsub i}s that build our chain. For this, we 
repurpose context-free items in order to track the entirety of the production that each nonterminal 
applies to produce the following nonterminal in the chain, with the bullet allowing us to pinpoint 
\emph{where} in the production the next nonterminal is located. With this definition, we can describe
rightmost derivations in an equivalent way as the authors do, with the particular advantage that we 
can perform structural induction on these chains, making it well-suited for theorem proving with
Isabelle unlike the authors' index-based definition of the decomposition.

\begin{example}\label{ex:rm_chain}
By our definition of rightmost chains, we would write \eqref{WSH rm chain} as~\footnote{Note that we 
once more omit most concatenation operators, replacing them by juxtaposition.}
\begin{multline*}
\isa{P\ {\isasymturnstile}\ {\isacharbrackleft}{\kern0pt}Nt\ S{\isacharprime}{\kern0pt}{\isacharbrackright}{\kern0pt}\ {\isasymmidarrow}{\isacharbrackleft}{\kern0pt}} X_{n-1}\ \isa{{\isasymrightarrow}\ {\isasymalpha}\isactrlsub n\ {\isasymcdot}\ Nt\ X\isactrlsub n\ {\isacharhash}{\kern0pt}\ {\isasymbeta}\isactrlsub n{\isacharbrackright}{\kern0pt}\ {\isacharhash}{\kern0pt}\ {\isacharbrackleft}{\kern0pt}} X_{n-2}\ \isa{{\isasymrightarrow}}\ 
\alpha_{n-1}\ \isa{{\isasymcdot}}\
  Nt\ X_{n-1}\ \#\ \beta_{n-1}]\\ 
  \isa{{\isacharhash}{\kern0pt}\ {\isasymdots}\ {\isacharhash}{\kern0pt}\ {\isacharbrackleft}{\kern0pt}{\isacharbrackleft}{\kern0pt}S{\isacharprime}{\kern0pt}\ {\isasymrightarrow}\ {\isasymalpha}\isactrlsub {\isadigit{1}}\ {\isasymcdot}\ Nt\ X\isactrlsub {\isadigit{1}}\ {\isacharhash}{\kern0pt}\ {\isasymbeta}\isactrlsub {\isadigit{1}}{\isacharbrackright}{\kern0pt}{\isacharbrackright}{\kern0pt}{\isasymrightarrow}r{\isacharasterisk}{\kern0pt}\ {\isasymalpha}\isactrlsub {\isadigit{1}}\ {\isasymdots}\ {\isasymalpha}\isactrlsub n\ {\isacharat}{\kern0pt}\ Nt\ X\isactrlsub n\ {\isacharhash}{\kern0pt}\ map\ Tm\ {\isacharparenleft}{\kern0pt}v\isactrlsub n{\isasymdots}v\isactrlsub {\isadigit{1}}{\isacharparenright}{\kern0pt}}.
\end{multline*}
\end{example}

We will now show the equivalence between rightmost chains and rightmost derivations.

\begin{lemma}\label{rm_chain_imp_derivers}
If \isa{\isaconst{P}\ {\isasymturnstile}\ \isafree{{\isasymalpha}}\ {\isasymmidarrow}\isafree{{\isasymrho}}{\isasymrightarrow}r{\isacharasterisk}{\kern0pt}\ \isafree{{\isasymbeta}}}, then \isa{\isaconst{P}\ {\isasymturnstile}\ \isafree{{\isasymalpha}}\ {\isasymRightarrow}r{\isacharasterisk}{\kern0pt}\ \isafree{{\isasymbeta}}}
\begin{proof}
By rule induction on the rightmost chain.
\end{proof}
\end{lemma}

\begin{lemma}\label{derivern_singleton_imp_produced}
If \isa{\isaconst{P}\ {\isasymturnstile}\ {\isacharbrackleft}{\kern0pt}\isaconst{Nt}\ \isafree{A}{\isacharbrackright}{\kern0pt}\ {\isasymRightarrow}r{\isacharparenleft}{\kern0pt}\isaconst{Suc}\ \isafree{n}{\isacharparenright}{\kern0pt}\ \isafree{{\isasymalpha}}\ {\isacharat}{\kern0pt}\ \isaconst{Nt}\ \isafree{X}\ {\isacharhash}{\kern0pt}\ \isafree{{\isasymbeta}}}, there exists a step in the rightmost derivation 
where \isa{{\isasymalpha}} was fully derived, \isa{X} was produced, and the string to the right of \isa{X} derives \isa{{\isasymbeta}}. More 
formally: there exist \isa{m\ {\isacharcolon}{\kern0pt}{\isacharcolon}{\kern0pt}\ nat}, \isa{{\isasymalpha}{\isacharprime}{\kern0pt}{\isacharcomma}{\kern0pt}\ {\isasymalpha}{\isacharprime}{\kern0pt}{\isacharprime}{\kern0pt}{\isacharcomma}{\kern0pt}\ {\isasymbeta}{\isacharprime}{\kern0pt}\ {\isacharcolon}{\kern0pt}{\isacharcolon}{\kern0pt}\ syms}, \isa{B\ {\isacharcolon}{\kern0pt}{\isacharcolon}{\kern0pt}\ {\isacharprime}{\kern0pt}n} and \isa{v\ {\isacharcolon}{\kern0pt}{\isacharcolon}{\kern0pt}\ {\isacharprime}{\kern0pt}t\ list} such that
\begin{gather*}
\isa{\isafree{m}\ {\isacharless}{\kern0pt}\ \isaconst{Suc}\ \isafree{n}}\\
\isa{\isaconst{P}\ {\isasymturnstile}\ {\isacharbrackleft}{\kern0pt}\isaconst{Nt}\ \isafree{A}{\isacharbrackright}{\kern0pt}\ {\isasymRightarrow}r{\isacharparenleft}{\kern0pt}\isafree{m}{\isacharparenright}{\kern0pt}\ \isafree{{\isasymalpha}{\isacharprime}{\kern0pt}}\ {\isacharat}{\kern0pt}\ \isaconst{Nt}\ \isafree{B}\ {\isacharhash}{\kern0pt}\ \isaconst{map}\ \isaconst{Tm}\ \isafree{v}}\\
\isa{\isaconst{P}\ {\isasymturnstile}\ \isafree{{\isasymalpha}{\isacharprime}{\kern0pt}}\ {\isacharat}{\kern0pt}\ \isaconst{Nt}\ \isafree{B}\ {\isacharhash}{\kern0pt}\ \isaconst{map}\ \isaconst{Tm}\ \isafree{v}\ {\isasymRightarrow}r\ \isafree{{\isasymalpha}{\isacharprime}{\kern0pt}}\ {\isacharat}{\kern0pt}\ \isafree{{\isasymalpha}{\isacharprime}{\kern0pt}{\isacharprime}{\kern0pt}}\ {\isacharat}{\kern0pt}\ \isaconst{Nt}\ \isafree{X}\ {\isacharhash}{\kern0pt}\ \isafree{{\isasymbeta}{\isacharprime}{\kern0pt}}\ {\isacharat}{\kern0pt}\ \isaconst{map}\ \isaconst{Tm}\ \isafree{v}}\\
\isa{\isafree{{\isasymalpha}}\ {\isacharequal}{\kern0pt}\ \isafree{{\isasymalpha}{\isacharprime}{\kern0pt}}\ {\isacharat}{\kern0pt}\ \isafree{{\isasymalpha}{\isacharprime}{\kern0pt}{\isacharprime}{\kern0pt}}}\\
\isa{\isaconst{P}\ {\isasymturnstile}\ \isafree{{\isasymbeta}{\isacharprime}{\kern0pt}}\ {\isacharat}{\kern0pt}\ \isaconst{map}\ \isaconst{Tm}\ \isafree{v}\ {\isasymRightarrow}r{\isacharasterisk}{\kern0pt}\ \isafree{{\isasymbeta}}}.
\end{gather*}
\begin{proof}
We begin the proof by strong induction on \isa{\isaconst{Suc}\ \isafree{n}} for arbitrary \isa{{\isasymalpha}} and \isa{{\isasymbeta}}. We now distinguish
the two usual cases of \isa{n}:

If \isa{n\ {\isacharequal}{\kern0pt}\ {\isadigit{0}}}, the implication holds for \isa{\isafree{m}\ {\isacharequal}{\kern0pt}\ {\isadigit{0}}}, \isa{\isafree{{\isasymalpha}{\isacharprime}{\kern0pt}}\ {\isacharequal}{\kern0pt}\ {\isacharbrackleft}{\kern0pt}{\isacharbrackright}{\kern0pt}}, \isa{\isafree{v}\ {\isacharequal}{\kern0pt}\ {\isacharbrackleft}{\kern0pt}{\isacharbrackright}{\kern0pt}}, 
  \isa{\isafree{{\isasymalpha}{\isacharprime}{\kern0pt}{\isacharprime}{\kern0pt}}\ {\isacharequal}{\kern0pt}\ \isafree{{\isasymalpha}}}, \isa{\isafree{{\isasymbeta}{\isacharprime}{\kern0pt}}\ {\isacharequal}{\kern0pt}\ \isafree{{\isasymbeta}}} and \isa{\isafree{A}\ {\isacharequal}{\kern0pt}\ \isafree{B}}. 

For the case where \isa{\isafree{n}\ {\isacharequal}{\kern0pt}\ \isaconst{Suc}\ \isafree{k}} for some \isa{k}, the derivation is of the form
\begin{multline}\label{prodd_stepn}
\isa{P\ {\isasymturnstile}\ {\isacharbrackleft}{\kern0pt}Nt\ A{\isacharbrackright}{\kern0pt}\ {\isasymRightarrow}r{\isacharparenleft}{\kern0pt}n{\isacharparenright}{\kern0pt}\ {\isasymalpha}{\isacharprime}{\kern0pt}\ {\isacharat}{\kern0pt}\ Nt\ B\ {\isacharhash}{\kern0pt}\ map\ Tm\ v\ {\isasymRightarrow}r\ {\isasymalpha}{\isacharprime}{\kern0pt}\ {\isacharat}{\kern0pt}\ {\isasymgamma}\ {\isacharat}{\kern0pt}\ map\ Tm\ v}\\ 
\isa{{\isacharequal}{\kern0pt}\ {\isasymalpha}\ {\isacharat}{\kern0pt}\ Nt\ X\ {\isacharhash}{\kern0pt}\ {\isasymbeta}}
\end{multline}

We now distinguish two further cases: \isa{X} is in \isa{{\isasymgamma}}, meaning \isa{{\isasymalpha}{\isacharprime}{\kern0pt}} is a prefix of \isa{{\isasymalpha}}, or 
\isa{X} is in \isa{{\isasymalpha}{\isacharprime}{\kern0pt}} and not in \isa{{\isasymgamma}}, meaning \isa{{\isasymalpha}} is a prefix of \isa{{\isasymalpha}{\isacharprime}{\kern0pt}}.

If \isa{X} is in \isa{{\isasymgamma}} and \isa{{\isasymalpha}{\isacharprime}{\kern0pt}} is a prefix of \isa{{\isasymalpha}}, there exist \isa{{\isasymdelta}{\isacharcomma}{\kern0pt}\ {\isasymzeta}\ {\isacharcolon}{\kern0pt}{\isacharcolon}{\kern0pt}\ syms} such that \isa{{\isasymgamma}\ {\isacharequal}{\kern0pt}\ {\isasymdelta}\ {\isacharat}{\kern0pt}\ Nt\ X\ {\isacharhash}{\kern0pt}\ {\isasymzeta}},
\isa{{\isasymalpha}\ {\isacharequal}{\kern0pt}\ {\isasymalpha}{\isacharprime}{\kern0pt}\ {\isacharat}{\kern0pt}\ {\isasymdelta}} and \isa{{\isasymbeta}\ {\isacharequal}{\kern0pt}\ {\isasymzeta}\ {\isacharat}{\kern0pt}\ map\ Tm\ v}. The implication then holds for \isa{m\ {\isacharequal}{\kern0pt}\ n}.

If \isa{X} is in \isa{{\isasymalpha}{\isacharprime}{\kern0pt}} and not in \isa{{\isasymgamma}}, there exist \isa{{\isasymdelta}{\isacharcomma}{\kern0pt}\ {\isasymzeta}\ {\isacharcolon}{\kern0pt}{\isacharcolon}{\kern0pt}\ syms} such that \isa{{\isasymalpha}{\isacharprime}{\kern0pt}\ {\isacharequal}{\kern0pt}\ {\isasymdelta}\ {\isacharat}{\kern0pt}\ Nt\ X\ {\isacharhash}{\kern0pt}\ {\isasymzeta}}. From
\eqref{prodd_stepn} we then get
\[ \isa{\isaconst{P}\ {\isasymturnstile}\ {\isacharbrackleft}{\kern0pt}\isaconst{Nt}\ \isafree{A}{\isacharbrackright}{\kern0pt}\ {\isasymRightarrow}r{\isacharparenleft}{\kern0pt}\isafree{n}{\isacharparenright}{\kern0pt}\ \isafree{{\isasymdelta}}\ {\isacharat}{\kern0pt}\ \isaconst{Nt}\ \isafree{X}\ {\isacharhash}{\kern0pt}\ \isafree{{\isasymzeta}}\ {\isacharat}{\kern0pt}\ \isaconst{Nt}\ \isafree{B}\ {\isacharhash}{\kern0pt}\ \isaconst{map}\ \isaconst{Tm}\ \isafree{v}}. \]
Furthermore, since \isa{Suc\ k\ {\isacharequal}{\kern0pt}\ n\ {\isacharless}{\kern0pt}\ Suc\ n}, and due to the fact that our induction hypothesis holds for 
arbitrary \isa{{\isasymalpha}} and \isa{{\isasymbeta}}, we can apply it for \isa{{\isasymdelta}} and \isa{{\isacharparenleft}{\kern0pt}{\isasymzeta}\ {\isacharat}{\kern0pt}\ Nt\ B\ {\isacharhash}{\kern0pt}\ map\ Tm\ v{\isacharparenright}{\kern0pt}}, by which the implication 
holds.
\end{proof}
\end{lemma}

\begin{lemma}\label{derivern_Suc_singleton_imp_rm_chain}
If \isa{\isaconst{P}\ {\isasymturnstile}\ {\isacharbrackleft}{\kern0pt}\isaconst{Nt}\ \isafree{A}{\isacharbrackright}{\kern0pt}\ {\isasymRightarrow}r{\isacharparenleft}{\kern0pt}\isaconst{Suc}\ \isafree{n}{\isacharparenright}{\kern0pt}\ \isafree{{\isasymalpha}}\ {\isacharat}{\kern0pt}\ \isaconst{Nt}\ \isafree{X}\ {\isacharhash}{\kern0pt}\ \isaconst{map}\ \isaconst{Tm}\ \isafree{v}}, then there exists a rightmost chain of the 
form 
\[ \isa{\isaconst{P}\ {\isasymturnstile}\ {\isacharbrackleft}{\kern0pt}\isaconst{Nt}\ \isafree{A}{\isacharbrackright}{\kern0pt}\ {\isasymmidarrow}{\isacharbrackleft}{\kern0pt}\isafree{B}\ {\isasymrightarrow}\ \isafree{{\isasymalpha}{\isacharprime}{\kern0pt}}\ {\isasymcdot}\ \isaconst{Nt}\ \isafree{X}\ {\isacharhash}{\kern0pt}\ \isafree{{\isasymbeta}}{\isacharbrackright}{\kern0pt}\ {\isacharhash}{\kern0pt}\ \isafree{{\isasymrho}}{\isasymrightarrow}r{\isacharasterisk}{\kern0pt}\ \isafree{{\isasymalpha}}\ {\isacharat}{\kern0pt}\ \isaconst{Nt}\ \isafree{X}\ {\isacharhash}{\kern0pt}\ \isaconst{map}\ \isaconst{Tm}\ \isafree{v}}. \]
\begin{proof}
We use strong induction on \isa{\isaconst{Suc}\ \isafree{n}} for arbitrary \isa{{\isasymalpha}{\isacharcomma}{\kern0pt}\ X{\isacharcomma}{\kern0pt}} and \isa{v}.

In the base case, \isa{n\ {\isacharequal}{\kern0pt}\ {\isadigit{0}}} implies 
\[ \isa{\isaconst{P}\ {\isasymturnstile}\ {\isacharbrackleft}{\kern0pt}\isaconst{Nt}\ \isafree{A}{\isacharbrackright}{\kern0pt}\ {\isasymmidarrow}{\isacharbrackleft}{\kern0pt}{\isacharbrackleft}{\kern0pt}\isafree{A}\ {\isasymrightarrow}\ \isafree{{\isasymalpha}}\ {\isasymcdot}\ \isaconst{Nt}\ \isafree{X}\ {\isacharhash}{\kern0pt}\ \isaconst{map}\ \isaconst{Tm}\ \isafree{v}{\isacharbrackright}{\kern0pt}{\isacharbrackright}{\kern0pt}{\isasymrightarrow}r{\isacharasterisk}{\kern0pt}\ \isafree{{\isasymalpha}}\ {\isacharat}{\kern0pt}\ \isaconst{Nt}\ \isafree{X}\ {\isacharhash}{\kern0pt}\ \isaconst{map}\ \isaconst{Tm}\ \isafree{v}}. \]

Otherwise, let \isa{n\ {\isacharequal}{\kern0pt}\ Suc\ m} for some \isa{m}. From 
\[ \isa{\isaconst{P}\ {\isasymturnstile}\ {\isacharbrackleft}{\kern0pt}\isaconst{Nt}\ \isafree{A}{\isacharbrackright}{\kern0pt}\ {\isasymRightarrow}r{\isacharparenleft}{\kern0pt}\isaconst{Suc}\ \isafree{n}{\isacharparenright}{\kern0pt}\ \isafree{{\isasymalpha}}\ {\isacharat}{\kern0pt}\ \isaconst{Nt}\ \isafree{X}\ {\isacharhash}{\kern0pt}\ \isaconst{map}\ \isaconst{Tm}\ \isafree{v}}, \]
there exist \isa{{\isasymbeta}{\isacharcomma}{\kern0pt}\ B{\isacharcomma}{\kern0pt}\ u{\isacharcomma}{\kern0pt}} and \isa{{\isasymgamma}} such that 
\begin{gather}
\isa{P\ {\isasymturnstile}\ {\isacharbrackleft}{\kern0pt}Nt\ A{\isacharbrackright}{\kern0pt}\ {\isasymRightarrow}r{\isacharparenleft}{\kern0pt}Suc\ m{\isacharparenright}{\kern0pt}\ {\isasymbeta}\ {\isacharat}{\kern0pt}\ Nt\ B\ {\isacharhash}{\kern0pt}\ map\ Tm\ u\ {\isasymRightarrow}r\ {\isasymbeta}\ {\isacharat}{\kern0pt}\ {\isasymgamma}\ {\isacharat}{\kern0pt}\ map\ Tm\ u}\label{der_rm.Suc_steps}\\
\intertext{and}
\isa{\isafree{{\isasymbeta}}\ {\isacharat}{\kern0pt}\ \isafree{{\isasymgamma}}\ {\isacharat}{\kern0pt}\ \isaconst{map}\ \isaconst{Tm}\ \isafree{u}\ {\isacharequal}{\kern0pt}\ \isafree{{\isasymalpha}}\ {\isacharat}{\kern0pt}\ \isaconst{Nt}\ \isafree{X}\ {\isacharhash}{\kern0pt}\ \isaconst{map}\ \isaconst{Tm}\ \isafree{v}}.\label{der_rm.bgu}
\end{gather}
Since \isa{Suc\ m\ {\isacharequal}{\kern0pt}\ n\ {\isacharless}{\kern0pt}\ Suc\ n}, by the induction hypothesis, \eqref{der_rm.Suc_steps} implies the
existence of some \isa{{\isasymrho}} such that 
\begin{equation}\label{der_rm.ih}
\isa{\isaconst{P}\ {\isasymturnstile}\ {\isacharbrackleft}{\kern0pt}\isaconst{Nt}\ \isafree{A}{\isacharbrackright}{\kern0pt}\ {\isasymmidarrow}\isafree{{\isasymrho}}{\isasymrightarrow}r{\isacharasterisk}{\kern0pt}\ \isafree{{\isasymbeta}}\ {\isacharat}{\kern0pt}\ \isaconst{Nt}\ \isafree{B}\ {\isacharhash}{\kern0pt}\ \isaconst{map}\ \isaconst{Tm}\ \isafree{u}}.
\end{equation}
Since \isa{B} derives \isa{{\isasymgamma}} in the last step of the derivation, we now distinguish two more cases:
either \isa{X} was produced by \isa{B} in this final derivation step, meaning \isa{\isaconst{Nt}\ \isafree{X}\ {\isasymin}\ \isaconst{set}\ \isafree{{\isasymgamma}}}, 
or, if \isa{B} did not produce \isa{X}, \isa{X} was already in the sentential form before the final step.

If \isa{\isaconst{Nt}\ \isafree{X}\ {\isasymin}\ \isaconst{set}\ \isafree{{\isasymgamma}}}, with \eqref{der_rm.bgu} we know there 
exist \isa{{\isasymdelta}\ {\isacharcolon}{\kern0pt}{\isacharcolon}{\kern0pt}\ syms} and \isa{w\ {\isacharcolon}{\kern0pt}{\isacharcolon}{\kern0pt}\ {\isacharprime}{\kern0pt}t\ list} such that \isa{\isafree{{\isasymgamma}}\ {\isacharequal}{\kern0pt}\ \isafree{{\isasymdelta}}\ {\isacharat}{\kern0pt}\ \isaconst{Nt}\ \isafree{X}\ {\isacharhash}{\kern0pt}\ \isaconst{map}\ \isaconst{Tm}\ \isafree{w}} and \isa{\isafree{w}\ {\isacharat}{\kern0pt}\ \isafree{u}\ {\isacharequal}{\kern0pt}\ \isafree{v}}. 
With \eqref{der_rm.ih} this implies
\begin{multline}\label{der_rm.True}
\isa{P\ {\isasymturnstile}\ {\isacharbrackleft}{\kern0pt}Nt\ A{\isacharbrackright}{\kern0pt}\ {\isasymmidarrow}{\isacharbrackleft}{\kern0pt}B\ {\isasymrightarrow}\ {\isasymdelta}\ {\isasymcdot}\ Nt\ X\ {\isacharhash}{\kern0pt}\ map\ Tm\ w{\isacharbrackright}{\kern0pt}\ {\isacharhash}{\kern0pt}\ {\isasymrho}{\isasymrightarrow}r{\isacharasterisk}{\kern0pt}}\\ 
  \isa{\isafree{{\isasymbeta}}\ {\isacharat}{\kern0pt}\ \isafree{{\isasymdelta}}\ {\isacharat}{\kern0pt}\ \isaconst{Nt}\ \isafree{X}\ {\isacharhash}{\kern0pt}\ \isaconst{map}\ \isaconst{Tm}\ {\isacharparenleft}{\kern0pt}\isafree{w}\ {\isacharat}{\kern0pt}\ \isafree{u}{\isacharparenright}{\kern0pt}\ {\isacharequal}{\kern0pt}\ \isafree{{\isasymbeta}}\ {\isacharat}{\kern0pt}\ \isafree{{\isasymgamma}}\ {\isacharat}{\kern0pt}\ \isaconst{map}\ \isaconst{Tm}\ \isafree{u}}
\end{multline}
By \eqref{der_rm.bgu}, this is exactly the chain we were aiming to construct, completing the first 
case.

If, on the other hand, \isa{\isaconst{Nt}\ \isafree{X}\ {\isasymnotin}\ \isaconst{set}\ \isafree{{\isasymgamma}}}, \eqref{der_rm.bgu} implies the existence of 
\isa{{\isasymdelta}\ {\isacharcolon}{\kern0pt}{\isacharcolon}{\kern0pt}\ syms} and \isa{y{\isacharcomma}{\kern0pt}\ z\ {\isacharcolon}{\kern0pt}{\isacharcolon}{\kern0pt}\ {\isacharprime}{\kern0pt}t\ list} such that
\begin{gather}
\isa{{\isasymbeta}\ {\isacharequal}{\kern0pt}\ {\isasymdelta}\ {\isacharat}{\kern0pt}\ Nt\ X\ {\isacharhash}{\kern0pt}\ map\ Tm\ y}\label{der_rm.b_dec}\\
\isa{{\isasymgamma}\ {\isacharequal}{\kern0pt}\ map\ Tm\ z}\label{der_rm.g_tms}\\
\isa{v\ {\isacharequal}{\kern0pt}\ y\ {\isacharat}{\kern0pt}\ z\ {\isacharat}{\kern0pt}\ u}.\label{der_rm.yzu}
\end{gather}
By Lemma~\ref{derivern_singleton_imp_produced}, \isa{n\ {\isacharequal}{\kern0pt}\ Suc\ m} implies the existence of some sentential 
form \isa{\isafree{{\isasymalpha}{\isacharprime}{\kern0pt}}\ {\isacharat}{\kern0pt}\ \isaconst{Nt}\ \isafree{C}\ {\isacharhash}{\kern0pt}\ \isaconst{map}\ \isaconst{Tm}\ \isafree{w}} such that
\begin{gather} 
\begin{multlined}\label{der_rm.prodd(1)}
\isa{P\ {\isasymturnstile}\ {\isacharbrackleft}{\kern0pt}Nt\ A{\isacharbrackright}{\kern0pt}\ {\isasymRightarrow}r{\isacharparenleft}{\kern0pt}k{\isacharparenright}{\kern0pt}\ {\isasymalpha}{\isacharprime}{\kern0pt}\ {\isacharat}{\kern0pt}\ Nt\ C\ {\isacharhash}{\kern0pt}\ map\ Tm\ w}\\
  \isa{{\isasymRightarrow}r\ {\isasymalpha}{\isacharprime}{\kern0pt}\ {\isacharat}{\kern0pt}\ {\isasymalpha}{\isacharprime}{\kern0pt}{\isacharprime}{\kern0pt}\ {\isacharat}{\kern0pt}\ Nt\ X\ {\isacharhash}{\kern0pt}\ {\isasymbeta}{\isacharprime}{\kern0pt}\ {\isacharat}{\kern0pt}\ map\ Tm\ w}
\end{multlined}\\
\intertext{and}
\isa{P\ {\isasymturnstile}\ {\isasymbeta}{\isacharprime}{\kern0pt}\ {\isacharat}{\kern0pt}\ map\ Tm\ w\ {\isasymRightarrow}r{\isacharasterisk}{\kern0pt}\ map\ Tm\ y\ {\isacharat}{\kern0pt}\ Nt\ B\ {\isacharhash}{\kern0pt}\ map\ Tm\ u}\label{der_rm.prodd(2)}
\end{gather}
for \isa{{\isasymdelta}\ {\isacharequal}{\kern0pt}\ {\isasymalpha}{\isacharprime}{\kern0pt}\ {\isacharat}{\kern0pt}\ {\isasymalpha}{\isacharprime}{\kern0pt}{\isacharprime}{\kern0pt}} and \isa{k\ {\isacharless}{\kern0pt}\ Suc\ m}. Moreover, with \eqref{der_rm.Suc_steps}, \eqref{der_rm.g_tms}, 
and \eqref{der_rm.yzu}, we have 
\begin{equation}\label{der_rm.suffix_derivers_v}
\isa{\isaconst{P}\ {\isasymturnstile}\ \isafree{{\isasymbeta}{\isacharprime}{\kern0pt}}\ {\isacharat}{\kern0pt}\ \isaconst{map}\ \isaconst{Tm}\ \isafree{w}\ {\isasymRightarrow}r{\isacharasterisk}{\kern0pt}\ \isaconst{map}\ \isaconst{Tm}\ \isafree{v}}
\end{equation}
From \eqref{der_rm.bgu},\eqref{der_rm.b_dec}, \eqref{der_rm.g_tms}, and \eqref{der_rm.yzu} we also 
have \isa{\isafree{{\isasymdelta}}\ {\isacharequal}{\kern0pt}\ \isafree{{\isasymalpha}}}. Since our induction hypothesis only holds for a nonzero number of steps, we
need to distinguish cases on the \isa{k} steps in \eqref{der_rm.prodd(1)}.

If \isa{k\ {\isacharequal}{\kern0pt}\ {\isadigit{0}}}, \eqref{der_rm.prodd(1)} implies that \isa{{\isasymalpha}{\isacharprime}{\kern0pt}\ {\isacharequal}{\kern0pt}\ {\isacharbrackleft}{\kern0pt}{\isacharbrackright}{\kern0pt}\ {\isacharequal}{\kern0pt}\ w} and \isa{A\ {\isacharequal}{\kern0pt}\ C}. This in turn implies
that \isa{{\isasymdelta}\ {\isacharequal}{\kern0pt}\ {\isasymalpha}{\isacharprime}{\kern0pt}{\isacharprime}{\kern0pt}}. With \eqref{der_rm.suffix_derivers_v} and substituting \isa{{\isasymdelta}} for \isa{{\isasymalpha}}, this implies 
\[ \isa{\isaconst{P}\ {\isasymturnstile}\ {\isacharbrackleft}{\kern0pt}\isaconst{Nt}\ \isafree{A}{\isacharbrackright}{\kern0pt}\ {\isasymmidarrow}{\isacharbrackleft}{\kern0pt}{\isacharbrackleft}{\kern0pt}\isafree{A}\ {\isasymrightarrow}\ \isafree{{\isasymalpha}{\isacharprime}{\kern0pt}{\isacharprime}{\kern0pt}}\ {\isasymcdot}\ \isaconst{Nt}\ \isafree{X}\ {\isacharhash}{\kern0pt}\ \isafree{{\isasymbeta}{\isacharprime}{\kern0pt}}{\isacharbrackright}{\kern0pt}{\isacharbrackright}{\kern0pt}{\isasymrightarrow}r{\isacharasterisk}{\kern0pt}\ \isafree{{\isasymalpha}}\ {\isacharat}{\kern0pt}\ \isaconst{Nt}\ \isafree{X}\ {\isacharhash}{\kern0pt}\ \isaconst{map}\ \isaconst{Tm}\ \isafree{v}}. \]

If, on the other hand, \isa{k\ {\isacharequal}{\kern0pt}\ Suc\ j} for some \isa{j}, we can apply the induction hypothesis for \isa{C}
with \eqref{der_rm.prodd(1)}, i.e., there exists some \isa{{\isasymrho}} such that \isa{{\isacharbrackleft}{\kern0pt}\isaconst{Nt}\ \isafree{A}{\isacharbrackright}{\kern0pt}} reaches
\isa{\isafree{{\isasymalpha}{\isacharprime}{\kern0pt}}\ {\isacharat}{\kern0pt}\ \isaconst{Nt}\ \isafree{C}\ {\isacharhash}{\kern0pt}\ \isaconst{map}\ \isaconst{Tm}\ \isafree{w}} with rightmost chain \isa{{\isasymrho}}~\footnote{The induction hypothesis actually
gives us more precise information about \isa{{\isasymrho}}, but we will not need it for our proof.}.
We can then use \eqref{der_rm.prodd(1)}, \eqref{der_rm.suffix_derivers_v}, and \isa{{\isasymalpha}\ {\isacharequal}{\kern0pt}\ {\isasymdelta}\ {\isacharequal}{\kern0pt}\ {\isasymalpha}{\isacharprime}{\kern0pt}\ {\isacharat}{\kern0pt}\ {\isasymalpha}{\isacharprime}{\kern0pt}{\isacharprime}{\kern0pt}}
to show that 
\[ \isa{\isaconst{P}\ {\isasymturnstile}\ {\isacharbrackleft}{\kern0pt}\isaconst{Nt}\ \isafree{A}{\isacharbrackright}{\kern0pt}\ {\isasymmidarrow}{\isacharbrackleft}{\kern0pt}\isafree{C}\ {\isasymrightarrow}\ \isafree{{\isasymalpha}{\isacharprime}{\kern0pt}{\isacharprime}{\kern0pt}}\ {\isasymcdot}\ \isaconst{Nt}\ \isafree{X}\ {\isacharhash}{\kern0pt}\ \isafree{{\isasymbeta}{\isacharprime}{\kern0pt}}{\isacharbrackright}{\kern0pt}\ {\isacharhash}{\kern0pt}\ \isafree{{\isasymrho}}{\isasymrightarrow}r{\isacharasterisk}{\kern0pt}\ \isafree{{\isasymalpha}}\ {\isacharat}{\kern0pt}\ \isaconst{Nt}\ \isafree{X}\ {\isacharhash}{\kern0pt}\ \isaconst{map}\ \isaconst{Tm}\ \isafree{v}} \]
holds, finishing the proof.
\end{proof}
\end{lemma}

Wilhelm et al.~\cite[p. 107]{Wilhelm} furthermore claim that if the IPDA with configuration
\isa{{\isacharparenleft}{\kern0pt}{\isacharbrackleft}{\kern0pt}\isafree{A}\ {\isasymrightarrow}\ \isafree{{\isasymalpha}}\ {\isasymcdot}\ \isafree{{\isasymbeta}}{\isacharbrackright}{\kern0pt}\ {\isacharhash}{\kern0pt}\ \isafree{{\isasymrho}{\isacharprime}{\kern0pt}}{\isacharcomma}{\kern0pt}\ \isafree{w}{\isacharparenright}{\kern0pt}} reaches the final configuration \isa{{\isacharparenleft}{\kern0pt}{\isacharbrackleft}{\kern0pt}{\isacharbrackleft}{\kern0pt}\isaconst{S{\isacharprime}{\kern0pt}}\ {\isasymrightarrow}\ \ {\isasymcdot}\ {\isacharbrackleft}{\kern0pt}\isaconst{Nt}\ \isaconst{S}{\isacharbrackright}{\kern0pt}\ {\isacharbrackright}{\kern0pt}{\isacharbrackright}{\kern0pt}{\isacharcomma}{\kern0pt}\ {\isacharbrackleft}{\kern0pt}{\isacharbrackright}{\kern0pt}{\isacharparenright}{\kern0pt}},
then \isa{{\isasymrho}{\isacharprime}{\kern0pt}} is of the form
\[ \isa{{\isasymrho}{\isacharprime}{\kern0pt}\ {\isacharequal}{\kern0pt}\ {\isacharbrackleft}{\kern0pt}} X_{n-1}\ \isa{{\isasymrightarrow}\ {\isasymalpha}\isactrlsub n\ {\isasymcdot}\ Nt\ X\isactrlsub n\ {\isacharhash}{\kern0pt}\ {\isasymbeta}\isactrlsub n{\isacharbrackright}{\kern0pt}\ {\isacharhash}{\kern0pt}\ {\isasymdots}\ {\isacharhash}{\kern0pt}\ {\isacharbrackleft}{\kern0pt}{\isacharbrackleft}{\kern0pt}S{\isacharprime}{\kern0pt}\ {\isasymrightarrow}\ {\isasymalpha}\isactrlsub {\isadigit{1}}\ {\isasymcdot}\ Nt\ X\isactrlsub {\isadigit{1}}\ {\isacharhash}{\kern0pt}\ {\isasymbeta}\isactrlsub {\isadigit{1}}{\isacharbrackright}{\kern0pt}{\isacharbrackright}{\kern0pt}} \]
for some \isa{n\ {\isasymge}\ {\isadigit{0}}} and \isa{X\isactrlsub n\ {\isacharequal}{\kern0pt}\ A}~\footnote{We have adapted the original claim to our own notation for 
the sake of consistency and clarity.} It is worth noting that this structure of \isa{{\isasymrho}{\isacharprime}{\kern0pt}} is essentially
the same as that of the item list in a rightmost chain (cf. Example~\ref{ex:rm_chain}). Therefore,
if some \isa{{\isasymsigma}\ {\isacharcolon}{\kern0pt}{\isacharcolon}{\kern0pt}\ item\ list} is part of some rightmost chain, we will be able to derive the same 
structure that Wilhelm et al.\ describe. We will now work towards proving that IPDA stacks reaching 
a final configuration have a stack that corresponds to some rightmost chain.%
\end{isamarkuptext}\isamarkuptrue%
%
\isadelimproof
%
\endisadelimproof
%
\isatagproof
%
\endisatagproof
{\isafoldproof}%
%
\isadelimproof
%
\endisadelimproof
%
\begin{isamarkuptext}%
\begin{lemma}\label{reaches_final_imp_last_is_init_or_final}
If \isa{{\isacharparenleft}{\kern0pt}{\isacharbrackleft}{\kern0pt}\isafree{A}\ {\isasymrightarrow}\ \isafree{{\isasymalpha}}\ {\isasymcdot}\ \isafree{{\isasymbeta}}{\isacharbrackright}{\kern0pt}\ {\isacharhash}{\kern0pt}\ \isafree{{\isasymrho}}{\isacharcomma}{\kern0pt}\ \isafree{w}{\isacharparenright}{\kern0pt}\ {\isasymturnstile}I{\isacharasterisk}{\kern0pt}\ {\isacharparenleft}{\kern0pt}{\isacharbrackleft}{\kern0pt}{\isacharbrackleft}{\kern0pt}\isaconst{S{\isacharprime}{\kern0pt}}\ {\isasymrightarrow}\ {\isacharbrackleft}{\kern0pt}\isaconst{Nt}\ \isaconst{S}{\isacharbrackright}{\kern0pt}\ {\isasymcdot}\ {\isacharbrackright}{\kern0pt}{\isacharbrackright}{\kern0pt}{\isacharcomma}{\kern0pt}\ {\isacharbrackleft}{\kern0pt}{\isacharbrackright}{\kern0pt}{\isacharparenright}{\kern0pt}}, then the last element in
\isa{{\isacharbrackleft}{\kern0pt}\isafree{A}\ {\isasymrightarrow}\ \isafree{{\isasymalpha}}\ {\isasymcdot}\ \isafree{{\isasymbeta}}{\isacharbrackright}{\kern0pt}\ {\isacharhash}{\kern0pt}\ \isafree{{\isasymrho}}} is either \isa{{\isacharbrackleft}{\kern0pt}\isaconst{S{\isacharprime}{\kern0pt}}\ {\isasymrightarrow}\ \ {\isasymcdot}\ {\isacharbrackleft}{\kern0pt}\isaconst{Nt}\ \isaconst{S{\isacharprime}{\kern0pt}}{\isacharbrackright}{\kern0pt}\ {\isacharbrackright}{\kern0pt}} or \isa{\mbox{{\isacharbrackleft}{\kern0pt}\isaconst{S{\isacharprime}{\kern0pt}}\ {\isasymrightarrow}\ {\isacharbrackleft}{\kern0pt}\isaconst{Nt}\ \isaconst{S}{\isacharbrackright}{\kern0pt}\ {\isasymcdot}\ {\isacharbrackright}{\kern0pt}}}.
\begin{proof}
By backwards induction on the length of the computation, making a case distinction on whether
the first step is shifting, expanding, or reducing in the transitive.
\end{proof}
\end{lemma}

\begin{lemma}\label{step_reaches_final_imp_S}
If \isa{{\isacharparenleft}{\kern0pt}{\isacharbrackleft}{\kern0pt}\isafree{A}\ {\isasymrightarrow}\ \isafree{{\isasymalpha}}\ {\isasymcdot}\ \isafree{{\isasymbeta}}{\isacharbrackright}{\kern0pt}\ {\isacharhash}{\kern0pt}\ \isafree{{\isasymrho}}\ {\isacharat}{\kern0pt}\ \isafree{{\isasymsigma}}{\isacharcomma}{\kern0pt}\ \isafree{u}{\isacharparenright}{\kern0pt}\ {\isasymturnstile}I\ {\isacharparenleft}{\kern0pt}{\isacharbrackleft}{\kern0pt}\isaconst{S{\isacharprime}{\kern0pt}}\ {\isasymrightarrow}\ {\isacharbrackleft}{\kern0pt}\isaconst{Nt}\ \isaconst{S}{\isacharbrackright}{\kern0pt}\ {\isasymcdot}\ {\isacharbrackright}{\kern0pt}\ {\isacharhash}{\kern0pt}\ \isafree{{\isasymsigma}}{\isacharcomma}{\kern0pt}\ \isafree{v}{\isacharparenright}{\kern0pt}},
then \isa{{\isacharbrackleft}{\kern0pt}A\ {\isasymrightarrow}\ {\isasymalpha}\ {\isasymcdot}\ {\isasymbeta}{\isacharbrackright}{\kern0pt}\ {\isacharhash}{\kern0pt}\ {\isasymrho}\ {\isacharequal}{\kern0pt}\ {\isacharbrackleft}{\kern0pt}{\isacharbrackleft}{\kern0pt}S\ {\isasymrightarrow}\ {\isasymalpha}\ {\isasymcdot}\ {\isacharbrackright}{\kern0pt}{\isacharcomma}{\kern0pt}\ {\isacharbrackleft}{\kern0pt}S{\isacharprime}{\kern0pt}\ {\isasymrightarrow}\ {\isasymcdot}\ {\isacharbrackleft}{\kern0pt}Nt\ S{\isacharbrackright}{\kern0pt}{\isacharbrackright}{\kern0pt}{\isacharbrackright}{\kern0pt}}.
\begin{proof}
By case distinction on the three types of transition.
\end{proof}
\end{lemma}

\begin{lemma}\label{rm_chain_Cons_imp_prod_rightmost}
If \isa{\isaconst{P}\ {\isasymturnstile}\ \isafree{{\isasymalpha}\isactrlsub {\isadigit{0}}}\ {\isasymmidarrow}{\isacharbrackleft}{\kern0pt}\isafree{A}\ {\isasymrightarrow}\ \isafree{{\isasymalpha}}\ {\isasymcdot}\ \isaconst{Nt}\ \isafree{B}\ {\isacharhash}{\kern0pt}\ \isafree{{\isasymbeta}}{\isacharbrackright}{\kern0pt}\ {\isacharhash}{\kern0pt}\ \isafree{{\isasymrho}}{\isasymrightarrow}r{\isacharasterisk}{\kern0pt}\ \isafree{{\isasymgamma}}}, there exist \isa{{\isasymdelta}\ {\isacharcolon}{\kern0pt}{\isacharcolon}{\kern0pt}\ syms} and \isa{u{\isacharcomma}{\kern0pt}\ v{\isacharcomma}{\kern0pt}\ w\ {\isacharcolon}{\kern0pt}{\isacharcolon}{\kern0pt}\ {\isacharprime}{\kern0pt}t\ list} 
such that \isa{\isafree{{\isasymgamma}}\ {\isacharequal}{\kern0pt}\ \isafree{{\isasymdelta}}\ {\isacharat}{\kern0pt}\ \isaconst{Nt}\ \isafree{B}\ {\isacharhash}{\kern0pt}\ \isaconst{map}\ \isaconst{Tm}\ \isafree{w}}, \isa{\isaconst{P}\ {\isasymturnstile}\ \isafree{{\isasymbeta}}\ {\isasymRightarrow}r{\isacharasterisk}{\kern0pt}\ \isaconst{map}\ \isaconst{Tm}\ \isafree{u}}, and \isa{\isafree{w}\ {\isacharequal}{\kern0pt}\ \isafree{u}\ {\isacharat}{\kern0pt}\ \isafree{v}}.
\begin{proof}
Trivial by rule inversion.
\end{proof}
\end{lemma}

\begin{lemma}\label{rm_chain_second_produces_hd}
If \isa{\isaconst{P{\isacharprime}{\kern0pt}}\ {\isasymturnstile}\ \isafree{{\isasymalpha}\isactrlsub {\isadigit{0}}}\ {\isasymmidarrow}{\isacharbrackleft}{\kern0pt}\isafree{A}\ {\isasymrightarrow}\ \isafree{{\isasymalpha}}\ {\isasymcdot}\ \isaconst{Nt}\ \isafree{B}\ {\isacharhash}{\kern0pt}\ \isafree{{\isasymbeta}}{\isacharbrackright}{\kern0pt}\ {\isacharhash}{\kern0pt}\ \isafree{i}\ {\isacharhash}{\kern0pt}\ \isafree{{\isasymrho}}{\isasymrightarrow}r{\isacharasterisk}{\kern0pt}\ \isafree{{\isasymgamma}}},
then there exist \isa{X{\isacharcomma}{\kern0pt}\ {\isasymalpha}{\isacharprime}{\kern0pt}{\isacharcomma}{\kern0pt}} and \isa{{\isasymbeta}{\isacharprime}{\kern0pt}} such that \isa{i\ {\isacharequal}{\kern0pt}\ {\isacharbrackleft}{\kern0pt}X\ {\isasymrightarrow}\ {\isasymalpha}{\isacharprime}{\kern0pt}\ {\isasymcdot}\ Nt\ A\ {\isacharhash}{\kern0pt}\ {\isasymbeta}{\isacharprime}{\kern0pt}{\isacharbrackright}{\kern0pt}}.
\begin{proof}
By rule inversion, we know there exist \isa{{\isasymalpha}{\isacharprime}{\kern0pt}\ {\isacharcolon}{\kern0pt}{\isacharcolon}{\kern0pt}\ syms} and \isa{v{\isacharcomma}{\kern0pt}\ u\ {\isacharcolon}{\kern0pt}{\isacharcolon}{\kern0pt}\ {\isacharprime}{\kern0pt}t\ list} where
\begin{subequations}
\begin{gather}
\isa{\isafree{{\isasymgamma}}\ {\isacharequal}{\kern0pt}\ \isafree{{\isasymalpha}{\isacharprime}{\kern0pt}}\ {\isacharat}{\kern0pt}\ \isafree{{\isasymalpha}}\ {\isacharat}{\kern0pt}\ \isaconst{Nt}\ \isafree{B}\ {\isacharhash}{\kern0pt}\ \isaconst{map}\ \isaconst{Tm}\ \isafree{u}\ {\isacharat}{\kern0pt}\ \isaconst{map}\ \isaconst{Tm}\ \isafree{v}}\\
\isa{\isaconst{P{\isacharprime}{\kern0pt}}\ {\isasymturnstile}\ \isafree{{\isasymalpha}\isactrlsub {\isadigit{0}}}\ {\isasymmidarrow}\isafree{i}\ {\isacharhash}{\kern0pt}\ \isafree{{\isasymrho}}{\isasymrightarrow}r{\isacharasterisk}{\kern0pt}\ \isafree{{\isasymalpha}{\isacharprime}{\kern0pt}}\ {\isacharat}{\kern0pt}\ \isaconst{Nt}\ \isafree{A}\ {\isacharhash}{\kern0pt}\ \isaconst{map}\ \isaconst{Tm}\ \isafree{v}}\label{rmc_snd_hd.step}\\
\isa{\isaconst{P{\isacharprime}{\kern0pt}}\ {\isasymturnstile}\ \isafree{{\isasymalpha}{\isacharprime}{\kern0pt}}\ {\isacharat}{\kern0pt}\ \isaconst{Nt}\ \isafree{A}\ {\isacharhash}{\kern0pt}\ \isaconst{map}\ \isaconst{Tm}\ \isafree{v}\ {\isasymRightarrow}r\ \isafree{{\isasymalpha}{\isacharprime}{\kern0pt}}\ {\isacharat}{\kern0pt}\ \isafree{{\isasymalpha}}\ {\isacharat}{\kern0pt}\ \isaconst{Nt}\ \isafree{B}\ {\isacharhash}{\kern0pt}\ \isafree{{\isasymbeta}}\ {\isacharat}{\kern0pt}\ \isaconst{map}\ \isaconst{Tm}\ \isafree{v}}\\
\isa{\isaconst{P{\isacharprime}{\kern0pt}}\ {\isasymturnstile}\ \isafree{{\isasymbeta}}\ {\isasymRightarrow}r{\isacharasterisk}{\kern0pt}\ \isaconst{map}\ \isaconst{Tm}\ \isafree{u}}
\end{gather}
\end{subequations}
The implication then follows from by a second rule inversion, this time on \eqref{rmc_snd_hd.step}.
\end{proof}
\end{lemma}

\begin{lemma}\label{ipda_reaches_final_imp_rm_chain}
If \isa{{\isacharparenleft}{\kern0pt}{\isacharbrackleft}{\kern0pt}\isafree{A}\ {\isasymrightarrow}\ \isafree{{\isasymalpha}}\ {\isasymcdot}\ \isafree{{\isasymbeta}}{\isacharbrackright}{\kern0pt}\ {\isacharhash}{\kern0pt}\ \isafree{{\isasymrho}}{\isacharcomma}{\kern0pt}\ \isafree{w}{\isacharparenright}{\kern0pt}\ {\isasymturnstile}I{\isacharasterisk}{\kern0pt}\ {\isacharparenleft}{\kern0pt}{\isacharbrackleft}{\kern0pt}{\isacharbrackleft}{\kern0pt}\isaconst{S{\isacharprime}{\kern0pt}}\ {\isasymrightarrow}\ {\isacharbrackleft}{\kern0pt}\isaconst{Nt}\ \isaconst{S}{\isacharbrackright}{\kern0pt}\ {\isasymcdot}\ {\isacharbrackright}{\kern0pt}{\isacharbrackright}{\kern0pt}{\isacharcomma}{\kern0pt}\ {\isacharbrackleft}{\kern0pt}{\isacharbrackright}{\kern0pt}{\isacharparenright}{\kern0pt}}, then either \mbox{\isa{\isafree{{\isasymrho}}\ {\isacharequal}{\kern0pt}\ {\isacharbrackleft}{\kern0pt}{\isacharbrackright}{\kern0pt}}},
or there exist \isa{{\isasymsigma}\ {\isacharcolon}{\kern0pt}{\isacharcolon}{\kern0pt}\ item\ list}, \isa{X\ {\isacharcolon}{\kern0pt}{\isacharcolon}{\kern0pt}\ {\isacharprime}{\kern0pt}n} and \isa{{\isasymalpha}{\isacharprime}{\kern0pt}{\isacharcomma}{\kern0pt}\ {\isasymbeta}{\isacharprime}{\kern0pt}{\isacharcomma}{\kern0pt}\ {\isasymgamma}\ {\isacharcolon}{\kern0pt}{\isacharcolon}{\kern0pt}\ syms} such that
\[ \isa{\isafree{{\isasymrho}}\ {\isacharequal}{\kern0pt}\ {\isacharbrackleft}{\kern0pt}\isafree{X}\ {\isasymrightarrow}\ \isafree{{\isasymalpha}{\isacharprime}{\kern0pt}}\ {\isasymcdot}\ \isaconst{Nt}\ \isafree{A}\ {\isacharhash}{\kern0pt}\ \isafree{{\isasymbeta}{\isacharprime}{\kern0pt}}{\isacharbrackright}{\kern0pt}\ {\isacharhash}{\kern0pt}\ \isafree{{\isasymsigma}}} \text{ and } \isa{\isaconst{P{\isacharprime}{\kern0pt}}\ {\isasymturnstile}\ {\isacharbrackleft}{\kern0pt}\isaconst{Nt}\ \isaconst{S{\isacharprime}{\kern0pt}}{\isacharbrackright}{\kern0pt}\ {\isasymmidarrow}\isafree{{\isasymrho}}{\isasymrightarrow}r{\isacharasterisk}{\kern0pt}\ \isafree{{\isasymgamma}}}. \]
\begin{proof}
We perform backwards induction on the length of the computation of \isa{\isaconst{I\isactrlsub G}} for arbitrary 
\isa{A{\isacharcomma}{\kern0pt}\ {\isasymalpha}{\isacharcomma}{\kern0pt}\ {\isasymbeta}{\isacharcomma}{\kern0pt}\ {\isasymrho}}, and \isa{w}.

The reflexive case implies directly that \isa{{\isasymrho}\ {\isacharequal}{\kern0pt}\ {\isacharbrackleft}{\kern0pt}{\isacharbrackright}{\kern0pt}}.

For the transitive case, the computation is of the form
\begin{equation}\label{ipda_rmc.step}
\isa{{\isacharparenleft}{\kern0pt}{\isacharbrackleft}{\kern0pt}A\ {\isasymrightarrow}\ {\isasymalpha}\ {\isasymcdot}\ {\isasymbeta}{\isacharbrackright}{\kern0pt}\ {\isacharhash}{\kern0pt}\ {\isasymrho}{\isacharcomma}{\kern0pt}\ w{\isacharparenright}{\kern0pt}\ {\isasymturnstile}I\ {\isacharparenleft}{\kern0pt}{\isacharbrackleft}{\kern0pt}B\ {\isasymrightarrow}\ {\isasymgamma}\ {\isasymcdot}\ {\isasymdelta}{\isacharbrackright}{\kern0pt}\ {\isacharhash}{\kern0pt}\ {\isasymtau}{\isacharcomma}{\kern0pt}\ v{\isacharparenright}{\kern0pt}\ {\isasymturnstile}I{\isacharasterisk}{\kern0pt}\ {\isacharparenleft}{\kern0pt}{\isacharbrackleft}{\kern0pt}}\isa{{\isacharbrackleft}{\kern0pt}\isaconst{S{\isacharprime}{\kern0pt}}\ {\isasymrightarrow}\ {\isacharbrackleft}{\kern0pt}\isaconst{Nt}\ \isaconst{S}{\isacharbrackright}{\kern0pt}\ {\isasymcdot}\ {\isacharbrackright}{\kern0pt}}\isa{{\isacharbrackright}{\kern0pt}{\isacharcomma}{\kern0pt}\ {\isacharbrackleft}{\kern0pt}{\isacharbrackright}{\kern0pt}{\isacharparenright}{\kern0pt}}
\end{equation}
for some \isa{B}, \isa{{\isasymgamma}}, \isa{{\isasymdelta}}, \isa{{\isasymtau}} and \isa{v}. This is due to the fact that in all three types of transition, 
\isa{I\isactrlsub G} replaces a nonzero amount of topmost stack symbols by a nonempty list, meaning that a step can 
never lead to an empty stack. We can now apply the induction hypothesis on the shorter computation 
starting on \isa{{\isacharparenleft}{\kern0pt}{\isacharbrackleft}{\kern0pt}\isafree{B}\ {\isasymrightarrow}\ \isafree{{\isasymgamma}}\ {\isasymcdot}\ \isafree{{\isasymdelta}}{\isacharbrackright}{\kern0pt}\ {\isacharhash}{\kern0pt}\ \isafree{{\isasymtau}}{\isacharcomma}{\kern0pt}\ \isafree{v}{\isacharparenright}{\kern0pt}}, meaning \isa{{\isasymtau}} is either empty, or it is in some rightmost
chain as we already described.

If \isa{{\isasymtau}\ {\isacharequal}{\kern0pt}\ {\isacharbrackleft}{\kern0pt}{\isacharbrackright}{\kern0pt}}, by Lemma~\ref{reaches_final_imp_last_is_init_or_final} we know that \isa{{\isacharbrackleft}{\kern0pt}\isafree{B}\ {\isasymrightarrow}\ \isafree{{\isasymgamma}}\ {\isasymcdot}\ \isafree{{\isasymdelta}}{\isacharbrackright}{\kern0pt}}
is either the initial or final state. We can prove this item cannot be the initial state by 
contradiction using \eqref{ipda_rmc.step}. Therefore, \isa{{\isacharbrackleft}{\kern0pt}\isafree{B}\ {\isasymrightarrow}\ \isafree{{\isasymgamma}}\ {\isasymcdot}\ \isafree{{\isasymdelta}}{\isacharbrackright}{\kern0pt}} must be the final state 
\isa{{\isacharbrackleft}{\kern0pt}\isaconst{S{\isacharprime}{\kern0pt}}\ {\isasymrightarrow}\ {\isacharbrackleft}{\kern0pt}\isaconst{Nt}\ \isaconst{S}{\isacharbrackright}{\kern0pt}\ {\isasymcdot}\ {\isacharbrackright}{\kern0pt}}, implying that \isa{{\isacharbrackleft}{\kern0pt}\isafree{A}\ {\isasymrightarrow}\ \isafree{{\isasymalpha}}\ {\isasymcdot}\ \isafree{{\isasymbeta}}{\isacharbrackright}{\kern0pt}\ {\isacharequal}{\kern0pt}\ {\isacharbrackleft}{\kern0pt}\isaconst{S}\ {\isasymrightarrow}\ \isafree{{\isasymalpha}}\ {\isasymcdot}\ {\isacharbrackright}{\kern0pt}} and 
\isa{{\isasymrho}\ {\isacharequal}{\kern0pt}\ {\isacharbrackleft}{\kern0pt}S{\isacharprime}{\kern0pt}\ {\isasymrightarrow}\ {\isasymcdot}\ {\isacharbrackleft}{\kern0pt}Nt\ S{\isacharbrackright}{\kern0pt}{\isacharbrackright}{\kern0pt}\ {\isacharhash}{\kern0pt}\ {\isacharbrackleft}{\kern0pt}{\isacharbrackright}{\kern0pt}} by Lemma~\ref{step_reaches_final_imp_S} for \isa{{\isasymsigma}\ {\isacharequal}{\kern0pt}\ {\isacharbrackleft}{\kern0pt}{\isacharbrackright}{\kern0pt}}. 
Therefore, \isa{{\isasymrho}} fulfills our claim.

For the second case, \isa{{\isasymtau}} has the structure
\begin{gather}\label{ipda_rm.ih}
\isa{\isafree{{\isasymtau}}\ {\isacharequal}{\kern0pt}\ {\isacharbrackleft}{\kern0pt}\isafree{X}\ {\isasymrightarrow}\ \isafree{{\isasymalpha}{\isacharprime}{\kern0pt}}\ {\isasymcdot}\ \isaconst{Nt}\ \isafree{B}\ {\isacharhash}{\kern0pt}\ \isafree{{\isasymbeta}{\isacharprime}{\kern0pt}}{\isacharbrackright}{\kern0pt}\ {\isacharhash}{\kern0pt}\ \isafree{{\isasymsigma}}}\\
\isa{\isaconst{P{\isacharprime}{\kern0pt}}\ {\isasymturnstile}\ {\isacharbrackleft}{\kern0pt}\isaconst{Nt}\ \isaconst{S{\isacharprime}{\kern0pt}}{\isacharbrackright}{\kern0pt}\ {\isasymmidarrow}\isafree{{\isasymtau}}{\isasymrightarrow}r{\isacharasterisk}{\kern0pt}\ \isafree{{\isasymzeta}}}
\end{gather}
for some \isa{X\ {\isacharcolon}{\kern0pt}{\isacharcolon}{\kern0pt}\ {\isacharprime}{\kern0pt}n}, \isa{{\isasymalpha}{\isacharprime}{\kern0pt}{\isacharcomma}{\kern0pt}\ {\isasymbeta}{\isacharcomma}{\kern0pt}\ {\isasymzeta}\ {\isacharcolon}{\kern0pt}{\isacharcolon}{\kern0pt}\ syms} and \isa{{\isasymsigma}\ {\isacharcolon}{\kern0pt}{\isacharcolon}{\kern0pt}\ item\ list}. We can now perform a case distinction 
on the step \isa{{\isacharparenleft}{\kern0pt}{\isacharbrackleft}{\kern0pt}\isafree{A}\ {\isasymrightarrow}\ \isafree{{\isasymalpha}}\ {\isasymcdot}\ \isafree{{\isasymbeta}}{\isacharbrackright}{\kern0pt}\ {\isacharhash}{\kern0pt}\ \isafree{{\isasymrho}}{\isacharcomma}{\kern0pt}\ \isafree{w}{\isacharparenright}{\kern0pt}\ {\isasymturnstile}I\ {\isacharparenleft}{\kern0pt}{\isacharbrackleft}{\kern0pt}\isafree{B}\ {\isasymrightarrow}\ \isafree{{\isasymgamma}}\ {\isasymcdot}\ \isafree{{\isasymdelta}}{\isacharbrackright}{\kern0pt}\ {\isacharhash}{\kern0pt}\ \isafree{{\isasymtau}}{\isacharcomma}{\kern0pt}\ \isafree{v}{\isacharparenright}{\kern0pt}}.

If the transition is shifting, we have \isa{\isafree{A}\ {\isacharequal}{\kern0pt}\ \isafree{B}} and \isa{\isafree{{\isasymrho}}\ {\isacharequal}{\kern0pt}\ \isafree{{\isasymtau}}}. The implication holds by 
\eqref{ipda_rm.ih}.

If the transition is reducing, we know that
\begin{equation}\label{ipda_rm.r}
\isa{{\isacharparenleft}{\kern0pt}{\isacharbrackleft}{\kern0pt}\isafree{A}\ {\isasymrightarrow}\ \isafree{{\isasymalpha}}\ {\isasymcdot}\ \isafree{{\isasymbeta}}{\isacharbrackright}{\kern0pt}\ {\isacharhash}{\kern0pt}\ \isafree{{\isasymrho}}{\isacharcomma}{\kern0pt}\ \isafree{w}{\isacharparenright}{\kern0pt}\ {\isacharequal}{\kern0pt}\ {\isacharparenleft}{\kern0pt}{\isacharbrackleft}{\kern0pt}\isafree{A}\ {\isasymrightarrow}\ \isafree{{\isasymalpha}}\ {\isasymcdot}\ {\isacharbrackright}{\kern0pt}\ {\isacharhash}{\kern0pt}\ {\isacharbrackleft}{\kern0pt}\isafree{B}\ {\isasymrightarrow}\ \isafree{{\isasymtheta}}\ {\isasymcdot}\ \isaconst{Nt}\ \isafree{A}\ {\isacharhash}{\kern0pt}\ \isafree{{\isasymdelta}}{\isacharbrackright}{\kern0pt}\ {\isacharhash}{\kern0pt}\ \isafree{{\isasymtau}}{\isacharcomma}{\kern0pt}\ \isafree{w}{\isacharparenright}{\kern0pt}},
\end{equation}
for some \isa{{\isasymtheta}} where \isa{\isafree{{\isasymgamma}}\ {\isacharequal}{\kern0pt}\ \isafree{{\isasymtheta}}\ {\isacharat}{\kern0pt}\ {\isacharbrackleft}{\kern0pt}\isaconst{Nt}\ \isafree{A}{\isacharbrackright}{\kern0pt}}, meaning that \isa{{\isacharparenleft}{\kern0pt}B{\isacharcomma}{\kern0pt}\ {\isasymtheta}\ {\isacharat}{\kern0pt}\ Nt\ A\ {\isacharhash}{\kern0pt}\ {\isasymdelta}{\isacharparenright}{\kern0pt}} is in \isa{P{\isacharprime}{\kern0pt}}. By
Lemma~\ref{rm_chain_Cons_imp_prod_rightmost}, \eqref{ipda_rm.ih} implies that \isa{{\isasymzeta}} is of the form
\[ \isa{\isafree{{\isasymzeta}}\ {\isacharequal}{\kern0pt}\ \isafree{{\isasymzeta}{\isacharprime}{\kern0pt}}\ {\isacharat}{\kern0pt}\ \isaconst{Nt}\ \isafree{B}\ {\isacharhash}{\kern0pt}\ \isaconst{map}\ \isaconst{Tm}\ \isafree{u}} \]
for some \isa{{\isasymzeta}{\isacharprime}{\kern0pt}} and \isa{u}. Moreover, since \isa{G{\isacharprime}{\kern0pt}} is reduced, there exists a \isa{v\ {\isacharcolon}{\kern0pt}{\isacharcolon}{\kern0pt}\ {\isacharprime}{\kern0pt}t\ list} that the string \isa{{\isasymdelta}} can derive.  
With the fact that \isa{B} produces \isa{{\isasymtheta}\ {\isacharat}{\kern0pt}\ Nt\ A\ {\isacharhash}{\kern0pt}\ {\isasymdelta}}, we can extend the chain in \eqref{ipda_rm.ih} by the 
item \isa{{\isacharbrackleft}{\kern0pt}\isafree{B}\ {\isasymrightarrow}\ \isafree{{\isasymtheta}}\ {\isasymcdot}\ \isaconst{Nt}\ \isafree{A}\ {\isacharhash}{\kern0pt}\ \isafree{{\isasymdelta}}{\isacharbrackright}{\kern0pt}}, i.e.
\begin{equation*}
\isa{\isaconst{P{\isacharprime}{\kern0pt}}\ {\isasymturnstile}\ {\isacharbrackleft}{\kern0pt}\isaconst{Nt}\ \isaconst{S{\isacharprime}{\kern0pt}}{\isacharbrackright}{\kern0pt}\ {\isasymmidarrow}{\isacharbrackleft}{\kern0pt}\isafree{B}\ {\isasymrightarrow}\ \isafree{{\isasymtheta}}\ {\isasymcdot}\ \isaconst{Nt}\ \isafree{A}\ {\isacharhash}{\kern0pt}\ \isafree{{\isasymdelta}}{\isacharbrackright}{\kern0pt}\ {\isacharhash}{\kern0pt}\ \isafree{{\isasymtau}}{\isasymrightarrow}r{\isacharasterisk}{\kern0pt}\ \isafree{{\isasymzeta}{\isacharprime}{\kern0pt}}\ {\isacharat}{\kern0pt}\ \isafree{{\isasymtheta}}\ {\isacharat}{\kern0pt}\ \isaconst{Nt}\ \isafree{A}\ {\isacharhash}{\kern0pt}\ \isaconst{map}\ \isaconst{Tm}\ {\isacharparenleft}{\kern0pt}\isafree{v}\ {\isacharat}{\kern0pt}\ \isafree{u}{\isacharparenright}{\kern0pt}}
\end{equation*}
The implication therefore holds with \eqref{ipda_rm.r}.

If the transition is expanding, we have
\begin{equation}\label{ipda_rm.e}
 \isa{{\isacharparenleft}{\kern0pt}{\isacharbrackleft}{\kern0pt}B\ {\isasymrightarrow}\ {\isasymgamma}\ {\isasymcdot}\ {\isasymdelta}{\isacharbrackright}{\kern0pt}\ {\isacharhash}{\kern0pt}\ {\isasymtau}{\isacharcomma}{\kern0pt}\ v{\isacharparenright}{\kern0pt}\ {\isacharequal}{\kern0pt}\ {\isacharparenleft}{\kern0pt}{\isacharbrackleft}{\kern0pt}B\ {\isasymrightarrow}\ {\isacharbrackleft}{\kern0pt}{\isacharbrackright}{\kern0pt}\ {\isasymcdot}\ {\isasymdelta}{\isacharbrackright}{\kern0pt}\ {\isacharhash}{\kern0pt}\ {\isacharbrackleft}{\kern0pt}A\ {\isasymrightarrow}\ {\isasymalpha}\ {\isasymcdot}\ {\isasymbeta}{\isacharbrackright}{\kern0pt}\ {\isacharhash}{\kern0pt}\ {\isasymrho}{\isacharcomma}{\kern0pt}\ w{\isacharparenright}{\kern0pt}}.
\end{equation}
If \isa{{\isasymrho}\ {\isacharequal}{\kern0pt}\ {\isacharbrackleft}{\kern0pt}{\isacharbrackright}{\kern0pt}}, the implication holds directly. Otherwise, if \isa{{\isasymrho}\ {\isacharequal}{\kern0pt}\ i\ {\isacharhash}{\kern0pt}\ {\isasymsigma}} for some \isa{i} and \isa{{\isasymsigma}}, 
from \eqref{ipda_rm.ih} and \eqref{ipda_rm.e} we get
\[ \isa{\isaconst{P{\isacharprime}{\kern0pt}}\ {\isasymturnstile}\ {\isacharbrackleft}{\kern0pt}\isaconst{Nt}\ \isaconst{S{\isacharprime}{\kern0pt}}{\isacharbrackright}{\kern0pt}\ {\isasymmidarrow}{\isacharbrackleft}{\kern0pt}\isafree{A}\ {\isasymrightarrow}\ \isafree{{\isasymalpha}}\ {\isasymcdot}\ \isafree{{\isasymbeta}}{\isacharbrackright}{\kern0pt}\ {\isacharhash}{\kern0pt}\ \isafree{i}\ {\isacharhash}{\kern0pt}\ \isafree{{\isasymsigma}}{\isasymrightarrow}r{\isacharasterisk}{\kern0pt}\ \isafree{{\isasymzeta}}}. \]
Finally, by inversion of the step rule for rightmost chains, this implies the existence of some \isa{{\isasymeta}} 
such that
\[ \isa{\isaconst{P{\isacharprime}{\kern0pt}}\ {\isasymturnstile}\ {\isacharbrackleft}{\kern0pt}\isaconst{Nt}\ \isaconst{S{\isacharprime}{\kern0pt}}{\isacharbrackright}{\kern0pt}\ {\isasymmidarrow}\isafree{{\isasymrho}}{\isasymrightarrow}r{\isacharasterisk}{\kern0pt}\ \isafree{{\isasymeta}}},\] 
and by Lemma~\ref{rm_chain_second_produces_hd}, \isa{i} is of the form \isa{\isafree{i}\ {\isacharequal}{\kern0pt}\ {\isacharbrackleft}{\kern0pt}\isafree{Z}\ {\isasymrightarrow}\ \isafree{{\isasymgamma}{\isacharprime}{\kern0pt}}\ {\isasymcdot}\ \isaconst{Nt}\ \isafree{A}\ {\isacharhash}{\kern0pt}\ \isafree{{\isasymdelta}{\isacharprime}{\kern0pt}}{\isacharbrackright}{\kern0pt}} 
for some \isa{Z{\isacharcomma}{\kern0pt}\ {\isasymgamma}{\isacharprime}{\kern0pt}}, and \isa{{\isasymdelta}{\isacharprime}{\kern0pt}}. Together with the existence of \isa{{\isasymeta}}, this completes the proof.
\end{proof}
\end{lemma}

We have now formalized the notion of rightmost chains, proved their existence is equivalent to that
of generic rightmost derivations, and proved that every IPDA stack that reaches an accepting state 
has such a chain of items on the stack. With these chains, we are now able to describe a very 
well-behaved relation between the nonterminals directly to the right of the bullet and the 
nonterminal to the left of the arrow of the preceding item: for any two neighboring items
in the chain's item list 
\[ \isa{{\isasymdots}\ {\isacharhash}{\kern0pt}\ {\isacharbrackleft}{\kern0pt}A\ {\isasymrightarrow}\ {\isasymalpha}\ {\isasymcdot}\ {\isasymbeta}{\isacharbrackright}{\kern0pt}\ {\isacharhash}{\kern0pt}\ {\isacharbrackleft}{\kern0pt}B\ {\isasymrightarrow}\ {\isasymgamma}\ {\isasymcdot}\ {\isasymdelta}{\isacharbrackright}{\kern0pt}\ {\isacharhash}{\kern0pt}\ {\isasymdots}} \]
we know that \isa{\isaconst{hd}\ \isafree{{\isasymdelta}}\ {\isacharequal}{\kern0pt}\ \isaconst{Nt}\ \isafree{A}}, meaning the LHS nonterminal of each item is produced in the 
successor item. Wilhelm et al.\ use these chains on multiple occasions only informally. By formalizing
them, we will be able in the future to induct on them more rigorously, which will be of utmost 
importance in the coming section.%
\end{isamarkuptext}\isamarkuptrue%
%
\isadelimdocument
%
\endisadelimdocument
%
\isatagdocument
%
\isamarkupchapter{The Characteristic Finite Automaton and the Canonical \isa{LR{\isacharparenleft}{\kern0pt}{\isadigit{0}}{\isacharparenright}{\kern0pt}} Automaton%
}
\isamarkuptrue%
%
\endisatagdocument
{\isafolddocument}%
%
\isadelimdocument
%
\endisadelimdocument
%
\begin{isamarkuptext}%
In this section, we will define the finite automata that the \isa{LR{\isacharparenleft}{\kern0pt}{\isadigit{0}}{\isacharparenright}{\kern0pt}} parser will be based on, 
and show certain properties that are fundamental to later prove the main theorems of this thesis.%
\end{isamarkuptext}\isamarkuptrue%
%
\isadelimdocument
%
\endisadelimdocument
%
\isatagdocument
%
\isamarkupsection{The Characteristic Finite Automaton%
}
\isamarkuptrue%
%
\endisatagdocument
{\isafolddocument}%
%
\isadelimdocument
%
\endisadelimdocument
%
\begin{isamarkuptext}%
As we discussed in the previous section, the IPDA nondeterministically selects productions when
performing expanding transitions, while shifting and reducing transitions are deterministic. Our goal
for the parser is to explore all
expansions in parallel, and selecting productions only for reduction. To achieve this, we will
first define an automaton, in this case an NFA, that we will call the
\concept{characteristic finite automaton} to \isa{G}. We base our implementation of finite automata on
the formalization thereof by Paulson~\cite{Paulson}.

\begin{definition}[Characteristic finite automaton]
The characteristic finite automaton to \isa{G} is the \isa{{\isacharparenleft}{\kern0pt}{\isacharparenleft}{\kern0pt}\isatfree{{\isacharprime}{\kern0pt}n}{\isacharcomma}{\kern0pt}\ \isatfree{{\isacharprime}{\kern0pt}t}{\isacharparenright}{\kern0pt}\ \isatconst{sym}{\isacharcomma}{\kern0pt}\ {\isacharparenleft}{\kern0pt}\isatfree{{\isacharprime}{\kern0pt}n}{\isacharcomma}{\kern0pt}\ \isatfree{{\isacharprime}{\kern0pt}t}{\isacharparenright}{\kern0pt}\ \isatconst{item}{\isacharparenright}{\kern0pt}\ \isatconst{nfa}}:
\begin{multline*}
  \isa{\ensuremath{\mathrm{char}(G)}} = \isa{{\isasymlparr}states} = \isa{\ensuremath{\mathrm{It}_{\isaconst{G{\isacharprime}{\kern0pt}}}}},\ \isa{init\ {\isacharequal}{\kern0pt}\ {\isacharbraceleft}{\kern0pt}{\isacharbrackleft}{\kern0pt}S{\isacharprime}{\kern0pt}\ {\isasymrightarrow}\ {\isacharbrackleft}{\kern0pt}{\isacharbrackright}{\kern0pt}\ {\isasymcdot}\ {\isacharbrackleft}{\kern0pt}Nt\ S{\isacharbrackright}{\kern0pt}{\isacharbrackright}{\kern0pt}{\isacharbraceright}{\kern0pt}{\isacharcomma}{\kern0pt}\ final}
    = \isa{\isaconst{completes}\ \ensuremath{\mathrm{It}_{\isaconst{G{\isacharprime}{\kern0pt}}}}} \isa{{\isacharcomma}{\kern0pt}\ nxt\ {\isacharequal}{\kern0pt}\ {\isasymDelta}\isactrlsub c{\isacharcomma}{\kern0pt}}\\
  \isa{eps\ {\isacharequal}{\kern0pt}\ {\isasymE}\isactrlsub c{\isasymrparr}}
\end{multline*}
with
% Fix style?
%
\begin{isabelle}%
\isafree{{\isasymDelta}\isactrlsub c}\ \isafree{q}\ \isafree{a}\ {\isacharequal}{\kern0pt}\isanewline
{\isacharparenleft}{\kern0pt}\textsf{case}\ \isafree{q}\ \textsf{of}\ {\isacharbrackleft}{\kern0pt}\isabound{X}\ {\isasymrightarrow}\ \isabound{{\isasymalpha}}\ {\isasymcdot}\ {\isacharbrackleft}{\kern0pt}{\isacharbrackright}{\kern0pt}{\isacharbrackright}{\kern0pt}\ {\isasymRightarrow}\ {\isasymemptyset}\isanewline
\isaindent{{\isacharparenleft}{\kern0pt}}{\isacharbar}{\kern0pt}\ {\isacharbrackleft}{\kern0pt}\isabound{X}\ {\isasymrightarrow}\ \isabound{{\isasymalpha}}\ {\isasymcdot}\ \isabound{Y}\ {\isacharhash}{\kern0pt}\ \isabound{{\isasymbeta}}{\isacharbrackright}{\kern0pt}\ {\isasymRightarrow}\isanewline
\isaindent{{\isacharparenleft}{\kern0pt}{\isacharbar}{\kern0pt}\ \ \ }\textsf{if}\ \isafree{a}\ {\isacharequal}{\kern0pt}\ \isabound{Y}\ {\isasymand}\ {\isacharparenleft}{\kern0pt}\isabound{X}{\isacharcomma}{\kern0pt}\ \isabound{{\isasymalpha}}\ {\isacharat}{\kern0pt}\ \isabound{Y}\ {\isacharhash}{\kern0pt}\ \isabound{{\isasymbeta}}{\isacharparenright}{\kern0pt}\ {\isasymin}\ \isaconst{P{\isacharprime}{\kern0pt}}\ \textsf{then}\ {\isacharbraceleft}{\kern0pt}{\isacharbrackleft}{\kern0pt}\isabound{X}\ {\isasymrightarrow}\ \isabound{{\isasymalpha}}\ {\isacharat}{\kern0pt}\ {\isacharbrackleft}{\kern0pt}\isabound{Y}{\isacharbrackright}{\kern0pt}\ {\isasymcdot}\ \isabound{{\isasymbeta}}{\isacharbrackright}{\kern0pt}{\isacharbraceright}{\kern0pt}\ \textsf{else}\ {\isasymemptyset}{\isacharparenright}{\kern0pt}%
\end{isabelle}
and
\begin{multline*}
\isa{{\isasymE}\isactrlsub c\ {\isacharequal}{\kern0pt}\ {\isacharbraceleft}{\kern0pt}} \isa{{\isacharparenleft}{\kern0pt}{\isacharbrackleft}{\kern0pt}\isafree{X}\ {\isasymrightarrow}\ \isafree{{\isasymalpha}}\ {\isasymcdot}\ \isaconst{Nt}\ \isafree{Y}\ {\isacharhash}{\kern0pt}\ \isafree{{\isasymbeta}}{\isacharbrackright}{\kern0pt}{\isacharcomma}{\kern0pt}\ {\isacharbrackleft}{\kern0pt}\isafree{Y}\ {\isasymrightarrow}\ \ {\isasymcdot}\ \isafree{{\isasymgamma}}\ {\isacharbrackright}{\kern0pt}{\isacharparenright}{\kern0pt}}\\
  \isa{{\isacharbar}{\kern0pt}\ X\ {\isasymalpha}\ Y\ {\isasymbeta}\ {\isasymgamma}{\isachardot}{\kern0pt}\ {\isacharparenleft}{\kern0pt}X{\isacharcomma}{\kern0pt}\ {\isasymalpha}\ {\isacharat}{\kern0pt}\ Nt\ Y\ {\isacharhash}{\kern0pt}\ {\isasymbeta}{\isacharparenright}{\kern0pt}\ {\isasymin}\ P{\isacharprime}{\kern0pt}\ {\isasymand}\ {\isacharparenleft}{\kern0pt}Y{\isacharcomma}{\kern0pt}\ {\isasymgamma}{\isacharparenright}{\kern0pt}\ {\isasymin}\ P{\isacharprime}{\kern0pt}{\isacharbraceright}{\kern0pt}}
\end{multline*}
\end{definition}

The characteristic finite automaton can therefore perform shifting and expanding transitions akin to 
those of the IPDA \isa{\isaconst{I\isactrlsub G}}, while shifting nonterminals corresponds to the reducing transitions of
\isa{\isaconst{I\isactrlsub G}}. \isa{\ensuremath{\mathrm{char}(G)}} can therefore reach an item in \isa{\ensuremath{\mathrm{It}_{\isaconst{G{\isacharprime}{\kern0pt}}}}} 
by reading the concatenation of the prefixes that \isa{\isaconst{I\isactrlsub G}} processed, i.e., the concatenation of
histories of the items on the stack of \isa{\isaconst{I\isactrlsub G}}, in order to reach this particular item, as
explained by Wilhelm et al.~\cite[p. 103]{Wilhelm}

Wilhelm et al.~\cite[p. 104]{Wilhelm} present a theorem stating the equivalence of three statements 
as one of the main results in the \isa{LR{\isacharparenleft}{\kern0pt}{\isadigit{0}}{\isacharparenright}{\kern0pt}} section of the book, Theorem 3.4.1, describing the relation 
between \isa{\ensuremath{\mathrm{char}(G)}}, rightmost derivations, and the IPDA. They make the following claim, which
we have adapted to match our own notation:

\begin{quote}
\textit{Let \isa{G} be a CFG and \isa{{\isasymgamma}\ {\isacharcolon}{\kern0pt}{\isacharcolon}{\kern0pt}\ syms}. The following three statements are equivalent:
\begin{enumerate}
\item There exists a computation \isa{{\isacharparenleft}{\kern0pt}{\isacharbrackleft}{\kern0pt}\isaconst{S{\isacharprime}{\kern0pt}}\ {\isasymrightarrow}\ \ {\isasymcdot}\ {\isacharbrackleft}{\kern0pt}\isaconst{Nt}\ \isaconst{S}{\isacharbrackright}{\kern0pt}\ {\isacharbrackright}{\kern0pt}{\isacharcomma}{\kern0pt}\ \isafree{{\isasymgamma}}{\isacharparenright}{\kern0pt}\ {\isasymturnstile}c{\isacharasterisk}{\kern0pt}\ {\isacharparenleft}{\kern0pt}{\isacharbrackleft}{\kern0pt}\isafree{A}\ {\isasymrightarrow}\ \isafree{{\isasymalpha}}\ {\isasymcdot}\ \isafree{{\isasymbeta}}{\isacharbrackright}{\kern0pt}{\isacharcomma}{\kern0pt}\ {\isacharbrackleft}{\kern0pt}{\isacharbrackright}{\kern0pt}{\isacharparenright}{\kern0pt}}
of the characteristic finite automaton \isa{\ensuremath{\mathrm{char}(G)}}.
\item There exists a rightmost derivation 
\begin{equation*}
\isa{P{\isacharprime}{\kern0pt}\ {\isasymturnstile}\ {\isacharbrackleft}{\kern0pt}Nt\ S{\isacharprime}{\kern0pt}{\isacharbrackright}{\kern0pt}\ {\isasymRightarrow}r{\isacharasterisk}{\kern0pt}\ {\isasymgamma}{\isacharprime}{\kern0pt}\ {\isacharat}{\kern0pt}\ Nt\ A\ {\isacharhash}{\kern0pt}\ map\ Tm\ w\ {\isasymRightarrow}r\ {\isasymgamma}{\isacharprime}{\kern0pt}\ {\isacharat}{\kern0pt}\ {\isasymalpha}\ {\isacharat}{\kern0pt}\ {\isasymbeta}\ {\isacharat}{\kern0pt}\ map\ Tm\ w}
\end{equation*}
with \isa{\isafree{{\isasymgamma}}\ {\isacharequal}{\kern0pt}\ \isafree{{\isasymgamma}{\isacharprime}{\kern0pt}}\ {\isacharat}{\kern0pt}\ \isafree{{\isasymalpha}}}.
\item There exists a computation \isa{{\isacharparenleft}{\kern0pt}{\isacharbrackleft}{\kern0pt}\isafree{A}\ {\isasymrightarrow}\ \isafree{{\isasymalpha}}\ {\isasymcdot}\ \isafree{{\isasymbeta}}{\isacharbrackright}{\kern0pt}\ {\isacharhash}{\kern0pt}\ \isafree{{\isasymrho}}{\isacharcomma}{\kern0pt}\ \isafree{w}{\isacharparenright}{\kern0pt}\ {\isasymturnstile}I{\isacharasterisk}{\kern0pt}\ \mbox{{\isacharparenleft}{\kern0pt}{\isacharbrackleft}{\kern0pt}{\isacharbrackleft}{\kern0pt}\isaconst{S{\isacharprime}{\kern0pt}}\ {\isasymrightarrow}\ {\isacharbrackleft}{\kern0pt}\isaconst{Nt}\ \isaconst{S}{\isacharbrackright}{\kern0pt}\ {\isasymcdot}\ {\isacharbrackright}{\kern0pt}{\isacharbrackright}{\kern0pt}{\isacharcomma}{\kern0pt}\ {\isacharbrackleft}{\kern0pt}{\isacharbrackright}{\kern0pt}{\isacharparenright}{\kern0pt}}}
of the IPDA \isa{\isaconst{I\isactrlsub G}} such that \isa{\isafree{{\isasymgamma}}\ {\isacharequal}{\kern0pt}\ \isaconst{hist}\ {\isacharparenleft}{\kern0pt}\isaconst{rev}\ \isafree{{\isasymrho}}{\isacharparenright}{\kern0pt}\ {\isacharat}{\kern0pt}\ \isafree{{\isasymalpha}}} holds.
\end{enumerate}}
\end{quote}

It is worth noting that we have modified the order of the statements, since in the book there
is a mismatch between the numbering of each statement in the theorem and in the proof thereof.

Wilhelm et al.\ prove the theorem by a circular proof of the form
$(1) \Longrightarrow (2) \Longrightarrow (3) \Longrightarrow (1)$. This 
theorem, however, is not a consequence of the chain of implications that was proved. Additionally,
there are certain typographic errors in the proofs which we will not address; we infer the author's
intention in such cases.

Implication $(1) \Longrightarrow (2)$ uses existing constants \isa{S{\isacharprime}{\kern0pt}} and \isa{S}, and fixes variables \isa{{\isasymgamma}} and
\isa{{\isacharbrackleft}{\kern0pt}\isafree{A}\ {\isasymrightarrow}\ \isafree{{\isasymalpha}}\ {\isasymcdot}\ \isafree{{\isasymbeta}}{\isacharbrackright}{\kern0pt}\ {\isasymin}\ \ensuremath{\mathrm{It}_{\isaconst{G{\isacharprime}{\kern0pt}}}}}. The conclusion then shows that there exist some \isa{{\isasymgamma}{\isacharprime}{\kern0pt}\ {\isacharcolon}{\kern0pt}{\isacharcolon}{\kern0pt}\ syms} and
\isa{w\ {\isacharcolon}{\kern0pt}{\isacharcolon}{\kern0pt}\ {\isacharprime}{\kern0pt}t\ list} for which the rightmost derivation exists and \isa{\isafree{{\isasymgamma}}\ {\isacharequal}{\kern0pt}\ \isafree{{\isasymgamma}{\isacharprime}{\kern0pt}}\ {\isacharat}{\kern0pt}\ \isafree{{\isasymalpha}}}. The proof itself
is flawed: the authors induct on the number of \isa{{\isasymepsilon}}-transitions instead, and the 
induction hypothesis, which is only established if the remaining input is \isa{{\isacharbrackleft}{\kern0pt}{\isacharbrackright}{\kern0pt}}, is
applied on a configuration with nonempty remaining input. Furthermore, the base case does not consider 
the scenario where \isa{\ensuremath{\mathrm{char}(G)}} reads \isa{Nt\ S} from the input before making any \isa{{\isasymepsilon}}-transitions. 
We will later prove the implication differently. 

Implication $(2) \Longrightarrow (3)$ fixes \isa{{\isasymgamma}{\isacharprime}{\kern0pt}{\isacharcomma}{\kern0pt}\ A{\isacharcomma}{\kern0pt}\ {\isasymalpha}{\isacharcomma}{\kern0pt}\ {\isasymbeta}{\isacharcomma}{\kern0pt}} and \isa{w}, and shows the existence of some \isa{{\isasymrho}}. Note
that the conclusion of $(1) \Longrightarrow (2)$ proved only an existentially quantified \isa{w}; this is an 
important distinction that we will address later.
In the statements \isa{\isafree{{\isasymgamma}}\ {\isacharequal}{\kern0pt}\ \isafree{{\isasymgamma}{\isacharprime}{\kern0pt}}\ {\isacharat}{\kern0pt}\ \isafree{{\isasymalpha}}} and \isa{\isafree{{\isasymgamma}}\ {\isacharequal}{\kern0pt}\ \isaconst{hist}\ {\isacharparenleft}{\kern0pt}\isaconst{rev}\ \isafree{{\isasymrho}}{\isacharparenright}{\kern0pt}\ {\isacharat}{\kern0pt}\ \isafree{{\isasymalpha}}}, variable \isa{{\isasymgamma}} is completely
meaningless; it bears no relation to the \isa{{\isasymgamma}} fixed in (1) since no \isa{{\isasymgamma}} is defined in (2) nor in (3). 
Therefore, this can be simplified to 
\isa{\isafree{{\isasymgamma}{\isacharprime}{\kern0pt}}\ {\isacharequal}{\kern0pt}\ \isaconst{hist}\ {\isacharparenleft}{\kern0pt}\isaconst{rev}\ \isafree{{\isasymrho}}{\isacharparenright}{\kern0pt}}. Furthermore, the proof for this implication assumes that 
\isa{\isaconst{P{\isacharprime}{\kern0pt}}\ {\isasymturnstile}\ \isafree{{\isasymbeta}}\ {\isasymRightarrow}r{\isacharasterisk}{\kern0pt}\ \isaconst{map}\ \isaconst{Tm}\ \isafree{v}}
for some \isa{v\ {\isacharcolon}{\kern0pt}{\isacharcolon}{\kern0pt}\ {\isacharprime}{\kern0pt}t\ list}. It is not clear what the statement intends to say, and this assumption
might need to be included in statement (2) for it to be correct, depending on what was meant 
originally: since \isa{w} is fixed, the statement of the existence of the IPDA computation seems to 
be stating that the computation holds for the same fixed \isa{w}, but this is not true; as we will see, 
the statement holds for input \isa{\isafree{v}\ {\isacharat}{\kern0pt}\ \isafree{w}}, where \isa{v} is again the word derived by \isa{{\isasymbeta}}. If \isa{w} is 
meant to be existentially quantified in the conclusion, this is not a problem, but this is not made
clear by the authors. The proof of the statement is also incorrect: it uses the book's informal
equivalent of our formalized rightmost chains, but states that one proves by induction that a 
configuration with topmost stack symbol $[X_{n-1}\ \to\ \alpha_n\ \cdot\ Nt\ A\ \#\ \beta_n]$
and input \isa{v\ {\isacharat}{\kern0pt}\ w} reaches the final configuration~\footnote{the authors naturally state the structure of
the entire stack, but only the topmost symbol is relevant for our argument.}. However, this 
configuration would first push the initial item \isa{{\isacharbrackleft}{\kern0pt}\isafree{A}\ {\isasymrightarrow}\ \ {\isasymcdot}\ \isafree{{\isasymalpha}}\ {\isacharat}{\kern0pt}\ \isafree{{\isasymbeta}}\ {\isacharbrackright}{\kern0pt}} onto the stack (this is the 
only possibility based on our assumptions about \isa{A}), and since we haven't made any assumptions 
regarding \isa{{\isasymalpha}}, we cannot make any further claims. It is therefore necessary for the topmost stack 
item to be \isa{{\isacharbrackleft}{\kern0pt}\isafree{A}\ {\isasymrightarrow}\ \isafree{{\isasymalpha}}\ {\isasymcdot}\ \isafree{{\isasymbeta}}{\isacharbrackright}{\kern0pt}} to be able to make any claims about the computation.

The proof of the final implication is fundamentally correct; it fixes \isa{{\isacharbrackleft}{\kern0pt}\isafree{A}\ {\isasymrightarrow}\ \isafree{{\isasymalpha}}\ {\isasymcdot}\ \isafree{{\isasymbeta}}{\isacharbrackright}{\kern0pt}} as 
usual, as well as \isa{{\isasymrho}} and \isa{w}. The conclusion does not contain any new variables since \isa{{\isasymgamma}} is just 
a shorthand for \isa{\isaconst{hist}\ {\isacharparenleft}{\kern0pt}\isaconst{rev}\ \isafree{{\isasymrho}}{\isacharparenright}{\kern0pt}\ {\isacharat}{\kern0pt}\ \isafree{{\isasymalpha}}} in practice. However, the main issue with this proof is,
as we previously mentioned, that the \isa{w} in the premise is not in the conclusion at all. In order 
to prove the equivalence of (2) and (3), we need the statement $(2) \Longrightarrow (3)$, which we have already
proved, and $(3) \Longrightarrow (2)$, which we prove through (1), i.e.,
$(3) \Longrightarrow (1) \Longrightarrow (2)$. However, since
the original \isa{w} in (3) is lost in (1), it is a new existentially quantified \isa{w} in (2). This means
that the statements (2) and (3) are not equivalent for the same \isa{w}, as it reads in the book, 
but for two different existentially quantified words \isa{u} and \isa{v}. However, the 
main results that we need from the theorem to formalize our parser still hold after addressing 
these errors. A corrected version Theorem 3.4.1:

\begin{theorem}[Equivalences between \isa{\ensuremath{\mathrm{char}(G)}}, rightmost derivations, and \isa{\isaconst{I\isactrlsub G}}]\label{char_derivers_ipda_iffs}
The following statements are equivalent:
\begin{enumerate}
\item There exists a computation \isa{{\isacharparenleft}{\kern0pt}{\isacharbrackleft}{\kern0pt}\isaconst{S{\isacharprime}{\kern0pt}}\ {\isasymrightarrow}\ \ {\isasymcdot}\ {\isacharbrackleft}{\kern0pt}\isaconst{Nt}\ \isaconst{S}{\isacharbrackright}{\kern0pt}\ {\isacharbrackright}{\kern0pt}{\isacharcomma}{\kern0pt}\ \isafree{{\isasymgamma}}\ {\isacharat}{\kern0pt}\ \isafree{{\isasymalpha}}{\isacharparenright}{\kern0pt}\ {\isasymturnstile}c{\isacharasterisk}{\kern0pt}\ {\isacharparenleft}{\kern0pt}{\isacharbrackleft}{\kern0pt}\isafree{A}\ {\isasymrightarrow}\ \isafree{{\isasymalpha}}\ {\isasymcdot}\ \isafree{{\isasymbeta}}{\isacharbrackright}{\kern0pt}{\isacharcomma}{\kern0pt}\ {\isacharbrackleft}{\kern0pt}{\isacharbrackright}{\kern0pt}{\isacharparenright}{\kern0pt}}
of the characteristic finite automaton \isa{\ensuremath{\mathrm{char}(G)}}.
\item There exists a rightmost derivation
\begin{equation*}
\isa{P{\isacharprime}{\kern0pt}\ {\isasymturnstile}\ {\isacharbrackleft}{\kern0pt}Nt\ S{\isacharprime}{\kern0pt}{\isacharbrackright}{\kern0pt}\ {\isasymRightarrow}r{\isacharasterisk}{\kern0pt}\ {\isasymgamma}\ {\isacharat}{\kern0pt}\ Nt\ A\ {\isacharhash}{\kern0pt}\ map\ Tm\ u\ {\isasymRightarrow}r\ {\isasymgamma}\ {\isacharat}{\kern0pt}\ {\isasymalpha}\ {\isacharat}{\kern0pt}\ {\isasymbeta}\ {\isacharat}{\kern0pt}\ map\ Tm\ u}.
\end{equation*}
for some \isa{u\ {\isacharcolon}{\kern0pt}{\isacharcolon}{\kern0pt}\ {\isacharprime}{\kern0pt}t\ list}.
\item There exists a computation \isa{{\isacharparenleft}{\kern0pt}{\isacharbrackleft}{\kern0pt}\isafree{A}\ {\isasymrightarrow}\ \isafree{{\isasymalpha}}\ {\isasymcdot}\ \isafree{{\isasymbeta}}{\isacharbrackright}{\kern0pt}\ {\isacharhash}{\kern0pt}\ \isafree{{\isasymrho}}{\isacharcomma}{\kern0pt}\ \isafree{v}{\isacharparenright}{\kern0pt}\ {\isasymturnstile}I{\isacharasterisk}{\kern0pt}\ {\isacharparenleft}{\kern0pt}{\isacharbrackleft}{\kern0pt}{\isacharbrackleft}{\kern0pt}\isaconst{S{\isacharprime}{\kern0pt}}\ {\isasymrightarrow}\ {\isacharbrackleft}{\kern0pt}\isaconst{Nt}\ \isaconst{S}{\isacharbrackright}{\kern0pt}\ {\isasymcdot}\ {\isacharbrackright}{\kern0pt}{\isacharbrackright}{\kern0pt}{\isacharcomma}{\kern0pt}\ {\isacharbrackleft}{\kern0pt}{\isacharbrackright}{\kern0pt}{\isacharparenright}{\kern0pt}}
of the IPDA \isa{\isaconst{I\isactrlsub G}} for some \isa{{\isasymrho}\ {\isacharcolon}{\kern0pt}{\isacharcolon}{\kern0pt}\ item\ list} and \isa{v\ {\isacharcolon}{\kern0pt}{\isacharcolon}{\kern0pt}\ {\isacharprime}{\kern0pt}t\ list} such that 
\isa{\mbox{\isaconst{hist}\ {\isacharparenleft}{\kern0pt}\isaconst{rev}\ \isafree{{\isasymrho}}{\isacharparenright}{\kern0pt}}\ {\isacharequal}{\kern0pt}\ \isafree{{\isasymgamma}}} holds.
\end{enumerate}
\end{theorem}

We will now work towards proving this theorem. We denote \isa{\ensuremath{\mathrm{char}(G)}} transition steps by \isa{{\isasymturnstile}c}, 
and as usual, the reflexive transitive closure by \isa{{\isasymturnstile}c{\isacharasterisk}{\kern0pt}}, and \isa{n}-length computations by \isa{{\isasymturnstile}c{\isacharparenleft}{\kern0pt}n{\isacharparenright}{\kern0pt}}. 
We first need two lemmas about \isa{\ensuremath{\mathrm{char}(G)}}:

\begin{lemma}\label{char_reachable_imp_substring}
If \isa{{\isacharparenleft}{\kern0pt}{\isacharbrackleft}{\kern0pt}\isaconst{S{\isacharprime}{\kern0pt}}\ {\isasymrightarrow}\ \ {\isasymcdot}\ {\isacharbrackleft}{\kern0pt}\isaconst{Nt}\ \isaconst{S}{\isacharbrackright}{\kern0pt}\ {\isacharbrackright}{\kern0pt}{\isacharcomma}{\kern0pt}\ \isafree{{\isasymgamma}}{\isacharparenright}{\kern0pt}\ {\isasymturnstile}c{\isacharasterisk}{\kern0pt}\ {\isacharparenleft}{\kern0pt}{\isacharbrackleft}{\kern0pt}\isafree{A}\ {\isasymrightarrow}\ \isafree{{\isasymalpha}}\ {\isasymcdot}\ \isafree{{\isasymbeta}}{\isacharbrackright}{\kern0pt}{\isacharcomma}{\kern0pt}\ \isafree{{\isasymdelta}}{\isacharparenright}{\kern0pt}}, there exists a \isa{{\isasymzeta}} such that 
\isa{\isafree{{\isasymgamma}}\ {\isacharequal}{\kern0pt}\ \isafree{{\isasymzeta}}\ {\isacharat}{\kern0pt}\ \isafree{{\isasymalpha}}\ {\isacharat}{\kern0pt}\ \isafree{{\isasymdelta}}}.
\begin{proof}
By induction on the length of the computation, making a case distinction on whether the last step is 
a read transition or an \isa{{\isasymepsilon}}-transition in the inductive step.
\end{proof} 
\end{lemma}

And now we begin our circular proof:

\begin{lemma}[Step 1]\label{char_imp_derivers}
If there exists a computation \isa{{\isacharparenleft}{\kern0pt}{\isacharbrackleft}{\kern0pt}\isaconst{S{\isacharprime}{\kern0pt}}\ {\isasymrightarrow}\ \ {\isasymcdot}\ {\isacharbrackleft}{\kern0pt}\isaconst{Nt}\ \isaconst{S}{\isacharbrackright}{\kern0pt}\ {\isacharbrackright}{\kern0pt}{\isacharcomma}{\kern0pt}\ \isafree{{\isasymgamma}}{\isacharparenright}{\kern0pt}\ {\isasymturnstile}c{\isacharasterisk}{\kern0pt}\ {\isacharparenleft}{\kern0pt}\mbox{{\isacharbrackleft}{\kern0pt}\isafree{A}\ {\isasymrightarrow}\ \isafree{{\isasymalpha}}\ {\isasymcdot}\ \isafree{{\isasymbeta}}{\isacharbrackright}{\kern0pt}}{\isacharcomma}{\kern0pt}\ {\isacharbrackleft}{\kern0pt}{\isacharbrackright}{\kern0pt}{\isacharparenright}{\kern0pt}} of 
\isa{\ensuremath{\mathrm{char}(G)}}, then there exist \isa{{\isasymgamma}{\isacharprime}{\kern0pt}\ {\isacharcolon}{\kern0pt}{\isacharcolon}{\kern0pt}\ syms} and \isa{w\ {\isacharcolon}{\kern0pt}{\isacharcolon}{\kern0pt}\ {\isacharprime}{\kern0pt}t\ list} for which there is a rightmost
derivation of the form \isa{P{\isacharprime}{\kern0pt}\ {\isasymturnstile}\ {\isacharbrackleft}{\kern0pt}Nt\ S{\isacharprime}{\kern0pt}{\isacharbrackright}{\kern0pt}\ {\isasymRightarrow}r{\isacharasterisk}{\kern0pt}\ {\isasymgamma}{\isacharprime}{\kern0pt}\ {\isacharat}{\kern0pt}\ Nt\ A\ {\isacharhash}{\kern0pt}\ map\ Tm\ w\ {\isasymRightarrow}\ {\isasymgamma}{\isacharprime}{\kern0pt}\ {\isacharat}{\kern0pt}\ {\isasymalpha}\ {\isacharat}{\kern0pt}\ {\isasymbeta}\ {\isacharat}{\kern0pt}\ map\ Tm\ w} with
\isa{\isafree{{\isasymgamma}}\ {\isacharequal}{\kern0pt}\ \isafree{{\isasymgamma}{\isacharprime}{\kern0pt}}\ {\isacharat}{\kern0pt}\ \isafree{{\isasymalpha}}}.
\begin{proof}
We proceed by induction over the length \isa{n} of the computation for arbitrary \isa{{\isasymgamma}{\isacharcomma}{\kern0pt}\ {\isasymalpha}{\isacharcomma}{\kern0pt}\ {\isasymbeta}{\isacharcomma}{\kern0pt}} and \isa{A}. 

If \isa{n\ {\isacharequal}{\kern0pt}\ {\isadigit{0}}} the implication holds trivially by reflexivity.

On the other hand, if \isa{n\ {\isacharequal}{\kern0pt}\ Suc\ m} for some \isa{m}, we perform a case distinction on the last step of
the computation.

If the last step is a read transition, it is of the form 
\isa{{\isacharparenleft}{\kern0pt}{\isacharbrackleft}{\kern0pt}\isafree{A}\ {\isasymrightarrow}\ \isafree{{\isasymalpha}{\isacharprime}{\kern0pt}}\ {\isasymcdot}\ \isafree{Y}\ {\isacharhash}{\kern0pt}\ \isafree{{\isasymbeta}}{\isacharbrackright}{\kern0pt}{\isacharcomma}{\kern0pt}\ {\isacharbrackleft}{\kern0pt}\isafree{Y}{\isacharbrackright}{\kern0pt}{\isacharparenright}{\kern0pt}\ {\isasymturnstile}c\ {\isacharparenleft}{\kern0pt}\mbox{{\isacharbrackleft}{\kern0pt}\isafree{A}\ {\isasymrightarrow}\ \isafree{{\isasymalpha}}\ {\isasymcdot}\ \isafree{{\isasymbeta}}{\isacharbrackright}{\kern0pt}}{\isacharcomma}{\kern0pt}\ {\isacharbrackleft}{\kern0pt}{\isacharbrackright}{\kern0pt}{\isacharparenright}{\kern0pt}}
for \isa{{\isasymalpha}{\isacharprime}{\kern0pt}} with \isa{\isafree{{\isasymalpha}}\ {\isacharequal}{\kern0pt}\ \isafree{{\isasymalpha}{\isacharprime}{\kern0pt}}\ {\isacharat}{\kern0pt}\ {\isacharbrackleft}{\kern0pt}\isafree{Y}{\isacharbrackright}{\kern0pt}}. With Lemma~\ref{char_reachable_imp_substring}, this implies 
that \isa{\isafree{{\isasymgamma}}\ {\isacharequal}{\kern0pt}\ \isafree{{\isasymdelta}}\ {\isacharat}{\kern0pt}\ \isafree{{\isasymalpha}{\isacharprime}{\kern0pt}}\ {\isacharat}{\kern0pt}\ {\isacharbrackleft}{\kern0pt}\isafree{Y}{\isacharbrackright}{\kern0pt}} for some \isa{{\isasymdelta}}, and therefore
\isa{{\isacharparenleft}{\kern0pt}{\isacharbrackleft}{\kern0pt}\isaconst{S{\isacharprime}{\kern0pt}}\ {\isasymrightarrow}\ \ {\isasymcdot}\ {\isacharbrackleft}{\kern0pt}\isaconst{Nt}\ \isaconst{S}{\isacharbrackright}{\kern0pt}\ {\isacharbrackright}{\kern0pt}{\isacharcomma}{\kern0pt}\ \isafree{{\isasymdelta}}\ {\isacharat}{\kern0pt}\ \isafree{{\isasymalpha}{\isacharprime}{\kern0pt}}{\isacharparenright}{\kern0pt}}\ \isa{{\isasymturnstile}c{\isacharparenleft}{\kern0pt}m{\isacharparenright}{\kern0pt}}\ \isa{{\isacharparenleft}{\kern0pt}{\isacharbrackleft}{\kern0pt}\isafree{A}\ {\isasymrightarrow}\ \isafree{{\isasymalpha}{\isacharprime}{\kern0pt}}\ {\isasymcdot}\ \isafree{Y}\ {\isacharhash}{\kern0pt}\ \isafree{{\isasymbeta}}{\isacharbrackright}{\kern0pt}{\isacharcomma}{\kern0pt}\ {\isacharbrackleft}{\kern0pt}{\isacharbrackright}{\kern0pt}{\isacharparenright}{\kern0pt}}.
By the induction hypothesis, this implies
\[ \isa{P{\isacharprime}{\kern0pt}\ {\isasymturnstile}\ {\isacharbrackleft}{\kern0pt}Nt\ S{\isacharprime}{\kern0pt}{\isacharbrackright}{\kern0pt}\ {\isasymRightarrow}r{\isacharasterisk}{\kern0pt}\ {\isasymgamma}{\isacharprime}{\kern0pt}\ {\isacharat}{\kern0pt}\ Nt\ A\ {\isacharhash}{\kern0pt}\ map\ Tm\ w\ {\isasymRightarrow}\ {\isasymdelta}\ {\isacharat}{\kern0pt}\ {\isasymalpha}{\isacharprime}{\kern0pt}\ {\isacharat}{\kern0pt}\ Y\ {\isacharhash}{\kern0pt}\ {\isasymbeta}\ {\isacharat}{\kern0pt}\ map\ Tm\ w}, \]
and the implication holds with \isa{\isafree{{\isasymalpha}}\ {\isacharequal}{\kern0pt}\ \isafree{{\isasymalpha}{\isacharprime}{\kern0pt}}\ {\isacharat}{\kern0pt}\ {\isacharbrackleft}{\kern0pt}\isafree{Y}{\isacharbrackright}{\kern0pt}}.

If the last step is an \isa{{\isasymepsilon}}-transition, it is of the form
\isa{{\isacharparenleft}{\kern0pt}{\isacharbrackleft}{\kern0pt}\isafree{X}\ {\isasymrightarrow}\ \isafree{{\isasymalpha}{\isacharprime}{\kern0pt}}\ {\isasymcdot}\ \isaconst{Nt}\ \isafree{A}\ {\isacharhash}{\kern0pt}\ \isafree{{\isasymbeta}{\isacharprime}{\kern0pt}}{\isacharbrackright}{\kern0pt}{\isacharcomma}{\kern0pt}\ {\isacharbrackleft}{\kern0pt}{\isacharbrackright}{\kern0pt}{\isacharparenright}{\kern0pt}\ {\isasymturnstile}c\ \mbox{{\isacharparenleft}{\kern0pt}{\isacharbrackleft}{\kern0pt}\isafree{A}\ {\isasymrightarrow}\ \isafree{{\isasymalpha}}\ {\isasymcdot}\ \isafree{{\isasymbeta}}{\isacharbrackright}{\kern0pt}{\isacharcomma}{\kern0pt}\ {\isacharbrackleft}{\kern0pt}{\isacharbrackright}{\kern0pt}{\isacharparenright}{\kern0pt}}},
which implies that \isa{\isafree{{\isasymalpha}}\ {\isacharequal}{\kern0pt}\ {\isacharbrackleft}{\kern0pt}{\isacharbrackright}{\kern0pt}}. By the IH we then have
\begin{equation}\label{char_der.ih}
\isa{P{\isacharprime}{\kern0pt}\ {\isasymturnstile}\ {\isacharbrackleft}{\kern0pt}Nt\ S{\isacharprime}{\kern0pt}{\isacharbrackright}{\kern0pt}\ {\isasymRightarrow}r{\isacharasterisk}{\kern0pt}\ {\isasymgamma}{\isacharprime}{\kern0pt}\ {\isacharat}{\kern0pt}\ Nt\ X\ {\isacharhash}{\kern0pt}\ map\ Tm\ w\ {\isasymRightarrow}r\ {\isasymgamma}{\isacharprime}{\kern0pt}\ {\isacharat}{\kern0pt}\ {\isasymalpha}{\isacharprime}{\kern0pt}\ {\isacharat}{\kern0pt}\ Nt\ A\ {\isacharhash}{\kern0pt}\ {\isasymbeta}{\isacharprime}{\kern0pt}\ {\isacharat}{\kern0pt}\ map\ Tm\ w}
\end{equation}
with \isa{\isafree{{\isasymgamma}}\ {\isacharequal}{\kern0pt}\ \isafree{{\isasymgamma}{\isacharprime}{\kern0pt}}\ {\isacharat}{\kern0pt}\ \isafree{{\isasymalpha}{\isacharprime}{\kern0pt}}}. Moreover, since \isa{G{\isacharprime}{\kern0pt}} is reduced, there exists a \isa{v\ {\isacharcolon}{\kern0pt}{\isacharcolon}{\kern0pt}\ {\isacharprime}{\kern0pt}t\ list} such that
\isa{\isaconst{P{\isacharprime}{\kern0pt}}\ {\isasymturnstile}\ \isafree{{\isasymbeta}{\isacharprime}{\kern0pt}}\ {\isasymRightarrow}r{\isacharasterisk}{\kern0pt}\ \isaconst{map}\ \isaconst{Tm}\ \isafree{v}}.
From this and \eqref{char_der.ih} we have
\begin{equation}
\isa{P{\isacharprime}{\kern0pt}\ {\isasymturnstile}\ {\isacharbrackleft}{\kern0pt}Nt\ S{\isacharprime}{\kern0pt}{\isacharbrackright}{\kern0pt}\ {\isasymRightarrow}r{\isacharasterisk}{\kern0pt}\ {\isasymgamma}\ {\isacharat}{\kern0pt}\ Nt\ A\ {\isacharhash}{\kern0pt}\ map\ Tm\ {\isacharparenleft}{\kern0pt}v{\isacharat}{\kern0pt}w{\isacharparenright}{\kern0pt}\ {\isasymRightarrow}r\ {\isasymgamma}\ {\isacharat}{\kern0pt}\ {\isasymalpha}\ {\isacharat}{\kern0pt}\ {\isasymbeta}\ {\isacharat}{\kern0pt}\ map\ Tm\ {\isacharparenleft}{\kern0pt}v{\isacharat}{\kern0pt}w{\isacharparenright}{\kern0pt}}.
\end{equation}
This completes the proof by the fact that \isa{\isafree{{\isasymalpha}}\ {\isacharequal}{\kern0pt}\ {\isacharbrackleft}{\kern0pt}{\isacharbrackright}{\kern0pt}}.
\end{proof}
\end{lemma}

We now move on towards proving the second step, for which we first show an auxiliary lemma.

\begin{lemma}\label{deriver_imp_IPDA_comp}
If \isa{\isaconst{P{\isacharprime}{\kern0pt}}\ {\isasymturnstile}\ {\isacharbrackleft}{\kern0pt}\isaconst{Nt}\ \isaconst{S{\isacharprime}{\kern0pt}}{\isacharbrackright}{\kern0pt}\ {\isasymRightarrow}r\ \isafree{{\isasymalpha}}\ {\isacharat}{\kern0pt}\ \isafree{{\isasymbeta}}} and \isa{\isaconst{P{\isacharprime}{\kern0pt}}\ {\isasymturnstile}\ \isafree{{\isasymbeta}}\ {\isasymRightarrow}{\isacharasterisk}{\kern0pt}\ \isaconst{map}\ \isaconst{Tm}\ \isafree{v}}, there exists a 
computation \isa{\mbox{{\isacharparenleft}{\kern0pt}{\isacharbrackleft}{\kern0pt}{\isacharbrackleft}{\kern0pt}\isaconst{S{\isacharprime}{\kern0pt}}\ {\isasymrightarrow}\ \isafree{{\isasymalpha}}\ {\isasymcdot}\ \isafree{{\isasymbeta}}{\isacharbrackright}{\kern0pt}{\isacharbrackright}{\kern0pt}{\isacharcomma}{\kern0pt}\ \isafree{v}{\isacharparenright}{\kern0pt}\ {\isasymturnstile}I{\isacharasterisk}{\kern0pt}\ {\isacharparenleft}{\kern0pt}{\isacharbrackleft}{\kern0pt}{\isacharbrackleft}{\kern0pt}\isaconst{S{\isacharprime}{\kern0pt}}\ {\isasymrightarrow}\ {\isacharbrackleft}{\kern0pt}\isaconst{Nt}\ \isaconst{S}{\isacharbrackright}{\kern0pt}\ {\isasymcdot}\ {\isacharbrackright}{\kern0pt}{\isacharbrackright}{\kern0pt}{\isacharcomma}{\kern0pt}\ {\isacharbrackleft}{\kern0pt}{\isacharbrackright}{\kern0pt}{\isacharparenright}{\kern0pt}}}.
\begin{proof}
Since the only possible production for \isa{S{\isacharprime}{\kern0pt}} is \isa{{\isacharparenleft}{\kern0pt}S{\isacharprime}{\kern0pt}{\isacharcomma}{\kern0pt}{\isacharbrackleft}{\kern0pt}Nt\ S{\isacharbrackright}{\kern0pt}{\isacharparenright}{\kern0pt}}, we have \isa{{\isacharbrackleft}{\kern0pt}\isaconst{Nt}\ \isaconst{S}{\isacharbrackright}{\kern0pt}\ {\isacharequal}{\kern0pt}\ \isafree{{\isasymalpha}}\ {\isacharat}{\kern0pt}\ \isafree{{\isasymbeta}}}.
The proof is then trivial by distinguishing the cases where \isa{{\isasymalpha}\ {\isacharequal}{\kern0pt}\ {\isacharbrackleft}{\kern0pt}Nt\ S{\isacharbrackright}{\kern0pt}} and \isa{{\isasymbeta}\ {\isacharequal}{\kern0pt}\ {\isacharbrackleft}{\kern0pt}{\isacharbrackright}{\kern0pt}}, and vice versa,
using Theorem~\ref{ipda.Lang_eq_Lang_G}.
\end{proof}
\end{lemma}

\begin{lemma}[Step 2]\label{derivers_imp_ipda}
If there exists a rightmost derivation 
\begin{equation*}
\isa{P{\isacharprime}{\kern0pt}\ {\isasymturnstile}\ {\isacharbrackleft}{\kern0pt}Nt\ S{\isacharprime}{\kern0pt}{\isacharbrackright}{\kern0pt}\ {\isasymRightarrow}r{\isacharasterisk}{\kern0pt}\ {\isasymgamma}\ {\isacharat}{\kern0pt}\ Nt\ A\ {\isacharhash}{\kern0pt}\ map\ Tm\ w\ {\isasymRightarrow}r\ {\isasymgamma}\ {\isacharat}{\kern0pt}\ {\isasymalpha}\ {\isacharat}{\kern0pt}\ {\isasymbeta}\ {\isacharat}{\kern0pt}\ map\ Tm\ w}
\end{equation*}
with \isa{\isaconst{P{\isacharprime}{\kern0pt}}\ {\isasymturnstile}\ \isafree{{\isasymbeta}}\ {\isasymRightarrow}{\isacharasterisk}{\kern0pt}\ \isaconst{map}\ \isaconst{Tm}\ \isafree{v}},
then there exists an \isa{\isaconst{IPDA}} computation
\isa{\mbox{{\isacharparenleft}{\kern0pt}{\isacharbrackleft}{\kern0pt}\isafree{A}\ {\isasymrightarrow}\ \isafree{{\isasymalpha}}\ {\isasymcdot}\ \isafree{{\isasymbeta}}{\isacharbrackright}{\kern0pt}\ {\isacharhash}{\kern0pt}\ \isafree{{\isasymrho}}{\isacharcomma}{\kern0pt}\ \isafree{v}\ {\isacharat}{\kern0pt}\ \isafree{w}{\isacharparenright}{\kern0pt}}\ {\isasymturnstile}I{\isacharasterisk}{\kern0pt}\ {\isacharparenleft}{\kern0pt}{\isacharbrackleft}{\kern0pt}{\isacharbrackleft}{\kern0pt}\isaconst{S{\isacharprime}{\kern0pt}}\ {\isasymrightarrow}\ {\isacharbrackleft}{\kern0pt}\isaconst{Nt}\ \isaconst{S}{\isacharbrackright}{\kern0pt}\ {\isasymcdot}\ {\isacharbrackright}{\kern0pt}{\isacharbrackright}{\kern0pt}{\isacharcomma}{\kern0pt}\ {\isacharbrackleft}{\kern0pt}{\isacharbrackright}{\kern0pt}{\isacharparenright}{\kern0pt}} with \isa{\isaconst{hist}\ {\isacharparenleft}{\kern0pt}\isaconst{rev}\ \isafree{{\isasymrho}}{\isacharparenright}{\kern0pt}\ {\isacharequal}{\kern0pt}\ \isafree{{\isasymgamma}}}.
\begin{proof}
We begin by distinguishing the cases where \isa{\isafree{n}\ {\isacharequal}{\kern0pt}\ {\isadigit{0}}} and \isa{{\isadigit{0}}\ {\isacharless}{\kern0pt}\ \isafree{n}}.

If \isa{\isafree{n}\ {\isacharequal}{\kern0pt}\ {\isadigit{0}}}, we have \isa{\isaconst{S{\isacharprime}{\kern0pt}}\ {\isacharequal}{\kern0pt}\ \isafree{A}}, \isa{{\isasymgamma}\ {\isacharequal}{\kern0pt}\ w\ {\isacharequal}{\kern0pt}\ {\isacharbrackleft}{\kern0pt}{\isacharbrackright}{\kern0pt}}, and \isa{\isafree{{\isasymalpha}}\ {\isacharat}{\kern0pt}\ \isafree{{\isasymbeta}}\ {\isacharequal}{\kern0pt}\ {\isacharbrackleft}{\kern0pt}\isaconst{Nt}\ \isaconst{S}{\isacharbrackright}{\kern0pt}}. The 
computation then exists by Lemma~\ref{deriver_imp_IPDA_comp}.

If \isa{{\isadigit{0}}\ {\isacharless}{\kern0pt}\ \isafree{n}}, Lemma~\ref{derivern_Suc_singleton_imp_rm_chain} implies the existence of some
\isa{{\isasymrho}} with \isa{\isaconst{P{\isacharprime}{\kern0pt}}\ {\isasymturnstile}\ {\isacharbrackleft}{\kern0pt}\isaconst{Nt}\ \isaconst{S{\isacharprime}{\kern0pt}}{\isacharbrackright}{\kern0pt}\ {\isasymmidarrow}\isafree{{\isasymrho}}{\isasymrightarrow}r{\isacharasterisk}{\kern0pt}\ \isafree{{\isasymgamma}}\ {\isacharat}{\kern0pt}\ \isaconst{Nt}\ \isafree{A}\ {\isacharhash}{\kern0pt}\ \isaconst{map}\ \isaconst{Tm}\ \isafree{w}}.
We can now induct on \isa{{\isasymrho}} for arbitrary \isa{{\isasymgamma}{\isacharcomma}{\kern0pt}\ {\isasymalpha}{\isacharcomma}{\kern0pt}\ {\isasymbeta}{\isacharcomma}{\kern0pt}\ A{\isacharcomma}{\kern0pt}\ w{\isacharcomma}{\kern0pt}} and \isa{v}, akin to the proof by 
Wilhelm et al.

If \isa{\isafree{{\isasymrho}}\ {\isacharequal}{\kern0pt}\ {\isacharbrackleft}{\kern0pt}{\isacharbrackright}{\kern0pt}}, we again have \isa{\isaconst{S{\isacharprime}{\kern0pt}}\ {\isacharequal}{\kern0pt}\ \isafree{A}}, \isa{{\isasymgamma}\ {\isacharequal}{\kern0pt}\ w\ {\isacharequal}{\kern0pt}\ {\isacharbrackleft}{\kern0pt}{\isacharbrackright}{\kern0pt}}, and \isa{\isafree{{\isasymalpha}}\ {\isacharat}{\kern0pt}\ \isafree{{\isasymbeta}}\ {\isacharequal}{\kern0pt}\ {\isacharbrackleft}{\kern0pt}\isaconst{Nt}\ \isaconst{S}{\isacharbrackright}{\kern0pt}}, and 
the proof is analogous to the case of \isa{\isafree{n}\ {\isacharequal}{\kern0pt}\ {\isadigit{0}}}.

If \isa{\isafree{{\isasymrho}}\ {\isacharequal}{\kern0pt}\ \isafree{i}\ {\isacharhash}{\kern0pt}\ \isafree{{\isasymsigma}}} for some \isa{i} and \isa{{\isasymsigma}}, by rule inversion we know that 
\begin{equation}\label{der_ipda.i}
\isa{\isafree{i}\ {\isacharequal}{\kern0pt}\ {\isacharbrackleft}{\kern0pt}\isafree{X}\ {\isasymrightarrow}\ \isafree{{\isasymalpha}{\isacharprime}{\kern0pt}}\ {\isasymcdot}\ \isaconst{Nt}\ \isafree{A}\ {\isacharhash}{\kern0pt}\ \isafree{{\isasymbeta}{\isacharprime}{\kern0pt}}{\isacharbrackright}{\kern0pt}}
\end{equation}
for some \isa{X}, \isa{{\isasymalpha}{\isacharprime}{\kern0pt}}, \isa{A} and \isa{{\isasymbeta}{\isacharprime}{\kern0pt}}. Furthermore, the chain is such that
\begin{gather*}
\isa{\isaconst{P{\isacharprime}{\kern0pt}}\ {\isasymturnstile}\ {\isacharbrackleft}{\kern0pt}\isaconst{Nt}\ \isaconst{S{\isacharprime}{\kern0pt}}{\isacharbrackright}{\kern0pt}\ {\isasymmidarrow}\isafree{{\isasymsigma}}{\isasymrightarrow}r{\isacharasterisk}{\kern0pt}\ \isafree{{\isasymalpha}{\isacharprime}{\kern0pt}{\isacharprime}{\kern0pt}}\ {\isacharat}{\kern0pt}\ \isaconst{Nt}\ \isafree{X}\ {\isacharhash}{\kern0pt}\ \isaconst{map}\ \isaconst{Tm}\ \isafree{v{\isacharprime}{\kern0pt}}}\\
\intertext{and}
\isa{\isaconst{P{\isacharprime}{\kern0pt}}\ {\isasymturnstile}\ \isafree{{\isasymalpha}{\isacharprime}{\kern0pt}{\isacharprime}{\kern0pt}}\ {\isacharat}{\kern0pt}\ \isaconst{Nt}\ \isafree{X}\ {\isacharhash}{\kern0pt}\ \isaconst{map}\ \isaconst{Tm}\ \isafree{v{\isacharprime}{\kern0pt}}\ {\isasymRightarrow}r\ \isafree{{\isasymalpha}{\isacharprime}{\kern0pt}{\isacharprime}{\kern0pt}}\ {\isacharat}{\kern0pt}\ \isafree{{\isasymalpha}{\isacharprime}{\kern0pt}}\ {\isacharat}{\kern0pt}\ \isaconst{Nt}\ \isafree{A}\ {\isacharhash}{\kern0pt}\ \isafree{{\isasymbeta}{\isacharprime}{\kern0pt}}\ {\isacharat}{\kern0pt}\ \isaconst{map}\ \isaconst{Tm}\ \isafree{v{\isacharprime}{\kern0pt}}}
\end{gather*}
with \isa{\isaconst{P{\isacharprime}{\kern0pt}}\ {\isasymturnstile}\ \isafree{{\isasymbeta}{\isacharprime}{\kern0pt}}\ {\isasymRightarrow}r{\isacharasterisk}{\kern0pt}\ \isaconst{map}\ \isaconst{Tm}\ \isafree{u}}, \isa{\isafree{u}\ {\isacharat}{\kern0pt}\ \isafree{v{\isacharprime}{\kern0pt}}\ {\isacharequal}{\kern0pt}\ \isafree{w}}, and \isa{\isafree{{\isasymalpha}{\isacharprime}{\kern0pt}{\isacharprime}{\kern0pt}}\ {\isacharat}{\kern0pt}\ \isafree{{\isasymalpha}{\isacharprime}{\kern0pt}}\ {\isacharequal}{\kern0pt}\ \isafree{{\isasymgamma}}}
for some \isa{{\isasymalpha}{\isacharprime}{\kern0pt}{\isacharprime}{\kern0pt}{\isacharcomma}{\kern0pt}\ u{\isacharcomma}{\kern0pt}} and \isa{v{\isacharprime}{\kern0pt}}. By the IH, all this implies the existence of some \isa{{\isasymtau}} where
\begin{equation*}
\isa{{\isacharparenleft}{\kern0pt}{\isacharbrackleft}{\kern0pt}\isafree{X}\ {\isasymrightarrow}\ \isafree{{\isasymalpha}{\isacharprime}{\kern0pt}}\ {\isacharat}{\kern0pt}\ {\isacharbrackleft}{\kern0pt}\isaconst{Nt}\ \isafree{A}{\isacharbrackright}{\kern0pt}\ {\isasymcdot}\ \isafree{{\isasymbeta}{\isacharprime}{\kern0pt}}{\isacharbrackright}{\kern0pt}\ {\isacharhash}{\kern0pt}\ \isafree{{\isasymtau}}{\isacharcomma}{\kern0pt}\ \isafree{u}\ {\isacharat}{\kern0pt}\ \isafree{v{\isacharprime}{\kern0pt}}{\isacharparenright}{\kern0pt}\ {\isasymturnstile}I{\isacharasterisk}{\kern0pt}\ {\isacharparenleft}{\kern0pt}{\isacharbrackleft}{\kern0pt}{\isacharbrackleft}{\kern0pt}\isaconst{S{\isacharprime}{\kern0pt}}\ {\isasymrightarrow}\ {\isacharbrackleft}{\kern0pt}\isaconst{Nt}\ \isaconst{S}{\isacharbrackright}{\kern0pt}\ {\isasymcdot}\ {\isacharbrackright}{\kern0pt}{\isacharbrackright}{\kern0pt}{\isacharcomma}{\kern0pt}\ {\isacharbrackleft}{\kern0pt}{\isacharbrackright}{\kern0pt}{\isacharparenright}{\kern0pt}}\\
\end{equation*}
and \isa{\isaconst{hist}\ {\isacharparenleft}{\kern0pt}\isaconst{rev}\ \isafree{{\isasymtau}}{\isacharparenright}{\kern0pt}\ {\isacharequal}{\kern0pt}\ \isafree{{\isasymalpha}{\isacharprime}{\kern0pt}{\isacharprime}{\kern0pt}}}.

Lastly, we can show with Lemma~\ref{derives_imp_completes} and our assumption that \isa{{\isasymbeta}} derives 
\isa{map\ Tm\ v} that the IPDA with configuration \isa{{\isacharparenleft}{\kern0pt}{\isacharbrackleft}{\kern0pt}\isafree{A}\ {\isasymrightarrow}\ \isafree{{\isasymalpha}}\ {\isasymcdot}\ \isafree{{\isasymbeta}}{\isacharbrackright}{\kern0pt}\ {\isacharhash}{\kern0pt}\ {\isacharbrackleft}{\kern0pt}\isafree{X}\ {\isasymrightarrow}\ \isafree{{\isasymalpha}{\isacharprime}{\kern0pt}}\ {\isasymcdot}\ \isaconst{Nt}\ \isafree{A}\ {\isacharhash}{\kern0pt}\ \isafree{{\isasymbeta}{\isacharprime}{\kern0pt}}{\isacharbrackright}{\kern0pt}\ {\isacharhash}{\kern0pt}\ \isafree{{\isasymtau}}{\isacharcomma}{\kern0pt}\ \isafree{v}\ {\isacharat}{\kern0pt}\ \isafree{w}{\isacharparenright}{\kern0pt}}
reaches \isa{\mbox{{\isacharparenleft}{\kern0pt}{\isacharbrackleft}{\kern0pt}\isafree{X}\ {\isasymrightarrow}\ \isafree{{\isasymalpha}{\isacharprime}{\kern0pt}}\ {\isacharat}{\kern0pt}\ {\isacharbrackleft}{\kern0pt}\isaconst{Nt}\ \isafree{A}{\isacharbrackright}{\kern0pt}\ {\isasymcdot}\ \isafree{{\isasymbeta}{\isacharprime}{\kern0pt}}{\isacharbrackright}{\kern0pt}\ {\isacharhash}{\kern0pt}\ \isafree{{\isasymtau}}{\isacharcomma}{\kern0pt}\ \isafree{u}\ {\isacharat}{\kern0pt}\ \isafree{v{\isacharprime}{\kern0pt}}{\isacharparenright}{\kern0pt}}}, completing the proof.
\end{proof}
\end{lemma}

It is worth noting that, as opposed to Wilhelm et al., we only assume 
\isa{\isaconst{P{\isacharprime}{\kern0pt}}\ {\isasymturnstile}\ \isafree{{\isasymbeta}}\ {\isasymRightarrow}{\isacharasterisk}{\kern0pt}\ \isaconst{map}\ \isaconst{Tm}\ \isafree{v}} for the above lemma, since specifying this as a rightmost 
derivation makes no difference. We will now prove further properties about \isa{\ensuremath{\mathrm{char}(G)}} which 
we will need for the third step.

\begin{lemma}\label{reaches_final_imp_completes}
If \isa{{\isacharparenleft}{\kern0pt}{\isacharbrackleft}{\kern0pt}\isafree{A}\ {\isasymrightarrow}\ \isafree{{\isasymalpha}}\ {\isasymcdot}\ \isafree{{\isasymbeta}}{\isacharbrackright}{\kern0pt}\ {\isacharhash}{\kern0pt}\ \isafree{{\isasymrho}}{\isacharcomma}{\kern0pt}\ \isafree{u}{\isacharparenright}{\kern0pt}}\ \isa{{\isasymturnstile}I{\isacharparenleft}{\kern0pt}n{\isacharparenright}{\kern0pt}}\ \isa{{\isacharparenleft}{\kern0pt}{\isacharbrackleft}{\kern0pt}{\isacharbrackleft}{\kern0pt}\isaconst{S{\isacharprime}{\kern0pt}}\ {\isasymrightarrow}\ {\isacharbrackleft}{\kern0pt}\isaconst{Nt}\ \isaconst{S}{\isacharbrackright}{\kern0pt}\ {\isasymcdot}\ {\isacharbrackright}{\kern0pt}{\isacharbrackright}{\kern0pt}{\isacharcomma}{\kern0pt}\ {\isacharbrackleft}{\kern0pt}{\isacharbrackright}{\kern0pt}{\isacharparenright}{\kern0pt}} holds, 
there exist a \isa{v\ {\isacharcolon}{\kern0pt}{\isacharcolon}{\kern0pt}\ {\isacharprime}{\kern0pt}t\ list} and \isa{i{\isacharcomma}{\kern0pt}\ j\ {\isacharcolon}{\kern0pt}{\isacharcolon}{\kern0pt}\ nat} such that
\begin{equation*} 
\isa{{\isacharparenleft}{\kern0pt}{\isacharbrackleft}{\kern0pt}A\ {\isasymrightarrow}\ {\isasymalpha}\ {\isasymcdot}\ {\isasymbeta}{\isacharbrackright}{\kern0pt}\ {\isacharhash}{\kern0pt}\ {\isasymrho}{\isacharcomma}{\kern0pt}\ u{\isacharparenright}{\kern0pt}\ {\isasymturnstile}I{\isacharparenleft}{\kern0pt}i{\isacharparenright}{\kern0pt}\ {\isacharparenleft}{\kern0pt}{\isacharbrackleft}{\kern0pt}A\ {\isasymrightarrow}\ {\isasymalpha}\ {\isacharat}{\kern0pt}\ {\isasymbeta}\ {\isasymcdot}\ {\isacharbrackleft}{\kern0pt}{\isacharbrackright}{\kern0pt}{\isacharbrackright}{\kern0pt}\ {\isacharhash}{\kern0pt}\ {\isasymrho}{\isacharcomma}{\kern0pt}\ v{\isacharparenright}{\kern0pt}\ {\isasymturnstile}I{\isacharparenleft}{\kern0pt}j{\isacharparenright}{\kern0pt}\ {\isacharparenleft}{\kern0pt}{\isacharbrackleft}{\kern0pt}}\isa{{\isacharbrackleft}{\kern0pt}\isaconst{S{\isacharprime}{\kern0pt}}\ {\isasymrightarrow}\ {\isacharbrackleft}{\kern0pt}\isaconst{Nt}\ \isaconst{S}{\isacharbrackright}{\kern0pt}\ {\isasymcdot}\ {\isacharbrackright}{\kern0pt}}\isa{{\isacharbrackright}{\kern0pt}{\isacharcomma}{\kern0pt}\ {\isacharbrackleft}{\kern0pt}{\isacharbrackright}{\kern0pt}{\isacharparenright}{\kern0pt}}
\end{equation*}
and \isa{\isafree{i}\ {\isacharplus}{\kern0pt}\ \isafree{j}\ {\isacharequal}{\kern0pt}\ \isafree{n}}.
\begin{proof}
We begin our proof by strong induction on \isa{n} for arbitrary \isa{A{\isacharcomma}{\kern0pt}\ u{\isacharcomma}{\kern0pt}\ {\isasymalpha}{\isacharcomma}{\kern0pt}\ {\isasymbeta}{\isacharcomma}{\kern0pt}} and \isa{{\isasymrho}}.

If \isa{n\ {\isacharequal}{\kern0pt}\ {\isadigit{0}}}, the implication is trivial.

Otherwise, if \isa{n\ {\isacharequal}{\kern0pt}\ Suc\ m} for some \isa{m}, we proceed with a case distinction on the first step of the 
computation. 

If the first step is a shifting or a reducing transition, the implication holds by the induction
hypothesis.

If the first step is an expanding transition, there exist \isa{Y\ {\isacharcolon}{\kern0pt}{\isacharcolon}{\kern0pt}\ {\isacharprime}{\kern0pt}n} and \isa{{\isasymgamma}\ {\isacharcolon}{\kern0pt}{\isacharcolon}{\kern0pt}\ syms} with
\isa{{\isasymbeta}\ {\isacharequal}{\kern0pt}\ Nt\ Y\ {\isacharhash}{\kern0pt}\ {\isasymgamma}} and
\begin{multline*}
\isa{{\isacharparenleft}{\kern0pt}{\isacharbrackleft}{\kern0pt}A\ {\isasymrightarrow}\ {\isasymalpha}\ {\isasymcdot}\ {\isasymbeta}{\isacharbrackright}{\kern0pt}\ {\isacharhash}{\kern0pt}\ {\isasymrho}{\isacharcomma}{\kern0pt}\ u{\isacharparenright}{\kern0pt}\ {\isasymturnstile}I\ {\isacharparenleft}{\kern0pt}{\isacharbrackleft}{\kern0pt}Y\ {\isasymrightarrow}\ {\isasymcdot}\ {\isasymgamma}{\isacharbrackright}{\kern0pt}\ {\isacharhash}{\kern0pt}\ {\isacharbrackleft}{\kern0pt}A\ {\isasymrightarrow}\ {\isasymalpha}\ {\isasymcdot}\ {\isasymbeta}{\isacharbrackright}{\kern0pt}\ {\isacharhash}{\kern0pt}\ {\isasymrho}{\isacharcomma}{\kern0pt}\ u{\isacharparenright}{\kern0pt}}\\
  \isa{{\isasymturnstile}I{\isacharparenleft}{\kern0pt}m{\isacharparenright}{\kern0pt}\ {\isacharparenleft}{\kern0pt}{\isacharbrackleft}{\kern0pt}}\isa{{\isacharbrackleft}{\kern0pt}\isaconst{S{\isacharprime}{\kern0pt}}\ {\isasymrightarrow}\ {\isacharbrackleft}{\kern0pt}\isaconst{Nt}\ \isaconst{S}{\isacharbrackright}{\kern0pt}\ {\isasymcdot}\ {\isacharbrackright}{\kern0pt}}\isa{{\isacharbrackright}{\kern0pt}{\isacharcomma}{\kern0pt}\ {\isacharbrackleft}{\kern0pt}{\isacharbrackright}{\kern0pt}{\isacharparenright}{\kern0pt}}.
\end{multline*}

By the IH, this implies the existence of some \isa{v\ {\isacharcolon}{\kern0pt}{\isacharcolon}{\kern0pt}\ {\isacharprime}{\kern0pt}t\ list} and \isa{i{\isacharcomma}{\kern0pt}\ j\ {\isacharcolon}{\kern0pt}{\isacharcolon}{\kern0pt}\ nat} with 
\begin{multline*}
\isa{{\isacharparenleft}{\kern0pt}{\isacharbrackleft}{\kern0pt}Y\ {\isasymrightarrow}\ \ {\isasymcdot}\ {\isasymgamma}{\isacharbrackright}{\kern0pt}\ {\isacharhash}{\kern0pt}\ {\isacharbrackleft}{\kern0pt}A\ {\isasymrightarrow}\ {\isasymalpha}\ {\isasymcdot}\ {\isasymbeta}{\isacharbrackright}{\kern0pt}\ {\isacharhash}{\kern0pt}\ {\isasymrho}{\isacharcomma}{\kern0pt}\ u{\isacharparenright}{\kern0pt}\ {\isasymturnstile}I{\isacharparenleft}{\kern0pt}i{\isacharparenright}{\kern0pt}\ {\isacharparenleft}{\kern0pt}{\isacharbrackleft}{\kern0pt}Y\ {\isasymrightarrow}\ {\isasymgamma}\ {\isasymcdot}\ {\isacharbrackright}{\kern0pt}\ {\isacharhash}{\kern0pt}\ {\isacharbrackleft}{\kern0pt}A\ {\isasymrightarrow}\ {\isasymalpha}\ {\isasymcdot}\ {\isasymbeta}{\isacharbrackright}{\kern0pt}\ {\isacharhash}{\kern0pt}\ {\isasymrho}{\isacharcomma}{\kern0pt}\ v{\isacharparenright}{\kern0pt}}\\
  \isa{{\isasymturnstile}I{\isacharparenleft}{\kern0pt}j{\isacharparenright}{\kern0pt}\ {\isacharparenleft}{\kern0pt}{\isacharbrackleft}{\kern0pt}}\isa{{\isacharbrackleft}{\kern0pt}\isaconst{S{\isacharprime}{\kern0pt}}\ {\isasymrightarrow}\ {\isacharbrackleft}{\kern0pt}\isaconst{Nt}\ \isaconst{S}{\isacharbrackright}{\kern0pt}\ {\isasymcdot}\ {\isacharbrackright}{\kern0pt}}\isa{{\isacharbrackright}{\kern0pt}{\isacharcomma}{\kern0pt}\ {\isacharbrackleft}{\kern0pt}{\isacharbrackright}{\kern0pt}{\isacharparenright}{\kern0pt}}
\end{multline*}
and \isa{i\ {\isacharplus}{\kern0pt}\ j\ {\isacharequal}{\kern0pt}\ m}. With the fact that \isa{{\isasymbeta}\ {\isacharequal}{\kern0pt}\ Nt\ Y\ {\isacharhash}{\kern0pt}\ {\isasymgamma}}, the first step of the computation
\[ \isa{{\isacharparenleft}{\kern0pt}{\isacharbrackleft}{\kern0pt}Y\ {\isasymrightarrow}\ {\isasymgamma}\ {\isasymcdot}\ {\isacharbrackright}{\kern0pt}\ {\isacharhash}{\kern0pt}\ {\isacharbrackleft}{\kern0pt}A\ {\isasymrightarrow}\ {\isasymalpha}\ {\isasymcdot}\ {\isasymbeta}{\isacharbrackright}{\kern0pt}\ {\isacharhash}{\kern0pt}\ {\isasymrho}{\isacharcomma}{\kern0pt}\ v{\isacharparenright}{\kern0pt}\ {\isasymturnstile}I{\isacharparenleft}{\kern0pt}j{\isacharparenright}{\kern0pt}\ {\isacharparenleft}{\kern0pt}{\isacharbrackleft}{\kern0pt}}\isa{{\isacharbrackleft}{\kern0pt}\isaconst{S{\isacharprime}{\kern0pt}}\ {\isasymrightarrow}\ {\isacharbrackleft}{\kern0pt}\isaconst{Nt}\ \isaconst{S}{\isacharbrackright}{\kern0pt}\ {\isasymcdot}\ {\isacharbrackright}{\kern0pt}}\isa{{\isacharbrackright}{\kern0pt}{\isacharcomma}{\kern0pt}\ {\isacharbrackleft}{\kern0pt}{\isacharbrackright}{\kern0pt}{\isacharparenright}{\kern0pt}} \]
is invariably a reducing transition, and since the resulting configuration reaches the accepting 
configuration in \isa{j\ {\isacharminus}{\kern0pt}\ {\isadigit{1}}\ {\isacharless}{\kern0pt}\ n} steps, we can use the IH again to finish the proof.
\end{proof}
\end{lemma}

\begin{lemma}\label{char_steps_consume}
If \isa{{\isacharparenleft}{\kern0pt}\isafree{A}{\isacharcomma}{\kern0pt}\ \isafree{{\isasymalpha}}\ {\isacharat}{\kern0pt}\ \isafree{{\isasymbeta}}\ {\isacharat}{\kern0pt}\ \isafree{{\isasymgamma}}{\isacharparenright}{\kern0pt}\ {\isasymin}\ \isaconst{P{\isacharprime}{\kern0pt}}}, \isa{\ensuremath{\mathrm{char}(G)}} computes
\[ \isa{{\isacharparenleft}{\kern0pt}{\isacharbrackleft}{\kern0pt}\isafree{A}\ {\isasymrightarrow}\ \isafree{{\isasymalpha}}\ {\isasymcdot}\ \isafree{{\isasymbeta}}\ {\isacharat}{\kern0pt}\ \isafree{{\isasymgamma}}{\isacharbrackright}{\kern0pt}{\isacharcomma}{\kern0pt}\ \isafree{{\isasymbeta}}\ {\isacharat}{\kern0pt}\ \isafree{{\isasymdelta}}{\isacharparenright}{\kern0pt}\ {\isasymturnstile}c{\isacharasterisk}{\kern0pt}\ {\isacharparenleft}{\kern0pt}{\isacharbrackleft}{\kern0pt}\isafree{A}\ {\isasymrightarrow}\ \isafree{{\isasymalpha}}\ {\isacharat}{\kern0pt}\ \isafree{{\isasymbeta}}\ {\isasymcdot}\ \isafree{{\isasymgamma}}{\isacharbrackright}{\kern0pt}{\isacharcomma}{\kern0pt}\ \isafree{{\isasymdelta}}{\isacharparenright}{\kern0pt}} \]
for any \isa{{\isasymdelta}}.
\begin{proof}
Trivial by induction on \isa{{\isasymbeta}} for arbitrary \isa{{\isasymalpha}}.
\end{proof}
\end{lemma}

And finally, we can prove the final step.

\begin{lemma}[Step 3]\label{ipda_imp_char}
If there exists a computation of \isa{I\isactrlsub G}
\[ \isa{{\isacharparenleft}{\kern0pt}{\isacharbrackleft}{\kern0pt}\isafree{A}\ {\isasymrightarrow}\ \isafree{{\isasymalpha}}\ {\isasymcdot}\ \isafree{{\isasymbeta}}{\isacharbrackright}{\kern0pt}\ {\isacharhash}{\kern0pt}\ \isafree{{\isasymrho}}{\isacharcomma}{\kern0pt}\ \isafree{w}{\isacharparenright}{\kern0pt}\ {\isasymturnstile}I{\isacharasterisk}{\kern0pt}\ {\isacharparenleft}{\kern0pt}{\isacharbrackleft}{\kern0pt}{\isacharbrackleft}{\kern0pt}\isaconst{S{\isacharprime}{\kern0pt}}\ {\isasymrightarrow}\ {\isacharbrackleft}{\kern0pt}\isaconst{Nt}\ \isaconst{S}{\isacharbrackright}{\kern0pt}\ {\isasymcdot}\ {\isacharbrackright}{\kern0pt}{\isacharbrackright}{\kern0pt}{\isacharcomma}{\kern0pt}\ {\isacharbrackleft}{\kern0pt}{\isacharbrackright}{\kern0pt}{\isacharparenright}{\kern0pt}}, \]
then for \isa{\ensuremath{\mathrm{char}(G)}} holds
\[ \isa{{\isacharparenleft}{\kern0pt}{\isacharbrackleft}{\kern0pt}\isaconst{S{\isacharprime}{\kern0pt}}\ {\isasymrightarrow}\ \ {\isasymcdot}\ {\isacharbrackleft}{\kern0pt}\isaconst{Nt}\ \isaconst{S}{\isacharbrackright}{\kern0pt}\ {\isacharbrackright}{\kern0pt}{\isacharcomma}{\kern0pt}\ \isaconst{hist}\ {\isacharparenleft}{\kern0pt}\isaconst{rev}\ \isafree{{\isasymrho}}{\isacharparenright}{\kern0pt}\ {\isacharat}{\kern0pt}\ \isafree{{\isasymalpha}}{\isacharparenright}{\kern0pt}\ {\isasymturnstile}c{\isacharasterisk}{\kern0pt}\ {\isacharparenleft}{\kern0pt}{\isacharbrackleft}{\kern0pt}\isafree{A}\ {\isasymrightarrow}\ \isafree{{\isasymalpha}}\ {\isasymcdot}\ \isafree{{\isasymbeta}}{\isacharbrackright}{\kern0pt}{\isacharcomma}{\kern0pt}\ {\isacharbrackleft}{\kern0pt}{\isacharbrackright}{\kern0pt}{\isacharparenright}{\kern0pt}}. \]
\begin{proof}
We induct on \isa{{\isasymrho}} for arbitrary \isa{A{\isacharcomma}{\kern0pt}\ {\isasymalpha}{\isacharcomma}{\kern0pt}\ {\isasymbeta}} and \isa{w}.

If \isa{{\isasymrho}\ {\isacharequal}{\kern0pt}\ {\isacharbrackleft}{\kern0pt}{\isacharbrackright}{\kern0pt}}, \isa{{\isacharbrackleft}{\kern0pt}\isafree{A}\ {\isasymrightarrow}\ \isafree{{\isasymalpha}}\ {\isasymcdot}\ \isafree{{\isasymbeta}}{\isacharbrackright}{\kern0pt}} is either the initial or the final state by 
Lemma~\ref{reaches_final_imp_last_is_init_or_final}. The implication can then be proved by 
distinguishing these two cases.

If \isa{{\isasymrho}\ {\isacharequal}{\kern0pt}\ i\ {\isacharhash}{\kern0pt}\ {\isasymsigma}}, we know by Lemma~\ref{ipda_reaches_final_imp_rm_chain} that
\begin{gather}\label{ipda_char.chain}
\isa{\isafree{i}\ {\isacharequal}{\kern0pt}\ {\isacharbrackleft}{\kern0pt}\isafree{X}\ {\isasymrightarrow}\ \isafree{{\isasymalpha}{\isacharprime}{\kern0pt}}\ {\isasymcdot}\ \isaconst{Nt}\ \isafree{A}\ {\isacharhash}{\kern0pt}\ \isafree{{\isasymbeta}{\isacharprime}{\kern0pt}}{\isacharbrackright}{\kern0pt}}\\
\isa{\isaconst{P{\isacharprime}{\kern0pt}}\ {\isasymturnstile}\ {\isacharbrackleft}{\kern0pt}\isaconst{Nt}\ \isaconst{S{\isacharprime}{\kern0pt}}{\isacharbrackright}{\kern0pt}\ {\isasymmidarrow}\isafree{i}\ {\isacharhash}{\kern0pt}\ \isafree{{\isasymsigma}}{\isasymrightarrow}r{\isacharasterisk}{\kern0pt}\ \isafree{{\isasymgamma}}}
\end{gather}
For some \isa{X\ {\isacharcolon}{\kern0pt}{\isacharcolon}{\kern0pt}\ {\isacharprime}{\kern0pt}n}, and \isa{{\isasymalpha}{\isacharprime}{\kern0pt}{\isacharcomma}{\kern0pt}\ {\isasymbeta}{\isacharcomma}{\kern0pt}\ {\isasymgamma}\ {\isacharcolon}{\kern0pt}{\isacharcolon}{\kern0pt}\ syms}. Furthermore, by Lemma~\ref{reaches_final_imp_completes}, 
our initial assumption of the accepting computation implies the existence of some \isa{v} with
\[ \isa{{\isacharparenleft}{\kern0pt}{\isacharbrackleft}{\kern0pt}\isafree{A}\ {\isasymrightarrow}\ \isafree{{\isasymalpha}}\ {\isacharat}{\kern0pt}\ \isafree{{\isasymbeta}}\ {\isasymcdot}\ {\isacharbrackright}{\kern0pt}\ {\isacharhash}{\kern0pt}\ \isafree{i}\ {\isacharhash}{\kern0pt}\ \isafree{{\isasymsigma}}{\isacharcomma}{\kern0pt}\ \isafree{v}{\isacharparenright}{\kern0pt}\ {\isasymturnstile}I{\isacharasterisk}{\kern0pt}\ {\isacharparenleft}{\kern0pt}{\isacharbrackleft}{\kern0pt}{\isacharbrackleft}{\kern0pt}\isaconst{S{\isacharprime}{\kern0pt}}\ {\isasymrightarrow}\ {\isacharbrackleft}{\kern0pt}\isaconst{Nt}\ \isaconst{S}{\isacharbrackright}{\kern0pt}\ {\isasymcdot}\ {\isacharbrackright}{\kern0pt}{\isacharbrackright}{\kern0pt}{\isacharcomma}{\kern0pt}\ {\isacharbrackleft}{\kern0pt}{\isacharbrackright}{\kern0pt}{\isacharparenright}{\kern0pt}} \]
Since in the LHS of this computation the topmost stack item is complete, the only possible first step
is a reducing transition, meaning 
\begin{multline}\label{ipda_char.red}
\isa{{\isacharparenleft}{\kern0pt}{\isacharbrackleft}{\kern0pt}\isafree{A}\ {\isasymrightarrow}\ \isafree{{\isasymalpha}}\ {\isacharat}{\kern0pt}\ \isafree{{\isasymbeta}}\ {\isasymcdot}\ {\isacharbrackright}{\kern0pt}\ {\isacharhash}{\kern0pt}\ \isafree{i}\ {\isacharhash}{\kern0pt}\ \isafree{{\isasymsigma}}{\isacharcomma}{\kern0pt}\ \isafree{v}{\isacharparenright}{\kern0pt}\ {\isasymturnstile}I\ {\isacharparenleft}{\kern0pt}{\isacharbrackleft}{\kern0pt}\isafree{X}\ {\isasymrightarrow}\ \isafree{{\isasymalpha}{\isacharprime}{\kern0pt}}\ {\isacharat}{\kern0pt}\ {\isacharbrackleft}{\kern0pt}\isaconst{Nt}\ \isafree{A}{\isacharbrackright}{\kern0pt}\ {\isasymcdot}\ \isafree{{\isasymbeta}{\isacharprime}{\kern0pt}}{\isacharbrackright}{\kern0pt}\ {\isacharhash}{\kern0pt}\ \isafree{{\isasymsigma}}{\isacharcomma}{\kern0pt}\ \isafree{v}{\isacharparenright}{\kern0pt}}\\
  \isa{{\isasymturnstile}I{\isacharasterisk}{\kern0pt}}\ \isa{{\isacharparenleft}{\kern0pt}{\isacharbrackleft}{\kern0pt}{\isacharbrackleft}{\kern0pt}\isaconst{S{\isacharprime}{\kern0pt}}\ {\isasymrightarrow}\ {\isacharbrackleft}{\kern0pt}\isaconst{Nt}\ \isaconst{S}{\isacharbrackright}{\kern0pt}\ {\isasymcdot}\ {\isacharbrackright}{\kern0pt}{\isacharbrackright}{\kern0pt}{\isacharcomma}{\kern0pt}\ {\isacharbrackleft}{\kern0pt}{\isacharbrackright}{\kern0pt}{\isacharparenright}{\kern0pt}}
\end{multline}
with \eqref{ipda_char.chain}. With the IH, and substituting \isa{\isafree{{\isasymrho}}\ {\isacharequal}{\kern0pt}\ \isafree{i}\ {\isacharhash}{\kern0pt}\ \isafree{{\isasymsigma}}} and 
\eqref{ipda_char.chain}, we now have
\[ \isa{{\isacharparenleft}{\kern0pt}{\isacharbrackleft}{\kern0pt}\isaconst{S{\isacharprime}{\kern0pt}}\ {\isasymrightarrow}\ \ {\isasymcdot}\ {\isacharbrackleft}{\kern0pt}\isaconst{Nt}\ \isaconst{S}{\isacharbrackright}{\kern0pt}\ {\isacharbrackright}{\kern0pt}{\isacharcomma}{\kern0pt}\ \isaconst{hist}\ {\isacharparenleft}{\kern0pt}\isaconst{rev}\ \isafree{{\isasymrho}}{\isacharparenright}{\kern0pt}\ {\isacharat}{\kern0pt}\ {\isacharbrackleft}{\kern0pt}\isaconst{Nt}\ \isafree{A}{\isacharbrackright}{\kern0pt}{\isacharparenright}{\kern0pt}\ {\isasymturnstile}c{\isacharasterisk}{\kern0pt}\ {\isacharparenleft}{\kern0pt}{\isacharbrackleft}{\kern0pt}\isafree{X}\ {\isasymrightarrow}\ \isafree{{\isasymalpha}{\isacharprime}{\kern0pt}}\ {\isacharat}{\kern0pt}\ {\isacharbrackleft}{\kern0pt}\isaconst{Nt}\ \isafree{A}{\isacharbrackright}{\kern0pt}\ {\isasymcdot}\ \isafree{{\isasymbeta}{\isacharprime}{\kern0pt}}{\isacharbrackright}{\kern0pt}{\isacharcomma}{\kern0pt}\ {\isacharbrackleft}{\kern0pt}{\isacharbrackright}{\kern0pt}{\isacharparenright}{\kern0pt}}. \]
By a case distinction on the final step of this computation, we know that the initial configuration
reaches \isa{{\isacharparenleft}{\kern0pt}{\isacharbrackleft}{\kern0pt}\isafree{X}\ {\isasymrightarrow}\ \isafree{{\isasymalpha}{\isacharprime}{\kern0pt}}\ {\isasymcdot}\ \isaconst{Nt}\ \isafree{A}\ {\isacharhash}{\kern0pt}\ \isafree{{\isasymbeta}{\isacharprime}{\kern0pt}}{\isacharbrackright}{\kern0pt}{\isacharcomma}{\kern0pt}\ {\isacharbrackleft}{\kern0pt}\isaconst{Nt}\ \isafree{A}{\isacharbrackright}{\kern0pt}{\isacharparenright}{\kern0pt}} immediately before the final step. Since only the
prefix \isa{\isaconst{hist}\ {\isacharparenleft}{\kern0pt}\isaconst{rev}\ \isafree{{\isasymrho}}{\isacharparenright}{\kern0pt}} of the input is consumed, the computation is independent of the
remaining string. This means that
\begin{equation}\label{ipda_char.calc}
\isa{{\isacharparenleft}{\kern0pt}{\isacharbrackleft}{\kern0pt}\isaconst{S{\isacharprime}{\kern0pt}}\ {\isasymrightarrow}\ \ {\isasymcdot}\ {\isacharbrackleft}{\kern0pt}\isaconst{Nt}\ \isaconst{S}{\isacharbrackright}{\kern0pt}\ {\isacharbrackright}{\kern0pt}{\isacharcomma}{\kern0pt}\ \isaconst{hist}\ {\isacharparenleft}{\kern0pt}\isaconst{rev}\ \isafree{{\isasymrho}}{\isacharparenright}{\kern0pt}\ {\isacharat}{\kern0pt}\ \isafree{{\isasymalpha}}{\isacharparenright}{\kern0pt}\ {\isasymturnstile}c{\isacharasterisk}{\kern0pt}\ {\isacharparenleft}{\kern0pt}{\isacharbrackleft}{\kern0pt}\isafree{X}\ {\isasymrightarrow}\ \isafree{{\isasymalpha}{\isacharprime}{\kern0pt}}\ {\isasymcdot}\ \isaconst{Nt}\ \isafree{A}\ {\isacharhash}{\kern0pt}\ \isafree{{\isasymbeta}{\isacharprime}{\kern0pt}}{\isacharbrackright}{\kern0pt}{\isacharcomma}{\kern0pt}\ \isafree{{\isasymalpha}}{\isacharparenright}{\kern0pt}}
\end{equation}
also holds. Furthermore, \eqref{ipda_char.red} implies that \mbox{\isa{{\isacharparenleft}{\kern0pt}\isafree{A}{\isacharcomma}{\kern0pt}\ \isafree{{\isasymalpha}}\ {\isacharat}{\kern0pt}\ \isafree{{\isasymbeta}}{\isacharparenright}{\kern0pt}\ {\isasymin}\ \isaconst{P{\isacharprime}{\kern0pt}}}} 
by the definition of the \isa{I\isactrlsub G} transition relations. Therefore, \eqref{ipda_char.calc} can continue
with
\[ \isa{{\isasymdots}\ {\isasymturnstile}c\ {\isacharparenleft}{\kern0pt}{\isacharbrackleft}{\kern0pt}A\ {\isasymrightarrow}\ {\isacharbrackleft}{\kern0pt}{\isacharbrackright}{\kern0pt}\ {\isasymcdot}\ {\isasymalpha}\ {\isacharat}{\kern0pt}\ {\isasymbeta}{\isacharbrackright}{\kern0pt}{\isacharcomma}{\kern0pt}\ {\isasymalpha}{\isacharparenright}{\kern0pt}}. \]
The RHS can then reach our target configuration by Lemma~\ref{char_steps_consume}, completing the
proof. Along with Lemmas~\ref{char_imp_derivers} and \ref{derivers_imp_ipda}, this also completes 
the proof of Theorem~\ref{char_derivers_ipda_iffs}.
\end{proof}
\end{lemma}

With this, we have proved all three implications. It is worth noting that, for the last two
implications, our proofs follow the main argumentation of Wilhelm et al., since our original
motivation to introduce the concept of rightmost chains was precisely to formalize the notion of
decomposing rightmost derivations akin to what Wilhelm et al.\ describe, and they, like we did in our
formalization of rightmost chains, argue that IPDA stacks leading to accepting configurations have
a structure that closely resembles these decompositions, but this resemblance is not explicitly stated
in the book.

We now define the notion of \concept{reliable prefixes}:

\begin{definition}[Reliable prefix]
\isa{{\isasymgamma}} is a reliable prefix for item \isa{{\isacharbrackleft}{\kern0pt}\isafree{A}\ {\isasymrightarrow}\ \isafree{{\isasymalpha}}\ {\isasymcdot}\ \isafree{{\isasymbeta}}{\isacharbrackright}{\kern0pt}} (with Isabelle notation 
``\isa{\isaconst{reliable{\isacharunderscore}{\kern0pt}prefix}\ {\isacharbrackleft}{\kern0pt}\isafree{A}\ {\isasymrightarrow}\ \isafree{{\isasymalpha}}\ {\isasymcdot}\ \isafree{{\isasymbeta}}{\isacharbrackright}{\kern0pt}\ \isafree{{\isasymgamma}}}'') if there exists a rightmost derivation 
\begin{equation*}
\isa{P{\isacharprime}{\kern0pt}\ {\isasymturnstile}\ {\isacharbrackleft}{\kern0pt}Nt\ S{\isacharprime}{\kern0pt}{\isacharbrackright}{\kern0pt}\ {\isasymRightarrow}r{\isacharasterisk}{\kern0pt}\ {\isasymgamma}{\isacharprime}{\kern0pt}\ {\isacharat}{\kern0pt}\ Nt\ A\ {\isacharhash}{\kern0pt}\ map\ Tm\ w\ {\isasymRightarrow}r\ {\isasymgamma}{\isacharprime}{\kern0pt}\ {\isacharat}{\kern0pt}\ {\isasymalpha}\ {\isacharat}{\kern0pt}\ {\isasymbeta}\ {\isacharat}{\kern0pt}\ map\ Tm\ w}
\end{equation*}
such that \isa{\isafree{{\isasymgamma}}\ {\isacharequal}{\kern0pt}\ \isafree{{\isasymgamma}{\isacharprime}{\kern0pt}}\ {\isacharat}{\kern0pt}\ \isafree{{\isasymalpha}}}. Alternatively, we say that the item \isa{{\isacharbrackleft}{\kern0pt}\isafree{A}\ {\isasymrightarrow}\ \isafree{{\isasymalpha}}\ {\isasymcdot}\ \isafree{{\isasymbeta}}{\isacharbrackright}{\kern0pt}} is 
\concept{valid} for \isa{{\isasymgamma}}. We also define the set of all valid items for \isa{{\isasymgamma}}:
\[ \isa{\isaconst{valids}\ \isavar{{\isasymgamma}}\ {\isasymequiv}\ {\isacharbraceleft}{\kern0pt}\isabound{i}\ {\isacharbar}{\kern0pt}\ \isaconst{reliable{\isacharunderscore}{\kern0pt}prefix}\ \isabound{i}\ \isavar{{\isasymgamma}}{\isacharbraceright}{\kern0pt}}. \]
\end{definition}

\begin{theorem}[Equivalence of \isa{\ensuremath{\mathrm{char}(G)}} computations and reliable prefixes]\label{char_eq_reliable_prefix}
There exists a \isa{\ensuremath{\mathrm{char}(G)}} computation \isa{{\isacharparenleft}{\kern0pt}{\isacharbrackleft}{\kern0pt}\isaconst{S{\isacharprime}{\kern0pt}}\ {\isasymrightarrow}\ \ {\isasymcdot}\ {\isacharbrackleft}{\kern0pt}\isaconst{Nt}\ \isaconst{S}{\isacharbrackright}{\kern0pt}\ {\isacharbrackright}{\kern0pt}{\isacharcomma}{\kern0pt}\ \isafree{{\isasymgamma}}{\isacharparenright}{\kern0pt}\ {\isasymturnstile}c{\isacharasterisk}{\kern0pt}\ {\isacharparenleft}{\kern0pt}{\isacharbrackleft}{\kern0pt}\isafree{A}\ {\isasymrightarrow}\ \isafree{{\isasymalpha}}\ {\isasymcdot}\ \isafree{{\isasymbeta}}{\isacharbrackright}{\kern0pt}{\isacharcomma}{\kern0pt}\ {\isacharbrackleft}{\kern0pt}{\isacharbrackright}{\kern0pt}{\isacharparenright}{\kern0pt}}
if and only if \isa{{\isacharbrackleft}{\kern0pt}\isafree{A}\ {\isasymrightarrow}\ \isafree{{\isasymalpha}}\ {\isasymcdot}\ \isafree{{\isasymbeta}}{\isacharbrackright}{\kern0pt}} is valid for \isa{{\isasymgamma}}.
\begin{proof}
This is a consequence of Lemmas~\ref{char_imp_derivers}, \ref{derivers_imp_ipda}, and 
\ref{ipda_imp_char}, using the fact that since \isa{G{\isacharprime}{\kern0pt}} is reduced, any sentential form derived from 
\isa{\isaconst{S{\isacharprime}{\kern0pt}}}, or a substring thereof, derives some \isa{v\ {\isacharcolon}{\kern0pt}{\isacharcolon}{\kern0pt}\ {\isacharprime}{\kern0pt}t\ list}.
\end{proof}
\end{theorem}

By this Theorem, we get two more interesting results:

\begin{corollary}
\isa{\ensuremath{\mathrm{char}(G)}} accepts exactly the set of reliable prefixes to complete items.
\qed
\end{corollary}

\begin{lemma}\label{char_fa_nextl_is_valids}
For the \isa{nextl} function of \isa{\ensuremath{\mathrm{char}(G)}} holds \isa{\isaconst{nextl}\ {\isacharparenleft}{\kern0pt}\isaconst{init}\ \ensuremath{\mathrm{char}(G)}{\isacharparenright}{\kern0pt}\ \isavar{{\isasymgamma}}\ {\isacharequal}{\kern0pt}\ \isaconst{valids}\ \isavar{{\isasymgamma}}}.
\qed
\end{lemma}%
\end{isamarkuptext}\isamarkuptrue%
%
\isadelimdocument
%
\endisadelimdocument
%
\isatagdocument
%
\isamarkupsection{The Canonical \isa{LR{\isacharparenleft}{\kern0pt}{\isadigit{0}}{\isacharparenright}{\kern0pt}} Automaton%
}
\isamarkuptrue%
%
\endisatagdocument
{\isafolddocument}%
%
\isadelimdocument
%
\endisadelimdocument
%
\begin{isamarkuptext}%
Now that we have defined \isa{\ensuremath{\mathrm{char}(G)}} and proved several useful properties, we can finally
define a deterministic automaton that our parser can use. Since \isa{\ensuremath{\mathrm{char}(G)}} is an NFA, we 
define the \concept{canonical \isa{LR{\isacharparenleft}{\kern0pt}{\isadigit{0}}{\isacharparenright}{\kern0pt}} automaton} \isa{\ensuremath{LR_0(G)}} as the \isa{{\isacharparenleft}{\kern0pt}{\isacharparenleft}{\kern0pt}\isatfree{{\isacharprime}{\kern0pt}n}{\isacharcomma}{\kern0pt}\ \isatfree{{\isacharprime}{\kern0pt}t}{\isacharparenright}{\kern0pt}\ \isatconst{sym}{\isacharcomma}{\kern0pt}\ {\isacharparenleft}{\kern0pt}\isatfree{{\isacharprime}{\kern0pt}n}{\isacharcomma}{\kern0pt}\ \isatfree{{\isacharprime}{\kern0pt}t}{\isacharparenright}{\kern0pt}\ \isatconst{item}\ \isatconst{set}{\isacharparenright}{\kern0pt}\ \isatconst{dfa}} resulting from 
the powerset construction of \isa{\ensuremath{\mathrm{char}(G)}} restricted to reachable states. The automaton is once more defined using 
Paulson's theory~\cite{Paulson}. We now show some properties we will need in our parser.

\begin{corollary}\label{nextl_dfa_LR0_is_valids}
The state that \isa{\ensuremath{LR_0(G)}} reaches after reading input word \isa{{\isasymgamma}} is \isa{valids\ {\isasymgamma}}.
\begin{proof}
This is a consequence of Lemma~\ref{char_fa_nextl_is_valids} and the definition of the transition 
function of the DFA resulting from the powerset construction.
\end{proof}
\end{corollary}

\begin{lemma}\label{state_imp_valids}
For every state \isa{\isafree{q}\ {\isasymin}\ \isaconst{states}\ \ensuremath{LR_0(G)}}, there exists a \isa{{\isasymgamma}\ {\isacharcolon}{\kern0pt}{\isacharcolon}{\kern0pt}\ syms} such that 
\isa{\isafree{q}\ {\isacharequal}{\kern0pt}\ \isaconst{valids}\ \isafree{{\isasymgamma}}}.
\begin{proof}
Since we define the states of \isa{\ensuremath{LR_0(G)}} to be restricted to the reachable states resulting 
from the powerset construction of \isa{\ensuremath{\mathrm{char}(G)}}, we also know that there exists an input string \isa{{\isasymgamma}} 
such that \isa{\ensuremath{LR_0(G)}} is in state \isa{q} after reading \isa{{\isasymgamma}}. By Corollary~\ref{nextl_dfa_LR0_is_valids},
this implies \isa{\mbox{\isafree{q}\ {\isacharequal}{\kern0pt}\ \isaconst{valids}\ \isafree{{\isasymgamma}}}}.
\end{proof}
\end{lemma}

\begin{lemma}\label{char_fa_nxts_is_shifts}
The equality
\[ \isa{{\isacharparenleft}{\kern0pt}{\isasymUnion}\isactrlbsub \isabound{i}{\isasymin}\isafree{Q}\isactrlesub \ \isaconst{nxt}\ \ensuremath{\mathrm{char}(G)}\ \isabound{i}\ \isafree{A}{\isacharparenright}{\kern0pt}\ {\isacharequal}{\kern0pt}\ {\isacharbraceleft}{\kern0pt}{\isacharbrackleft}{\kern0pt}\isabound{X}\ {\isasymrightarrow}\ \isabound{{\isasymalpha}}\ {\isacharat}{\kern0pt}\ {\isacharbrackleft}{\kern0pt}\isafree{A}{\isacharbrackright}{\kern0pt}\ {\isasymcdot}\ \isabound{{\isasymbeta}}{\isacharbrackright}{\kern0pt}\ {\isacharbar}{\kern0pt}\ {\isacharbrackleft}{\kern0pt}\isabound{X}\ {\isasymrightarrow}\ \isabound{{\isasymalpha}}\ {\isasymcdot}\ \isafree{A}\ {\isacharhash}{\kern0pt}\ \isabound{{\isasymbeta}}{\isacharbrackright}{\kern0pt}\ {\isasymin}\ \isafree{Q}{\isacharbraceright}{\kern0pt}}. \] 
holds for any \isa{\isafree{Q}\ {\isasymsubseteq}\ \ensuremath{\mathrm{It}_{\isaconst{G{\isacharprime}{\kern0pt}}}}}.
\begin{proof}
We fix some \isa{\isafree{i}\ {\isasymin}\ {\isacharparenleft}{\kern0pt}{\isasymUnion}\isactrlbsub \isabound{i}{\isasymin}\isafree{Q}\isactrlesub \ \isaconst{nxt}\ \ensuremath{\mathrm{char}(G)}\ \isabound{i}\ \isafree{A}{\isacharparenright}{\kern0pt}}, and we need to show that it is in the set 
on the RHS of the equation. Since \isa{i} is in the union, there exists a \isa{j} such that
\isa{\isafree{i}\ {\isasymin}\ \isaconst{nxt}\ \ensuremath{\mathrm{char}(G)}\ \isafree{j}\ \isafree{A}}. Since the \isa{nxt} relation is only defined for incomplete items, 
\isa{j} must be of the form \isa{{\isacharbrackleft}{\kern0pt}\isafree{X}\ {\isasymrightarrow}\ \isafree{{\isasymalpha}}\ {\isasymcdot}\ \isafree{B}\ {\isacharhash}{\kern0pt}\ \isafree{{\isasymbeta}}{\isacharbrackright}{\kern0pt}}. 
\isa{\isafree{i}\ {\isasymin}\ {\isacharbraceleft}{\kern0pt}{\isacharbrackleft}{\kern0pt}\isabound{X}\ {\isasymrightarrow}\ \isabound{{\isasymalpha}}\ {\isacharat}{\kern0pt}\ {\isacharbrackleft}{\kern0pt}\isafree{A}{\isacharbrackright}{\kern0pt}\ {\isasymcdot}\ \isabound{{\isasymbeta}}{\isacharbrackright}{\kern0pt}\ {\isacharbar}{\kern0pt}\ {\isacharbrackleft}{\kern0pt}\isabound{X}\ {\isasymrightarrow}\ \isabound{{\isasymalpha}}\ {\isasymcdot}\ \isafree{A}\ {\isacharhash}{\kern0pt}\ \isabound{{\isasymbeta}}{\isacharbrackright}{\kern0pt}\ {\isasymin}\ \isafree{Q}{\isacharbraceright}{\kern0pt}} then follows trivially by the 
definition of the \isa{nxt} transition function.

For the converse, there exists some \isa{{\isacharbrackleft}{\kern0pt}\isafree{X}\ {\isasymrightarrow}\ \isafree{{\isasymalpha}}\ {\isasymcdot}\ \isafree{A}\ {\isacharhash}{\kern0pt}\ \isafree{{\isasymbeta}}{\isacharbrackright}{\kern0pt}\ {\isasymin}\ \isafree{Q}} for which 
\isa{\isafree{i}\ {\isacharequal}{\kern0pt}\ {\isacharbrackleft}{\kern0pt}\isafree{X}\ {\isasymrightarrow}\ \isafree{{\isasymalpha}}\ {\isacharat}{\kern0pt}\ {\isacharbrackleft}{\kern0pt}\isafree{A}{\isacharbrackright}{\kern0pt}\ {\isasymcdot}\ \isafree{{\isasymbeta}}{\isacharbrackright}{\kern0pt}}. With \isa{\isafree{Q}\ {\isasymsubseteq}\ \ensuremath{\mathrm{It}_{\isaconst{G{\isacharprime}{\kern0pt}}}}}, this implies \isa{\isafree{i}\ {\isasymin}\ \ensuremath{\mathrm{It}_{\isaconst{G{\isacharprime}{\kern0pt}}}}},
and \isa{\isafree{i}\ {\isasymin}\ {\isacharparenleft}{\kern0pt}{\isasymUnion}\isactrlbsub \isabound{i}{\isasymin}\isafree{Q}\isactrlesub \ \isaconst{nxt}\ \ensuremath{\mathrm{char}(G)}\ \isabound{i}\ \isafree{A}{\isacharparenright}{\kern0pt}} follows directly by definition of \isa{nxt}.
\end{proof}
\end{lemma}

\begin{lemma}\label{eps_reliable_preserved}
If \isa{i} is valid for some \isa{{\isasymgamma}\ {\isacharcolon}{\kern0pt}{\isacharcolon}{\kern0pt}\ syms}, and for some item \isa{k} 
\[ \isa{{\isacharparenleft}{\kern0pt}\isafree{i}{\isacharcomma}{\kern0pt}\ \isafree{k}{\isacharparenright}{\kern0pt}\ {\isasymin}\ {\isacharparenleft}{\kern0pt}\isaconst{eps}\ \ensuremath{\mathrm{char}(G)}{\isacharparenright}{\kern0pt}\isactrlsup {\isacharasterisk}{\kern0pt}} \] 
holds, then \isa{k} is also valid for \isa{{\isasymgamma}}.
\begin{proof}
The proof is by backward induction on the number of \isa{{\isasymepsilon}}-transition steps using 
Theorem~\ref{char_eq_reliable_prefix}.
\end{proof}
\end{lemma}

\begin{lemma}\label{dfa_LR0_nxt_is_epsclo_of_shift}
For any state of \isa{\ensuremath{LR_0(G)}} \isa{\isafree{Q}} and a symbol \isa{Y}, 
\isa{\isaconst{nxt}\ \ensuremath{LR_0(G)}\ \isafree{Q}\ \isafree{Y}} is equivalent to the set
\[ \isa{\isaconst{epsclo}\ {\isacharbraceleft}{\kern0pt}{\isacharbrackleft}{\kern0pt}\isabound{X}\ {\isasymrightarrow}\ \isabound{{\isasymalpha}}\ {\isacharat}{\kern0pt}\ {\isacharbrackleft}{\kern0pt}\isafree{Y}{\isacharbrackright}{\kern0pt}\ {\isasymcdot}\ \isabound{{\isasymbeta}}{\isacharbrackright}{\kern0pt}\ {\isacharbar}{\kern0pt}\ {\isacharbrackleft}{\kern0pt}\isabound{X}\ {\isasymrightarrow}\ \isabound{{\isasymalpha}}\ {\isasymcdot}\ \isafree{Y}\ {\isacharhash}{\kern0pt}\ \isabound{{\isasymbeta}}{\isacharbrackright}{\kern0pt}\ {\isasymin}\ \isafree{Q}{\isacharbraceright}{\kern0pt}} \]
where \isa{\isaconst{epsclo}\ \isafree{A}} denotes the \isa{{\isasymepsilon}}-closure of the set \isa{A} for \isa{\ensuremath{\mathrm{char}(G)}}.
\begin{proof}
Follows from the definition of the \isa{nxt} function of the DFA resulting from the powerset construction,
as well as Lemma~\ref{char_fa_nxts_is_shifts}.
\end{proof}
\end{lemma}

We now prove a lemma that will be once again essential to numerous proofs in the coming section, but
was not part of the work of Wilhelm et al.

\begin{lemma}\label{nxt_dfa_LR0_shift_is_valids_app}
If \isa{\isaconst{valids}\ \isafree{{\isasymgamma}}\ {\isasymin}\ \isaconst{states}\ \ensuremath{LR_0(G)}}, then \isa{\isaconst{nxt}\ \ensuremath{LR_0(G)}\ {\isacharparenleft}{\kern0pt}\isaconst{valids}\ \isavar{{\isasymgamma}}{\isacharparenright}{\kern0pt}\ \isavar{X}\ {\isacharequal}{\kern0pt}\ \isaconst{valids}\ {\isacharparenleft}{\kern0pt}\isavar{{\isasymgamma}}\ {\isacharat}{\kern0pt}\ {\isacharbrackleft}{\kern0pt}\isavar{X}{\isacharbrackright}{\kern0pt}{\isacharparenright}{\kern0pt}}.
\begin{proof}
By Lemma~\ref{dfa_LR0_nxt_is_epsclo_of_shift}, \isa{\isaconst{nxt}\ \ensuremath{LR_0(G)}\ {\isacharparenleft}{\kern0pt}\isaconst{valids}\ \isafree{{\isasymgamma}}{\isacharparenright}{\kern0pt}\ \isafree{X}} is equivalent to 
the set
\[ \isa{\isaconst{epsclo}\ {\isacharbraceleft}{\kern0pt}{\isacharbrackleft}{\kern0pt}\isabound{A}\ {\isasymrightarrow}\ \isabound{{\isasymalpha}}\ {\isacharat}{\kern0pt}\ {\isacharbrackleft}{\kern0pt}\isafree{X}{\isacharbrackright}{\kern0pt}\ {\isasymcdot}\ \isabound{{\isasymbeta}}{\isacharbrackright}{\kern0pt}\ {\isacharbar}{\kern0pt}\ {\isacharbrackleft}{\kern0pt}\isabound{A}\ {\isasymrightarrow}\ \isabound{{\isasymalpha}}\ {\isasymcdot}\ \isafree{X}\ {\isacharhash}{\kern0pt}\ \isabound{{\isasymbeta}}{\isacharbrackright}{\kern0pt}\ {\isasymin}\ \isaconst{valids}\ \isafree{{\isasymgamma}}{\isacharbraceright}{\kern0pt}}. \]
We will abbreviate this set as \isa{{\isasymE}}.

First, we will show that \isa{{\isasymE}} is a subset of \isa{\isaconst{valids}\ {\isacharparenleft}{\kern0pt}\isafree{{\isasymgamma}}\ {\isacharat}{\kern0pt}\ {\isacharbrackleft}{\kern0pt}\isafree{X}{\isacharbrackright}{\kern0pt}{\isacharparenright}{\kern0pt}} by 
Lemma~\ref{eps_reliable_preserved}.

We now show that \isa{\isaconst{valids}\ {\isacharparenleft}{\kern0pt}\isafree{{\isasymgamma}}\ {\isacharat}{\kern0pt}\ {\isacharbrackleft}{\kern0pt}\isafree{X}{\isacharbrackright}{\kern0pt}{\isacharparenright}{\kern0pt}} is a subset of \isa{{\isasymE}} to complete the proof.

We assume \isa{\isafree{i}\ {\isasymin}\ \isaconst{valids}\ {\isacharparenleft}{\kern0pt}\isafree{{\isasymgamma}}\ {\isacharat}{\kern0pt}\ {\isacharbrackleft}{\kern0pt}\isafree{X}{\isacharbrackright}{\kern0pt}{\isacharparenright}{\kern0pt}}, which means that either
\isa{i} is a complete item, or an incomplete item of the form \isa{{\isacharbrackleft}{\kern0pt}\isafree{A}\ {\isasymrightarrow}\ \isafree{{\isasymalpha}}\ {\isacharat}{\kern0pt}\ {\isacharbrackleft}{\kern0pt}\isafree{X}{\isacharbrackright}{\kern0pt}\ {\isasymcdot}\ \isafree{{\isasymbeta}}{\isacharbrackright}{\kern0pt}}.
We now distinguish these two cases.

If \isa{i} is complete, we know it is a reachable state of \isa{\ensuremath{\mathrm{char}(G)}}. We show by induction on the 
length of the \isa{\ensuremath{\mathrm{char}(G)}} computation from the initial state to \isa{i} that there exists an item 
\isa{{\isacharbrackleft}{\kern0pt}\isafree{Y}\ {\isasymrightarrow}\ \isafree{{\isasymdelta}}\ {\isacharat}{\kern0pt}\ {\isacharbrackleft}{\kern0pt}\isafree{X}{\isacharbrackright}{\kern0pt}\ {\isasymcdot}\ \isafree{{\isasymzeta}}{\isacharbrackright}{\kern0pt}} such that
\[ \isa{{\isacharparenleft}{\kern0pt}{\isacharbrackleft}{\kern0pt}S{\isacharprime}{\kern0pt}\ {\isasymrightarrow}\ {\isacharbrackleft}{\kern0pt}{\isacharbrackright}{\kern0pt}\ {\isasymcdot}\ {\isacharbrackleft}{\kern0pt}Nt\ S{\isacharbrackright}{\kern0pt}{\isacharbrackright}{\kern0pt}{\isacharcomma}{\kern0pt}\ {\isasymgamma}\ {\isacharat}{\kern0pt}\ {\isacharbrackleft}{\kern0pt}X{\isacharbrackright}{\kern0pt}{\isacharparenright}{\kern0pt}\ {\isasymturnstile}c{\isacharasterisk}{\kern0pt}\ {\isacharparenleft}{\kern0pt}{\isacharbrackleft}{\kern0pt}Y\ {\isasymrightarrow}\ {\isasymdelta}\ {\isacharat}{\kern0pt}\ {\isacharbrackleft}{\kern0pt}X{\isacharbrackright}{\kern0pt}\ {\isasymcdot}\ {\isasymzeta}{\isacharbrackright}{\kern0pt}{\isacharcomma}{\kern0pt}\ {\isacharbrackleft}{\kern0pt}{\isacharbrackright}{\kern0pt}{\isacharparenright}{\kern0pt}\ {\isasymturnstile}c{\isacharasterisk}{\kern0pt}\ {\isacharparenleft}{\kern0pt}i{\isacharcomma}{\kern0pt}\ {\isacharbrackleft}{\kern0pt}{\isacharbrackright}{\kern0pt}{\isacharparenright}{\kern0pt}}. \]
We can then show that \isa{\isafree{i}\ {\isasymin}\ \isafree{{\isasymE}}} by Theorem~\ref{char_eq_reliable_prefix} and the fact
that \isa{{\isacharparenleft}{\kern0pt}{\isacharbrackleft}{\kern0pt}\isafree{Y}\ {\isasymrightarrow}\ \isafree{{\isasymdelta}}\ {\isacharat}{\kern0pt}\ {\isacharbrackleft}{\kern0pt}\isafree{X}{\isacharbrackright}{\kern0pt}\ {\isasymcdot}\ \isafree{{\isasymzeta}}{\isacharbrackright}{\kern0pt}{\isacharcomma}{\kern0pt}\ \isafree{i}{\isacharparenright}{\kern0pt}\ {\isasymin}\ {\isacharparenleft}{\kern0pt}\isaconst{eps}\ \ensuremath{\mathrm{char}(G)}{\isacharparenright}{\kern0pt}\isactrlsup {\isacharasterisk}{\kern0pt}}, using Lemma~\ref{eps_reliable_preserved}
once again.
\end{proof}
\end{lemma}

As is the case with any DFA resulting from the powerset construction of an NFA, the canonical 
\isa{LR{\isacharparenleft}{\kern0pt}{\isadigit{0}}{\isacharparenright}{\kern0pt}} automaton effectively navigates all possible \isa{\ensuremath{\mathrm{char}(G)}} computations simultaneously,
thereby achieving deterministic behavior. Furthermore, Lemmas~\ref{state_imp_valids} and
\ref{nxt_dfa_LR0_shift_is_valids_app} tell us that all \isa{\ensuremath{LR_0(G)}} states are sets of valid items,
and if the current items are valid for \isa{{\isasymgamma}}, passing a symbol \isa{X} to the transition function of
\isa{\ensuremath{LR_0(G)}} yields the set of items valid for \isa{{\isasymgamma}\ {\isacharat}{\kern0pt}\ {\isacharbrackleft}{\kern0pt}X{\isacharbrackright}{\kern0pt}}. Hence, we can use the canonical \isa{LR{\isacharparenleft}{\kern0pt}{\isadigit{0}}{\isacharparenright}{\kern0pt}}
automaton to determine in
parallel which items one could possibly be processing based on what has already been derived. This
also shows why we can use \isa{\ensuremath{LR_0(G)}} in our parser in order to find a rightmost derivation for the
input: building a reliable prefix from left to right starting from the left of the input is the
opposite of applying derivations on the rightmost nonterminal, which build the word from right to
left. Thus, we are able to reconstruct the rightmost derivation backwards, with the presence of
complete items allowing us to detect positions for reduction as we will now see in the definition
of our parser.%
\end{isamarkuptext}\isamarkuptrue%
%
\isadelimdocument
%
\endisadelimdocument
%
\isatagdocument
%
\isamarkupchapter{The \isa{LR{\isacharparenleft}{\kern0pt}{\isadigit{0}}{\isacharparenright}{\kern0pt}} Parser%
}
\isamarkuptrue%
%
\endisatagdocument
{\isafolddocument}%
%
\isadelimdocument
%
\endisadelimdocument
%
\isadelimproof
%
\endisadelimproof
%
\isatagproof
%
\endisatagproof
{\isafoldproof}%
%
\isadelimproof
%
\endisadelimproof
%
\isadelimdocument
%
\endisadelimdocument
%
\isatagdocument
%
\isamarkupsection{Definition%
}
\isamarkuptrue%
%
\endisatagdocument
{\isafolddocument}%
%
\isadelimdocument
%
\endisadelimdocument
%
\begin{isamarkuptext}%
\begin{definition}[\isa{LR{\isacharparenleft}{\kern0pt}{\isadigit{0}}{\isacharparenright}{\kern0pt}} parser]
Let \isa{\isafree{{\isasymDelta}\isactrlsub G}} be the transition relation of the canonical \isa{LR{\isacharparenleft}{\kern0pt}{\isadigit{0}}{\isacharparenright}{\kern0pt}} automaton \isa{\ensuremath{LR_0(G)}}, 
\isa{q\isactrlsub {\isadigit{0}}} the initial state of \isa{\ensuremath{LR_0(G)}}, \isa{Q} the set of states of \isa{\ensuremath{LR_0(G)}}, and \isa{f} the singleton 
set \isa{{\isacharbraceleft}{\kern0pt}{\isacharbrackleft}{\kern0pt}\isaconst{S{\isacharprime}{\kern0pt}}\ {\isasymrightarrow}\ {\isasymcdot}\ {\isacharbrackright}{\kern0pt}{\isacharbraceright}{\kern0pt}}.
The \concept{\isa{LR{\isacharparenleft}{\kern0pt}{\isadigit{0}}{\isacharparenright}{\kern0pt}} parser} to the CFG \isa{G} is the \isa{{\isacharparenleft}{\kern0pt}{\isacharparenleft}{\kern0pt}\isatfree{{\isacharprime}{\kern0pt}n}{\isacharcomma}{\kern0pt}\ \isatfree{{\isacharprime}{\kern0pt}t}{\isacharparenright}{\kern0pt}\ \isatconst{item}\ \isatconst{set}{\isacharcomma}{\kern0pt}\ \isatfree{{\isacharprime}{\kern0pt}t}{\isacharparenright}{\kern0pt}\ \isatconst{gpda}}
\[ \isa{\ensuremath{P_0(G)}} = \isa{{\isasymlparr}\isaconst{states}\ {\isacharequal}{\kern0pt}\ \isafree{Q}\ {\isasymunion}\ {\isacharbraceleft}{\kern0pt}\isafree{f}{\isacharbraceright}{\kern0pt}{\isacharcomma}{\kern0pt}\ \isaconst{init}\ {\isacharequal}{\kern0pt}\ \isafree{q\isactrlsub {\isadigit{0}}}{\isacharcomma}{\kern0pt}\ \isaconst{final}\ {\isacharequal}{\kern0pt}\ {\isacharbraceleft}{\kern0pt}\isafree{f}{\isacharbraceright}{\kern0pt}{\isacharcomma}{\kern0pt}\ \isaconst{nxt}\ {\isacharequal}{\kern0pt}\ \isafree{{\isasymDelta}\isactrlsub {\isadigit{0}}}{\isacharcomma}{\kern0pt}\ \isaconst{eps}\ {\isacharequal}{\kern0pt}\ \isafree{{\isasymE}\isactrlsub {\isadigit{0}}}{\isasymrparr}}. \]
\end{definition}

Similarly to the IPDA, we define three types of transitions, two of which are in the
transition relation \isa{eps}:
\begin{enumerate}
\item A \concept{reading} transition takes place when the parser reads an input symbol \isa{a} and pushes 
the successor state \isa{\isafree{{\isasymDelta}\isactrlsub G}\ \isafree{q}\ {\isacharparenleft}{\kern0pt}\isaconst{Tm}\ \isafree{a}{\isacharparenright}{\kern0pt}} onto the stack, where \isa{q} is the current state. This 
transition takes place only if the successor state is not empty. Therefore, we define our \isa{nxt} 
relation as
\[ \isa{\isafree{{\isasymDelta}\isactrlsub {\isadigit{0}}}\ {\isacharequal}{\kern0pt}\ {\isacharbraceleft}{\kern0pt}{\isacharparenleft}{\kern0pt}{\isacharbrackleft}{\kern0pt}\isabound{q}{\isacharbrackright}{\kern0pt}{\isacharcomma}{\kern0pt}\ \isabound{a}{\isacharcomma}{\kern0pt}\ {\isacharbrackleft}{\kern0pt}\isafree{{\isasymDelta}\isactrlsub G}\ \isabound{q}\ {\isacharparenleft}{\kern0pt}\isaconst{Tm}\ \isabound{a}{\isacharparenright}{\kern0pt}{\isacharcomma}{\kern0pt}\ \isabound{q}{\isacharbrackright}{\kern0pt}{\isacharparenright}{\kern0pt}\ {\isacharbar}{\kern0pt}\ \isabound{q}\ {\isasymin}\ \isafree{Q}\ {\isasymand}\ \isafree{{\isasymDelta}\isactrlsub G}\ \isabound{q}\ {\isacharparenleft}{\kern0pt}\isaconst{Tm}\ \isabound{a}{\isacharparenright}{\kern0pt}\ {\isasymnoteq}\ {\isasymemptyset}{\isacharbraceright}{\kern0pt}}. \] 
\item A \concept{reducing} transition occurs when the current state contains a complete
item \isa{{\isacharbrackleft}{\kern0pt}\isafree{X}\ {\isasymrightarrow}\ \isafree{{\isasymalpha}}\ {\isasymcdot}\ {\isacharbrackright}{\kern0pt}}. If this is the case, the parser first pops the first \isa{{\isacharbar}{\kern0pt}\isafree{{\isasymalpha}}{\isacharbar}{\kern0pt}} 
items off the stack, and pushes the successor state under \isa{X} onto the stack. In other words, the 
parser replaces the topmost stack list \isa{{\isacharbrackleft}{\kern0pt}q\isactrlsub {\isadigit{1}}{\isacharcomma}{\kern0pt}\ q\isactrlsub {\isadigit{2}}{\isacharcomma}{\kern0pt}\ {\isasymdots}{\isacharcomma}{\kern0pt}} \isa{\ensuremath{\isafree{q}_{{\isacharbar}{\kern0pt}\isafree{{\isasymalpha}}{\isacharbar}{\kern0pt}}}} \isa{{\isacharcomma}{\kern0pt}\ q{\isacharbrackright}{\kern0pt}} by 
\isa{{\isacharbrackleft}{\kern0pt}\isafree{{\isasymDelta}\isactrlsub G}\ \isafree{q}\ {\isacharparenleft}{\kern0pt}\isaconst{Nt}\ \isafree{X}{\isacharparenright}{\kern0pt}{\isacharcomma}{\kern0pt}\ \isafree{q}{\isacharbrackright}{\kern0pt}}. Formally we define these transitions by the set
\begin{multline*}
  \isa{{\isacharbraceleft}{\kern0pt}let\ q\ {\isacharequal}{\kern0pt}\ last\ {\isacharparenleft}{\kern0pt}q\isactrlsub n{\isacharhash}{\kern0pt}qs{\isacharparenright}{\kern0pt}\ in\ {\isacharparenleft}{\kern0pt}q\isactrlsub n\ {\isacharhash}{\kern0pt}\ qs{\isacharcomma}{\kern0pt}\ {\isasymDelta}\isactrlsub G\ q\ {\isacharparenleft}{\kern0pt}Nt\ X{\isacharparenright}{\kern0pt}\ {\isacharhash}{\kern0pt}\ {\isacharbrackleft}{\kern0pt}q{\isacharbrackright}{\kern0pt}{\isacharparenright}{\kern0pt}}\\
      \isa{{\isacharbar}{\kern0pt}\ q\isactrlsub n\ qs\ X\ {\isasymalpha}{\isachardot}{\kern0pt}\ set\ {\isacharparenleft}{\kern0pt}q\isactrlsub n{\isacharhash}{\kern0pt}qs{\isacharparenright}{\kern0pt}\ {\isasymsubseteq}\ Q\ {\isasymand}\ {\isacharbrackleft}{\kern0pt}X\ {\isasymrightarrow}\ {\isasymalpha}\ {\isasymcdot}\ {\isacharbrackright}{\kern0pt}\ {\isasymin}\ q\isactrlsub n}\\ 
      \isa{{\isasymand}} \isa{{\isacharbar}{\kern0pt}\isafree{{\isasymalpha}}{\isacharbar}{\kern0pt}\ {\isacharequal}{\kern0pt}\ {\isacharbar}{\kern0pt}\isafree{qs}{\isacharbar}{\kern0pt}\ {\isasymand}\ \isafree{X}\ {\isasymnoteq}\ \isaconst{S{\isacharprime}{\kern0pt}}}\isa{{\isacharbraceright}{\kern0pt}{\isachardot}{\kern0pt}}
\end{multline*}
It is worth noting that, if \isa{{\isacharbar}{\kern0pt}\isafree{{\isasymalpha}}{\isacharbar}{\kern0pt}\ {\isacharequal}{\kern0pt}\ {\isadigit{0}}}, no items are popped off the stack, meaning that
the stack simply is extended by the successor of the topmost state.
\item Lastly, \concept{finishing} transitions reduce the complete item \isa{{\isacharbrackleft}{\kern0pt}\isaconst{S{\isacharprime}{\kern0pt}}\ {\isasymrightarrow}\ {\isacharbrackleft}{\kern0pt}\isaconst{Nt}\ \isaconst{S}{\isacharbrackright}{\kern0pt}\ {\isasymcdot}\ {\isacharbrackright}{\kern0pt}}. 
This transition signals that the parser has successfully reduced the processed input to the start
symbol, and it does so by reducing the topmost state to the final state \isa{\isafree{f}} if the
second to last state is \isa{q\isactrlsub {\isadigit{0}}}.
Finish transitions therefore correspond to the set 
\[ \isa{{\isacharbraceleft}{\kern0pt}{\isacharparenleft}{\kern0pt}{\isacharbrackleft}{\kern0pt}\isabound{q}{\isacharcomma}{\kern0pt}\ \isafree{q\isactrlsub {\isadigit{0}}}{\isacharbrackright}{\kern0pt}{\isacharcomma}{\kern0pt}\ {\isacharbrackleft}{\kern0pt}\isafree{f}{\isacharbrackright}{\kern0pt}{\isacharparenright}{\kern0pt}\ {\isacharbar}{\kern0pt}\ \isabound{q}\ {\isasymin}\ \isafree{Q}\ {\isasymand}\ {\isacharbrackleft}{\kern0pt}\isaconst{S{\isacharprime}{\kern0pt}}\ {\isasymrightarrow}\ {\isacharbrackleft}{\kern0pt}\isaconst{Nt}\ \isaconst{S}{\isacharbrackright}{\kern0pt}\ {\isasymcdot}\ {\isacharbrackright}{\kern0pt}\ {\isasymin}\ \isabound{q}{\isacharbraceright}{\kern0pt}}. \]
\end{enumerate}
The set \isa{{\isasymE}\isactrlsub {\isadigit{0}}} is therefore defined as the union of the sets for reduce and finish transitions. 

Our definition of \isa{\ensuremath{P_0(G)}} is very similar to that of Wilhelm et al., except that we, as we did when
defining \isa{I\isactrlsub G}, restrict the elements of the transition relations to elements of \isa{Q} explicitly
to overcome the same problem we had in the IPDA section. In this case, however, this has one more
benefit: since \isa{Q} is actually a subset of the states of \isa{\ensuremath{P_0(G)}}, by restricting our transition functions
to \isa{Q} only, instead of \isa{\isafree{Q}\ {\isasymunion}\ {\isacharbraceleft}{\kern0pt}\isafree{f}{\isacharbraceright}{\kern0pt}}, our transition relations are more well-defined, since we
guarantee that every state that gets passed to the \isa{\ensuremath{LR_0(G)}} transition function is a state in 
\isa{\ensuremath{LR_0(G)}}. A second detail we added is the condition \isa{\isafree{X}\ {\isasymnoteq}\ \isaconst{S{\isacharprime}{\kern0pt}}} in the reduce transition 
relation. This could lead to nondeterministic behavior, since the finish transition condition 
\isa{{\isacharbrackleft}{\kern0pt}\isaconst{S{\isacharprime}{\kern0pt}}\ {\isasymrightarrow}\ {\isacharbrackleft}{\kern0pt}\isaconst{Nt}\ \isaconst{S}{\isacharbrackright}{\kern0pt}\ {\isasymcdot}\ {\isacharbrackright}{\kern0pt}\ {\isasymin}\ \isafree{q}} also implies the reduce transition condition 
\isa{{\isacharbrackleft}{\kern0pt}\isafree{X}\ {\isasymrightarrow}\ \isafree{{\isasymalpha}}\ {\isasymcdot}\ {\isacharbrackright}{\kern0pt}\ {\isasymin}\ \isafree{q}}. Restricting reduce transitions to items fulfilling this additional
property causes reduce and finish transitions to be mutually exclusive, avoiding this problem.

Another minor distinction is that Wilhelm et al.\ define the final state of the parser as an 
unspecified fresh state not in \isa{Q}. To achieve this, we simply fixed \isa{f} as the singleton set we 
described above because the item \isa{{\isacharbrackleft}{\kern0pt}S{\isacharprime}{\kern0pt}\ {\isasymrightarrow}\ {\isasymcdot}\ {\isacharbrackright}{\kern0pt}} is not in \isa{\ensuremath{\mathrm{It}_{\isaconst{G{\isacharprime}{\kern0pt}}}}}, and therefore not in \isa{Q}, 
by the definition of \isa{G{\isacharprime}{\kern0pt}}. The choice of \isa{f}, however, is irrelevant as long as it is not an element
of \isa{Q}.%
\end{isamarkuptext}\isamarkuptrue%
%
\isadelimdocument
%
\endisadelimdocument
%
\isatagdocument
%
\isamarkupsubsection{Finiteness of the Transition Relations of \isa{\ensuremath{P_0(G)}}\label{sec:finite}%
}
\isamarkuptrue%
%
\endisatagdocument
{\isafolddocument}%
%
\isadelimdocument
%
\endisadelimdocument
%
\begin{isamarkuptext}%
We mentioned when defining GPDAs that we do not assume the transition relation to be finite.
This is only because the first GPDA we defined, i.e., the IPDA, only acts as a theoretical construct 
that aids us in proving useful and interesting properties in order to achieve our end goal of a
deterministic parsing algorithm. This is not the case for the \isa{LR{\isacharparenleft}{\kern0pt}{\isadigit{0}}{\isacharparenright}{\kern0pt}} parser, since this is precisely
the machine we want to use for syntactic analysis. In order to define a parser that is usable
in practice, its transition relation must be executable, i.e., we need to be able to define a
function that computes our parser's transitions. If one defines an automaton as having a transition function,
as is the case for Paulson's DFAs, for example, executability is guaranteed by construction, but the
same does not hold in general when transitions are defined as a relation, which is precisely the
situation of our parser definition.

One possibility of making such a relation executable is through tables: if one has an (n+1)-ary
relation \isa{R}, one can trivially define a function that maps each tuple \isa{{\isacharparenleft}{\kern0pt}x\isactrlsub {\isadigit{1}}{\isacharcomma}{\kern0pt}\ x\isactrlsub {\isadigit{2}}{\isacharcomma}{\kern0pt}\ {\isasymdots}{\isacharcomma}{\kern0pt}\ x\isactrlsub n{\isacharparenright}{\kern0pt}} to a list
\isa{ys} such that for every \isa{y} in \isa{ys} holds \isa{{\isacharparenleft}{\kern0pt}x\isactrlsub {\isadigit{1}}{\isacharcomma}{\kern0pt}\ x\isactrlsub {\isadigit{2}}{\isacharcomma}{\kern0pt}\ {\isasymdots}{\isacharcomma}{\kern0pt}\ x\isactrlsub n{\isacharcomma}{\kern0pt}\ y{\isacharparenright}{\kern0pt}\ {\isasymin}\ R} by iterating through the
set. It is easy to see, however, that this method is only possible if \isa{R} is finite. We will now
prove the finiteness of the transition relations of \isa{\ensuremath{P_0(G)}}.

\begin{proposition}\label{finite_nxt_P0}
The \isa{\isaconst{nxt}} relation of \isa{\ensuremath{P_0(G)}} is finite.
\begin{proof}
Let \isa{\isafree{f}\ {\isacharparenleft}{\kern0pt}\isafree{q}{\isacharcomma}{\kern0pt}\ \isafree{a}{\isacharparenright}{\kern0pt}\ {\isacharequal}{\kern0pt}\ {\isacharparenleft}{\kern0pt}{\isacharbrackleft}{\kern0pt}\isafree{q}{\isacharbrackright}{\kern0pt}{\isacharcomma}{\kern0pt}\ \isafree{a}{\isacharcomma}{\kern0pt}\ {\isacharbrackleft}{\kern0pt}\isaconst{nxt}\ \ensuremath{LR_0(G)}\ \isafree{q}\ {\isacharparenleft}{\kern0pt}\isaconst{Tm}\ \isafree{a}{\isacharparenright}{\kern0pt}{\isacharcomma}{\kern0pt}\ \isafree{q}{\isacharbrackright}{\kern0pt}{\isacharparenright}{\kern0pt}} be a mapping from \isa{item\ set\ {\isasymtimes}\ {\isacharprime}{\kern0pt}t} to
\isa{item\ set\ list\ {\isasymtimes}\ {\isacharprime}{\kern0pt}t\ {\isasymtimes}\ item\ set\ list}. Furthermore, let \isa{T} be the set 
\[ \isa{\isafree{T}\ {\isacharequal}{\kern0pt}\ {\isacharbraceleft}{\kern0pt}{\isacharparenleft}{\kern0pt}\isabound{q}{\isacharcomma}{\kern0pt}\ \isabound{a}{\isacharparenright}{\kern0pt}\ {\isacharbar}{\kern0pt}\ \isabound{q}\ {\isasymin}\ \isaconst{states}\ \ensuremath{LR_0(G)}\ {\isasymand}\ \isaconst{nxt}\ \ensuremath{LR_0(G)}\ \isabound{q}\ {\isacharparenleft}{\kern0pt}\isaconst{Tm}\ \isabound{a}{\isacharparenright}{\kern0pt}\ {\isasymnoteq}\ {\isasymemptyset}{\isacharbraceright}{\kern0pt}}. \]
It is easy to see that \isa{nxt} is a subset of \isa{\isafree{f}\ {\isacharbackquote}{\kern0pt}\ \isafree{T}}. Thus, it suffices to prove that \isa{T} is 
finite.

\isa{T} is the union of the sets \isa{{\isacharbraceleft}{\kern0pt}{\isacharparenleft}{\kern0pt}\isafree{q}{\isacharcomma}{\kern0pt}\ \isabound{a}{\isacharparenright}{\kern0pt}\ {\isacharbar}{\kern0pt}\ \isaconst{nxt}\ \ensuremath{LR_0(G)}\ \isafree{q}\ {\isacharparenleft}{\kern0pt}\isaconst{Tm}\ \isabound{a}{\isacharparenright}{\kern0pt}\ {\isasymnoteq}\ {\isasymemptyset}{\isacharbraceright}{\kern0pt}} for each
state \isa{q} of \isa{\ensuremath{LR_0(G)}}. We can then prove that each of these sets is finite by proving that for
any item set \isa{q}, the set \isa{{\isacharbraceleft}{\kern0pt}\isabound{X}\ {\isacharbar}{\kern0pt}\ \isaconst{nxt}\ \ensuremath{LR_0(G)}\ \isafree{q}\ \isabound{X}\ {\isasymnoteq}\ {\isasymemptyset}{\isacharbraceright}{\kern0pt}} is finite. This is because if for a 
symbol \isa{X} holds \isa{\isaconst{nxt}\ \ensuremath{LR_0(G)}\ \isafree{q}\ \isafree{X}\ {\isasymnoteq}\ {\isasymemptyset}}, there must exist an item in the \isa{{\isasymepsilon}}-closure of 
\isa{\isafree{q}\ {\isasyminter}\ \ensuremath{\mathrm{It}_{\isaconst{G{\isacharprime}{\kern0pt}}}}} of the form \isa{{\isacharbrackleft}{\kern0pt}\isafree{A}\ {\isasymrightarrow}\ \isafree{{\isasymalpha}}\ {\isasymcdot}\ \isafree{X}\ {\isacharhash}{\kern0pt}\ \isafree{{\isasymbeta}}{\isacharbrackright}{\kern0pt}}. Since the \isa{{\isasymepsilon}}-closure of \isa{\ensuremath{\mathrm{char}(G)}}
is finite for any set, the proof is complete.
\end{proof}
\end{proposition}

\begin{lemma}\label{finite_lists_complete_length_eq}
The set 
\begin{multline*}
\isa{L\ {\isacharequal}{\kern0pt}\ {\isacharbraceleft}{\kern0pt}q\ {\isacharhash}{\kern0pt}\ qs{\isacharbar}{\kern0pt}\ q\ qs\ X\ {\isasymalpha}{\isachardot}{\kern0pt}\ set\ {\isacharparenleft}{\kern0pt}q\ {\isacharhash}{\kern0pt}\ qs{\isacharparenright}{\kern0pt}\ {\isasymsubseteq}\ dfa{\isachardot}{\kern0pt}states\ LR\isactrlsub {\isadigit{0}}}\\ 
 \isa{{\isasymand}\ {\isacharbrackleft}{\kern0pt}X\ {\isasymrightarrow}\ {\isasymalpha}\ {\isasymcdot}\ {\isacharbrackright}{\kern0pt}\ {\isasymin}\ q\ {\isasymand}\ length\ {\isasymalpha}\ {\isacharequal}{\kern0pt}\ length\ qs\ {\isasymand}\ X\ {\isasymnoteq}\ S{\isacharprime}{\kern0pt}{\isacharbraceright}{\kern0pt}}
\end{multline*}
is finite.
\begin{proof}
We know \isa{L} is a subset of \isa{\isafree{f}\ {\isacharbackquote}{\kern0pt}\ \isafree{F}}, where
\begin{multline*}
F = \isa{f\ {\isacharbackquote}{\kern0pt}\ {\isacharbraceleft}{\kern0pt}{\isacharparenleft}{\kern0pt}q{\isacharcomma}{\kern0pt}\ qs{\isacharcomma}{\kern0pt}\ X{\isacharcomma}{\kern0pt}\ {\isasymalpha}{\isacharparenright}{\kern0pt}{\isacharbar}{\kern0pt}\ q\ qs\ X\ {\isasymalpha}{\isachardot}{\kern0pt}\ set\ {\isacharparenleft}{\kern0pt}q\ {\isacharhash}{\kern0pt}\ qs{\isacharparenright}{\kern0pt}\ {\isasymsubseteq}\ dfa{\isachardot}{\kern0pt}states\ LR\isactrlsub {\isadigit{0}}\ {\isasymand}\ {\isacharparenleft}{\kern0pt}{\isacharbrackleft}{\kern0pt}X\ {\isasymrightarrow}\ {\isasymalpha}\ {\isasymcdot}\ {\isacharbrackleft}{\kern0pt}{\isacharbrackright}{\kern0pt}{\isacharbrackright}{\kern0pt}\ {\isasymin}\ q{\isacharparenright}{\kern0pt}}\\
  \isa{{\isasymand}\ length\ {\isasymalpha}\ {\isacharequal}{\kern0pt}\ length\ qs\ {\isasymand}\ X\ {\isasymnoteq}\ S{\isacharprime}{\kern0pt}{\isacharbraceright}{\kern0pt}}
\end{multline*}
and \isa{\isafree{f}\ {\isacharparenleft}{\kern0pt}\isafree{q}{\isacharcomma}{\kern0pt}\ \isafree{qs}{\isacharcomma}{\kern0pt}\ \isafree{X}{\isacharcomma}{\kern0pt}\ \isafree{{\isasymalpha}}{\isacharparenright}{\kern0pt}\ {\isacharequal}{\kern0pt}\ \isafree{q}\ {\isacharhash}{\kern0pt}\ \isafree{qs}}. Moreover, it is easy to see that \isa{F} is finite: 
for a tuple \isa{{\isacharparenleft}{\kern0pt}\isafree{q}{\isacharcomma}{\kern0pt}\ \isafree{qs}{\isacharcomma}{\kern0pt}\ \isafree{X}{\isacharcomma}{\kern0pt}\ \isafree{{\isasymalpha}}{\isacharparenright}{\kern0pt}\ {\isasymin}\ \isafree{F}}, there are finitely many possibilities for \isa{q}, since it must
be a subset of \isa{\ensuremath{\mathrm{It}_{\isaconst{G{\isacharprime}{\kern0pt}}}}}. Furthermore, there are finitely many possible \isa{qs}, since such a list can
only be in a tuple of \isa{F} if there exists some complete item \isa{{\isacharbrackleft}{\kern0pt}\isafree{Y}\ {\isasymrightarrow}\ \isafree{{\isasymbeta}}\ {\isasymcdot}\ {\isacharbrackright}{\kern0pt}\ {\isasymin}\ \ensuremath{\mathrm{It}_{\isaconst{G{\isacharprime}{\kern0pt}}}}} with
\isa{{\isacharbar}{\kern0pt}\isafree{qs}{\isacharbar}{\kern0pt}\ {\isacharequal}{\kern0pt}\ {\isacharbar}{\kern0pt}\isafree{{\isasymbeta}}{\isacharbar}{\kern0pt}}. Since \isa{\isaconst{set}\ \isafree{qs}\ {\isasymsubseteq}\ \ensuremath{\mathrm{It}_{\isaconst{G{\isacharprime}{\kern0pt}}}}}, this means there are finitely many
\isa{qs} that can be in a tuple in \isa{F}. Finally, for \isa{X} and \isa{{\isasymalpha}} must hold \isa{{\isacharbrackleft}{\kern0pt}\isafree{X}\ {\isasymrightarrow}\ \isafree{{\isasymalpha}}\ {\isasymcdot}\ {\isacharbrackright}{\kern0pt}\ {\isasymin}\ \ensuremath{\mathrm{It}_{\isaconst{G{\isacharprime}{\kern0pt}}}}},
and we know there are only finitely many complete items in \isa{\ensuremath{\mathrm{It}_{\isaconst{G{\isacharprime}{\kern0pt}}}}} since \isa{\ensuremath{\mathrm{It}_{\isaconst{G{\isacharprime}{\kern0pt}}}}}
itself is finite.

Thus, \isa{L} is a subset of the image of a finite set, meaning it is also finite.
\end{proof}
\end{lemma}

\begin{lemma}\label{finite_eps_P0_reduce}
The set of reduce transitions \isa{R} of \isa{\ensuremath{P_0(G)}} is finite.
\begin{proof}
Consider set \isa{L} from Lemma~\ref{finite_lists_complete_length_eq}, which this Lemma tells us is finite.
It can easily be shown that \isa{\isafree{R}\ {\isasymsubseteq}\ \isafree{L}\ {\isasymtimes}\ {\isacharbraceleft}{\kern0pt}{\isacharbrackleft}{\kern0pt}\isabound{p}{\isacharcomma}{\kern0pt}\ \isabound{q}{\isacharbrackright}{\kern0pt}\ {\isacharbar}{\kern0pt}\ \isabound{p}\ {\isasymin}\ \isaconst{states}\ \ensuremath{LR_0(G)}\ {\isasymand}\ \isabound{q}\ {\isasymin}\ \isaconst{states}\ \ensuremath{LR_0(G)}{\isacharbraceright}{\kern0pt}}.
Moreover, the set \isa{{\isacharbraceleft}{\kern0pt}{\isacharbrackleft}{\kern0pt}\isabound{p}{\isacharcomma}{\kern0pt}\ \isabound{q}{\isacharbrackright}{\kern0pt}\ {\isacharbar}{\kern0pt}\ \isabound{p}\ {\isasymin}\ \isaconst{states}\ \ensuremath{LR_0(G)}\ {\isasymand}\ \isabound{q}\ {\isasymin}\ \isaconst{states}\ \ensuremath{LR_0(G)}{\isacharbraceright}{\kern0pt}} can trivially be
shown to be finite through a bijection with \isa{\isaconst{states}\ \ensuremath{LR_0(G)}\ {\isasymtimes}\ \isaconst{states}\ \ensuremath{LR_0(G)}},
completing the proof.
\end{proof}
\end{lemma}

\begin{proposition}\label{finite_eps_P0}
The \isa{\isaconst{eps}} relation of \isa{\ensuremath{P_0(G)}} is finite.
\begin{proof}
By definition, \isa{\isaconst{eps}} is the union of the reduce and finish relation sets. 

The set of finish transitions is trivially a subset of \isa{\isafree{f}\ {\isacharbackquote}{\kern0pt}\ \isaconst{states}\ \ensuremath{LR_0(G)}} 
with \isa{\isafree{f}\ \isafree{q}\ {\isacharequal}{\kern0pt}\ {\isacharparenleft}{\kern0pt}{\isacharbrackleft}{\kern0pt}\isafree{q}{\isacharcomma}{\kern0pt}\ \isaconst{init}\ \ensuremath{LR_0(G)}{\isacharbrackright}{\kern0pt}{\isacharcomma}{\kern0pt}\ {\isacharbrackleft}{\kern0pt}{\isacharbraceleft}{\kern0pt}{\isacharbrackleft}{\kern0pt}\isaconst{S{\isacharprime}{\kern0pt}}\ {\isasymrightarrow}\ {\isasymcdot}\ {\isacharbrackright}{\kern0pt}{\isacharbraceright}{\kern0pt}{\isacharbrackright}{\kern0pt}{\isacharparenright}{\kern0pt}}. The finiteness of \isa{\isaconst{eps}} then follows by
the finiteness of the states of \isa{\ensuremath{LR_0(G)}} and Lemma~\ref{finite_eps_P0_reduce}.
\end{proof}
\end{proposition}

With Propositions~\ref{finite_nxt_P0} and \ref{finite_eps_P0} we have shown a necessary condition
for the development of an executable transition function of the \isa{LR{\isacharparenleft}{\kern0pt}{\isadigit{0}}{\isacharparenright}{\kern0pt}} parser using the method we
described in the beginning of Section~\ref{sec:finite}.%
\end{isamarkuptext}\isamarkuptrue%
%
\isadelimdocument
%
\endisadelimdocument
%
\isatagdocument
%
\isamarkupsection{\isa{LR{\isacharparenleft}{\kern0pt}{\isadigit{0}}{\isacharparenright}{\kern0pt}}-Adequate and Inadequate States%
}
\isamarkuptrue%
%
\endisatagdocument
{\isafolddocument}%
%
\isadelimdocument
%
\endisadelimdocument
%
\begin{isamarkuptext}%
Even though our modified \isa{\ensuremath{P_0(G)}} definition circumvents the issue between reduce and finish 
transitions we described before, nondeterminism is still a problem in many other cases. We 
define two types of \concept{conflicts} that can be present in parser states which can lead to 
nondeterministic behavior, as defined by Wilhelm et al.~\cite[p. 110]{Wilhelm}:

\begin{definition}[Shift-reduce conflicts, reduce-reduce conflicts, and \isa{LR{\isacharparenleft}{\kern0pt}{\isadigit{0}}{\isacharparenright}{\kern0pt}} inadequacy]
Let \isa{q} be a state of \isa{\ensuremath{P_0(G)}}. We say \isa{q} has a \concept{shift-reduce conflict} if it allows for 
\isa{\ensuremath{P_0(G)}} to make both a shift and a reduce transition.

Moreover, \isa{q} is said to have a \concept{reduce-reduce conflict} if it allows performing reducing
transitions for two distinct productions.

If \isa{q} has either of these conflicts, it is \concept{\isa{LR{\isacharparenleft}{\kern0pt}{\isadigit{0}}{\isacharparenright}{\kern0pt}}-inadequate}. Otherwise, it is 
\concept{\isa{LR{\isacharparenleft}{\kern0pt}{\isadigit{0}}{\isacharparenright}{\kern0pt}}-adequate}.
\end{definition}

It is worth noting that in this context, we consider finishing transitions a special case of reducing
transitions. Therefore, for a state to permit a reducing transition for a production \isa{{\isacharparenleft}{\kern0pt}\isafree{A}{\isacharcomma}{\kern0pt}\ \isafree{{\isasymalpha}}{\isacharparenright}{\kern0pt}},
it suffices for it to contain the complete item \isa{{\isacharbrackleft}{\kern0pt}\isafree{A}\ {\isasymrightarrow}\ \isafree{{\isasymalpha}}\ {\isasymcdot}\ {\isacharbrackright}{\kern0pt}}.

Because of this, we formally define \isa{q} to have a reduce-reduce conflict if 
\isa{\mbox{{\isadigit{1}}\ {\isacharless}{\kern0pt}\ {\isacharbar}{\kern0pt}\isaconst{completes}\ \isafree{q}{\isacharbar}{\kern0pt}}}. Meanwhile, there is a shift-reduce conflict if there exist two 
tuples \isa{{\isacharparenleft}{\kern0pt}{\isacharbrackleft}{\kern0pt}\isafree{q}{\isacharbrackright}{\kern0pt}{\isacharcomma}{\kern0pt}\ \isafree{a}{\isacharcomma}{\kern0pt}\ \isafree{rs\isactrlsub {\isadigit{1}}}{\isacharparenright}{\kern0pt}\ {\isasymin}\ \isaconst{nxt}\ \ensuremath{P_0(G)}} and \isa{{\isacharparenleft}{\kern0pt}\isafree{q}\ {\isacharhash}{\kern0pt}\ \isafree{qs}{\isacharcomma}{\kern0pt}\ \isafree{rs\isactrlsub {\isadigit{2}}}{\isacharparenright}{\kern0pt}\ {\isasymin}\ \isaconst{eps}\ \ensuremath{P_0(G)}}. This is
equivalent to \isa{q} containing at least one complete item, and an incomplete item of the form 
\isa{{\isacharbrackleft}{\kern0pt}\isafree{X}\ {\isasymrightarrow}\ \isafree{{\isasymalpha}}\ {\isasymcdot}\ \isaconst{Tm}\ \isafree{a}\ {\isacharhash}{\kern0pt}\ \isafree{{\isasymbeta}}{\isacharbrackright}{\kern0pt}}.

We will now show a lemma that Wilhelm et al.\ present without proof, and which we will need in the
future to show the final theorem of the \isa{LR{\isacharparenleft}{\kern0pt}{\isadigit{0}}{\isacharparenright}{\kern0pt}} section in the book. First, we prove an auxiliary
lemma:

\begin{lemma}\label{derivern_non_word_imp_hd_eps_reachable}
If \isa{\isaconst{P{\isacharprime}{\kern0pt}}\ {\isasymturnstile}\ {\isacharbrackleft}{\kern0pt}\isaconst{Nt}\ \isafree{A}{\isacharbrackright}{\kern0pt}\ {\isasymRightarrow}r{\isacharparenleft}{\kern0pt}\isaconst{Suc}\ \isafree{n}{\isacharparenright}{\kern0pt}\ \isafree{Y}\ {\isacharhash}{\kern0pt}\ \isafree{{\isasymgamma}}} and \isa{Y\ {\isacharhash}{\kern0pt}\ {\isasymgamma}} is not a terminal word, then there exist
\isa{X\ {\isacharcolon}{\kern0pt}{\isacharcolon}{\kern0pt}\ {\isacharprime}{\kern0pt}n} and \isa{{\isasymalpha}{\isacharcomma}{\kern0pt}\ {\isasymbeta}\ {\isacharcolon}{\kern0pt}{\isacharcolon}{\kern0pt}\ syms} such that \mbox{\isa{{\isacharbrackleft}{\kern0pt}\isafree{A}\ {\isasymrightarrow}\ \ {\isasymcdot}\ \isafree{{\isasymalpha}}\ {\isacharbrackright}{\kern0pt}\ {\isasymin}\ \ensuremath{\mathrm{It}_{\isaconst{G{\isacharprime}{\kern0pt}}}}}} and 
\isa{{\isacharparenleft}{\kern0pt}{\isacharbrackleft}{\kern0pt}\isafree{A}\ {\isasymrightarrow}\ \ {\isasymcdot}\ \isafree{{\isasymalpha}}\ {\isacharbrackright}{\kern0pt}{\isacharcomma}{\kern0pt}\ {\isacharbrackleft}{\kern0pt}\isafree{X}\ {\isasymrightarrow}\ \ {\isasymcdot}\ \isafree{Y}\ {\isacharhash}{\kern0pt}\ \isafree{{\isasymbeta}}\ {\isacharbrackright}{\kern0pt}{\isacharparenright}{\kern0pt}\ {\isasymin}\ {\isacharparenleft}{\kern0pt}\isaconst{eps}\ \ensuremath{\mathrm{char}(G)}{\isacharparenright}{\kern0pt}\isactrlsup {\isacharasterisk}{\kern0pt}}.
\begin{proof}
We induct on \isa{n} for arbitrary \isa{Y} and \isa{{\isasymgamma}}.

If \isa{n\ {\isacharequal}{\kern0pt}\ {\isadigit{0}}}, then \isa{{\isacharparenleft}{\kern0pt}\isafree{A}{\isacharcomma}{\kern0pt}\ \isafree{Y}\ {\isacharhash}{\kern0pt}\ \isafree{{\isasymgamma}}{\isacharparenright}{\kern0pt}\ {\isasymin}\ \isaconst{P{\isacharprime}{\kern0pt}}}, and the implication holds trivially by reflexivity.

If \isa{n\ {\isachargreater}{\kern0pt}\ {\isadigit{0}}}, we consider the final step in the derivation. The sentential form that derives 
\isa{\isafree{Y}\ {\isacharhash}{\kern0pt}\ \isafree{{\isasymgamma}}} in the last step must be nonempty and contain at least one nonterminal symbol. Let
\isa{\isafree{X}\ {\isacharhash}{\kern0pt}\ \isafree{{\isasymdelta}}} be this sentential form. By the induction hypothesis, there exist
\isa{{\isasymalpha}{\isacharcomma}{\kern0pt}\ {\isasymbeta}\ {\isacharcolon}{\kern0pt}{\isacharcolon}{\kern0pt}\ syms} and \isa{W\ {\isacharcolon}{\kern0pt}{\isacharcolon}{\kern0pt}\ {\isacharprime}{\kern0pt}n} such that \mbox{\isa{{\isacharbrackleft}{\kern0pt}\isafree{A}\ {\isasymrightarrow}\ \ {\isasymcdot}\ \isafree{{\isasymalpha}}\ {\isacharbrackright}{\kern0pt}\ {\isasymin}\ \ensuremath{\mathrm{It}_{\isaconst{G{\isacharprime}{\kern0pt}}}}}} and 
\begin{equation}\label{der_hd_eps.ih}
\isa{{\isacharparenleft}{\kern0pt}{\isacharbrackleft}{\kern0pt}\isafree{A}\ {\isasymrightarrow}\ \ {\isasymcdot}\ \isafree{{\isasymalpha}}\ {\isacharbrackright}{\kern0pt}{\isacharcomma}{\kern0pt}\ {\isacharbrackleft}{\kern0pt}\isafree{W}\ {\isasymrightarrow}\ \ {\isasymcdot}\ \isafree{X}\ {\isacharhash}{\kern0pt}\ \isafree{{\isasymbeta}}\ {\isacharbrackright}{\kern0pt}{\isacharparenright}{\kern0pt}\ {\isasymin}\ {\isacharparenleft}{\kern0pt}\isaconst{eps}\ \ensuremath{\mathrm{char}(G)}{\isacharparenright}{\kern0pt}\isactrlsup {\isacharasterisk}{\kern0pt}}. 
\end{equation}
We can now distinguish cases based on whether \isa{{\isasymdelta}} consists exclusively of terminals.

If \isa{{\isasymdelta}} contains no nonterminals, there exists an \isa{X{\isacharprime}{\kern0pt}\ {\isacharcolon}{\kern0pt}{\isacharcolon}{\kern0pt}\ {\isacharprime}{\kern0pt}n} such that \isa{\isafree{X}\ {\isacharequal}{\kern0pt}\ \isaconst{Nt}\ \isafree{X{\isacharprime}{\kern0pt}}}. Furthermore, 
there also exists a production \isa{{\isacharparenleft}{\kern0pt}\isafree{X{\isacharprime}{\kern0pt}}{\isacharcomma}{\kern0pt}\ \isafree{{\isasymchi}}{\isacharparenright}{\kern0pt}\ {\isasymin}\ \isaconst{P{\isacharprime}{\kern0pt}}} which was applied in the last derivation 
step was applied, meaning \isa{\isafree{{\isasymchi}}\ {\isacharat}{\kern0pt}\ \isafree{{\isasymdelta}}\ {\isacharequal}{\kern0pt}\ \isafree{Y}\ {\isacharhash}{\kern0pt}\ \isafree{{\isasymgamma}}}. This implies that \isa{{\isasymchi}} cannot be the empty word; 
if this were the case, \isa{\isafree{Y}\ {\isacharhash}{\kern0pt}\ \isafree{{\isasymgamma}}} would contain no nonterminals, which contradicts our assumption 
that it is not a terminal word. Therefore, \isa{\isafree{{\isasymchi}}\ {\isacharequal}{\kern0pt}\ \isafree{Y}\ {\isacharhash}{\kern0pt}\ \isafree{{\isasymzeta}}} for some \isa{{\isasymzeta}\ {\isacharcolon}{\kern0pt}{\isacharcolon}{\kern0pt}\ syms}. From all this 
follows \isa{{\isacharparenleft}{\kern0pt}{\isacharbrackleft}{\kern0pt}\isafree{W}\ {\isasymrightarrow}\ \ {\isasymcdot}\ \isaconst{Nt}\ \isafree{X{\isacharprime}{\kern0pt}}\ {\isacharhash}{\kern0pt}\ \isafree{{\isasymbeta}}\ {\isacharbrackright}{\kern0pt}{\isacharcomma}{\kern0pt}\ {\isacharbrackleft}{\kern0pt}\isafree{X{\isacharprime}{\kern0pt}}\ {\isasymrightarrow}\ \ {\isasymcdot}\ \isafree{Y}\ {\isacharhash}{\kern0pt}\ \isafree{{\isasymzeta}}\ {\isacharbrackright}{\kern0pt}{\isacharparenright}{\kern0pt}\ {\isasymin}\ \mbox{\isaconst{eps}\ \ensuremath{\mathrm{char}(G)}}}.
The implication then holds by \eqref{der_hd_eps.ih}.

On the other hand, if \isa{{\isasymdelta}} contains at least a nonterminal, we know that \isa{\isafree{X}\ {\isacharequal}{\kern0pt}\ \isafree{Y}} since there
exists a nonterminal to the right of \isa{X}, meaning that after the final rightmost derivation step,
it remains unchanged. This completes the proof by \eqref{der_hd_eps.ih}.
\end{proof}
\end{lemma}

\begin{lemma}\label{LR0_adequate_cases}
If a state \isa{q} is \isa{LR{\isacharparenleft}{\kern0pt}{\isadigit{0}}{\isacharparenright}{\kern0pt}}-adequate, one of the following holds:
\begin{enumerate}
\item \isa{\isaconst{completes}\ \isafree{q}\ {\isacharequal}{\kern0pt}\ {\isasymemptyset}}
\item \isa{\isafree{q}\ {\isacharequal}{\kern0pt}\ {\isacharbraceleft}{\kern0pt}{\isacharbrackleft}{\kern0pt}\isafree{A}\ {\isasymrightarrow}\ \isafree{{\isasymalpha}}\ {\isasymcdot}\ {\isacharbrackright}{\kern0pt}{\isacharbraceright}{\kern0pt}} for some \isa{A} and \isa{{\isasymalpha}}.
\item \isa{completes\ q\ {\isacharequal}{\kern0pt}\ {\isacharbraceleft}{\kern0pt}{\isacharbrackleft}{\kern0pt}A\ {\isasymrightarrow}\ {\isasymcdot}\ {\isacharbrackright}{\kern0pt}{\isacharbraceright}{\kern0pt}} and for every incomplete item \isa{i\ {\isasymin}\ q}, 
there exist \isa{X{\isacharcomma}{\kern0pt}\ Y\ {\isacharcolon}{\kern0pt}{\isacharcolon}{\kern0pt}\ {\isacharprime}{\kern0pt}n} and \isa{{\isasymalpha}{\isacharcomma}{\kern0pt}\ {\isasymbeta}\ {\isacharcolon}{\kern0pt}{\isacharcolon}{\kern0pt}\ syms} such that \isa{\isafree{i}\ {\isacharequal}{\kern0pt}\ {\isacharbrackleft}{\kern0pt}\isafree{X}\ {\isasymrightarrow}\ \isafree{{\isasymalpha}}\ {\isasymcdot}\ \isaconst{Nt}\ \isafree{Y}\ {\isacharhash}{\kern0pt}\ \isafree{{\isasymbeta}}{\isacharbrackright}{\kern0pt}} 
and every rightmost derivation of a word \isa{w} 
\[ \isa{P{\isacharprime}{\kern0pt}\ {\isasymturnstile}\ {\isacharbrackleft}{\kern0pt}Nt\ Y{\isacharbrackright}{\kern0pt}\ {\isasymRightarrow}r{\isacharasterisk}{\kern0pt}\ {\isasymgamma}\ {\isasymRightarrow}r\ map\ Tm\ w} \]
implies \isa{\isafree{{\isasymgamma}}\ {\isacharequal}{\kern0pt}\ \isaconst{Nt}\ \isafree{A}\ {\isacharhash}{\kern0pt}\ \isaconst{map}\ \isaconst{Tm}\ \isafree{w}}.
\end{enumerate}
\begin{proof}
Since \isa{q} is a state of \isa{\ensuremath{P_0(G)}}, it is either the final state \isa{{\isacharbraceleft}{\kern0pt}{\isacharbrackleft}{\kern0pt}\isaconst{S{\isacharprime}{\kern0pt}}\ {\isasymrightarrow}\ {\isasymcdot}\ {\isacharbrackright}{\kern0pt}{\isacharbraceright}{\kern0pt}} or a state of 
\isa{\ensuremath{LR_0(G)}}. If \isa{\isafree{q}\ {\isacharequal}{\kern0pt}\ {\isacharbraceleft}{\kern0pt}{\isacharbrackleft}{\kern0pt}\isaconst{S{\isacharprime}{\kern0pt}}\ {\isasymrightarrow}\ {\isasymcdot}\ {\isacharbrackright}{\kern0pt}{\isacharbraceright}{\kern0pt}}, it is trivially \isa{LR{\isacharparenleft}{\kern0pt}{\isadigit{0}}{\isacharparenright}{\kern0pt}}-adequate and fulfills (2).

Now, we assume \isa{q} is a state of \isa{\ensuremath{LR_0(G)}}. Since \isa{q} is \isa{LR{\isacharparenleft}{\kern0pt}{\isadigit{0}}{\isacharparenright}{\kern0pt}}-adequate, we can distinguish two cases:
either \isa{q} has no complete items, or exactly one complete item. If \isa{q} has no complete items, it 
fulfills condition (1).

Otherwise, 
\begin{equation}\label{ad_cases.compq}
\isa{\isaconst{completes}\ \isafree{q}\ {\isacharequal}{\kern0pt}\ {\isacharbraceleft}{\kern0pt}{\isacharbrackleft}{\kern0pt}\isafree{A}\ {\isasymrightarrow}\ \isafree{{\isasymalpha}}\ {\isasymcdot}\ {\isacharbrackright}{\kern0pt}{\isacharbraceright}{\kern0pt}}
\end{equation} 
for some \isa{A} and \isa{{\isasymalpha}}. We now consider another two cases on whether \isa{q} has any incomplete items;
if it has none, it fulfills condition (2).

Otherwise, let \isa{{\isacharbrackleft}{\kern0pt}\isafree{X}\ {\isasymrightarrow}\ \isafree{{\isasymalpha}{\isacharprime}{\kern0pt}}\ {\isasymcdot}\ \isafree{Z}\ {\isacharhash}{\kern0pt}\ \isafree{{\isasymbeta}}{\isacharbrackright}{\kern0pt}\ {\isasymin}\ \isafree{q}}. We know that \isa{Z} must be a nonterminal, since if it
were some terminal symbol \isa{Tm\ a}, the reading transition 
\isa{\mbox{{\isacharparenleft}{\kern0pt}{\isacharbrackleft}{\kern0pt}\isafree{q}{\isacharbrackright}{\kern0pt}{\isacharcomma}{\kern0pt}\ \isafree{a}{\isacharcomma}{\kern0pt}\ {\isacharbrackleft}{\kern0pt}\isaconst{nxt}\ \ensuremath{LR_0(G)}\ \isafree{q}\ {\isacharparenleft}{\kern0pt}\isaconst{Tm}\ \isafree{a}{\isacharparenright}{\kern0pt}{\isacharcomma}{\kern0pt}\ \isafree{q}{\isacharbrackright}{\kern0pt}{\isacharparenright}{\kern0pt}}} would be possible, and since \isa{q} also contains a 
complete item, a shift-reduce conflict would be present, contradicting the premise that \isa{q} is 
\isa{LR{\isacharparenleft}{\kern0pt}{\isadigit{0}}{\isacharparenright}{\kern0pt}}-adequate. Thus, let \isa{\isafree{Z}\ {\isacharequal}{\kern0pt}\ \isaconst{Nt}\ \isafree{Y}}.

Since \isa{G{\isacharprime}{\kern0pt}} is reduced, we know that there must exist a rightmost derivation of a terminal word \isa{w}
of the form \isa{P{\isacharprime}{\kern0pt}\ {\isasymturnstile}\ {\isacharbrackleft}{\kern0pt}Nt\ Y{\isacharbrackright}{\kern0pt}\ {\isasymRightarrow}r{\isacharasterisk}{\kern0pt}\ {\isasymgamma}\ {\isasymRightarrow}r\ map\ Tm\ w} for some \isa{{\isasymgamma}\ {\isacharcolon}{\kern0pt}{\isacharcolon}{\kern0pt}\ syms}. To complete our proof, 
we need to show that \isa{\mbox{\isafree{{\isasymgamma}}\ {\isacharequal}{\kern0pt}\ \isaconst{Nt}\ \isafree{A}\ {\isacharhash}{\kern0pt}\ \isaconst{map}\ \isaconst{Tm}\ \isafree{w}}} and \isa{\isafree{{\isasymalpha}}\ {\isacharequal}{\kern0pt}\ {\isacharbrackleft}{\kern0pt}{\isacharbrackright}{\kern0pt}}. We begin by doing a 
case distinction on the reflexive transitive closure of the derivation from \isa{Y} to \isa{{\isasymgamma}}:

In the reflexive case, \isa{Y} derives \isa{\isaconst{map}\ \isaconst{Tm}\ \isafree{w}} in a single step, meaning 
\isa{{\isacharparenleft}{\kern0pt}\isafree{Y}{\isacharcomma}{\kern0pt}\ \isaconst{map}\ \isaconst{Tm}\ \isafree{w}{\isacharparenright}{\kern0pt}\ {\isasymin}\ \isaconst{P{\isacharprime}{\kern0pt}}}.
Therefore, \isa{{\isacharparenleft}{\kern0pt}{\isacharbrackleft}{\kern0pt}\isafree{X}\ {\isasymrightarrow}\ \isafree{{\isasymalpha}{\isacharprime}{\kern0pt}}\ {\isasymcdot}\ \isaconst{Nt}\ \isafree{Y}\ {\isacharhash}{\kern0pt}\ \isafree{{\isasymbeta}}{\isacharbrackright}{\kern0pt}{\isacharcomma}{\kern0pt}\ {\isacharbrackleft}{\kern0pt}\isafree{Y}\ {\isasymrightarrow}\ \ {\isasymcdot}\ \isaconst{map}\ \isaconst{Tm}\ \isafree{w}\ {\isacharbrackright}{\kern0pt}{\isacharparenright}{\kern0pt}} is in the \isa{{\isasymepsilon}}-transition relation of 
\isa{\ensuremath{\mathrm{char}(G)}}; this and \isa{{\isacharbrackleft}{\kern0pt}\isafree{X}\ {\isasymrightarrow}\ \isafree{{\isasymalpha}{\isacharprime}{\kern0pt}}\ {\isasymcdot}\ \isaconst{Nt}\ \isafree{Y}\ {\isacharhash}{\kern0pt}\ \isafree{{\isasymbeta}}{\isacharbrackright}{\kern0pt}\ {\isasymin}\ \isafree{q}} imply that 
\isa{\mbox{{\isacharbrackleft}{\kern0pt}\isafree{Y}\ {\isasymrightarrow}\ \ {\isasymcdot}\ \isaconst{map}\ \isaconst{Tm}\ \isafree{w}\ {\isacharbrackright}{\kern0pt}\ {\isasymin}\ \isafree{q}}} as well. We can now perform a case distinction on \isa{w}: if 
\isa{\isafree{w}\ {\isacharequal}{\kern0pt}\ {\isacharbrackleft}{\kern0pt}{\isacharbrackright}{\kern0pt}}, the implication holds since \isa{{\isacharbrackleft}{\kern0pt}\isafree{Y}\ {\isasymrightarrow}\ \ {\isasymcdot}\ \isaconst{map}\ \isaconst{Tm}\ \isafree{w}\ {\isacharbrackright}{\kern0pt}\ {\isacharequal}{\kern0pt}\ {\isacharbrackleft}{\kern0pt}\isafree{A}\ {\isasymrightarrow}\ \isafree{{\isasymalpha}}\ {\isasymcdot}\ {\isacharbrackright}{\kern0pt}} and
\isa{{\isacharbrackleft}{\kern0pt}\isaconst{Nt}\ \isafree{Y}{\isacharbrackright}{\kern0pt}\ {\isacharequal}{\kern0pt}\ \isafree{{\isasymgamma}}}. If on the other hand \isa{w} is nonempty, \isa{q} has a shift-reduce conflict, which
contradicts its \isa{LR{\isacharparenleft}{\kern0pt}{\isadigit{0}}{\isacharparenright}{\kern0pt}}-adequacy. This completes the reflexive case, and we can now move on to the 
other case.

For the transitive case, we know \isa{Y} derives \isa{{\isasymgamma}} in \isa{\isaconst{Suc}\ \isafree{n}} steps for some \isa{n\ {\isacharcolon}{\kern0pt}{\isacharcolon}{\kern0pt}\ nat}. 
Since \isa{{\isasymgamma}} derives \isa{map\ Tm\ w} in the final step, it must be nonempty and contain at least one 
nonterminal. Therefore, let 
\begin{equation}\label{ad_cases.gWd}
\isa{\isafree{{\isasymgamma}}\ {\isacharequal}{\kern0pt}\ \isafree{W}\ {\isacharhash}{\kern0pt}\ \isafree{{\isasymdelta}}}
\end{equation}
 for some \isa{W\ {\isacharcolon}{\kern0pt}{\isacharcolon}{\kern0pt}\ sym} and \isa{{\isasymdelta}\ {\isacharcolon}{\kern0pt}{\isacharcolon}{\kern0pt}\ syms}. From all this, 
Lemma~\ref{derivern_non_word_imp_hd_eps_reachable} tells us there exist \isa{{\isasymgamma}{\isacharprime}{\kern0pt}{\isacharcomma}{\kern0pt}\ {\isasymzeta}\ {\isacharcolon}{\kern0pt}{\isacharcolon}{\kern0pt}\ syms} and 
\isa{B\ {\isacharcolon}{\kern0pt}{\isacharcolon}{\kern0pt}\ {\isacharprime}{\kern0pt}n} such that \isa{{\isacharbrackleft}{\kern0pt}\isafree{Y}\ {\isasymrightarrow}\ \ {\isasymcdot}\ \isafree{{\isasymgamma}{\isacharprime}{\kern0pt}}\ {\isacharbrackright}{\kern0pt}\ {\isasymin}\ \ensuremath{\mathrm{It}_{\isaconst{G{\isacharprime}{\kern0pt}}}}} and
\begin{equation}\label{ad_cases.YB_eps} 
\isa{{\isacharparenleft}{\kern0pt}{\isacharbrackleft}{\kern0pt}\isafree{Y}\ {\isasymrightarrow}\ \ {\isasymcdot}\ \isafree{{\isasymgamma}{\isacharprime}{\kern0pt}}\ {\isacharbrackright}{\kern0pt}{\isacharcomma}{\kern0pt}\ {\isacharbrackleft}{\kern0pt}\isafree{B}\ {\isasymrightarrow}\ \ {\isasymcdot}\ \isafree{W}\ {\isacharhash}{\kern0pt}\ \isafree{{\isasymzeta}}\ {\isacharbrackright}{\kern0pt}{\isacharparenright}{\kern0pt}\ {\isasymin}\ {\isacharparenleft}{\kern0pt}\isaconst{eps}\ \ensuremath{\mathrm{char}(G)}{\isacharparenright}{\kern0pt}\isactrlsup {\isacharasterisk}{\kern0pt}}. 
\end{equation}
Moreover, by contradiction we know that \isa{W} must be a nonterminal: if this is not the case, 
\isa{\isafree{W}\ {\isacharequal}{\kern0pt}\ \isaconst{Tm}\ \isafree{a}} for some \isa{a\ {\isacharcolon}{\kern0pt}{\isacharcolon}{\kern0pt}\ {\isacharprime}{\kern0pt}t}. By transitivity, this along with 
\isa{{\isacharbrackleft}{\kern0pt}\isafree{X}\ {\isasymrightarrow}\ \isafree{{\isasymalpha}{\isacharprime}{\kern0pt}}\ {\isasymcdot}\ \isafree{Z}\ {\isacharhash}{\kern0pt}\ \isafree{{\isasymbeta}}{\isacharbrackright}{\kern0pt}\ {\isasymin}\ \isafree{q}}, \isa{\isafree{Z}\ {\isacharequal}{\kern0pt}\ \isaconst{Nt}\ \isafree{Y}}, and \eqref{ad_cases.YB_eps} implies that 
\isa{{\isacharbrackleft}{\kern0pt}\isafree{B}\ {\isasymrightarrow}\ \ {\isasymcdot}\ \isaconst{Tm}\ \isafree{a}\ {\isacharhash}{\kern0pt}\ \isafree{{\isasymzeta}}\ {\isacharbrackright}{\kern0pt}\ {\isasymin}\ \isafree{q}}, permitting a reading transition on state \isa{q}. Since 
\isa{\isaconst{completes}\ \isafree{q}} is nonempty, such a reading transition would make \isa{q} \isa{LR{\isacharparenleft}{\kern0pt}{\isadigit{0}}{\isacharparenright}{\kern0pt}}-inadequate, 
leading to a contradiction to our original assumption. Therefore, let \isa{\isafree{W}\ {\isacharequal}{\kern0pt}\ \isaconst{Nt}\ \isafree{C}} for some
\isa{C\ {\isacharcolon}{\kern0pt}{\isacharcolon}{\kern0pt}\ {\isacharprime}{\kern0pt}n}. Then the last step of the rightmost derivation must consist of some production 
\isa{{\isacharparenleft}{\kern0pt}\isafree{C}{\isacharcomma}{\kern0pt}\ \isaconst{map}\ \isaconst{Tm}\ \isafree{u}{\isacharparenright}{\kern0pt}\ {\isasymin}\ \isaconst{P{\isacharprime}{\kern0pt}}} for some terminal word \isa{u} with 
\begin{equation}\label{ad_cases.udw}
\isa{\isaconst{map}\ \isaconst{Tm}\ \isafree{u}\ {\isacharat}{\kern0pt}\ \isafree{{\isasymdelta}}\ {\isacharequal}{\kern0pt}\ \isaconst{map}\ \isaconst{Tm}\ \isafree{w}}.
\end{equation}
This implies that \isa{{\isacharparenleft}{\kern0pt}{\isacharbrackleft}{\kern0pt}\isafree{B}\ {\isasymrightarrow}\ \ {\isasymcdot}\ \isafree{W}\ {\isacharhash}{\kern0pt}\ \isafree{{\isasymzeta}}\ {\isacharbrackright}{\kern0pt}{\isacharcomma}{\kern0pt}\ {\isacharbrackleft}{\kern0pt}\isafree{C}\ {\isasymrightarrow}\ \ {\isasymcdot}\ \isaconst{map}\ \isaconst{Tm}\ \isafree{u}\ {\isacharbrackright}{\kern0pt}{\isacharparenright}{\kern0pt}\ {\isasymin}\ \isaconst{eps}\ \ensuremath{\mathrm{char}(G)}}, and 
once again with \isa{{\isacharbrackleft}{\kern0pt}\isafree{X}\ {\isasymrightarrow}\ \isafree{{\isasymalpha}{\isacharprime}{\kern0pt}}\ {\isasymcdot}\ \isafree{Z}\ {\isacharhash}{\kern0pt}\ \isafree{{\isasymbeta}}{\isacharbrackright}{\kern0pt}\ {\isasymin}\ \isafree{q}}, \isa{\isafree{Z}\ {\isacharequal}{\kern0pt}\ \isaconst{Nt}\ \isafree{Y}}, and \eqref{ad_cases.YB_eps}, 
this implies that \isa{{\isacharparenleft}{\kern0pt}{\isacharbrackleft}{\kern0pt}\isafree{X}\ {\isasymrightarrow}\ \isafree{{\isasymalpha}{\isacharprime}{\kern0pt}}\ {\isasymcdot}\ \isafree{Z}\ {\isacharhash}{\kern0pt}\ \isafree{{\isasymbeta}}{\isacharbrackright}{\kern0pt}{\isacharcomma}{\kern0pt}\ {\isacharbrackleft}{\kern0pt}\isafree{C}\ {\isasymrightarrow}\ \ {\isasymcdot}\ \isaconst{map}\ \isaconst{Tm}\ \isafree{u}\ {\isacharbrackright}{\kern0pt}{\isacharparenright}{\kern0pt}\ {\isasymin}\ {\isacharparenleft}{\kern0pt}\isaconst{eps}\ \ensuremath{\mathrm{char}(G)}{\isacharparenright}{\kern0pt}\isactrlsup {\isacharasterisk}{\kern0pt}} by transitivity.
This means that the item \isa{{\isacharbrackleft}{\kern0pt}\isafree{C}\ {\isasymrightarrow}\ \ {\isasymcdot}\ \isaconst{map}\ \isaconst{Tm}\ \isafree{u}\ {\isacharbrackright}{\kern0pt}} is also in \isa{q}. This and \eqref{ad_cases.compq} 
imply 
\begin{equation}\label{C_eq_comp}
\isa{{\isacharbrackleft}{\kern0pt}C\ {\isasymrightarrow}\ {\isacharbrackleft}{\kern0pt}{\isacharbrackright}{\kern0pt}\ {\isasymcdot}\ map\ Tm\ u{\isacharbrackright}{\kern0pt}\ {\isacharequal}{\kern0pt}\ {\isacharbrackleft}{\kern0pt}A\ {\isasymrightarrow}\ {\isasymalpha}\ {\isasymcdot}\ {\isacharbrackleft}{\kern0pt}{\isacharbrackright}{\kern0pt}{\isacharbrackright}{\kern0pt}}.
\end{equation}
This follows from the fact that \isa{{\isacharbrackleft}{\kern0pt}\isafree{C}\ {\isasymrightarrow}\ \ {\isasymcdot}\ \isaconst{map}\ \isaconst{Tm}\ \isafree{u}\ {\isacharbrackright}{\kern0pt}} can only be an incomplete item if \isa{u} is 
nonempty, but this would allow a reading transition to take place. As we already showed, this would 
lead to a contradiction. Therefore, by \eqref{ad_cases.gWd}, \eqref{ad_cases.udw}, and 
\eqref{C_eq_comp}, the criteria for condition (3) of our original claim are fulfilled. 
\end{proof}
\end{lemma}%
\end{isamarkuptext}\isamarkuptrue%
%
\isadelimdocument
%
\endisadelimdocument
%
\isatagdocument
%
\isamarkupsection{\isa{LR{\isacharparenleft}{\kern0pt}k{\isacharparenright}{\kern0pt}} Grammars%
}
\isamarkuptrue%
%
\endisatagdocument
{\isafolddocument}%
%
\isadelimdocument
%
\endisadelimdocument
%
\begin{isamarkuptext}%
\begin{definition}[\isa{LR{\isacharparenleft}{\kern0pt}k{\isacharparenright}{\kern0pt}} Grammar]
\isa{G{\isacharprime}{\kern0pt}} is an \concept{\isa{LR{\isacharparenleft}{\kern0pt}k{\isacharparenright}{\kern0pt}} grammar} if for any \isa{{\isasymalpha}{\isacharcomma}{\kern0pt}\ {\isasymbeta}{\isacharcomma}{\kern0pt}\ {\isasymgamma}\ {\isacharcolon}{\kern0pt}{\isacharcolon}{\kern0pt}\ syms}, \isa{X{\isacharcomma}{\kern0pt}\ Y\ {\isacharcolon}{\kern0pt}{\isacharcolon}{\kern0pt}\ {\isacharprime}{\kern0pt}n}, and 
\isa{w{\isacharcomma}{\kern0pt}\ x{\isacharcomma}{\kern0pt}\ y\ {\isacharcolon}{\kern0pt}{\isacharcolon}{\kern0pt}\ {\isacharprime}{\kern0pt}t\ list},
\begin{gather*}
\begin{multlined}
\isa{P{\isacharprime}{\kern0pt}\ {\isasymturnstile}\ {\isacharbrackleft}{\kern0pt}Nt\ S{\isacharprime}{\kern0pt}{\isacharbrackright}{\kern0pt}\ {\isasymRightarrow}r{\isacharasterisk}{\kern0pt}\ {\isasymalpha}\ {\isacharat}{\kern0pt}\ Nt\ X\ {\isacharhash}{\kern0pt}\ map\ Tm\ w\ {\isasymRightarrow}r\ {\isasymalpha}\ {\isacharat}{\kern0pt}\ {\isasymbeta}\ {\isacharat}{\kern0pt}\ map\ Tm\ w}
\end{multlined}
\intertext{and}
\begin{multlined}
\isa{P{\isacharprime}{\kern0pt}\ {\isasymturnstile}\ {\isacharbrackleft}{\kern0pt}Nt\ S{\isacharprime}{\kern0pt}{\isacharbrackright}{\kern0pt}\ {\isasymRightarrow}r{\isacharasterisk}{\kern0pt}\ {\isasymgamma}\ {\isacharat}{\kern0pt}\ Nt\ Y\ {\isacharhash}{\kern0pt}\ map\ Tm\ x\ {\isasymRightarrow}r\ {\isasymalpha}\ {\isacharat}{\kern0pt}\ {\isasymbeta}\ {\isacharat}{\kern0pt}\ map\ Tm\ y}
\end{multlined}
\end{gather*}
and \isa{\isafree{w}{\isacharbar}{\kern0pt}\isactrlbsub \isafree{k}\isactrlesub \ {\isacharequal}{\kern0pt}\ \isafree{y}{\isacharbar}{\kern0pt}\isactrlbsub \isafree{k}\isactrlesub }\\[4\jot] 
implies \isa{\isafree{{\isasymalpha}}\ {\isacharequal}{\kern0pt}\ \isafree{{\isasymgamma}}} \isa{\isafree{X}\ {\isacharequal}{\kern0pt}\ \isafree{Y}}, and \isa{\isafree{x}\ {\isacharequal}{\kern0pt}\ \isafree{y}}.
\end{definition}

It is worth noting that our definition, in accordance with that of Wilhelm et 
al.~\cite[p. 113]{Wilhelm}, is restricted to the extension of a CFG. We will discuss a more 
generalized definition applicable to arbitrary CFGs in a later section, but the original
definition suffices to prove the main theorem for the section on \isa{LR{\isacharparenleft}{\kern0pt}{\isadigit{0}}{\isacharparenright}{\kern0pt}} parsing presented by 
Wilhelm et al.~\cite[p. 114]{Wilhelm}, since the theorem relates the \isa{LR{\isacharparenleft}{\kern0pt}k{\isacharparenright}{\kern0pt}} condition only to the
parser, which we have also restricted to extended grammars. Before this theorem, we show some auxiliary 
lemmas.

\begin{lemma}\label{derivern_Suc_substring_reliable}
If a rightmost derivation \isa{\isaconst{P{\isacharprime}{\kern0pt}}\ {\isasymturnstile}\ {\isacharbrackleft}{\kern0pt}\isaconst{Nt}\ \isaconst{S{\isacharprime}{\kern0pt}}{\isacharbrackright}{\kern0pt}\ {\isasymRightarrow}r{\isacharparenleft}{\kern0pt}\isaconst{Suc}\ \isafree{n}{\isacharparenright}{\kern0pt}\ \isafree{{\isasymalpha}}\ {\isacharat}{\kern0pt}\ \isafree{{\isasymbeta}}\ {\isacharat}{\kern0pt}\ \isaconst{Nt}\ \isafree{X}\ {\isacharhash}{\kern0pt}\ \mbox{\isaconst{map}\ \isaconst{Tm}\ \isafree{w}}} exists,
there also exists an incomplete item \isa{\mbox{{\isacharbrackleft}{\kern0pt}\isafree{A}\ {\isasymrightarrow}\ \isafree{{\isasymalpha}{\isacharprime}{\kern0pt}}\ {\isasymcdot}\ \isafree{Y}\ {\isacharhash}{\kern0pt}\ \isafree{{\isasymbeta}{\isacharprime}{\kern0pt}}{\isacharbrackright}{\kern0pt}}} valid for \isa{{\isasymalpha}} with
\begin{gather*}
\isa{\isaconst{P{\isacharprime}{\kern0pt}}\ {\isasymturnstile}\ {\isacharbrackleft}{\kern0pt}\isaconst{Nt}\ \isaconst{S{\isacharprime}{\kern0pt}}{\isacharbrackright}{\kern0pt}\ {\isasymRightarrow}r{\isacharasterisk}{\kern0pt}\ \isafree{{\isasymalpha}}\ {\isacharat}{\kern0pt}\ \isafree{Y}\ {\isacharhash}{\kern0pt}\ \isafree{{\isasymbeta}{\isacharprime}{\kern0pt}}\ {\isacharat}{\kern0pt}\ \isaconst{map}\ \isaconst{Tm}\ \isafree{w{\isacharprime}{\kern0pt}}}\\
\isa{\isaconst{P{\isacharprime}{\kern0pt}}\ {\isasymturnstile}\ \isafree{Y}\ {\isacharhash}{\kern0pt}\ \isafree{{\isasymbeta}{\isacharprime}{\kern0pt}}\ {\isacharat}{\kern0pt}\ \isaconst{map}\ \isaconst{Tm}\ \isafree{w{\isacharprime}{\kern0pt}}\ {\isasymRightarrow}r{\isacharasterisk}{\kern0pt}\ \isafree{{\isasymbeta}}\ {\isacharat}{\kern0pt}\ \isaconst{Nt}\ \isafree{X}\ {\isacharhash}{\kern0pt}\ \isaconst{map}\ \isaconst{Tm}\ \isafree{w}}
\end{gather*}
\begin{proof}
Since the rightmost derivation has length greater than zero, we know there exists a \isa{{\isasymrho}} where
\isa{\isaconst{P{\isacharprime}{\kern0pt}}\ {\isasymturnstile}\ {\isacharbrackleft}{\kern0pt}\isaconst{Nt}\ \isaconst{S{\isacharprime}{\kern0pt}}{\isacharbrackright}{\kern0pt}\ {\isasymmidarrow}\isafree{{\isasymrho}}{\isasymrightarrow}r{\isacharasterisk}{\kern0pt}\ \isafree{{\isasymalpha}}\ {\isacharat}{\kern0pt}\ \isafree{{\isasymbeta}}\ {\isacharat}{\kern0pt}\ \isaconst{Nt}\ \isafree{X}\ {\isacharhash}{\kern0pt}\ \isaconst{map}\ \isaconst{Tm}\ \isafree{w}}. The length of the derivation also 
implies \isa{{\isacharbrackleft}{\kern0pt}\isaconst{Nt}\ \isaconst{S{\isacharprime}{\kern0pt}}{\isacharbrackright}{\kern0pt}\ {\isasymnoteq}\ \isafree{{\isasymalpha}}\ {\isacharat}{\kern0pt}\ \isafree{{\isasymbeta}}\ {\isacharat}{\kern0pt}\ \isaconst{Nt}\ \isafree{X}\ {\isacharhash}{\kern0pt}\ \isaconst{map}\ \isaconst{Tm}\ \isafree{w}}. We now use this inequality to induct on \isa{{\isasymrho}} 
for arbitrary \isa{X{\isacharcomma}{\kern0pt}\ {\isasymalpha}{\isacharcomma}{\kern0pt}\ {\isasymbeta}} and \isa{w}. 

The case \isa{{\isasymrho}\ {\isacharequal}{\kern0pt}\ {\isacharbrackleft}{\kern0pt}{\isacharbrackright}{\kern0pt}} cannot hold, since it would contradict the inequality.

If \isa{{\isasymrho}\ {\isacharequal}{\kern0pt}\ i\ {\isacharhash}{\kern0pt}\ {\isasymsigma}} for some \isa{i} and \isa{{\isasymsigma}}, performing rule inversion on the rightmost chain implies the
existence of some \isa{A\ {\isacharcolon}{\kern0pt}{\isacharcolon}{\kern0pt}\ {\isacharprime}{\kern0pt}n}, \isa{{\isasymalpha}{\isacharprime}{\kern0pt}{\isacharcomma}{\kern0pt}\ {\isasymalpha}{\isacharprime}{\kern0pt}{\isacharprime}{\kern0pt}{\isacharcomma}{\kern0pt}\ {\isasymbeta}{\isacharprime}{\kern0pt}\ {\isacharcolon}{\kern0pt}{\isacharcolon}{\kern0pt}\ syms}, and \isa{u{\isacharcomma}{\kern0pt}\ v\ {\isacharcolon}{\kern0pt}{\isacharcolon}{\kern0pt}\ {\isacharprime}{\kern0pt}t\ list} such that
\begin{subequations}
\begin{gather}
\isa{\isafree{i}\ {\isacharequal}{\kern0pt}\ {\isacharbrackleft}{\kern0pt}\isafree{A}\ {\isasymrightarrow}\ \isafree{{\isasymalpha}{\isacharprime}{\kern0pt}{\isacharprime}{\kern0pt}}\ {\isasymcdot}\ \isaconst{Nt}\ \isafree{X}\ {\isacharhash}{\kern0pt}\ \isafree{{\isasymbeta}{\isacharprime}{\kern0pt}}{\isacharbrackright}{\kern0pt}}\label{der_sucn.start}\\
\isa{\isafree{{\isasymalpha}}\ {\isacharat}{\kern0pt}\ \isafree{{\isasymbeta}}\ {\isacharat}{\kern0pt}\ \isaconst{Nt}\ \isafree{X}\ {\isacharhash}{\kern0pt}\ \isaconst{map}\ \isaconst{Tm}\ \isafree{w}\ {\isacharequal}{\kern0pt}\ \isafree{{\isasymalpha}{\isacharprime}{\kern0pt}}\ {\isacharat}{\kern0pt}\ \isafree{{\isasymalpha}{\isacharprime}{\kern0pt}{\isacharprime}{\kern0pt}}\ {\isacharat}{\kern0pt}\ \isaconst{Nt}\ \isafree{X}\ {\isacharhash}{\kern0pt}\ \isaconst{map}\ \isaconst{Tm}\ \isafree{u}\ {\isacharat}{\kern0pt}\ \isaconst{map}\ \isaconst{Tm}\ \isafree{v}}\label{der_sucn.eq}\\
\isa{\isaconst{P{\isacharprime}{\kern0pt}}\ {\isasymturnstile}\ {\isacharbrackleft}{\kern0pt}\isaconst{Nt}\ \isaconst{S{\isacharprime}{\kern0pt}}{\isacharbrackright}{\kern0pt}\ {\isasymmidarrow}\isafree{{\isasymsigma}}{\isasymrightarrow}r{\isacharasterisk}{\kern0pt}\ \isafree{{\isasymalpha}{\isacharprime}{\kern0pt}}\ {\isacharat}{\kern0pt}\ \isaconst{Nt}\ \isafree{A}\ {\isacharhash}{\kern0pt}\ \isaconst{map}\ \isaconst{Tm}\ \isafree{v}}\label{der_sucn.sig}\\
\isa{\isaconst{P{\isacharprime}{\kern0pt}}\ {\isasymturnstile}\ \isafree{{\isasymalpha}{\isacharprime}{\kern0pt}}\ {\isacharat}{\kern0pt}\ \isaconst{Nt}\ \isafree{A}\ {\isacharhash}{\kern0pt}\ \isaconst{map}\ \isaconst{Tm}\ \isafree{v}\ {\isasymRightarrow}r\ \isafree{{\isasymalpha}{\isacharprime}{\kern0pt}}\ {\isacharat}{\kern0pt}\ \isafree{{\isasymalpha}{\isacharprime}{\kern0pt}{\isacharprime}{\kern0pt}}\ {\isacharat}{\kern0pt}\ \isaconst{Nt}\ \isafree{X}\ {\isacharhash}{\kern0pt}\ \isafree{{\isasymbeta}{\isacharprime}{\kern0pt}}\ {\isacharat}{\kern0pt}\ \isaconst{map}\ \isaconst{Tm}\ \isafree{v}}\label{der_sucn.deriver}\\
\isa{\isaconst{P{\isacharprime}{\kern0pt}}\ {\isasymturnstile}\ \isafree{{\isasymbeta}{\isacharprime}{\kern0pt}}\ {\isasymRightarrow}r{\isacharasterisk}{\kern0pt}\ \isaconst{map}\ \isaconst{Tm}\ \isafree{u}}\label{der_sucn.end}
\end{gather}
\end{subequations}
We can now use \eqref{der_sucn.sig} in a second rule inversion.

The reflexive case implies \isa{{\isasymalpha}\ {\isacharequal}{\kern0pt}\ {\isasymbeta}\ {\isacharequal}{\kern0pt}\ {\isacharbrackleft}{\kern0pt}{\isacharbrackright}{\kern0pt}}, \isa{X\ {\isacharequal}{\kern0pt}\ S} and \isa{w\ {\isacharequal}{\kern0pt}\ {\isacharbrackleft}{\kern0pt}{\isacharbrackright}{\kern0pt}}, and the implication holds for item 
\isa{{\isacharbrackleft}{\kern0pt}S{\isacharprime}{\kern0pt}\ {\isasymrightarrow}\ {\isacharbrackleft}{\kern0pt}{\isacharbrackright}{\kern0pt}\ {\isasymcdot}\ Nt\ S\ {\isacharhash}{\kern0pt}\ {\isacharbrackleft}{\kern0pt}{\isacharbrackright}{\kern0pt}{\isacharbrackright}{\kern0pt}}.

In the \isa{step} case, we can distinguish two cases based on \eqref{der_sucn.eq}: either \isa{{\isasymalpha}} is a 
prefix of \isa{{\isasymalpha}{\isacharprime}{\kern0pt}}, or vice versa. 

If \isa{{\isasymalpha}} is a prefix of \isa{{\isasymalpha}{\isacharprime}{\kern0pt}}, there exists some \isa{{\isasymgamma}} where \isa{\isafree{{\isasymalpha}{\isacharprime}{\kern0pt}}\ {\isacharequal}{\kern0pt}\ \isafree{{\isasymalpha}}\ {\isacharat}{\kern0pt}\ \isafree{{\isasymgamma}}} and 
\isa{\isafree{{\isasymbeta}}\ {\isacharequal}{\kern0pt}\ \isafree{{\isasymgamma}}\ {\isacharat}{\kern0pt}\ \isafree{{\isasymalpha}{\isacharprime}{\kern0pt}{\isacharprime}{\kern0pt}}}. Therefore, this is the case where a previous, shorter chain already derived 
our prefix \isa{{\isasymalpha}}, meaning we will work towards using our induction hypothesis. To achieve this, we can 
use equations \eqref{der_sucn.start} - \eqref{der_sucn.end} to show that 
\isa{\isaconst{P{\isacharprime}{\kern0pt}}\ {\isasymturnstile}\ \isafree{{\isasymgamma}}\ {\isacharat}{\kern0pt}\ \isaconst{Nt}\ \isafree{A}\ {\isacharhash}{\kern0pt}\ \isaconst{map}\ \isaconst{Tm}\ \isafree{v}\ {\isasymRightarrow}r{\isacharasterisk}{\kern0pt}\ \isafree{{\isasymbeta}}\ {\isacharat}{\kern0pt}\ \isaconst{Nt}\ \isafree{X}\ {\isacharhash}{\kern0pt}\ \isaconst{map}\ \isaconst{Tm}\ \isafree{w}}.
Furthermore, with our \isa{step} assumptions for \eqref{der_sucn.sig}, 
\isa{{\isacharbrackleft}{\kern0pt}\isaconst{Nt}\ \isaconst{S{\isacharprime}{\kern0pt}}{\isacharbrackright}{\kern0pt}\ {\isasymnoteq}\ \isafree{{\isasymalpha}}\ {\isacharat}{\kern0pt}\ \isafree{{\isasymgamma}}\ {\isacharat}{\kern0pt}\ \isaconst{Nt}\ \isafree{A}\ {\isacharhash}{\kern0pt}\ \isaconst{map}\ \isaconst{Tm}\ \isafree{v}} must hold. With all this we can show the claim holds 
by the IH. 

On the other hand, if \isa{{\isasymalpha}{\isacharprime}{\kern0pt}} is a prefix of \isa{{\isasymalpha}} there exists some \isa{{\isasymgamma}} with
\isa{\isafree{{\isasymalpha}}\ {\isacharequal}{\kern0pt}\ \isafree{{\isasymalpha}{\isacharprime}{\kern0pt}}\ {\isacharat}{\kern0pt}\ \isafree{{\isasymgamma}}} and \isa{\isafree{{\isasymalpha}{\isacharprime}{\kern0pt}{\isacharprime}{\kern0pt}}\ {\isacharequal}{\kern0pt}\ \isafree{{\isasymgamma}}\ {\isacharat}{\kern0pt}\ \isafree{{\isasymbeta}}}. Furthermore, by \eqref{der_sucn.deriver} and 
\eqref{der_sucn.eq} we have 
\begin{equation}\label{der_sucn.A_deriver}
\isa{\isaconst{P{\isacharprime}{\kern0pt}}\ {\isasymturnstile}\ \isafree{{\isasymalpha}{\isacharprime}{\kern0pt}}\ {\isacharat}{\kern0pt}\ \isaconst{Nt}\ \isafree{A}\ {\isacharhash}{\kern0pt}\ \isaconst{map}\ \isaconst{Tm}\ \isafree{v}\ {\isasymRightarrow}r\ \isafree{{\isasymalpha}}\ {\isacharat}{\kern0pt}\ \isafree{{\isasymbeta}}\ {\isacharat}{\kern0pt}\ \isaconst{Nt}\ \isafree{X}\ {\isacharhash}{\kern0pt}\ \isafree{{\isasymbeta}{\isacharprime}{\kern0pt}}\ {\isacharat}{\kern0pt}\ \isaconst{map}\ \isaconst{Tm}\ \isafree{v}}.
\end{equation}
All this implies that \isa{{\isacharbrackleft}{\kern0pt}\isafree{A}\ {\isasymrightarrow}\ \isafree{{\isasymgamma}}\ {\isasymcdot}\ \isafree{{\isasymbeta}}\ {\isacharat}{\kern0pt}\ \isaconst{Nt}\ \isafree{X}\ {\isacharhash}{\kern0pt}\ \isafree{{\isasymbeta}{\isacharprime}{\kern0pt}}{\isacharbrackright}{\kern0pt}} is valid for \isa{{\isasymalpha}}. Moreover,
from \eqref{der_sucn.end} we have that 
\isa{\isaconst{P{\isacharprime}{\kern0pt}}\ {\isasymturnstile}\ \isafree{{\isasymbeta}}\ {\isacharat}{\kern0pt}\ \isaconst{Nt}\ \isafree{X}\ {\isacharhash}{\kern0pt}\ \isafree{{\isasymbeta}{\isacharprime}{\kern0pt}}\ {\isacharat}{\kern0pt}\ \isaconst{map}\ \isaconst{Tm}\ \isafree{v}\ {\isasymRightarrow}r{\isacharasterisk}{\kern0pt}\ \isafree{{\isasymbeta}}\ {\isacharat}{\kern0pt}\ \isaconst{Nt}\ \isafree{X}\ {\isacharhash}{\kern0pt}\ \isaconst{map}\ \isaconst{Tm}\ \isafree{w}}, again by \eqref{der_sucn.eq}.
Finally, Lemma~\ref{rm_chain_imp_derivers} with \eqref{der_sucn.sig}, followed by
\eqref{der_sucn.A_deriver} by transitivity, imply that the claim holds for item 
\isa{{\isacharbrackleft}{\kern0pt}\isafree{A}\ {\isasymrightarrow}\ \isafree{{\isasymgamma}}\ {\isasymcdot}\ \isafree{{\isasymbeta}}\ {\isacharat}{\kern0pt}\ \isaconst{Nt}\ \isafree{X}\ {\isacharhash}{\kern0pt}\ \isafree{{\isasymbeta}{\isacharprime}{\kern0pt}}{\isacharbrackright}{\kern0pt}}.
\end{proof}
\end{lemma}

\begin{lemma}\label{derivers_substring_reliable}
If a rightmost derivation \isa{\isaconst{P{\isacharprime}{\kern0pt}}\ {\isasymturnstile}\ {\isacharbrackleft}{\kern0pt}\isaconst{Nt}\ \isaconst{S{\isacharprime}{\kern0pt}}{\isacharbrackright}{\kern0pt}\ {\isasymRightarrow}r{\isacharasterisk}{\kern0pt}\ \isafree{{\isasymalpha}}\ {\isacharat}{\kern0pt}\ \isafree{{\isasymbeta}}\ {\isacharat}{\kern0pt}\ \isaconst{Nt}\ \isafree{X}\ {\isacharhash}{\kern0pt}\ \isaconst{map}\ \isaconst{Tm}\ \isafree{w}}
exists, there also exists an incomplete item valid for \isa{{\isasymalpha}}.
\begin{proof}
The proof is by a case distinction on the derivation length. If it is \isa{{\isadigit{0}}}, \isa{{\isasymalpha}} must be an empty 
string. Since \isa{{\isacharbrackleft}{\kern0pt}\isaconst{S{\isacharprime}{\kern0pt}}\ {\isasymrightarrow}\ \ {\isasymcdot}\ {\isacharbrackleft}{\kern0pt}\isaconst{Nt}\ \isaconst{S}{\isacharbrackright}{\kern0pt}\ {\isacharbrackright}{\kern0pt}} is valid for \isa{{\isacharbrackleft}{\kern0pt}{\isacharbrackright}{\kern0pt}}, the implication holds.

If the length is greater than \isa{{\isadigit{0}}}, the claim follows directly from 
Lemma~\ref{derivern_Suc_substring_reliable}.
\end{proof}
\end{lemma}

\begin{lemma}\label{prefix_comp_unique_imp_eps}
If the set ``\isa{\isaconst{valids}\ \isafree{{\isasymalpha}}}'' is \isa{LR{\isacharparenleft}{\kern0pt}{\isadigit{0}}{\isacharparenright}{\kern0pt}}-adequate and contains item \isa{\mbox{{\isacharbrackleft}{\kern0pt}\isafree{X}\ {\isasymrightarrow}\ {\isasymcdot}\ {\isacharbrackright}{\kern0pt}}}, then 
\begin{equation*}
\isa{P{\isacharprime}{\kern0pt}\ {\isasymturnstile}\ {\isacharbrackleft}{\kern0pt}Nt\ S{\isacharprime}{\kern0pt}{\isacharbrackright}{\kern0pt}\ {\isasymRightarrow}r{\isacharasterisk}{\kern0pt}\ {\isasymalpha}\ {\isacharat}{\kern0pt}\ {\isasymalpha}{\isacharprime}{\kern0pt}\ {\isacharat}{\kern0pt}\ Nt\ Y\ {\isacharhash}{\kern0pt}\ map\ Tm\ x\ {\isasymRightarrow}r\ {\isasymalpha}\ {\isacharat}{\kern0pt}\ map\ Tm\ y}
\end{equation*}
implies \isa{\isafree{{\isasymalpha}{\isacharprime}{\kern0pt}}\ {\isacharequal}{\kern0pt}\ {\isacharbrackleft}{\kern0pt}{\isacharbrackright}{\kern0pt}\ {\isasymand}\ \isafree{Y}\ {\isacharequal}{\kern0pt}\ \isafree{X}\ {\isasymand}\ \isafree{x}\ {\isacharequal}{\kern0pt}\ \isafree{y}}.
\begin{proof}
We first distinguish cases on the length of the derivation of \isa{{\isasymalpha}\ {\isacharat}{\kern0pt}\ {\isasymalpha}{\isacharprime}{\kern0pt}\ {\isacharat}{\kern0pt}\ Nt\ Y\ {\isacharhash}{\kern0pt}\ map\ Tm\ x}.

If the length is \isa{{\isadigit{0}}}, the implication holds trivially by reflexivity.

If the length is nonzero, by Lemma~\ref{derivern_Suc_substring_reliable} there exists an item 
\isa{\mbox{{\isacharbrackleft}{\kern0pt}\isafree{A}\ {\isasymrightarrow}\ \isafree{{\isasymbeta}}\ {\isasymcdot}\ \isafree{Z{\isacharprime}{\kern0pt}}\ {\isacharhash}{\kern0pt}\ \isafree{{\isasymgamma}}{\isacharbrackright}{\kern0pt}}} in \isa{\isaconst{valids}\ \isafree{{\isasymalpha}}} such that for some \isa{z} holds
\begin{equation}\label{pref_eps.deriver}
 \isa{P{\isacharprime}{\kern0pt}\ {\isasymturnstile}\ {\isacharbrackleft}{\kern0pt}Nt\ S{\isacharprime}{\kern0pt}{\isacharbrackright}{\kern0pt}\ {\isasymRightarrow}r{\isacharasterisk}{\kern0pt}\ {\isasymalpha}\ {\isacharat}{\kern0pt}\ Z{\isacharprime}{\kern0pt}\ {\isacharhash}{\kern0pt}\ {\isasymgamma}\ {\isacharat}{\kern0pt}\ map\ Tm\ z\ {\isasymRightarrow}r{\isacharasterisk}{\kern0pt}\ {\isasymalpha}{\isacharprime}{\kern0pt}\ {\isacharat}{\kern0pt}\ Nt\ Y\ {\isacharhash}{\kern0pt}\ map\ Tm\ x}.
\end{equation}
By case 3 of Lemma~\ref{LR0_adequate_cases}, we know there exists some nonterminal \isa{Z{\isacharprime}{\kern0pt}{\isacharprime}{\kern0pt}} with
\isa{\isafree{Z{\isacharprime}{\kern0pt}}\ {\isacharequal}{\kern0pt}\ \isaconst{Nt}\ \isafree{Z{\isacharprime}{\kern0pt}{\isacharprime}{\kern0pt}}}. This, the final step in the rightmost derivation in our assumption, and 
\eqref{pref_eps.deriver} imply that the derivation
\begin{equation*}
\isa{P{\isacharprime}{\kern0pt}\ {\isasymturnstile}\ Nt\ Z{\isacharprime}{\kern0pt}{\isacharprime}{\kern0pt}\ {\isacharhash}{\kern0pt}\ {\isasymgamma}\ {\isacharat}{\kern0pt}\ map\ Tm\ z\ {\isasymRightarrow}r{\isacharasterisk}{\kern0pt}\ {\isasymalpha}{\isacharprime}{\kern0pt}\ {\isacharat}{\kern0pt}\ Nt\ Y\ {\isacharhash}{\kern0pt}\ map\ Tm\ x\ {\isasymRightarrow}r\ map\ Tm\ y}.
\end{equation*}
can be decomposed into
\begin{equation*}
\isa{P{\isacharprime}{\kern0pt}\ {\isasymturnstile}\ Nt\ Z{\isacharprime}{\kern0pt}{\isacharprime}{\kern0pt}\ {\isacharhash}{\kern0pt}\ {\isasymgamma}\ {\isacharat}{\kern0pt}\ map\ Tm\ z\ {\isasymRightarrow}r{\isacharasterisk}{\kern0pt}\ Nt\ Z{\isacharprime}{\kern0pt}{\isacharprime}{\kern0pt}\ {\isacharhash}{\kern0pt}\ map\ Tm\ v\ {\isasymRightarrow}r{\isacharasterisk}{\kern0pt}\ {\isasymgamma}{\isacharprime}{\kern0pt}\ {\isacharat}{\kern0pt}\ map\ Tm\ v\ {\isasymRightarrow}r\ map\ Tm\ y}
\end{equation*}
for some \isa{{\isasymgamma}{\isacharprime}{\kern0pt}\ {\isacharcolon}{\kern0pt}{\isacharcolon}{\kern0pt}\ syms} and \isa{v\ {\isacharcolon}{\kern0pt}{\isacharcolon}{\kern0pt}\ {\isacharprime}{\kern0pt}t\ list} with \isa{\isafree{{\isasymgamma}{\isacharprime}{\kern0pt}}\ {\isacharat}{\kern0pt}\ \isaconst{map}\ \isaconst{Tm}\ \isafree{v}\ {\isacharequal}{\kern0pt}\ \isafree{{\isasymalpha}{\isacharprime}{\kern0pt}}\ {\isacharat}{\kern0pt}\ \isaconst{Nt}\ \isafree{Y}\ {\isacharhash}{\kern0pt}\ \isaconst{map}\ \isaconst{Tm}\ \isafree{x}}.
This decomposition in turn implies the existence of some \isa{u\ {\isacharcolon}{\kern0pt}{\isacharcolon}{\kern0pt}\ {\isacharprime}{\kern0pt}t\ list} such that 
\isa{\isafree{y}\ {\isacharequal}{\kern0pt}\ \isafree{u}\ {\isacharat}{\kern0pt}\ \isafree{v}} and
\[ \isa{P{\isacharprime}{\kern0pt}\ {\isasymturnstile}\ {\isacharbrackleft}{\kern0pt}Nt\ Z{\isacharprime}{\kern0pt}{\isacharprime}{\kern0pt}{\isacharbrackright}{\kern0pt}\ {\isasymRightarrow}r{\isacharasterisk}{\kern0pt}\ {\isasymgamma}{\isacharprime}{\kern0pt}\ {\isasymRightarrow}r\ map\ Tm\ u}.  \]
By case 3 of Lemma~\ref{LR0_adequate_cases}, this implies that \isa{\isafree{{\isasymgamma}{\isacharprime}{\kern0pt}}\ {\isacharequal}{\kern0pt}\ \isaconst{Nt}\ \isafree{X}\ {\isacharhash}{\kern0pt}\ \isaconst{map}\ \isaconst{Tm}\ \isafree{u}}.
This and the decomposition of the derivation imply \isa{{\isasymalpha}{\isacharprime}{\kern0pt}\ {\isacharat}{\kern0pt}\ Nt\ Y\ {\isacharhash}{\kern0pt}\ map\ Tm\ x\ {\isacharequal}{\kern0pt}\ Nt\ X\ {\isacharhash}{\kern0pt}\ map\ Tm\ y}.
From this follows \isa{\isafree{{\isasymalpha}{\isacharprime}{\kern0pt}}\ {\isacharequal}{\kern0pt}\ {\isacharbrackleft}{\kern0pt}{\isacharbrackright}{\kern0pt}\ {\isasymand}\ \isafree{Y}\ {\isacharequal}{\kern0pt}\ \isafree{X}\ {\isasymand}\ \isafree{x}\ {\isacharequal}{\kern0pt}\ \isafree{y}}, completing the proof.
\end{proof}
\end{lemma}

\begin{lemma}\label{is_LR_wlogI}
In order to prove that \isa{G{\isacharprime}{\kern0pt}} is an \isa{LR{\isacharparenleft}{\kern0pt}k{\isacharparenright}{\kern0pt}} grammar, it suffices to prove the following property for
arbitrary \isa{{\isasymalpha}{\isacharcomma}{\kern0pt}\ {\isasymbeta}{\isacharcomma}{\kern0pt}\ {\isasymgamma}{\isacharcomma}{\kern0pt}\ {\isasymdelta}\ {\isacharcolon}{\kern0pt}{\isacharcolon}{\kern0pt}\ syms}, \isa{X{\isacharcomma}{\kern0pt}\ Y\ {\isacharcolon}{\kern0pt}{\isacharcolon}{\kern0pt}\ {\isacharprime}{\kern0pt}n}, and \isa{w{\isacharcomma}{\kern0pt}\ x{\isacharcomma}{\kern0pt}\ y\ {\isacharcolon}{\kern0pt}{\isacharcolon}{\kern0pt}\ {\isacharprime}{\kern0pt}t\ list}:

The conjunction of the following statements
\begin{gather*}
\begin{multlined}
\isa{P{\isacharprime}{\kern0pt}\ {\isasymturnstile}\ {\isacharbrackleft}{\kern0pt}Nt\ S{\isacharprime}{\kern0pt}{\isacharbrackright}{\kern0pt}\ {\isasymRightarrow}r{\isacharasterisk}{\kern0pt}\ {\isasymalpha}\ {\isacharat}{\kern0pt}\ Nt\ X\ {\isacharhash}{\kern0pt}\ map\ Tm\ w\ {\isasymRightarrow}r\ {\isasymalpha}\ {\isacharat}{\kern0pt}\ {\isasymbeta}\ {\isacharat}{\kern0pt}\ map\ Tm\ w},
\end{multlined}\\
\begin{multlined}
\isa{P{\isacharprime}{\kern0pt}\ {\isasymturnstile}\ {\isacharbrackleft}{\kern0pt}Nt\ S{\isacharprime}{\kern0pt}{\isacharbrackright}{\kern0pt}\ {\isasymRightarrow}r{\isacharasterisk}{\kern0pt}\ {\isasymgamma}\ {\isacharat}{\kern0pt}\ Nt\ Y\ {\isacharhash}{\kern0pt}\ map\ Tm\ x\ {\isasymRightarrow}r\ {\isasymalpha}\ {\isacharat}{\kern0pt}\ {\isasymbeta}\ {\isacharat}{\kern0pt}\ map\ Tm\ y},
\end{multlined}\\
\isa{\isafree{{\isasymgamma}}\ {\isacharat}{\kern0pt}\ \isafree{{\isasymdelta}}\ {\isacharat}{\kern0pt}\ \isaconst{map}\ \isaconst{Tm}\ \isafree{x}\ {\isacharequal}{\kern0pt}\ \isafree{{\isasymalpha}}\ {\isacharat}{\kern0pt}\ \isafree{{\isasymbeta}}\ {\isacharat}{\kern0pt}\ \isaconst{map}\ \isaconst{Tm}\ \isafree{y}},\\
\isa{{\isacharbar}{\kern0pt}\isafree{{\isasymalpha}}\ {\isacharat}{\kern0pt}\ \isafree{{\isasymbeta}}{\isacharbar}{\kern0pt}\ {\isasymle}\ {\isacharbar}{\kern0pt}\isafree{{\isasymgamma}}\ {\isacharat}{\kern0pt}\ \isafree{{\isasymdelta}}{\isacharbar}{\kern0pt}},\\
\isa{\isafree{w}{\isacharbar}{\kern0pt}\isactrlbsub \isafree{k}\isactrlesub \ {\isacharequal}{\kern0pt}\ \isafree{y}{\isacharbar}{\kern0pt}\isactrlbsub \isafree{k}\isactrlesub }
\end{gather*}
implies \isa{\isafree{{\isasymalpha}}\ {\isacharequal}{\kern0pt}\ \isafree{{\isasymgamma}}}, \isa{\isafree{X}\ {\isacharequal}{\kern0pt}\ \isafree{Y}}, and \isa{\isafree{x}\ {\isacharequal}{\kern0pt}\ \isafree{y}}.
\begin{proof}
We begin by assuming that our implication holds, i.e., the conjunction of the statements above 
indeed implies \isa{\isafree{{\isasymalpha}}\ {\isacharequal}{\kern0pt}\ \isafree{{\isasymgamma}}\ {\isasymand}\ \isafree{X}\ {\isacharequal}{\kern0pt}\ \isafree{Y}\ {\isasymand}\ \isafree{x}\ {\isacharequal}{\kern0pt}\ \isafree{y}}. Let \isa{I} be this implication. Now it must be shown 
that \isa{I} implies the \isa{LR{\isacharparenleft}{\kern0pt}k{\isacharparenright}{\kern0pt}} condition for \isa{G{\isacharprime}{\kern0pt}}. We therefore fix two rightmost derivations
\begin{gather}
\isa{P{\isacharprime}{\kern0pt}\ {\isasymturnstile}\ {\isacharbrackleft}{\kern0pt}Nt\ S{\isacharprime}{\kern0pt}{\isacharbrackright}{\kern0pt}\ {\isasymRightarrow}r{\isacharasterisk}{\kern0pt}\ {\isasymalpha}\ {\isacharat}{\kern0pt}\ Nt\ X\ {\isacharhash}{\kern0pt}\ map\ Tm\ w\ {\isasymRightarrow}r\ {\isasymalpha}\ {\isacharat}{\kern0pt}\ {\isasymbeta}\ {\isacharat}{\kern0pt}\ map\ Tm\ w}\label{LR_wlog.X}
\intertext{and}
\isa{P{\isacharprime}{\kern0pt}\ {\isasymturnstile}\ {\isacharbrackleft}{\kern0pt}Nt\ S{\isacharprime}{\kern0pt}{\isacharbrackright}{\kern0pt}\ {\isasymRightarrow}r{\isacharasterisk}{\kern0pt}\ {\isasymgamma}\ {\isacharat}{\kern0pt}\ Nt\ Y\ {\isacharhash}{\kern0pt}\ map\ Tm\ x\ {\isasymRightarrow}r\ {\isasymalpha}\ {\isacharat}{\kern0pt}\ {\isasymbeta}\ {\isacharat}{\kern0pt}\ map\ Tm\ y}\label{LR_wlog.Y}
\intertext{such that}
\isa{\isafree{w}{\isacharbar}{\kern0pt}\isactrlbsub \isafree{k}\isactrlesub \ {\isacharequal}{\kern0pt}\ \isafree{y}{\isacharbar}{\kern0pt}\isactrlbsub \isafree{k}\isactrlesub }\label{LR_wlog.wkyk},
\end{gather}
and our goal is to show that \isa{\isafree{{\isasymalpha}}\ {\isacharequal}{\kern0pt}\ \isafree{{\isasymgamma}}\ {\isasymand}\ \isafree{X}\ {\isacharequal}{\kern0pt}\ \isafree{Y}\ {\isasymand}\ \isafree{x}\ {\isacharequal}{\kern0pt}\ \isafree{y}} holds.

By \eqref{LR_wlog.Y}, there exists a production \isa{{\isacharparenleft}{\kern0pt}Y{\isacharcomma}{\kern0pt}\ {\isasymdelta}{\isacharparenright}{\kern0pt}} that was applied in the last step, meaning 
\begin{equation}\label{LR_wlog.deriver}
\isa{\isafree{{\isasymalpha}}\ {\isacharat}{\kern0pt}\ \isafree{{\isasymbeta}}\ {\isacharat}{\kern0pt}\ \isaconst{map}\ \isaconst{Tm}\ \isafree{y}\ {\isacharequal}{\kern0pt}\ \isafree{{\isasymgamma}}\ {\isacharat}{\kern0pt}\ \isafree{{\isasymdelta}}\ {\isacharat}{\kern0pt}\ \isaconst{map}\ \isaconst{Tm}\ \isafree{x}}.
\end{equation}
We now distinguish two cases: \isa{{\isacharbar}{\kern0pt}\isafree{{\isasymalpha}}\ {\isacharat}{\kern0pt}\ \isafree{{\isasymbeta}}{\isacharbar}{\kern0pt}\ {\isasymle}\ {\isacharbar}{\kern0pt}\isafree{{\isasymgamma}}\ {\isacharat}{\kern0pt}\ \isafree{{\isasymdelta}}{\isacharbar}{\kern0pt}} and 
\isa{{\isacharbar}{\kern0pt}\isafree{{\isasymgamma}}\ {\isacharat}{\kern0pt}\ \isafree{{\isasymdelta}}{\isacharbar}{\kern0pt}\ {\isasymle}\ {\isacharbar}{\kern0pt}\isafree{{\isasymalpha}}\ {\isacharat}{\kern0pt}\ \isafree{{\isasymbeta}}{\isacharbar}{\kern0pt}}.

If \isa{{\isacharbar}{\kern0pt}\isafree{{\isasymalpha}}\ {\isacharat}{\kern0pt}\ \isafree{{\isasymbeta}}{\isacharbar}{\kern0pt}\ {\isasymle}\ {\isacharbar}{\kern0pt}\isafree{{\isasymgamma}}\ {\isacharat}{\kern0pt}\ \isafree{{\isasymdelta}}{\isacharbar}{\kern0pt}}, \isa{\isafree{{\isasymalpha}}\ {\isacharequal}{\kern0pt}\ \isafree{{\isasymgamma}}\ {\isasymand}\ \isafree{X}\ {\isacharequal}{\kern0pt}\ \isafree{Y}\ {\isasymand}\ \isafree{x}\ {\isacharequal}{\kern0pt}\ \isafree{y}} follows by \isa{I} and
equations \eqref{LR_wlog.X} - \eqref{LR_wlog.deriver}.

If \isa{{\isacharbar}{\kern0pt}\isafree{{\isasymgamma}}\ {\isacharat}{\kern0pt}\ \isafree{{\isasymdelta}}{\isacharbar}{\kern0pt}\ {\isasymle}\ {\isacharbar}{\kern0pt}\isafree{{\isasymalpha}}\ {\isacharat}{\kern0pt}\ \isafree{{\isasymbeta}}{\isacharbar}{\kern0pt}}, by \eqref{LR_wlog.deriver} we know that there exists 
some \isa{z\ {\isacharcolon}{\kern0pt}{\isacharcolon}{\kern0pt}\ {\isacharprime}{\kern0pt}t\ list} such that 
\begin{equation}\label{LR_wlog.ab_eq}
\isa{\isafree{{\isasymalpha}}\ {\isacharat}{\kern0pt}\ \isafree{{\isasymbeta}}\ {\isacharequal}{\kern0pt}\ \isafree{{\isasymgamma}}\ {\isacharat}{\kern0pt}\ \isafree{{\isasymdelta}}\ {\isacharat}{\kern0pt}\ \isaconst{map}\ \isaconst{Tm}\ \isafree{z}} \text{ and } \isa{\isafree{x}\ {\isacharequal}{\kern0pt}\ \isafree{z}\ {\isacharat}{\kern0pt}\ \isafree{y}}.
\end{equation} 
Furthermore, this and \eqref{LR_wlog.wkyk} imply that \isa{\isafree{x}{\isacharbar}{\kern0pt}\isactrlbsub \isafree{k}\isactrlesub \ {\isacharequal}{\kern0pt}\ {\isacharparenleft}{\kern0pt}\isafree{z}\ {\isacharat}{\kern0pt}\ \isafree{w}{\isacharparenright}{\kern0pt}{\isacharbar}{\kern0pt}\isactrlbsub \isafree{k}\isactrlesub }. With this 
we can finally show \isa{\isafree{{\isasymalpha}}\ {\isacharequal}{\kern0pt}\ \isafree{{\isasymgamma}}\ {\isasymand}\ \isafree{X}\ {\isacharequal}{\kern0pt}\ \isafree{Y}\ {\isasymand}\ \isafree{x}\ {\isacharequal}{\kern0pt}\ \isafree{y}} using \isa{I} after substituting
\eqref{LR_wlog.ab_eq} into the final step of \eqref{LR_wlog.X}, and \eqref{LR_wlog.deriver} into
the final step of \eqref{LR_wlog.Y}.
\end{proof}
\end{lemma}

This lemma tells us that when proving the \isa{LR{\isacharparenleft}{\kern0pt}k{\isacharparenright}{\kern0pt}} condition holds for \isa{G{\isacharprime}{\kern0pt}}, we can assume without 
loss of generality that \isa{{\isasymalpha}\ {\isacharat}{\kern0pt}\ {\isasymbeta}} is a prefix of \isa{{\isasymgamma}\ {\isacharat}{\kern0pt}\ {\isasymdelta}}, where \isa{{\isasymdelta}} is the handle of 
\isa{{\isasymalpha}\ {\isacharat}{\kern0pt}\ {\isasymbeta}\ {\isacharat}{\kern0pt}\ map\ Tm\ y}. In other words, we can assume WLOG that the handle of \isa{{\isasymalpha}\ {\isacharat}{\kern0pt}\ {\isasymbeta}\ {\isacharat}{\kern0pt}\ map\ Tm\ y} is not 
contained in \isa{{\isasymalpha}\ {\isacharat}{\kern0pt}\ {\isasymbeta}}, except in the case of equality. This assumption will greatly simplify the proof 
of the following theorem.

\begin{theorem}\label{is_LR0_iff_no_LR0_inadequates}
\isa{G{\isacharprime}{\kern0pt}} is an \isa{LR{\isacharparenleft}{\kern0pt}{\isadigit{0}}{\isacharparenright}{\kern0pt}} grammar if and only if \isa{\ensuremath{LR_0(G)}} has no \isa{LR{\isacharparenleft}{\kern0pt}{\isadigit{0}}{\isacharparenright}{\kern0pt}}-inadequate states.
\begin{proof}
For direction ``\isa{{\isasymLongrightarrow}}'', we prove the implication by contradiction.
Assume \isa{G{\isacharprime}{\kern0pt}} is \isa{LR{\isacharparenleft}{\kern0pt}{\isadigit{0}}{\isacharparenright}{\kern0pt}} and \isa{\ensuremath{LR_0(G)}} has an \isa{LR{\isacharparenleft}{\kern0pt}{\isadigit{0}}{\isacharparenright}{\kern0pt}}-inadequate state \isa{q} meaning it either has a 
shift-reduce or a reduce-reduce conflict. 

If \isa{q} has a shift-reduce conflict, there exist two items \isa{{\isacharbrackleft}{\kern0pt}\isafree{X}\ {\isasymrightarrow}\ \isafree{{\isasymbeta}}\ {\isasymcdot}\ {\isacharbrackright}{\kern0pt}} and 
\isa{\mbox{{\isacharbrackleft}{\kern0pt}\isafree{Y}\ {\isasymrightarrow}\ \isafree{{\isasymdelta}}\ {\isasymcdot}\ \isaconst{Tm}\ \isafree{a}\ {\isacharhash}{\kern0pt}\ \isafree{{\isasymalpha}}{\isacharbrackright}{\kern0pt}}} in \isa{q}. By Lemma~\ref{state_imp_valids}, we know there exists some 
\isa{{\isasymgamma}\ {\isacharcolon}{\kern0pt}{\isacharcolon}{\kern0pt}\ syms} that is a reliable prefix for both of these states. Therefore, there exist rightmost
derivations of the form
\begin{gather}
  \isa{P{\isacharprime}{\kern0pt}\ {\isasymturnstile}\ {\isacharbrackleft}{\kern0pt}Nt\ S{\isacharprime}{\kern0pt}{\isacharbrackright}{\kern0pt}\ {\isasymRightarrow}r{\isacharasterisk}{\kern0pt}\ {\isasymgamma}{\isacharprime}{\kern0pt}\ {\isacharat}{\kern0pt}\ Nt\ X\ {\isacharhash}{\kern0pt}\ map\ Tm\ w\ {\isasymRightarrow}r\ {\isasymgamma}{\isacharprime}{\kern0pt}\ {\isacharat}{\kern0pt}\ {\isasymbeta}\ {\isacharat}{\kern0pt}\ map\ Tm\ w}\label{LR0_iff.X_derivers1}
\intertext{and}
  \isa{P{\isacharprime}{\kern0pt}\ {\isasymturnstile}\ {\isacharbrackleft}{\kern0pt}Nt\ S{\isacharprime}{\kern0pt}{\isacharbrackright}{\kern0pt}\ {\isasymRightarrow}r{\isacharasterisk}{\kern0pt}\ {\isasymnu}\ {\isacharat}{\kern0pt}\ Nt\ Y\ {\isacharhash}{\kern0pt}\ map\ Tm\ y\ {\isasymRightarrow}r\ {\isasymnu}\ {\isacharat}{\kern0pt}\ {\isasymdelta}\ {\isacharat}{\kern0pt}\ Tm\ a\ {\isacharhash}{\kern0pt}\ {\isasymalpha}\ {\isacharat}{\kern0pt}\ map\ Tm\ y}\label{LR0_iff.Y_derivers1}
\intertext{such that}
\isa{{\isasymgamma}\ {\isacharequal}{\kern0pt}\ {\isasymgamma}{\isacharprime}{\kern0pt}\ {\isacharat}{\kern0pt}\ {\isasymbeta}\ {\isacharequal}{\kern0pt}\ {\isasymnu}\ {\isacharat}{\kern0pt}\ {\isasymdelta}}\label{LR0_iff.g_rp}
\end{gather}
Since \isa{G{\isacharprime}{\kern0pt}} is reduced, there exists some \isa{v\ {\isacharcolon}{\kern0pt}{\isacharcolon}{\kern0pt}\ {\isacharprime}{\kern0pt}t\ list} derived by \isa{{\isasymalpha}}. 
If the derivation of this \isa{v} has length \isa{{\isadigit{0}}}, i.e., \isa{{\isasymalpha}\ {\isacharequal}{\kern0pt}\ map\ Tm\ v}, we get a contradiction, since
the \isa{LR{\isacharparenleft}{\kern0pt}{\isadigit{0}}{\isacharparenright}{\kern0pt}} condition requires $\isa{y} \overset{!}{=} \isa{a\ {\isacharhash}{\kern0pt}\ v\ {\isacharat}{\kern0pt}\ y}$, which clearly cannot hold.

On the other hand, if the rightmost derivation of \isa{v} has nonzero length, it is of the form 
\[ \isa{P{\isacharprime}{\kern0pt}\ {\isasymturnstile}\ {\isasymalpha}\ {\isasymRightarrow}r{\isacharasterisk}{\kern0pt}\ map\ Tm\ u\ {\isacharat}{\kern0pt}\ Nt\ Z\ {\isacharhash}{\kern0pt}\ map\ Tm\ x\ {\isasymRightarrow}r\ map\ Tm\ {\isacharparenleft}{\kern0pt}u\ {\isacharat}{\kern0pt}\ v{\isacharprime}{\kern0pt}\ {\isacharat}{\kern0pt}\ x{\isacharparenright}{\kern0pt}\ {\isacharequal}{\kern0pt}\ map\ Tm\ v}. \]
With \eqref{LR0_iff.Y_derivers1} and \eqref{LR0_iff.g_rp}, this implies
\begin{multline*}
\isa{P{\isacharprime}{\kern0pt}\ {\isasymturnstile}\ {\isacharbrackleft}{\kern0pt}Nt\ S{\isacharprime}{\kern0pt}{\isacharbrackright}{\kern0pt}\ {\isasymRightarrow}r{\isacharasterisk}{\kern0pt}\ {\isasymnu}\ {\isacharat}{\kern0pt}\ {\isasymdelta}\ {\isacharat}{\kern0pt}\ Tm\ a\ {\isacharhash}{\kern0pt}\ map\ Tm\ u\ {\isacharat}{\kern0pt}\ Nt\ Z\ {\isacharhash}{\kern0pt}\ map\ Tm\ {\isacharparenleft}{\kern0pt}x\ {\isacharat}{\kern0pt}\ y{\isacharparenright}{\kern0pt}}\\
  \isa{{\isasymRightarrow}r\ {\isasymgamma}{\isacharprime}{\kern0pt}\ {\isacharat}{\kern0pt}\ {\isasymbeta}\ {\isacharat}{\kern0pt}\ Tm\ a\ {\isacharhash}{\kern0pt}\ map\ Tm\ {\isacharparenleft}{\kern0pt}u\ {\isacharat}{\kern0pt}\ v{\isacharprime}{\kern0pt}\ {\isacharat}{\kern0pt}\ x\ {\isacharat}{\kern0pt}\ y{\isacharparenright}{\kern0pt}}.
\end{multline*}
This is again a contradiction, since by \eqref{LR0_iff.X_derivers1} the \isa{LR{\isacharparenleft}{\kern0pt}{\isadigit{0}}{\isacharparenright}{\kern0pt}} condition requires 
$\isa{x\ {\isacharat}{\kern0pt}\ y} \overset{!}{=} \isa{a\ {\isacharhash}{\kern0pt}\ u\ {\isacharat}{\kern0pt}\ v{\isacharprime}{\kern0pt}\ {\isacharat}{\kern0pt}\ x\ {\isacharat}{\kern0pt}\ y}$. This finishes the case where state \isa{q} has a 
shift-reduce conflict.

In the case where \isa{q} has a reduce-reduce conflict, there exist two distinct complete items 
\isa{{\isacharbrackleft}{\kern0pt}\isafree{X}\ {\isasymrightarrow}\ \isafree{{\isasymbeta}}\ {\isasymcdot}\ {\isacharbrackright}{\kern0pt}} and \isa{{\isacharbrackleft}{\kern0pt}\isafree{Y}\ {\isasymrightarrow}\ \isafree{{\isasymdelta}}\ {\isasymcdot}\ {\isacharbrackright}{\kern0pt}} in \isa{q}. Once again by Lemma~\ref{state_imp_valids},
there exists a \isa{{\isasymgamma}} for which both items are valid. Therefore, there exist rightmost derivations of 
the form
\begin{gather*}
\isa{P{\isacharprime}{\kern0pt}\ {\isasymturnstile}\ {\isacharbrackleft}{\kern0pt}Nt\ S{\isacharprime}{\kern0pt}{\isacharbrackright}{\kern0pt}\ {\isasymRightarrow}r{\isacharasterisk}{\kern0pt}\ {\isasymgamma}{\isacharprime}{\kern0pt}\ {\isacharat}{\kern0pt}\ Nt\ X\ {\isacharhash}{\kern0pt}\ map\ Tm\ w\ {\isasymRightarrow}r\ {\isasymgamma}{\isacharprime}{\kern0pt}\ {\isacharat}{\kern0pt}\ {\isasymbeta}\ {\isacharat}{\kern0pt}\ map\ Tm\ w}
\intertext{and}
\isa{P{\isacharprime}{\kern0pt}\ {\isasymturnstile}\ {\isacharbrackleft}{\kern0pt}Nt\ S{\isacharprime}{\kern0pt}{\isacharbrackright}{\kern0pt}\ {\isasymRightarrow}r{\isacharasterisk}{\kern0pt}\ {\isasymnu}\ {\isacharat}{\kern0pt}\ Nt\ Y\ {\isacharhash}{\kern0pt}\ map\ Tm\ y\ {\isasymRightarrow}r\ {\isasymnu}\ {\isacharat}{\kern0pt}\ {\isasymdelta}\ {\isacharat}{\kern0pt}\ map\ Tm\ y}
\end{gather*}
with \isa{{\isasymgamma}\ {\isacharequal}{\kern0pt}\ {\isasymgamma}{\isacharprime}{\kern0pt}\ {\isacharat}{\kern0pt}\ {\isasymbeta}\ {\isacharequal}{\kern0pt}\ {\isasymnu}\ {\isacharat}{\kern0pt}\ {\isasymdelta}}. Since the two items are distinct, there are two cases:
\begin{itemize}
\item \isa{{\isasymbeta}\ {\isasymnoteq}\ {\isasymdelta}}, implying \isa{{\isasymgamma}{\isacharprime}{\kern0pt}\ {\isasymnoteq}\ {\isasymnu}}; or
\item \isa{X\ {\isasymnoteq}\ Y}
\end{itemize}
Either case contradicts the \isa{LR{\isacharparenleft}{\kern0pt}{\isadigit{0}}{\isacharparenright}{\kern0pt}} condition, thus proving the reduce-reduce case.

We now move on to direction ``\isa{{\isasymLongleftarrow}}''. We assume the canonical \isa{LR{\isacharparenleft}{\kern0pt}{\isadigit{0}}{\isacharparenright}{\kern0pt}} automaton has no 
\isa{LR{\isacharparenleft}{\kern0pt}{\isadigit{0}}{\isacharparenright}{\kern0pt}}-inadequate states, must prove that the \isa{LR{\isacharparenleft}{\kern0pt}{\isadigit{0}}{\isacharparenright}{\kern0pt}} condition holds for \isa{G{\isacharprime}{\kern0pt}}. By 
Lemma~\ref{is_LR_wlogI}, consider rightmost derivations
\begin{subequations}
\begin{gather}
\isa{P{\isacharprime}{\kern0pt}\ {\isasymturnstile}\ {\isacharbrackleft}{\kern0pt}Nt\ S{\isacharprime}{\kern0pt}{\isacharbrackright}{\kern0pt}\ {\isasymRightarrow}r{\isacharasterisk}{\kern0pt}\ {\isasymalpha}\ {\isacharat}{\kern0pt}\ Nt\ X\ {\isacharhash}{\kern0pt}\ map\ Tm\ w\ {\isasymRightarrow}r\ {\isasymalpha}\ {\isacharat}{\kern0pt}\ {\isasymbeta}\ {\isacharat}{\kern0pt}\ map\ Tm\ w}\label{LR0_iff.X_derivers2}\\
\isa{P{\isacharprime}{\kern0pt}\ {\isasymturnstile}\ {\isacharbrackleft}{\kern0pt}Nt\ S{\isacharprime}{\kern0pt}{\isacharbrackright}{\kern0pt}\ {\isasymRightarrow}r{\isacharasterisk}{\kern0pt}\ {\isasymgamma}\ {\isacharat}{\kern0pt}\ Nt\ Y\ {\isacharhash}{\kern0pt}\ map\ Tm\ x\ {\isasymRightarrow}r\ {\isasymalpha}\ {\isacharat}{\kern0pt}\ {\isasymbeta}\ {\isacharat}{\kern0pt}\ map\ Tm\ y}\label{LR0_iff.Y_derivers2}\\
\intertext{such that the following hold}
\isa{\isafree{{\isasymgamma}}\ {\isacharat}{\kern0pt}\ \isafree{{\isasymdelta}}\ {\isacharat}{\kern0pt}\ \isaconst{map}\ \isaconst{Tm}\ \isafree{x}\ {\isacharequal}{\kern0pt}\ \isafree{{\isasymalpha}}\ {\isacharat}{\kern0pt}\ \isafree{{\isasymbeta}}\ {\isacharat}{\kern0pt}\ \isaconst{map}\ \isaconst{Tm}\ \isafree{y}}\label{LR0_iff.ydx}\\
\isa{{\isacharbar}{\kern0pt}\isafree{{\isasymalpha}}\ {\isacharat}{\kern0pt}\ \isafree{{\isasymbeta}}{\isacharbar}{\kern0pt}\ {\isasymle}\ {\isacharbar}{\kern0pt}\isafree{{\isasymgamma}}\ {\isacharat}{\kern0pt}\ \isafree{{\isasymdelta}}{\isacharbar}{\kern0pt}}\label{LR0_iff.leq}
\end{gather}
\end{subequations}
We have to show \isa{\isafree{{\isasymalpha}}\ {\isacharequal}{\kern0pt}\ \isafree{{\isasymgamma}}}, \isa{\isafree{X}\ {\isacharequal}{\kern0pt}\ \isafree{Y}}, and \isa{\isafree{x}\ {\isacharequal}{\kern0pt}\ \isafree{y}}. We omit the premise 
\isa{\isafree{w}{\isacharbar}{\kern0pt}\isactrlbsub {\isadigit{0}}\isactrlesub \ {\isacharequal}{\kern0pt}\ \isafree{y}{\isacharbar}{\kern0pt}\isactrlbsub {\isadigit{0}}\isactrlesub } since it is trivially true for any two lists.

Let \isa{p} be the state that \isa{\ensuremath{LR_0(G)}} reaches after reading input word \isa{\isafree{{\isasymalpha}}\ {\isacharat}{\kern0pt}\ \isafree{{\isasymbeta}}}. By 
Corollary~\ref{nextl_dfa_LR0_is_valids}, we know that this state is the set \isa{\isaconst{valids}\ {\isacharparenleft}{\kern0pt}\isafree{{\isasymalpha}}\ {\isacharat}{\kern0pt}\ \isafree{{\isasymbeta}}{\isacharparenright}{\kern0pt}}. 
With \eqref{LR0_iff.X_derivers2}, this implies that \isa{\mbox{{\isacharbrackleft}{\kern0pt}\isafree{X}\ {\isasymrightarrow}\ \isafree{{\isasymbeta}}\ {\isasymcdot}\ {\isacharbrackright}{\kern0pt}}\ {\isasymin}\ \isafree{p}}. Since by our initial 
assumption, \isa{p} is \isa{LR{\isacharparenleft}{\kern0pt}{\isadigit{0}}{\isacharparenright}{\kern0pt}}-adequate, we can now consider the three cases that
Lemma~\ref{LR0_adequate_cases} gives us.

If \isa{p} consists of a single complete item, \isa{\isafree{p}\ {\isacharequal}{\kern0pt}\ {\isacharbraceleft}{\kern0pt}{\isacharbrackleft}{\kern0pt}\isafree{X}\ {\isasymrightarrow}\ \isafree{{\isasymbeta}}\ {\isasymcdot}\ {\isacharbrackright}{\kern0pt}{\isacharbraceright}{\kern0pt}} must hold. Moreover, 
\eqref{LR0_iff.ydx} and \eqref{LR0_iff.leq} imply the existence of some \isa{z\ {\isacharcolon}{\kern0pt}{\isacharcolon}{\kern0pt}\ {\isacharprime}{\kern0pt}t} list such that
\begin{equation}\label{LR0_iff.z_defs}
\isa{\isafree{{\isasymgamma}}\ {\isacharat}{\kern0pt}\ \isafree{{\isasymdelta}}\ {\isacharequal}{\kern0pt}\ \isafree{{\isasymalpha}}\ {\isacharat}{\kern0pt}\ \isafree{{\isasymbeta}}\ {\isacharat}{\kern0pt}\ \isaconst{map}\ \isaconst{Tm}\ \isafree{z}} \text{ and } \isa{\isafree{z}\ {\isacharat}{\kern0pt}\ \isafree{x}\ {\isacharequal}{\kern0pt}\ \isafree{y}}.
\end{equation}

We can now distinguish two more cases based on whether \isa{{\isasymalpha}\ {\isacharat}{\kern0pt}\ {\isasymbeta}} is a prefix of \isa{{\isasymgamma}}.

If this is true, there exist \isa{u{\isacharcomma}{\kern0pt}\ v\ {\isacharcolon}{\kern0pt}{\isacharcolon}{\kern0pt}\ {\isacharprime}{\kern0pt}t\ list} such that \isa{{\isasymgamma}\ {\isacharequal}{\kern0pt}\ {\isasymalpha}\ {\isacharat}{\kern0pt}\ {\isasymbeta}\ {\isacharat}{\kern0pt}\ map\ Tm\ u}, \isa{{\isasymdelta}\ {\isacharequal}{\kern0pt}\ map\ Tm\ v}, and
\isa{z\ {\isacharequal}{\kern0pt}\ u\ {\isacharat}{\kern0pt}\ v}. Since from \eqref{LR0_iff.X_derivers2} we have 
\[ \isa{\isaconst{P{\isacharprime}{\kern0pt}}\ {\isasymturnstile}\ {\isacharbrackleft}{\kern0pt}\isaconst{Nt}\ \isaconst{S{\isacharprime}{\kern0pt}}{\isacharbrackright}{\kern0pt}\ {\isasymRightarrow}r{\isacharasterisk}{\kern0pt}\ \isafree{{\isasymgamma}}\ {\isacharat}{\kern0pt}\ \isaconst{Nt}\ \isafree{Y}\ {\isacharhash}{\kern0pt}\ \isaconst{map}\ \isaconst{Tm}\ \isafree{x}}, \] 
we can substitute \isa{\isafree{{\isasymgamma}}\ {\isacharequal}{\kern0pt}\ \isafree{{\isasymalpha}}\ {\isacharat}{\kern0pt}\ \isafree{{\isasymbeta}}\ {\isacharat}{\kern0pt}\ \isaconst{map}\ \isaconst{Tm}\ \isafree{u}}. By Lemma~\ref{derivers_substring_reliable} this implies the 
existence of some incomplete item in \isa{\isaconst{valids}\ {\isacharparenleft}{\kern0pt}\isafree{{\isasymalpha}}\ {\isacharat}{\kern0pt}\ \isafree{{\isasymbeta}}{\isacharparenright}{\kern0pt}}, contradicting the fact that
\isa{\isafree{p}\ {\isacharequal}{\kern0pt}\ {\isacharbraceleft}{\kern0pt}{\isacharbrackleft}{\kern0pt}\isafree{X}\ {\isasymrightarrow}\ \isafree{{\isasymbeta}}\ {\isasymcdot}\ {\isacharbrackright}{\kern0pt}{\isacharbraceright}{\kern0pt}}.

Otherwise, if \isa{{\isasymalpha}\ {\isacharat}{\kern0pt}\ {\isasymbeta}} is not a substring of \isa{{\isasymgamma}}, \eqref{LR0_iff.z_defs} implies the existence of 
some \isa{{\isasymdelta}{\isacharprime}{\kern0pt}\ {\isacharcolon}{\kern0pt}{\isacharcolon}{\kern0pt}\ syms} such that \isa{\isafree{{\isasymgamma}}\ {\isacharat}{\kern0pt}\ \isafree{{\isasymdelta}{\isacharprime}{\kern0pt}}\ {\isacharequal}{\kern0pt}\ \isafree{{\isasymalpha}}\ {\isacharat}{\kern0pt}\ \isafree{{\isasymbeta}}} and \isa{\isafree{{\isasymdelta}}\ {\isacharequal}{\kern0pt}\ \isafree{{\isasymdelta}{\isacharprime}{\kern0pt}}\ {\isacharat}{\kern0pt}\ \isaconst{map}\ \isaconst{Tm}\ \isafree{z}}. This implies 
that item \isa{{\isacharbrackleft}{\kern0pt}\isafree{Y}\ {\isasymrightarrow}\ \isafree{{\isasymdelta}{\isacharprime}{\kern0pt}}\ {\isasymcdot}\ \isaconst{map}\ \isaconst{Tm}\ \isafree{z}{\isacharbrackright}{\kern0pt}} is valid for \isa{\isafree{{\isasymalpha}}\ {\isacharat}{\kern0pt}\ \isafree{{\isasymbeta}}} by \eqref{LR0_iff.Y_derivers2}, 
\eqref{LR0_iff.ydx}, and the fact that \isa{\isafree{z}\ {\isacharat}{\kern0pt}\ \isafree{x}\ {\isacharequal}{\kern0pt}\ \isafree{y}}, meaning it is in \isa{p}. This and 
\isa{\isafree{p}\ {\isacharequal}{\kern0pt}\ {\isacharbraceleft}{\kern0pt}{\isacharbrackleft}{\kern0pt}\isafree{X}\ {\isasymrightarrow}\ \isafree{{\isasymbeta}}\ {\isasymcdot}\ {\isacharbrackright}{\kern0pt}{\isacharbraceright}{\kern0pt}} in turn imply \isa{\isafree{{\isasymalpha}}\ {\isacharequal}{\kern0pt}\ \isafree{{\isasymgamma}}\ {\isasymand}\ \isafree{X}\ {\isacharequal}{\kern0pt}\ \isafree{Y}\ {\isasymand}\ \isafree{x}\ {\isacharequal}{\kern0pt}\ \isafree{y}}, meaning the implication holds. % clarify?
This completes the case where \isa{p} consists of a single complete item.

In the case where \isa{\isaconst{completes}\ \isafree{p}\ {\isacharequal}{\kern0pt}\ {\isacharbraceleft}{\kern0pt}{\isacharbrackleft}{\kern0pt}\isafree{A}\ {\isasymrightarrow}\ {\isasymcdot}\ {\isacharbrackright}{\kern0pt}{\isacharbraceright}{\kern0pt}}, the fact that \isa{{\isacharbrackleft}{\kern0pt}\isafree{X}\ {\isasymrightarrow}\ \isafree{{\isasymbeta}}\ {\isasymcdot}\ {\isacharbrackright}{\kern0pt}\ {\isasymin}\ \isafree{p}} forces
\isa{{\isacharbrackleft}{\kern0pt}\isafree{A}\ {\isasymrightarrow}\ {\isasymcdot}\ {\isacharbrackright}{\kern0pt}\ {\isacharequal}{\kern0pt}\ {\isacharbrackleft}{\kern0pt}\isafree{X}\ {\isasymrightarrow}\ \isafree{{\isasymbeta}}\ {\isasymcdot}\ {\isacharbrackright}{\kern0pt}}. From \eqref{LR0_iff.ydx} we can now distinguish cases based on
whether \isa{\isafree{{\isasymalpha}}\ {\isacharat}{\kern0pt}\ \isafree{{\isasymbeta}}} is a prefix of \isa{{\isasymgamma}}.

If this is the case, there exists some \isa{{\isasymzeta}} such that \isa{\isafree{{\isasymgamma}}\ {\isacharequal}{\kern0pt}\ \isafree{{\isasymalpha}}\ {\isacharat}{\kern0pt}\ \isafree{{\isasymbeta}}\ {\isacharat}{\kern0pt}\ \isafree{{\isasymzeta}}}, and we can rewrite 
\eqref{LR0_iff.Y_derivers2} as \isa{P{\isacharprime}{\kern0pt}\ {\isasymturnstile}\ {\isacharbrackleft}{\kern0pt}Nt\ S{\isacharprime}{\kern0pt}{\isacharbrackright}{\kern0pt}\ {\isasymRightarrow}r{\isacharasterisk}{\kern0pt}\ {\isacharparenleft}{\kern0pt}{\isasymalpha}\ {\isacharat}{\kern0pt}\ {\isasymbeta}{\isacharparenright}{\kern0pt}\ {\isacharat}{\kern0pt}\ {\isasymzeta}\ {\isacharat}{\kern0pt}\ Nt\ Y\ {\isacharhash}{\kern0pt}\ map\ Tm\ x\ {\isasymRightarrow}r\ {\isasymalpha}\ {\isacharat}{\kern0pt}\ {\isasymbeta}\ {\isacharat}{\kern0pt}\ map\ Tm\ y}.
By Lemma~\ref{prefix_comp_unique_imp_eps} this implies that \isa{\isafree{{\isasymzeta}}\ {\isacharequal}{\kern0pt}\ {\isacharbrackleft}{\kern0pt}{\isacharbrackright}{\kern0pt}\ {\isasymand}\ \isafree{X}\ {\isacharequal}{\kern0pt}\ \isafree{Y}\ {\isasymand}\ \isafree{x}\ {\isacharequal}{\kern0pt}\ \isafree{y}} holds,
which trivially implies \isa{{\isasymalpha}\ {\isacharequal}{\kern0pt}\ {\isasymgamma}} since \isa{{\isasymbeta}\ {\isacharequal}{\kern0pt}\ {\isacharbrackleft}{\kern0pt}{\isacharbrackright}{\kern0pt}}.

Lastly, if \isa{\isafree{{\isasymalpha}}\ {\isacharat}{\kern0pt}\ \isafree{{\isasymbeta}}} is not a prefix of \isa{{\isasymgamma}}, there exist \isa{{\isasymzeta}\ {\isacharcolon}{\kern0pt}{\isacharcolon}{\kern0pt}\ syms} and \isa{x{\isacharprime}{\kern0pt}\ {\isacharcolon}{\kern0pt}{\isacharcolon}{\kern0pt}\ {\isacharprime}{\kern0pt}t\ list} such 
that \isa{\isafree{{\isasymalpha}}\ {\isacharat}{\kern0pt}\ \isafree{{\isasymbeta}}\ {\isacharequal}{\kern0pt}\ \isafree{{\isasymgamma}}\ {\isacharat}{\kern0pt}\ \isafree{{\isasymzeta}}}, \isa{\isafree{{\isasymdelta}}\ {\isacharequal}{\kern0pt}\ \isafree{{\isasymzeta}}\ {\isacharat}{\kern0pt}\ \isaconst{map}\ \isaconst{Tm}\ \isafree{x{\isacharprime}{\kern0pt}}}, and \isa{\isafree{y}\ {\isacharequal}{\kern0pt}\ \isafree{x{\isacharprime}{\kern0pt}}\ {\isacharat}{\kern0pt}\ \isafree{x}}. We can then 
perform a final case distinction on \isa{x{\isacharprime}{\kern0pt}}, where the case \isa{\isafree{x{\isacharprime}{\kern0pt}}\ {\isacharequal}{\kern0pt}\ {\isacharbrackleft}{\kern0pt}{\isacharbrackright}{\kern0pt}} forces 
\isa{{\isacharbrackleft}{\kern0pt}\isafree{Y}\ {\isasymrightarrow}\ \isafree{{\isasymzeta}}\ {\isasymcdot}\ \isaconst{map}\ \isaconst{Tm}\ \isafree{x{\isacharprime}{\kern0pt}}{\isacharbrackright}{\kern0pt}\ {\isacharequal}{\kern0pt}\ {\isacharbrackleft}{\kern0pt}\isafree{X}\ {\isasymrightarrow}\ {\isasymcdot}\ {\isacharbrackright}{\kern0pt}}, and \isa{\isafree{x{\isacharprime}{\kern0pt}}\ {\isasymnoteq}\ {\isacharbrackleft}{\kern0pt}{\isacharbrackright}{\kern0pt}} implies a contradiction by 
Lemma~\ref{LR0_adequate_cases}. This completes the proof of the Theorem.
\end{proof}
\end{theorem}

We have now proved the main theorem of the \isa{LR{\isacharparenleft}{\kern0pt}{\isadigit{0}}{\isacharparenright}{\kern0pt}} section of the book. Our proof 
of the first implication, 
``\textit{\isa{G{\isacharprime}{\kern0pt}} is \isa{LR{\isacharparenleft}{\kern0pt}{\isadigit{0}}{\isacharparenright}{\kern0pt}} \isa{{\isasymLongrightarrow}} \isa{\ensuremath{LR_0(G)}} has no \isa{LR{\isacharparenleft}{\kern0pt}{\isadigit{0}}{\isacharparenright}{\kern0pt}}-inadequate states}'', is 
quite similar to the one given by the authors. Their proof of the converse, however, is incorrect.
Recall the original rightmost derivations as stated in the proof:
\begin{subequations}
\begin{gather}
\isa{P{\isacharprime}{\kern0pt}\ {\isasymturnstile}\ {\isacharbrackleft}{\kern0pt}Nt\ S{\isacharprime}{\kern0pt}{\isacharbrackright}{\kern0pt}\ {\isasymRightarrow}r{\isacharasterisk}{\kern0pt}\ {\isasymalpha}\ {\isacharat}{\kern0pt}\ Nt\ X\ {\isacharhash}{\kern0pt}\ map\ Tm\ w\ {\isasymRightarrow}r\ {\isasymalpha}\ {\isacharat}{\kern0pt}\ {\isasymbeta}\ {\isacharat}{\kern0pt}\ map\ Tm\ w}\label{iff_disc1}\\
\isa{P{\isacharprime}{\kern0pt}\ {\isasymturnstile}\ {\isacharbrackleft}{\kern0pt}Nt\ S{\isacharprime}{\kern0pt}{\isacharbrackright}{\kern0pt}\ {\isasymRightarrow}r{\isacharasterisk}{\kern0pt}\ {\isasymgamma}\ {\isacharat}{\kern0pt}\ Nt\ Y\ {\isacharhash}{\kern0pt}\ map\ Tm\ x\ {\isasymRightarrow}r\ {\isasymalpha}\ {\isacharat}{\kern0pt}\ {\isasymbeta}\ {\isacharat}{\kern0pt}\ map\ Tm\ y}\label{iff_disc2}
\end{gather}
\end{subequations}
Instead of distinguishing on the 3 cases provided by Lemma~\ref{LR0_adequate_cases} like we did,
they distinguish cases on whether \isa{{\isasymbeta}} is empty. In both cases, they seem to implicitly assume that 
the final step in \eqref{iff_disc2} implies that there exists a valid item for \isa{\isafree{{\isasymalpha}}\ {\isacharat}{\kern0pt}\ \isafree{{\isasymbeta}}} of
the form \isa{{\isacharbrackleft}{\kern0pt}\isafree{Y}\ {\isasymrightarrow}\ \isafree{{\isasymzeta}}\ {\isasymcdot}\ \isafree{{\isasymeta}}{\isacharbrackright}{\kern0pt}}.

This does not hold in general; consider the case where the handle of \isa{\isafree{{\isasymalpha}}\ {\isacharat}{\kern0pt}\ \isafree{{\isasymbeta}}\ {\isacharat}{\kern0pt}\ \isaconst{map}\ \isaconst{Tm}\ \isafree{y}}
is some \isa{{\isasymdelta}} with \isa{\isafree{{\isasymalpha}}\ {\isacharequal}{\kern0pt}\ \isafree{{\isasymgamma}}\ {\isacharat}{\kern0pt}\ \isafree{{\isasymdelta}}}, and \isa{\isaconst{map}\ \isaconst{Tm}\ \isafree{x}\ {\isacharequal}{\kern0pt}\ \isafree{{\isasymbeta}}\ {\isacharat}{\kern0pt}\ \isaconst{map}\ \isaconst{Tm}\ \isafree{y}}. In such a case, the item 
\isa{{\isacharbrackleft}{\kern0pt}\isafree{Y}\ {\isasymrightarrow}\ \isafree{{\isasymdelta}}\ {\isasymcdot}\ {\isacharbrackright}{\kern0pt}} would be valid for \isa{{\isasymalpha}}, which does not imply
\isa{{\isacharbrackleft}{\kern0pt}\isafree{Y}\ {\isasymrightarrow}\ \isafree{{\isasymdelta}}\ {\isasymcdot}\ {\isacharbrackright}{\kern0pt}\ {\isasymin}\ \isaconst{valids}\ {\isacharparenleft}{\kern0pt}\isafree{{\isasymalpha}}\ {\isacharat}{\kern0pt}\ \isafree{{\isasymbeta}}{\isacharparenright}{\kern0pt}} in general. The same issue can arise in other cases. In
particular, if \isa{\isafree{{\isasymalpha}}\ {\isacharat}{\kern0pt}\ \isafree{{\isasymbeta}}} is a proper prefix of \isa{{\isasymgamma}}, one cannot use \eqref{iff_disc2} to infer
the existence of any item of the form \isa{{\isacharbrackleft}{\kern0pt}\isafree{Y}\ {\isasymrightarrow}\ \isafree{{\isasymzeta}}\ {\isasymcdot}\ \isafree{{\isasymeta}}{\isacharbrackright}{\kern0pt}} that is valid for \isa{{\isasymalpha}\ {\isacharat}{\kern0pt}\ {\isasymbeta}} by the
definition of reliable prefixes: if there exists a nonempty string \isa{{\isasymtheta}} with \isa{\isafree{{\isasymgamma}}\ {\isacharequal}{\kern0pt}\ \isafree{{\isasymalpha}}\ {\isacharat}{\kern0pt}\ \isafree{{\isasymbeta}}\ {\isacharat}{\kern0pt}\ \isafree{{\isasymtheta}}},
\eqref{iff_disc2} can only yield an item of the form \isa{{\isacharbrackleft}{\kern0pt}\isafree{Y}\ {\isasymrightarrow}\ \isafree{{\isasymzeta}}\ {\isasymcdot}\ \isafree{{\isasymeta}}{\isacharbrackright}{\kern0pt}} valid for
the string \isa{\isafree{{\isasymalpha}}\ {\isacharat}{\kern0pt}\ \isafree{{\isasymbeta}}\ {\isacharat}{\kern0pt}\ \isafree{{\isasymtheta}}\ {\isacharat}{\kern0pt}\ \isafree{{\isasymzeta}}}. For this reason, our argument for the case where 
\isa{p} consists of a single complete item and \isa{\isafree{{\isasymalpha}}\ {\isacharat}{\kern0pt}\ \isafree{{\isasymbeta}}} is a prefix of \isa{{\isasymgamma}} is based on
Lemma~\ref{derivers_substring_reliable}, which derives a contradiction independently of \isa{Y}. 

In general, the issues caused by approaching the proof through a case distinction on \isa{{\isasymbeta}} are mostly
caused by the fact that if \isa{\isafree{{\isasymgamma}}\ {\isacharat}{\kern0pt}\ \isafree{{\isasymdelta}}\ {\isacharat}{\kern0pt}\ \isaconst{map}\ \isaconst{Tm}\ \isafree{x}\ {\isacharequal}{\kern0pt}\ \isafree{{\isasymalpha}}\ {\isacharat}{\kern0pt}\ \isafree{{\isasymbeta}}\ {\isacharat}{\kern0pt}\ \isaconst{map}\ \isaconst{Tm}\ \isafree{y}} where \isa{{\isasymdelta}} is
the handle of \isa{{\isasymalpha}\ {\isacharat}{\kern0pt}\ {\isasymbeta}\ {\isacharat}{\kern0pt}\ map\ Tm\ y} in \eqref{iff_disc2}, there are far too many cases to consider
as to how the lists are split, and the cases where \isa{{\isasymgamma}\ {\isacharat}{\kern0pt}\ {\isasymdelta}} is a proper prefix of \isa{\isafree{{\isasymalpha}}\ {\isacharat}{\kern0pt}\ \isafree{{\isasymbeta}}} are 
particularly problematic. This was the motivation towards assuming more well-behaved rightmost 
derivations WLOG, where it is guaranteed that either \isa{\isafree{{\isasymalpha}}\ {\isacharat}{\kern0pt}\ \isafree{{\isasymbeta}}} was already previously derived, 
allowing us to use Lemma~\ref{derivers_substring_reliable}, or a prefix of \isa{{\isasymdelta}} is a suffix of \isa{{\isasymalpha}\ {\isacharat}{\kern0pt}\ {\isasymbeta}}, 
in which case we are guaranteed that an item corresponding to production \isa{{\isacharparenleft}{\kern0pt}\isafree{Y}{\isacharcomma}{\kern0pt}\ \isafree{{\isasymdelta}}{\isacharparenright}{\kern0pt}} is valid 
for \isa{\isafree{{\isasymalpha}}\ {\isacharat}{\kern0pt}\ \isafree{{\isasymbeta}}}, allowing us to argue about \isa{p}.%
\end{isamarkuptext}\isamarkuptrue%
%
\isadelimdocument
%
\endisadelimdocument
%
\isatagdocument
%
\isamarkupsubsection{Preservation of the \isa{LR{\isacharparenleft}{\kern0pt}k{\isacharparenright}{\kern0pt}} Condition in Extended Grammars%
}
\isamarkuptrue%
%
\endisatagdocument
{\isafolddocument}%
%
\isadelimdocument
%
\endisadelimdocument
%
\begin{isamarkuptext}%
There is another aspect related to \isa{LR{\isacharparenleft}{\kern0pt}k{\isacharparenright}{\kern0pt}} grammars worth discussing, namely the definition of the 
\isa{LR{\isacharparenleft}{\kern0pt}k{\isacharparenright}{\kern0pt}} condition itself. As we stated before, our definition, which is based on that of Wilhelm et 
al., is restricted to the extension of a CFG. Since extending a grammar preserves several properties, 
such as the language and quality of being reduced, as we showed in Theorem~\ref{Lang_preserved} and
Lemma~\ref{G'_reduced}, one could intuitively assume that the extension has no effect on the original
grammar's satisfaction of the \isa{LR{\isacharparenleft}{\kern0pt}k{\isacharparenright}{\kern0pt}} condition. In our formalization of \isa{LR{\isacharparenleft}{\kern0pt}k{\isacharparenright}{\kern0pt}}, we remained working
within the context of our extended grammars, which fixes \isa{G{\isacharprime}{\kern0pt}} as the extension of the original CFG \isa{G}.
Since our grammars are fixed, we defined the \isa{LR{\isacharparenleft}{\kern0pt}k{\isacharparenright}{\kern0pt}} condition per Wilhelm et al.\ as the predicate 
\isa{is{\isacharunderscore}{\kern0pt}LR\ {\isacharcolon}{\kern0pt}{\isacharcolon}{\kern0pt}\ nat\ {\isasymRightarrow}\ bool}:
\begin{quote}
%
\begin{isabelle}%
\isaconst{is{\isacharunderscore}{\kern0pt}LR}\ \isavar{k}\ {\isasymequiv}\isanewline
{\isasymforall}\isabound{{\isasymalpha}}\ \isabound{X}\ \isabound{{\isasymbeta}}\ \isabound{w}\ \isabound{{\isasymgamma}}\ \isabound{Y}\ \isabound{x}\ \isabound{y}{\isachardot}{\kern0pt}\isanewline
\isaindent{\ \ \ }\isaconst{P{\isacharprime}{\kern0pt}}\ {\isasymturnstile}\ {\isacharbrackleft}{\kern0pt}\isaconst{Nt}\ \isaconst{S{\isacharprime}{\kern0pt}}{\isacharbrackright}{\kern0pt}\ {\isasymRightarrow}r{\isacharasterisk}{\kern0pt}\ \isabound{{\isasymalpha}}\ {\isacharat}{\kern0pt}\ \isaconst{Nt}\ \isabound{X}\ {\isacharhash}{\kern0pt}\ \isaconst{map}\ \isaconst{Tm}\ \isabound{w}\ {\isasymand}\isanewline
\isaindent{\ \ \ }\isaconst{P{\isacharprime}{\kern0pt}}\ {\isasymturnstile}\ \isabound{{\isasymalpha}}\ {\isacharat}{\kern0pt}\ \isaconst{Nt}\ \isabound{X}\ {\isacharhash}{\kern0pt}\ \isaconst{map}\ \isaconst{Tm}\ \isabound{w}\ {\isasymRightarrow}r\ \isabound{{\isasymalpha}}\ {\isacharat}{\kern0pt}\ \isabound{{\isasymbeta}}\ {\isacharat}{\kern0pt}\ \isaconst{map}\ \isaconst{Tm}\ \isabound{w}\ {\isasymand}\isanewline
\isaindent{\ \ \ }\isaconst{P{\isacharprime}{\kern0pt}}\ {\isasymturnstile}\ {\isacharbrackleft}{\kern0pt}\isaconst{Nt}\ \isaconst{S{\isacharprime}{\kern0pt}}{\isacharbrackright}{\kern0pt}\ {\isasymRightarrow}r{\isacharasterisk}{\kern0pt}\ \isabound{{\isasymgamma}}\ {\isacharat}{\kern0pt}\ \isaconst{Nt}\ \isabound{Y}\ {\isacharhash}{\kern0pt}\ \isaconst{map}\ \isaconst{Tm}\ \isabound{x}\ {\isasymand}\isanewline
\isaindent{\ \ \ }\isaconst{P{\isacharprime}{\kern0pt}}\ {\isasymturnstile}\ \isabound{{\isasymgamma}}\ {\isacharat}{\kern0pt}\ \isaconst{Nt}\ \isabound{Y}\ {\isacharhash}{\kern0pt}\ \isaconst{map}\ \isaconst{Tm}\ \isabound{x}\ {\isasymRightarrow}r\ \isabound{{\isasymalpha}}\ {\isacharat}{\kern0pt}\ \isabound{{\isasymbeta}}\ {\isacharat}{\kern0pt}\ \isaconst{map}\ \isaconst{Tm}\ \isabound{y}\ {\isasymand}\ \isabound{w}{\isacharbar}{\kern0pt}\isactrlbsub \isavar{k}\isactrlesub \ {\isacharequal}{\kern0pt}\ \isabound{y}{\isacharbar}{\kern0pt}\isactrlbsub \isavar{k}\isactrlesub \ {\isasymlongrightarrow}\isanewline
\isaindent{\ \ \ }\isabound{{\isasymalpha}}\ {\isacharequal}{\kern0pt}\ \isabound{{\isasymgamma}}\ {\isasymand}\ \isabound{X}\ {\isacharequal}{\kern0pt}\ \isabound{Y}\ {\isasymand}\ \isabound{x}\ {\isacharequal}{\kern0pt}\ \isabound{y}%
\end{isabelle}
\end{quote}
Consider now the more general \isa{is{\isacharunderscore}{\kern0pt}LR{\isacharprime}{\kern0pt}\ {\isacharcolon}{\kern0pt}{\isacharcolon}{\kern0pt}\ nat\ {\isasymRightarrow}\ {\isacharparenleft}{\kern0pt}{\isacharprime}{\kern0pt}a{\isacharcomma}{\kern0pt}\ {\isacharprime}{\kern0pt}b{\isacharparenright}{\kern0pt}\ Cfg\ {\isasymRightarrow}\ bool}:
\begin{quote}
%
\begin{isabelle}%
\isaconst{is{\isacharunderscore}{\kern0pt}LR{\isacharprime}{\kern0pt}}\ \isavar{k}\ \isavar{G}\ {\isasymequiv}\isanewline
{\isasymforall}\isabound{{\isasymalpha}}\ \isabound{X}\ \isabound{{\isasymbeta}}\ \isabound{w}\ \isabound{{\isasymgamma}}\ \isabound{Y}\ \isabound{x}\ \isabound{y}{\isachardot}{\kern0pt}\isanewline
\isaindent{\ \ \ }\isaconst{Prods}\ \isavar{G}\ {\isasymturnstile}\ {\isacharbrackleft}{\kern0pt}\isaconst{Nt}\ {\isacharparenleft}{\kern0pt}\isaconst{Start}\ \isavar{G}{\isacharparenright}{\kern0pt}{\isacharbrackright}{\kern0pt}\ {\isasymRightarrow}r{\isacharasterisk}{\kern0pt}\ \isabound{{\isasymalpha}}\ {\isacharat}{\kern0pt}\ \isaconst{Nt}\ \isabound{X}\ {\isacharhash}{\kern0pt}\ \isaconst{map}\ \isaconst{Tm}\ \isabound{w}\ {\isasymand}\isanewline
\isaindent{\ \ \ }\isaconst{Prods}\ \isavar{G}\ {\isasymturnstile}\ \isabound{{\isasymalpha}}\ {\isacharat}{\kern0pt}\ \isaconst{Nt}\ \isabound{X}\ {\isacharhash}{\kern0pt}\ \isaconst{map}\ \isaconst{Tm}\ \isabound{w}\ {\isasymRightarrow}r\ \isabound{{\isasymalpha}}\ {\isacharat}{\kern0pt}\ \isabound{{\isasymbeta}}\ {\isacharat}{\kern0pt}\ \isaconst{map}\ \isaconst{Tm}\ \isabound{w}\ {\isasymand}\isanewline
\isaindent{\ \ \ }\isaconst{Prods}\ \isavar{G}\ {\isasymturnstile}\ {\isacharbrackleft}{\kern0pt}\isaconst{Nt}\ {\isacharparenleft}{\kern0pt}\isaconst{Start}\ \isavar{G}{\isacharparenright}{\kern0pt}{\isacharbrackright}{\kern0pt}\ {\isasymRightarrow}r{\isacharasterisk}{\kern0pt}\ \isabound{{\isasymgamma}}\ {\isacharat}{\kern0pt}\ \isaconst{Nt}\ \isabound{Y}\ {\isacharhash}{\kern0pt}\ \isaconst{map}\ \isaconst{Tm}\ \isabound{x}\ {\isasymand}\isanewline
\isaindent{\ \ \ }\isaconst{Prods}\ \isavar{G}\ {\isasymturnstile}\ \isabound{{\isasymgamma}}\ {\isacharat}{\kern0pt}\ \isaconst{Nt}\ \isabound{Y}\ {\isacharhash}{\kern0pt}\ \isaconst{map}\ \isaconst{Tm}\ \isabound{x}\ {\isasymRightarrow}r\ \isabound{{\isasymalpha}}\ {\isacharat}{\kern0pt}\ \isabound{{\isasymbeta}}\ {\isacharat}{\kern0pt}\ \isaconst{map}\ \isaconst{Tm}\ \isabound{y}\ {\isasymand}\ \isabound{w}{\isacharbar}{\kern0pt}\isactrlbsub \isavar{k}\isactrlesub \ {\isacharequal}{\kern0pt}\ \isabound{y}{\isacharbar}{\kern0pt}\isactrlbsub \isavar{k}\isactrlesub \ {\isasymlongrightarrow}\isanewline
\isaindent{\ \ \ }\isabound{{\isasymalpha}}\ {\isacharequal}{\kern0pt}\ \isabound{{\isasymgamma}}\ {\isasymand}\ \isabound{X}\ {\isacharequal}{\kern0pt}\ \isabound{Y}\ {\isasymand}\ \isabound{x}\ {\isacharequal}{\kern0pt}\ \isabound{y}%
\end{isabelle}
\end{quote}
This generalized definition of \isa{LR{\isacharparenleft}{\kern0pt}k{\isacharparenright}{\kern0pt}} for arbitrary CFGs is not preserved by extension, i.e., for 
a CFG \isa{G} and its extension \isa{G{\isacharprime}{\kern0pt}}, the equality 
$\isa{\isaconst{is{\isacharunderscore}{\kern0pt}LR{\isacharprime}{\kern0pt}}\ \isafree{k}\ \isafree{G}} \overset{?}{=} \isa{\isaconst{is{\isacharunderscore}{\kern0pt}LR{\isacharprime}{\kern0pt}}\ \isafree{k}\ \isaconst{G{\isacharprime}{\kern0pt}}}$
does not hold in general. We will prove two necessary and sufficient conditions that show that for
grammars with certain properties, the extension does preserve the \isa{LR{\isacharparenleft}{\kern0pt}k{\isacharparenright}{\kern0pt}} condition.

One direction of the claim does hold for all \isa{k}:

\begin{lemma}\label{LR'k_G'_imp_LR'k_G}
Let \isa{G} be a CFG and \isa{G{\isacharprime}{\kern0pt}} its extension. For all \isa{k} holds
\[ \isa{\isaconst{is{\isacharunderscore}{\kern0pt}LR{\isacharprime}{\kern0pt}}\ \isavar{k}\ \isaconst{G{\isacharprime}{\kern0pt}}\ {\isasymLongrightarrow}\ \isaconst{is{\isacharunderscore}{\kern0pt}LR{\isacharprime}{\kern0pt}}\ \isavar{k}\ \isafree{G}} \]
\begin{proof}
Trivial using the fact that \isa{\isaconst{P}\ {\isasymsubseteq}\ \isaconst{P{\isacharprime}{\kern0pt}}}.
\end{proof}
\end{lemma}

We now show the converse for \isa{k\ {\isacharequal}{\kern0pt}\ {\isadigit{0}}} and \isa{k\ {\isachargreater}{\kern0pt}\ {\isadigit{0}}} separately. For the remainder of this section, let 
\isa{G} again be a reduced grammar with finite productions \isa{P}, start symbol \isa{S} and nonempty language, 
and \isa{G{\isacharprime}{\kern0pt}} its extension with start symbol \isa{S{\isacharprime}{\kern0pt}} and productions \isa{P{\isacharprime}{\kern0pt}}. 

\begin{lemma}[Preservation of \isa{LR{\isacharparenleft}{\kern0pt}{\isadigit{0}}{\isacharparenright}{\kern0pt}}]\label{LR'0_G_imp_LR'0_G'}
Assume \isa{G} is an \isa{LR{\isacharparenleft}{\kern0pt}{\isadigit{0}}{\isacharparenright}{\kern0pt}} grammar. Furthermore, assume there does not exist a left recursion of the
form \isa{\isaconst{P}\ {\isasymturnstile}\ {\isacharbrackleft}{\kern0pt}\isaconst{Nt}\ \isaconst{S}{\isacharbrackright}{\kern0pt}\ {\isasymRightarrow}r{\isacharparenleft}{\kern0pt}\isaconst{Suc}\ \isafree{n}{\isacharparenright}{\kern0pt}\ \isaconst{Nt}\ \isaconst{S}\ {\isacharhash}{\kern0pt}\ \isafree{{\isasymalpha}}}.
for any \isa{n\ {\isacharcolon}{\kern0pt}{\isacharcolon}{\kern0pt}\ nat} or \isa{{\isasymalpha}\ {\isacharcolon}{\kern0pt}{\isacharcolon}{\kern0pt}\ syms}. Then \isa{G{\isacharprime}{\kern0pt}} is an \isa{LR{\isacharparenleft}{\kern0pt}{\isadigit{0}}{\isacharparenright}{\kern0pt}} grammar.
\begin{proof}
Consider the rightmost derivations 
\begin{gather*}
\isa{P{\isacharprime}{\kern0pt}\ {\isasymturnstile}\ {\isacharbrackleft}{\kern0pt}Nt\ S{\isacharprime}{\kern0pt}{\isacharbrackright}{\kern0pt}\ {\isasymRightarrow}r{\isacharparenleft}{\kern0pt}n{\isacharparenright}{\kern0pt}\ {\isasymalpha}\ {\isacharat}{\kern0pt}\ Nt\ X\ {\isacharhash}{\kern0pt}\ map\ Tm\ w}
  \isa{{\isasymRightarrow}r\ {\isasymalpha}\ {\isacharat}{\kern0pt}\ {\isasymbeta}\ {\isacharat}{\kern0pt}\ map\ Tm\ w}\\
\isa{P{\isacharprime}{\kern0pt}\ {\isasymturnstile}\ {\isacharbrackleft}{\kern0pt}Nt\ S{\isacharprime}{\kern0pt}{\isacharbrackright}{\kern0pt}\ {\isasymRightarrow}r{\isacharparenleft}{\kern0pt}m{\isacharparenright}{\kern0pt}\ {\isasymgamma}\ {\isacharat}{\kern0pt}\ Nt\ Y\ {\isacharhash}{\kern0pt}\ map\ Tm\ x}
  \isa{{\isasymRightarrow}r\ {\isasymalpha}\ {\isacharat}{\kern0pt}\ {\isasymbeta}\ {\isacharat}{\kern0pt}\ map\ Tm\ y}
\end{gather*}
We have to prove that \isa{\isafree{{\isasymalpha}}\ {\isacharequal}{\kern0pt}\ \isafree{{\isasymgamma}}\ {\isasymand}\ \isafree{X}\ {\isacharequal}{\kern0pt}\ \isafree{Y}\ {\isasymand}\ \isafree{x}\ {\isacharequal}{\kern0pt}\ \isafree{y}} holds. We again omit the assumption 
\isa{\isafree{w}{\isacharbar}{\kern0pt}\isactrlbsub {\isadigit{0}}\isactrlesub \ {\isacharequal}{\kern0pt}\ \isafree{y}{\isacharbar}{\kern0pt}\isactrlbsub {\isadigit{0}}\isactrlesub } since it is trivially true.

We distinguish cases on the lengths of the derivations \isa{m} and \isa{n}.

If \isa{n\ {\isacharequal}{\kern0pt}\ m\ {\isacharequal}{\kern0pt}\ {\isadigit{0}}}, the implication holds trivially.

If \isa{n\ {\isacharequal}{\kern0pt}\ {\isadigit{0}}} and \isa{m\ {\isacharequal}{\kern0pt}\ Suc\ m{\isacharprime}{\kern0pt}} for some \isa{m{\isacharprime}{\kern0pt}}, we know from \isa{n\ {\isacharequal}{\kern0pt}\ {\isadigit{0}}} that \isa{{\isasymalpha}\ {\isacharequal}{\kern0pt}\ {\isacharbrackleft}{\kern0pt}{\isacharbrackright}{\kern0pt}}, \isa{{\isasymbeta}\ {\isacharequal}{\kern0pt}\ {\isacharbrackleft}{\kern0pt}Nt\ S{\isacharbrackright}{\kern0pt}}, and 
\isa{w\ {\isacharequal}{\kern0pt}\ {\isacharbrackleft}{\kern0pt}{\isacharbrackright}{\kern0pt}}.
With Lemma~\ref{G'_derivern_Suc_imp_G_derivern}, this implies 
\isa{\isaconst{P{\isacharprime}{\kern0pt}}\ {\isasymturnstile}\ {\isacharbrackleft}{\kern0pt}\isaconst{Nt}\ \isaconst{S{\isacharprime}{\kern0pt}}{\isacharbrackright}{\kern0pt}\ {\isasymRightarrow}r{\isacharparenleft}{\kern0pt}\isaconst{Suc}\ {\isacharparenleft}{\kern0pt}\isaconst{Suc}\ \isafree{m{\isacharprime}{\kern0pt}}{\isacharparenright}{\kern0pt}{\isacharparenright}{\kern0pt}\ \isaconst{Nt}\ \isaconst{S}\ {\isacharhash}{\kern0pt}\ \isaconst{map}\ \isaconst{Tm}\ \isafree{y}}, which contradicts our assumption about
the absence of this kind of left recursions for \isa{S} in \isa{G}.

The opposite case, \isa{n\ {\isacharequal}{\kern0pt}\ Suc\ n{\isacharprime}{\kern0pt}} for some \isa{n{\isacharprime}{\kern0pt}} and \isa{m\ {\isacharequal}{\kern0pt}\ {\isadigit{0}}}, is analogous.

Finally, the case where both derivations have nonzero length is trivial by
Lemma~\ref{G'_derivern_Suc_imp_G_derivern}.
\end{proof}
\end{lemma}

We will now show that a non-left-recursive start symbol through rightmost derivations is a necessary
condition as well.

\begin{lemma}\label{G_LR0_lrec_imp_not_G'_LR0}
If \isa{G} is an \isa{LR{\isacharparenleft}{\kern0pt}{\isadigit{0}}{\isacharparenright}{\kern0pt}} grammar with a left recursion of the form 
\isa{\isaconst{P}\ {\isasymturnstile}\ {\isacharbrackleft}{\kern0pt}\isaconst{Nt}\ \isaconst{S}{\isacharbrackright}{\kern0pt}\ {\isasymRightarrow}r{\isacharparenleft}{\kern0pt}\isaconst{Suc}\ \isafree{n}{\isacharparenright}{\kern0pt}\ \isaconst{Nt}\ \isaconst{S}\ {\isacharhash}{\kern0pt}\ \isafree{{\isasymalpha}}}, \isa{G{\isacharprime}{\kern0pt}} is not an \isa{LR{\isacharparenleft}{\kern0pt}{\isadigit{0}}{\isacharparenright}{\kern0pt}} grammar.
\begin{proof}
Since it is a reduced grammar, we know there exists some \isa{v\ {\isacharcolon}{\kern0pt}{\isacharcolon}{\kern0pt}\ {\isacharprime}{\kern0pt}t\ list} derived by \isa{{\isasymalpha}}, i.e., there 
also exists some \isa{k} with
\isa{\isaconst{P}\ {\isasymturnstile}\ {\isacharbrackleft}{\kern0pt}\isaconst{Nt}\ \isaconst{S}{\isacharbrackright}{\kern0pt}\ {\isasymRightarrow}r{\isacharparenleft}{\kern0pt}\isaconst{Suc}\ \isafree{k}{\isacharparenright}{\kern0pt}\ \isaconst{Nt}\ \isaconst{S}\ {\isacharhash}{\kern0pt}\ \isaconst{map}\ \isaconst{Tm}\ \isafree{v}}.
We can then decompose this derivation as 
\isa{P\ {\isasymturnstile}\ {\isacharbrackleft}{\kern0pt}Nt\ S{\isacharbrackright}{\kern0pt}\ {\isasymRightarrow}r{\isacharasterisk}{\kern0pt}\ {\isasymgamma}\ {\isacharat}{\kern0pt}\ Nt\ Y\ {\isacharhash}{\kern0pt}\ map\ Tm\ x\ {\isasymRightarrow}r\ Nt\ S\ {\isacharhash}{\kern0pt}\ map\ Tm\ v} for some \isa{{\isasymgamma}\ {\isacharcolon}{\kern0pt}{\isacharcolon}{\kern0pt}\ syms}, \isa{Y\ {\isacharcolon}{\kern0pt}{\isacharcolon}{\kern0pt}\ {\isacharprime}{\kern0pt}n}, and 
\isa{x\ {\isacharcolon}{\kern0pt}{\isacharcolon}{\kern0pt}\ {\isacharprime}{\kern0pt}t\ list}. Trivially, this implies \isa{Y\ {\isasymnoteq}\ S{\isacharprime}{\kern0pt}}.

Furthermore, we know \isa{G{\isacharprime}{\kern0pt}} has a derivation 
\isa{P{\isacharprime}{\kern0pt}\ {\isasymturnstile}\ {\isacharbrackleft}{\kern0pt}Nt\ S{\isacharprime}{\kern0pt}{\isacharbrackright}{\kern0pt}\ {\isasymRightarrow}r{\isacharasterisk}{\kern0pt}\ {\isasymalpha}\ {\isacharat}{\kern0pt}\ Nt\ X\ {\isacharhash}{\kern0pt}\ map\ Tm\ w\ {\isasymRightarrow}r\ {\isasymalpha}\ {\isacharat}{\kern0pt}\ {\isasymbeta}\ {\isacharat}{\kern0pt}\ map\ Tm\ w} with \isa{{\isasymalpha}\ {\isacharequal}{\kern0pt}\ {\isacharbrackleft}{\kern0pt}{\isacharbrackright}{\kern0pt}}, \isa{X\ {\isacharequal}{\kern0pt}\ S{\isacharprime}{\kern0pt}}, 
\isa{{\isasymbeta}\ {\isacharequal}{\kern0pt}\ {\isacharbrackleft}{\kern0pt}Nt\ S{\isacharbrackright}{\kern0pt}}, and \isa{w\ {\isacharequal}{\kern0pt}\ {\isacharbrackleft}{\kern0pt}{\isacharbrackright}{\kern0pt}} by reflexivity. We can therefore express the right sentential form 
\isa{Nt\ S\ {\isacharhash}{\kern0pt}\ map\ Tm\ v} as \isa{{\isasymalpha}\ {\isacharat}{\kern0pt}\ {\isasymbeta}\ {\isacharat}{\kern0pt}\ map\ Tm\ v}, and the \isa{LR{\isacharparenleft}{\kern0pt}{\isadigit{0}}{\isacharparenright}{\kern0pt}} condition is thus violated since \isa{X\ {\isasymnoteq}\ Y}.
\end{proof}
\end{lemma}

From these lemmas and Lemma~\ref{LR'k_G'_imp_LR'k_G} follows the first equivalence.

\begin{theorem}\label{G'_LR0_iff_G_LR0_lrec_free}
\isa{G{\isacharprime}{\kern0pt}} is an \isa{LR{\isacharparenleft}{\kern0pt}{\isadigit{0}}{\isacharparenright}{\kern0pt}} grammar if and only if \isa{G} is an \isa{LR{\isacharparenleft}{\kern0pt}{\isadigit{0}}{\isacharparenright}{\kern0pt}} grammar with no left recursions of the 
form \isa{\isaconst{P}\ {\isasymturnstile}\ {\isacharbrackleft}{\kern0pt}\isaconst{Nt}\ \isaconst{S}{\isacharbrackright}{\kern0pt}\ {\isasymRightarrow}r{\isacharparenleft}{\kern0pt}\isaconst{Suc}\ \isafree{n}{\isacharparenright}{\kern0pt}\ \isaconst{Nt}\ \isaconst{S}\ {\isacharhash}{\kern0pt}\ \isafree{{\isasymalpha}}} for any \isa{n} or \isa{{\isasymalpha}}.
\end{theorem}

We now move on to \isa{k\ {\isachargreater}{\kern0pt}\ {\isadigit{0}}}.

\begin{lemma}[Preservation of \isa{LR{\isacharparenleft}{\kern0pt}k{\isacharparenright}{\kern0pt}} for nonzero \isa{k}]\label{LR'_Suc_G_imp_LR'_Suc_G'}
Let \isa{G} be an \isa{LR{\isacharparenleft}{\kern0pt}k{\isacharparenright}{\kern0pt}} grammar for \isa{k\ {\isachargreater}{\kern0pt}\ {\isadigit{0}}}. Furthermore, assume there does not exist a cycle
\isa{\isaconst{P}\ {\isasymturnstile}\ {\isacharbrackleft}{\kern0pt}\isaconst{Nt}\ \isaconst{S}{\isacharbrackright}{\kern0pt}\ {\isasymRightarrow}r{\isacharparenleft}{\kern0pt}\isaconst{Suc}\ \isafree{n}{\isacharparenright}{\kern0pt}\ {\isacharbrackleft}{\kern0pt}\isaconst{Nt}\ \isaconst{S}{\isacharbrackright}{\kern0pt}} for any \isa{n\ {\isacharcolon}{\kern0pt}{\isacharcolon}{\kern0pt}\ nat}. Then \isa{G{\isacharprime}{\kern0pt}} is an \isa{LR{\isacharparenleft}{\kern0pt}k{\isacharparenright}{\kern0pt}} grammar
\begin{proof}
Consider the rightmost derivations 
\begin{gather*}
\isa{P{\isacharprime}{\kern0pt}\ {\isasymturnstile}\ {\isacharbrackleft}{\kern0pt}Nt\ S{\isacharprime}{\kern0pt}{\isacharbrackright}{\kern0pt}\ {\isasymRightarrow}r{\isacharparenleft}{\kern0pt}n{\isacharparenright}{\kern0pt}\ {\isasymalpha}\ {\isacharat}{\kern0pt}\ Nt\ X\ {\isacharhash}{\kern0pt}\ map\ Tm\ w}
  \isa{{\isasymRightarrow}r\ {\isasymalpha}\ {\isacharat}{\kern0pt}\ {\isasymbeta}\ {\isacharat}{\kern0pt}\ map\ Tm\ w}\\
\isa{P{\isacharprime}{\kern0pt}\ {\isasymturnstile}\ {\isacharbrackleft}{\kern0pt}Nt\ S{\isacharprime}{\kern0pt}{\isacharbrackright}{\kern0pt}\ {\isasymRightarrow}r{\isacharparenleft}{\kern0pt}m{\isacharparenright}{\kern0pt}\ {\isasymgamma}\ {\isacharat}{\kern0pt}\ Nt\ Y\ {\isacharhash}{\kern0pt}\ map\ Tm\ x}
  \isa{{\isasymRightarrow}r\ {\isasymalpha}\ {\isacharat}{\kern0pt}\ {\isasymbeta}\ {\isacharat}{\kern0pt}\ map\ Tm\ y}
\end{gather*}
with \isa{\isafree{w}{\isacharbar}{\kern0pt}\isactrlbsub \isafree{k}\isactrlesub \ {\isacharequal}{\kern0pt}\ \isafree{y}{\isacharbar}{\kern0pt}\isactrlbsub \isafree{k}\isactrlesub }
We have to prove that \isa{\isafree{{\isasymalpha}}\ {\isacharequal}{\kern0pt}\ \isafree{{\isasymgamma}}\ {\isasymand}\ \isafree{X}\ {\isacharequal}{\kern0pt}\ \isafree{Y}\ {\isasymand}\ \isafree{x}\ {\isacharequal}{\kern0pt}\ \isafree{y}} holds.
We will first focus on the case where \isa{n\ {\isacharequal}{\kern0pt}\ {\isadigit{0}}} and \isa{m\ {\isacharequal}{\kern0pt}\ Suc\ m{\isacharprime}{\kern0pt}} for some \isa{m{\isacharprime}{\kern0pt}}:

This case once again implies that 
\begin{equation}\label{LR_Suc_k}
\isa{\isafree{{\isasymalpha}}\ {\isacharequal}{\kern0pt}\ {\isacharbrackleft}{\kern0pt}{\isacharbrackright}{\kern0pt}}, \isa{\isafree{{\isasymbeta}}\ {\isacharequal}{\kern0pt}\ {\isacharbrackleft}{\kern0pt}\isaconst{Nt}\ \isaconst{S}{\isacharbrackright}{\kern0pt}}, \text{ and } \isa{\isafree{w}\ {\isacharequal}{\kern0pt}\ {\isacharbrackleft}{\kern0pt}{\isacharbrackright}{\kern0pt}},
\end{equation}
and we again have a left recursion of the form 
\isa{\isaconst{P{\isacharprime}{\kern0pt}}\ {\isasymturnstile}\ {\isacharbrackleft}{\kern0pt}\isaconst{Nt}\ \isaconst{S{\isacharprime}{\kern0pt}}{\isacharbrackright}{\kern0pt}\ {\isasymRightarrow}r{\isacharparenleft}{\kern0pt}\isaconst{Suc}\ {\isacharparenleft}{\kern0pt}\isaconst{Suc}\ \isafree{m{\isacharprime}{\kern0pt}}{\isacharparenright}{\kern0pt}{\isacharparenright}{\kern0pt}\ \isaconst{Nt}\ \isaconst{S}\ {\isacharhash}{\kern0pt}\ \isaconst{map}\ \isaconst{Tm}\ \isafree{y}}.
However, in this case, since \isa{{\isadigit{0}}\ {\isacharless}{\kern0pt}\ \isafree{k}} and \isa{\isafree{w}\ {\isacharequal}{\kern0pt}\ {\isacharbrackleft}{\kern0pt}{\isacharbrackright}{\kern0pt}}, \isa{\isafree{w}{\isacharbar}{\kern0pt}\isactrlbsub \isafree{k}\isactrlesub \ {\isacharequal}{\kern0pt}\ \isafree{y}{\isacharbar}{\kern0pt}\isactrlbsub \isafree{k}\isactrlesub } 
can only hold if \isa{y\ {\isacharequal}{\kern0pt}\ {\isacharbrackleft}{\kern0pt}{\isacharbrackright}{\kern0pt}}. This contradicts the assumption that \isa{G} is cycle-free.

The proof for the other cases is analogous to that of Lemma~\ref{LR'0_G_imp_LR'0_G'}.
\end{proof}
\end{lemma}

We will once again show that this condition is necessary.

\begin{lemma}\label{G_LR_Suc_cycle_imp_not_G'_LR_Suc}
If \isa{G} is an \isa{LR{\isacharparenleft}{\kern0pt}k{\isacharparenright}{\kern0pt}} grammar for \isa{k\ {\isachargreater}{\kern0pt}\ {\isadigit{0}}} with a cycle through its start symbol \isa{S} reachable
via rightmost derivations, \isa{G{\isacharprime}{\kern0pt}} is not an
\isa{LR{\isacharparenleft}{\kern0pt}k{\isacharparenright}{\kern0pt}} grammar.
\begin{proof}
A cycle through \isa{S} via rightmost derivations implies the existence of some \isa{n} such that
\isa{P\ {\isasymturnstile}\ {\isacharbrackleft}{\kern0pt}Nt\ S{\isacharbrackright}{\kern0pt}\ {\isasymRightarrow}r{\isacharparenleft}{\kern0pt}n{\isacharparenright}{\kern0pt}\ {\isasymgamma}\ {\isacharat}{\kern0pt}\ Nt\ Y\ {\isacharhash}{\kern0pt}\ map\ Tm\ x\ {\isasymRightarrow}r\ {\isacharbrackleft}{\kern0pt}Nt\ S{\isacharbrackright}{\kern0pt}} for some \isa{{\isasymgamma}\ {\isacharcolon}{\kern0pt}{\isacharcolon}{\kern0pt}\ syms}, \isa{Y\ {\isacharcolon}{\kern0pt}{\isacharcolon}{\kern0pt}\ {\isacharprime}{\kern0pt}n}, and 
\isa{x\ {\isacharcolon}{\kern0pt}{\isacharcolon}{\kern0pt}\ {\isacharprime}{\kern0pt}t\ list}. This again implies that \isa{Y\ {\isasymnoteq}\ S{\isacharprime}{\kern0pt}}, and the \isa{LR{\isacharparenleft}{\kern0pt}k{\isacharparenright}{\kern0pt}} condition is violated similarly 
to how it occurs in the proof of Lemma~\ref{G_LR0_lrec_imp_not_G'_LR0}, with one important 
distinction: in this case, the condition is violated because \isa{x} must be empty for 
\isa{{\isasymgamma}\ {\isacharat}{\kern0pt}\ Nt\ Y\ {\isacharhash}{\kern0pt}\ map\ Tm\ x} to derive \isa{{\isacharbrackleft}{\kern0pt}Nt\ S{\isacharbrackright}{\kern0pt}}. Therefore, \isa{{\isasymgamma}\ {\isacharat}{\kern0pt}\ Nt\ Y\ {\isacharhash}{\kern0pt}\ map\ Tm\ x} derives 
\isa{\isafree{{\isasymalpha}}\ {\isacharat}{\kern0pt}\ \isafree{{\isasymbeta}}\ {\isacharat}{\kern0pt}\ \isaconst{map}\ \isaconst{Tm}\ \isafree{x}} with \isa{{\isasymalpha}\ {\isacharequal}{\kern0pt}\ {\isasymbeta}\ {\isacharequal}{\kern0pt}\ map\ Tm\ x\ {\isacharequal}{\kern0pt}\ {\isacharbrackleft}{\kern0pt}{\isacharbrackright}{\kern0pt}} and thus the property 
\isa{\isafree{x}{\isacharbar}{\kern0pt}\isactrlbsub \isafree{k}\isactrlesub \ {\isacharequal}{\kern0pt}\ \isafree{w}{\isacharbar}{\kern0pt}\isactrlbsub \isafree{k}\isactrlesub } = \isa{{\isacharbrackleft}{\kern0pt}{\isacharbrackright}{\kern0pt}} holds. Thus, \isa{X\ {\isasymnoteq}\ Y} violates the \isa{LR{\isacharparenleft}{\kern0pt}k{\isacharparenright}{\kern0pt}} condition once again.
\end{proof}
\end{lemma}

Once more, from the above two lemmas and Lemma~\ref{LR'k_G'_imp_LR'k_G} follows the equality.

\begin{theorem}\label{G'_LR_Suc_iff_G_LR_Suc_cycle_free}
\isa{G{\isacharprime}{\kern0pt}} is an \isa{LR{\isacharparenleft}{\kern0pt}k{\isacharparenright}{\kern0pt}} grammar for \isa{k\ {\isachargreater}{\kern0pt}\ {\isadigit{0}}} if and only if \isa{G} has no cycles of the form 
\isa{\isaconst{P}\ {\isasymturnstile}\ {\isacharbrackleft}{\kern0pt}\isaconst{Nt}\ \isaconst{S}{\isacharbrackright}{\kern0pt}\ {\isasymRightarrow}r{\isacharparenleft}{\kern0pt}\isaconst{Suc}\ \isafree{n}{\isacharparenright}{\kern0pt}\ {\isacharbrackleft}{\kern0pt}\isaconst{Nt}\ \isaconst{S}{\isacharbrackright}{\kern0pt}} for any \isa{n}.
\end{theorem}

It is worth highlighting that in the CFG section of the book, Wilhelm et al.\
explain~\cite[p. 52]{Wilhelm}:
\begin{quote}
Context-free grammars that describe programming languages should be unambiguous. If this is the case, 
then there exist exactly one parse tree, one leftmost and one rightmost derivation for each 
syntactically correct program.
\end{quote}

This is of great importance to our result from Theorem~\ref{G'_LR_Suc_iff_G_LR_Suc_cycle_free}: a 
grammar with a cycle through its start symbol has infinitely many derivations for every word in the
language, i.e., it is ambiguous. We can therefore conclude, based on our results and the authors'
claims about the importance of grammars being unambiguous, that the precondition for the preservation
of the \isa{LR{\isacharparenleft}{\kern0pt}k{\isacharparenright}{\kern0pt}} condition for \isa{k\ {\isachargreater}{\kern0pt}\ {\isadigit{0}}}, i.e., the grammar being cycle-free under rightmost derivations, 
is essentially guaranteed to hold in practice. This is because the extension only causes the 
previously satisfied condition to be violated in grammars which the theory we are formalizing 
already deems inadequate to begin with. In fact, our result even shows that one could extend a
grammar with a cycle through its start symbol without introducing ambiguity, provided that this cycle
does not consist exclusively of rightmost derivations. The extension thus only affects a \emph{subset}
of grammars that are already dismissed as unsuitable.

For the case of \isa{k\ {\isacharequal}{\kern0pt}\ {\isadigit{0}}}, however, this does become problematic. Grammars with left recursions
are fairly standard, meaning that our precondition shows that only a subset of \isa{LR{\isacharparenleft}{\kern0pt}{\isadigit{0}}{\isacharparenright}{\kern0pt}} grammars that
are used in practice allows for the extension, namely those that lack a left-recursive start symbol
whose left recursion is reachable with rightmost derivations. Since our parser and its underlying
automata are defined exclusively for extended grammars, this means that it is possible for an
unextended grammar \isa{G} to fulfill the \isa{LR{\isacharparenleft}{\kern0pt}{\isadigit{0}}{\isacharparenright}{\kern0pt}} condition while \isa{\ensuremath{LR_0(G)}} contains \isa{LR{\isacharparenleft}{\kern0pt}{\isadigit{0}}{\isacharparenright}{\kern0pt}}-inadequate
states. Therefore, the \isa{LR{\isacharparenleft}{\kern0pt}{\isadigit{0}}{\isacharparenright}{\kern0pt}} parser as defined by Wilhelm et al.\ is not able to parse
an arbitrary \isa{LR{\isacharparenleft}{\kern0pt}{\isadigit{0}}{\isacharparenright}{\kern0pt}} grammar deterministically in general.

Theorem~\ref{G'_LR0_iff_G_LR0_lrec_free} also shows that one needs to be careful when treating the
\isa{LR{\isacharparenleft}{\kern0pt}{\isadigit{0}}{\isacharparenright}{\kern0pt}} condition as interchangeable for \isa{G} and \isa{G{\isacharprime}{\kern0pt}}. Wilhelm et al.\ do so 
themselves~\cite[p. 113]{Wilhelm}: after they introduce their definition of \isa{LR{\isacharparenleft}{\kern0pt}k{\isacharparenright}{\kern0pt}}, which is
defined only for an extended grammar, they immediately go on to show an example of an unextended
grammar that fulfills the (generalized) \isa{LR{\isacharparenleft}{\kern0pt}{\isadigit{0}}{\isacharparenright}{\kern0pt}} condition. This is not problematic for this
particular grammar, but given that left recursions are not uncommon in general, our results suggest
that it is sometimes necessary to differentiate explicitly between the fulfillment of the condition
by the unextended and extended grammars instead of treating them as equivalent.
 
With these results and our proof of Theorem~\ref{is_LR0_iff_no_LR0_inadequates}, we have completed
our goal of verifying the \isa{LR{\isacharparenleft}{\kern0pt}{\isadigit{0}}{\isacharparenright}{\kern0pt}} parsing theory as presented by Wilhelm et al. We will now address
the final topic we will discuss: the parser's language.%
\end{isamarkuptext}\isamarkuptrue%
%
\isadelimdocument
%
\endisadelimdocument
%
\isatagdocument
%
\isamarkupsection{Correctness of the \isa{LR{\isacharparenleft}{\kern0pt}{\isadigit{0}}{\isacharparenright}{\kern0pt}} Parser%
}
\isamarkuptrue%
%
\endisatagdocument
{\isafolddocument}%
%
\isadelimdocument
%
\endisadelimdocument
%
\begin{isamarkuptext}%
Now that we have finished formalizing the basic \isa{LR{\isacharparenleft}{\kern0pt}{\isadigit{0}}{\isacharparenright}{\kern0pt}} theory, the most interesting question
that remains open is whether \isa{LR{\isacharparenleft}{\kern0pt}{\isadigit{0}}{\isacharparenright}{\kern0pt}} parser we have defined is correct.
We will answer this question by proving that \isa{\ensuremath{P_0(G)}} accepts exactly \isa{L{\isacharparenleft}{\kern0pt}\isafree{G}{\isacharparenright}{\kern0pt}}. In this
section, we denote parser steps, their reflexive transitive closure, and \isa{n}-step computations by
\isa{{\isasymturnstile}P}, \isa{{\isasymturnstile}P{\isacharasterisk}{\kern0pt}}, and \isa{{\isasymturnstile}P{\isacharparenleft}{\kern0pt}n{\isacharparenright}{\kern0pt}} respectively.%
\end{isamarkuptext}\isamarkuptrue%
%
\isadelimdocument
%
\endisadelimdocument
%
\isatagdocument
%
\isamarkupsubsection{Stack Words: Proving Soundness%
}
\isamarkuptrue%
%
\endisatagdocument
{\isafolddocument}%
%
\isadelimdocument
%
\endisadelimdocument
%
\begin{isamarkuptext}%
We will begin our correctness proof by showing that our parser is sound, i.e., every word
accepted by our parser is a word in \isa{L{\isacharparenleft}{\kern0pt}\isafree{G}{\isacharparenright}{\kern0pt}}. We will show an invariant for \isa{\ensuremath{P_0(G)}} 
computations from which we will derive this property. Consider the following excerpt of Wilhelm
et al.\ on the parser's stack~\cite[p. 110]{Wilhelm}:
\begin{quote}
The construction of \isa{\ensuremath{LR_0(G)}} guarantees that for each non-initial and non-final state \isa{q} there
exists exactly one entry symbol under which the automaton can make a transition into \isa{q}. The
pushdown contents \isa{q\isactrlsub {\isadigit{0}}{\isacharcomma}{\kern0pt}\ {\isasymdots}{\isacharcomma}{\kern0pt}\ q\isactrlsub n} with \isa{\isafree{q\isactrlsub {\isadigit{0}}}\ {\isacharequal}{\kern0pt}\ \ensuremath{\isafree{q}_{\isafree{G}{\isacharcomma}{\kern0pt}{\isadigit{0}}}}}~\footnote{Recall that the authors define
the topmost stack symbol to be the rightmost one when writing out a list, which is the opposite of 
our convention.} [correspond] therefore to a uniquely determined word 
\mbox{\isa{{\isasymalpha}\ {\isacharequal}{\kern0pt}\ X\isactrlsub {\isadigit{1}}\ {\isasymdots}\ X\isactrlsub n\ {\isasymin}\ {\isacharparenleft}{\kern0pt}V\isactrlsub T\ {\isasymunion}\ V\isactrlsub N{\isacharparenright}{\kern0pt}\isactrlsup {\isacharasterisk}{\kern0pt}}}~\footnote{\isa{V\isactrlsub T} and \isa{V\isactrlsub N} denote the sets of the grammar's terminals and
nonterminals respectively.} for which
$\Delta_G(q_i, X_{i+1}) = q_{i+1}$ holds. This word \isa{{\isasymalpha}} is a reliable prefix, and \isa{q\isactrlsub n} is the
set of all items valid for \isa{{\isasymalpha}}.
\end{quote}

The authors do not develop this idea further. It is of particular interest, as we will soon see,
what the relation between the reliable prefix \isa{{\isasymalpha}} and the input that has been consumed at that point
is. We use this intuitive explanation of the inner workings of the parser's stack as inspiration to 
formalize what we call the parser's \concept{stack words}. For configurations of \isa{\ensuremath{P_0(G)}} \isa{c\isactrlsub {\isadigit{0}}} and
\isa{c\isactrlsub {\isadigit{1}}}, we write \isa{\isafree{{\isasymalpha}}\ {\isasymTurnstile}\ \isafree{c\isactrlsub {\isadigit{0}}}\ {\isasymturnstile}P{\isacharasterisk}{\kern0pt}\ \isafree{c\isactrlsub {\isadigit{1}}}} to mean that \isa{c\isactrlsub {\isadigit{0}}} reaches \isa{c\isactrlsub {\isadigit{1}}} with stack word \isa{{\isasymalpha}}. This is 
defined inductively with a
\concept{reflexive} rule:
\begin{gather*}
\isa{{\isacharbrackleft}{\kern0pt}{\isacharbrackright}{\kern0pt}\ {\isasymTurnstile}\ \isavar{c\isactrlsub {\isadigit{0}}}\ {\isasymturnstile}P{\isacharasterisk}{\kern0pt}\ \isavar{c\isactrlsub {\isadigit{0}}}}\\
\intertext{and a \concept{step} rule:}
\isa{\mbox{}\inferrule{\mbox{\isavar{{\isasymalpha}}\ {\isasymTurnstile}\ \isavar{c\isactrlsub {\isadigit{0}}}\ {\isasymturnstile}P{\isacharasterisk}{\kern0pt}\ {\isacharparenleft}{\kern0pt}\isavar{ps}\ {\isacharat}{\kern0pt}\ \isavar{q}\ {\isacharhash}{\kern0pt}\ \isavar{qs}{\isacharcomma}{\kern0pt}\ \isavar{u}{\isacharparenright}{\kern0pt}}\\\ \mbox{{\isacharparenleft}{\kern0pt}\isavar{ps}\ {\isacharat}{\kern0pt}\ \isavar{q}\ {\isacharhash}{\kern0pt}\ \isavar{qs}{\isacharcomma}{\kern0pt}\ \isavar{u}{\isacharparenright}{\kern0pt}\ {\isasymturnstile}P\ {\isacharparenleft}{\kern0pt}\isaconst{nxt}\ \ensuremath{LR_0(G)}\ \isavar{q}\ \isavar{X}\ {\isacharhash}{\kern0pt}\ \isavar{q}\ {\isacharhash}{\kern0pt}\ \isavar{qs}{\isacharcomma}{\kern0pt}\ \isavar{v}{\isacharparenright}{\kern0pt}}}{\mbox{\isavar{X}\ {\isacharhash}{\kern0pt}\ \isactrlbsub {\isacharbar}{\kern0pt}\isavar{ps}{\isacharbar}{\kern0pt}\isactrlesub {\isacharbar}{\kern0pt}\isavar{{\isasymalpha}}\ {\isasymTurnstile}\ \isavar{c\isactrlsub {\isadigit{0}}}\ {\isasymturnstile}P{\isacharasterisk}{\kern0pt}\ {\isacharparenleft}{\kern0pt}\isaconst{nxt}\ \ensuremath{LR_0(G)}\ \isavar{q}\ \isavar{X}\ {\isacharhash}{\kern0pt}\ \isavar{q}\ {\isacharhash}{\kern0pt}\ \isavar{qs}{\isacharcomma}{\kern0pt}\ \isavar{v}{\isacharparenright}{\kern0pt}}}}
\end{gather*}

With stack words, we are now able to track the reliable prefix that corresponds to the parser's stack
akin to what the authors describe: note that in both read and reduce transitions, our parser always
pushes \isa{\isaconst{nxt}\ \ensuremath{LR_0(G)}\ \isafree{q}\ \isafree{X}} on top of the stack, where \isa{q} is the second-topmost state after 
the transition was made. In reduce transitions, some topmost stack string is deleted and replaced by
the successor of the state that becomes the topmost state immediately after the deletion. 
This topmost string corresponds to the list \isa{ps} in the \isa{step} rule. In read transitions, the stack
is simply extended by the successor of the topmost state, so in this case our \isa{ps} is the empty list.
Since in 
both cases we remove \isa{{\isacharbar}{\kern0pt}\isafree{ps}{\isacharbar}{\kern0pt}} symbols from the stack, we do the exact same for the stack 
word, which is why we replace the old stack word \isa{{\isasymalpha}} by the suffix \isa{\isactrlbsub {\isacharbar}{\kern0pt}\isafree{ps}{\isacharbar}{\kern0pt}\isactrlesub {\isacharbar}{\kern0pt}\isafree{{\isasymalpha}}}. By 
pushing the symbol \isa{X} that gets passed to the \isa{nxt} function to compute the new topmost state, we
keep track of the symbols that the authors describe in the excerpt. It is worth remarking that our 
stack words are stored in reverse, once again, to make the inductive rules more fitting to the
Isabelle list datatype. It is also important to note that stack words as we defined them have a 
caveat: the stack word that is being tracked only truly corresponds to the parser's stack if 
for \isa{\isafree{{\isasymalpha}}\ {\isasymTurnstile}\ \isafree{c\isactrlsub {\isadigit{0}}}\ {\isasymturnstile}P{\isacharasterisk}{\kern0pt}\ \isafree{c\isactrlsub {\isadigit{1}}}}, the stack of \isa{c\isactrlsub {\isadigit{0}}} is empty, i.e., \isa{\isafree{c\isactrlsub {\isadigit{0}}}\ {\isacharequal}{\kern0pt}\ {\isacharparenleft}{\kern0pt}{\isacharbrackleft}{\kern0pt}\isafree{p}{\isacharbrackright}{\kern0pt}{\isacharcomma}{\kern0pt}\ \isafree{w}{\isacharparenright}{\kern0pt}} for some state 
\isa{p} and input \isa{w}. This limitation, however, is not problematic for us, since the starting 
configuration of the parser fulfills this precondition.

In order to prove the parser sound, we will first need to show some additional properties.

\begin{lemma}\label{inj_nxt_dfa_LR0_if_nonempty}
If \isa{\isaconst{nxt}\ \ensuremath{LR_0(G)}\ \isafree{q}\ \isafree{X}\ {\isacharequal}{\kern0pt}\ \isaconst{nxt}\ \ensuremath{LR_0(G)}\ \isafree{q}\ \isafree{Y}} and \isa{\isaconst{nxt}\ \ensuremath{LR_0(G)}\ \isafree{q}\ \isafree{X}\ {\isasymnoteq}\ {\isasymemptyset}}, then
\isa{\isafree{X}\ {\isacharequal}{\kern0pt}\ \isafree{Y}} holds.
\begin{proof}
With Lemma~\ref{char_fa_nxts_is_shifts} we can show that for any state \isa{p} of \isa{\ensuremath{LR_0(G)}} holds
\[ \isa{\isaconst{nxt}\ \ensuremath{LR_0(G)}\ \isafree{p}\ \isafree{X}\ {\isacharequal}{\kern0pt}\ \isaconst{epsclo}\ {\isacharbraceleft}{\kern0pt}{\isacharbrackleft}{\kern0pt}\isabound{X{\isacharprime}{\kern0pt}}\ {\isasymrightarrow}\ \isabound{{\isasymalpha}}\ {\isacharat}{\kern0pt}\ {\isacharbrackleft}{\kern0pt}\isafree{X}{\isacharbrackright}{\kern0pt}\ {\isasymcdot}\ \isabound{{\isasymbeta}}{\isacharbrackright}{\kern0pt}\ {\isacharbar}{\kern0pt}\ {\isacharbrackleft}{\kern0pt}\isabound{X{\isacharprime}{\kern0pt}}\ {\isasymrightarrow}\ \isabound{{\isasymalpha}}\ {\isasymcdot}\ \isafree{X}\ {\isacharhash}{\kern0pt}\ \isabound{{\isasymbeta}}{\isacharbrackright}{\kern0pt}\ {\isasymin}\ \isafree{p}{\isacharbraceright}{\kern0pt}}. \]
The claim then follows trivially.
\end{proof}
\end{lemma}

\begin{lemma}\label{P0_init_stack_word_length_inv}
If \isa{{\isacharparenleft}{\kern0pt}{\isacharbrackleft}{\kern0pt}\isaconst{init}\ \ensuremath{P_0(G)}{\isacharbrackright}{\kern0pt}{\isacharcomma}{\kern0pt}\ \isafree{u}{\isacharparenright}{\kern0pt}} reaches some configuration \isa{{\isacharparenleft}{\kern0pt}\isafree{qs}{\isacharcomma}{\kern0pt}\ \isafree{v}{\isacharparenright}{\kern0pt}} with stack word \isa{{\isasymalpha}}, 
then \isa{{\isacharbar}{\kern0pt}\isafree{qs}{\isacharbar}{\kern0pt}\ {\isacharequal}{\kern0pt}\ {\isacharbar}{\kern0pt}\isafree{{\isasymalpha}}{\isacharbar}{\kern0pt}\ {\isacharplus}{\kern0pt}\ {\isadigit{1}}} holds.
\begin{proof}
By rule induction on the stack word.
\end{proof}
\end{lemma}

\begin{lemma}\label{stack_word_imp_P0_steps}
If \isa{\isafree{{\isasymalpha}}\ {\isasymTurnstile}\ \isafree{c\isactrlsub {\isadigit{0}}}\ {\isasymturnstile}P{\isacharasterisk}{\kern0pt}\ \isafree{c\isactrlsub {\isadigit{1}}}}, then \isa{\isafree{c\isactrlsub {\isadigit{0}}}\ {\isasymturnstile}P{\isacharasterisk}{\kern0pt}\ \isafree{c\isactrlsub {\isadigit{1}}}}.
\begin{proof}
By rule induction on the stack word.
\end{proof}
\end{lemma}

\begin{lemma}\label{P0_steps_imp_stack_word}
If a computation \isa{\isafree{c\isactrlsub {\isadigit{0}}}\ {\isasymturnstile}P{\isacharasterisk}{\kern0pt}\ {\isacharparenleft}{\kern0pt}\isafree{q}\ {\isacharhash}{\kern0pt}\ \isafree{qs}{\isacharcomma}{\kern0pt}\ \isafree{w}{\isacharparenright}{\kern0pt}} exists for some state of \isa{\ensuremath{LR_0(G)}} \isa{q}, 
there exists a stack word \isa{{\isasymalpha}} with which \isa{c\isactrlsub {\isadigit{0}}} reaches \isa{{\isacharparenleft}{\kern0pt}q\ {\isacharhash}{\kern0pt}\ qs{\isacharcomma}{\kern0pt}\ w{\isacharparenright}{\kern0pt}}.
\begin{proof}
The proof is by induction on the length of the computation, with a case distinction on the 
transition type in the last step for the transitive case, and using the fact that after a 
finish transition, the topmost state on the stack is not a state of \isa{\ensuremath{LR_0(G)}}.
\end{proof}
\end{lemma}

\begin{lemma}\label{P0_nth_is_valids_of_nth_stack_word}
If \isa{\isafree{{\isasymalpha}}\ {\isasymTurnstile}\ {\isacharparenleft}{\kern0pt}{\isacharbrackleft}{\kern0pt}\isaconst{init}\ \ensuremath{P_0(G)}{\isacharbrackright}{\kern0pt}{\isacharcomma}{\kern0pt}\ \isafree{u}{\isacharparenright}{\kern0pt}\ {\isasymturnstile}P{\isacharasterisk}{\kern0pt}\ {\isacharparenleft}{\kern0pt}\isafree{q}\ {\isacharhash}{\kern0pt}\ \isafree{qs}{\isacharcomma}{\kern0pt}\ \isafree{v}{\isacharparenright}{\kern0pt}}, then 
\isa{{\isacharparenleft}{\kern0pt}\isafree{q}\ {\isacharhash}{\kern0pt}\ \isafree{qs}{\isacharparenright}{\kern0pt}\ensuremath{_{[\mathit{\isafree{n}}]}}\ {\isacharequal}{\kern0pt}\ \isaconst{valids}\ {\isacharparenleft}{\kern0pt}\isaconst{rev}\ \isactrlbsub \isafree{n}\isactrlesub {\isacharbar}{\kern0pt}\isafree{{\isasymalpha}}{\isacharparenright}{\kern0pt}} holds for $\isa{{\isadigit{0}}\ {\isasymle}\ n} < \isa{{\isacharbar}{\kern0pt}\isafree{q}\ {\isacharhash}{\kern0pt}\ \isafree{qs}{\isacharbar}{\kern0pt}}$.
\begin{proof}
We perform rule induction on the stack word for arbitrary \isa{q{\isacharcomma}{\kern0pt}\ qs{\isacharcomma}{\kern0pt}\ u{\isacharcomma}{\kern0pt}} and \isa{v}.

The reflexive case holds by showing that \isa{\isaconst{init}\ \ensuremath{P_0(G)}\ {\isacharequal}{\kern0pt}\ \isaconst{valids}\ {\isacharbrackleft}{\kern0pt}{\isacharbrackright}{\kern0pt}} using 
Lemma~\ref{nextl_dfa_LR0_is_valids}.

For the step case, we know the initial configuration reaches some configuration 
\isa{{\isacharparenleft}{\kern0pt}\isafree{ps}\ {\isacharat}{\kern0pt}\ \isafree{q}\ {\isacharhash}{\kern0pt}\ \isafree{qs}{\isacharcomma}{\kern0pt}\ \isafree{v}{\isacharparenright}{\kern0pt}} with stack word \isa{{\isasymalpha}}, and then the parser takes a step from this
configuration into \isa{{\isacharparenleft}{\kern0pt}\isaconst{nxt}\ \ensuremath{LR_0(G)}\ \isafree{q}\ \isafree{X}\ {\isacharhash}{\kern0pt}\ \isafree{q}\ {\isacharhash}{\kern0pt}\ \isafree{qs}{\isacharcomma}{\kern0pt}\ \isafree{w}{\isacharparenright}{\kern0pt}}. We need to show that 
for $\isa{{\isadigit{0}}\ {\isasymle}\ n} < \isa{{\isacharbar}{\kern0pt}\isaconst{nxt}\ \ensuremath{LR_0(G)}\ \isafree{q}\ \isafree{X}\ {\isacharhash}{\kern0pt}\ \isafree{q}\ {\isacharhash}{\kern0pt}\ \isafree{qs}{\isacharbar}{\kern0pt}}$, 
\isa{\mbox{{\isacharparenleft}{\kern0pt}\isaconst{nxt}\ \ensuremath{LR_0(G)}\ \isafree{q}\ \isafree{X}\ {\isacharhash}{\kern0pt}\ \isafree{q}\ {\isacharhash}{\kern0pt}\ \isafree{qs}{\isacharparenright}{\kern0pt}\ensuremath{_{[\mathit{\isafree{n}}]}}}\ {\isacharequal}{\kern0pt}\ \isaconst{valids}\ {\isacharparenleft}{\kern0pt}\isaconst{rev}\ \isactrlbsub \isafree{n}\isactrlesub {\isacharbar}{\kern0pt}{\isacharparenleft}{\kern0pt}\isafree{X}\ {\isacharhash}{\kern0pt}\ \isactrlbsub {\isacharbar}{\kern0pt}\isafree{ps}{\isacharbar}{\kern0pt}\isactrlesub {\isacharbar}{\kern0pt}\isafree{{\isasymalpha}}{\isacharparenright}{\kern0pt}{\isacharparenright}{\kern0pt}}
holds. We now distinguish cases on the type of transition in the final step.

In the reading transition case, we know the stack word after the step is exactly \isa{X\ {\isacharhash}{\kern0pt}\ {\isasymalpha}}.
Furthermore, by the induction hypothesis we know that for all $\isa{{\isadigit{0}}\ {\isasymle}\ n} < \isa{{\isacharbar}{\kern0pt}\isafree{q}\ {\isacharhash}{\kern0pt}\ \isafree{qs}{\isacharbar}{\kern0pt}}$
holds \isa{{\isacharparenleft}{\kern0pt}\isafree{q}\ {\isacharhash}{\kern0pt}\ \isafree{qs}{\isacharparenright}{\kern0pt}\ensuremath{_{[\mathit{\isafree{n}}]}}\ {\isacharequal}{\kern0pt}\ \isaconst{valids}\ {\isacharparenleft}{\kern0pt}\isaconst{rev}\ \isactrlbsub \isafree{n}\isactrlesub {\isacharbar}{\kern0pt}\isafree{{\isasymalpha}}{\isacharparenright}{\kern0pt}}. Thus, we know \isa{q} is exactly \isa{\isaconst{valids}\ \isafree{{\isasymalpha}}}.
Moreover, by Lemma~\ref{nxt_dfa_LR0_shift_is_valids_app}, \isa{\isaconst{nxt}\ \ensuremath{LR_0(G)}\ \isafree{q}\ \isafree{X}} is 
\isa{\mbox{\isaconst{valids}\ {\isacharparenleft}{\kern0pt}\isaconst{rev}\ {\isacharparenleft}{\kern0pt}\isafree{X}\ {\isacharhash}{\kern0pt}\ \isafree{{\isasymalpha}}{\isacharparenright}{\kern0pt}{\isacharparenright}{\kern0pt}}}. With all this, the implication holds.

If the transition is reducing, the induction hypothesis tells us that for all 
$\isa{{\isadigit{0}}\ {\isasymle}\ n} < \isa{{\isacharbar}{\kern0pt}\isafree{q}\ {\isacharhash}{\kern0pt}\ \isafree{qs}{\isacharbar}{\kern0pt}}$ holds 
\isa{{\isacharparenleft}{\kern0pt}\isafree{q}\ {\isacharhash}{\kern0pt}\ \isafree{qs}{\isacharparenright}{\kern0pt}\ensuremath{_{[\mathit{\isafree{n}}]}}\ {\isacharequal}{\kern0pt}\ \isaconst{valids}\ {\isacharparenleft}{\kern0pt}\isaconst{rev}\ \isactrlbsub {\isacharparenleft}{\kern0pt}{\isacharbar}{\kern0pt}\isafree{ps}{\isacharbar}{\kern0pt}\ {\isacharplus}{\kern0pt}\ \isafree{n}{\isacharparenright}{\kern0pt}\isactrlesub {\isacharbar}{\kern0pt}\isafree{{\isasymalpha}}{\isacharparenright}{\kern0pt}}, which implies that 
\isa{\isafree{q}\ {\isacharequal}{\kern0pt}\ \isaconst{valids}\ {\isacharparenleft}{\kern0pt}\isaconst{rev}\ \isactrlbsub {\isacharbar}{\kern0pt}\isafree{ps}{\isacharbar}{\kern0pt}\isactrlesub {\isacharbar}{\kern0pt}\isafree{{\isasymalpha}}{\isacharparenright}{\kern0pt}}. With this we can then use
Lemma~\ref{nxt_dfa_LR0_shift_is_valids_app} to show that the implication once again holds, 
finishing the proof.
\end{proof}
\end{lemma}

This lemma highlights the correspondence between the stack word and the stack that we described. In
particular, compare this to what Wilhelm et al.\ explain in the excerpt we quoted: ``This word \isa{{\isasymalpha}} is
a reliable prefix, and \isa{q\isactrlsub n} is the set of all items valid for \isa{{\isasymalpha}}.'', where \isa{q\isactrlsub n} refers to the topmost
stack state. Our lemma shows that this property is not only true for \isa{q\isactrlsub n}, but also for every
item on the stack and its corresponding suffix of the stack word.

We will now show the parser's soundness with the following invariant:

\begin{lemma}\label{P0_invariant}
If \isa{\isafree{{\isasymalpha}}\ {\isasymTurnstile}\ {\isacharparenleft}{\kern0pt}{\isacharbrackleft}{\kern0pt}\isaconst{init}\ \ensuremath{P_0(G)}{\isacharbrackright}{\kern0pt}{\isacharcomma}{\kern0pt}\ \isafree{u}\ {\isacharat}{\kern0pt}\ \isafree{v}{\isacharparenright}{\kern0pt}\ {\isasymturnstile}P{\isacharasterisk}{\kern0pt}\ {\isacharparenleft}{\kern0pt}\isafree{q}\ {\isacharhash}{\kern0pt}\ \isafree{qs}{\isacharcomma}{\kern0pt}\ \isafree{v}{\isacharparenright}{\kern0pt}} and \isa{\isafree{q}\ {\isasymnoteq}\ {\isasymemptyset}}, then 
  \isa{\isaconst{P{\isacharprime}{\kern0pt}}\ {\isasymturnstile}\ \isaconst{rev}\ \isafree{{\isasymalpha}}\ {\isasymRightarrow}r{\isacharasterisk}{\kern0pt}\ \isaconst{map}\ \isaconst{Tm}\ \isafree{u}}
\begin{proof}
We proceed by rule induction on the stack word for arbitrary \isa{q{\isacharcomma}{\kern0pt}\ qs{\isacharcomma}{\kern0pt}\ u{\isacharcomma}{\kern0pt}} and \isa{v}.

The reflexive case is trivial since \isa{\isafree{u}\ {\isacharat}{\kern0pt}\ \isafree{v}\ {\isacharequal}{\kern0pt}\ \isafree{v}} implies \isa{\isafree{u}\ {\isacharequal}{\kern0pt}\ {\isacharbrackleft}{\kern0pt}{\isacharbrackright}{\kern0pt}}.

In the step case, we know the starting configuration with input \isa{\isafree{u}\ {\isacharat}{\kern0pt}\ \isafree{v}} reaches some
configuration \isa{{\isacharparenleft}{\kern0pt}\isafree{ps}\ {\isacharat}{\kern0pt}\ \isafree{q}\ {\isacharhash}{\kern0pt}\ \isafree{qs}{\isacharcomma}{\kern0pt}\ \isafree{w}{\isacharparenright}{\kern0pt}} with some stack word \isa{{\isasymalpha}}. We will call this configuration
\isa{c\isactrlsub {\isadigit{1}}}. Furthermore, we also know that the successor configuration is
\isa{{\isacharparenleft}{\kern0pt}\mbox{\isaconst{nxt}\ \ensuremath{LR_0(G)}\ \isafree{q}\ \isafree{X}\ {\isacharhash}{\kern0pt}\ \isafree{q}\ {\isacharhash}{\kern0pt}\ \isafree{qs}}{\isacharcomma}{\kern0pt}\ \isafree{v}{\isacharparenright}{\kern0pt}} with \isa{\isaconst{nxt}\ \ensuremath{LR_0(G)}\ \isafree{q}\ \isafree{X}\ {\isasymnoteq}\ {\isasymemptyset}}. Let \isa{c\isactrlsub {\isadigit{2}}} be
this successor configuration. 

Our goal is to show that \isa{\isaconst{P{\isacharprime}{\kern0pt}}\ {\isasymturnstile}\ \isaconst{rev}\ {\isacharparenleft}{\kern0pt}\isafree{X}\ {\isacharhash}{\kern0pt}\ \isactrlbsub {\isacharbar}{\kern0pt}\isafree{ps}{\isacharbar}{\kern0pt}\isactrlesub {\isacharbar}{\kern0pt}\isafree{{\isasymalpha}}{\isacharparenright}{\kern0pt}\ {\isasymRightarrow}r{\isacharasterisk}{\kern0pt}\ \isaconst{map}\ \isaconst{Tm}\ \isafree{u}} holds. 

From the case assumptions, we know there exists some \isa{u{\isacharprime}{\kern0pt}} such that \isa{\isafree{u}\ {\isacharat}{\kern0pt}\ \isafree{v}\ {\isacharequal}{\kern0pt}\ \isafree{u{\isacharprime}{\kern0pt}}\ {\isacharat}{\kern0pt}\ \isafree{w}}. By
the induction hypothesis, this implies that \isa{\isaconst{P{\isacharprime}{\kern0pt}}\ {\isasymturnstile}\ \isaconst{rev}\ \isafree{{\isasymalpha}}\ {\isasymRightarrow}r{\isacharasterisk}{\kern0pt}\ \isaconst{map}\ \isaconst{Tm}\ \isafree{u{\isacharprime}{\kern0pt}}}. We can now 
distinguish cases based on the type of transition in the step \isa{\isafree{c\isactrlsub {\isadigit{1}}}\ {\isasymturnstile}P\ \isafree{c\isactrlsub {\isadigit{2}}}}.

If the step is a read transition, we know that \isa{ps} is the empty list, and \isa{\isafree{u}\ {\isacharequal}{\kern0pt}\ \isafree{u{\isacharprime}{\kern0pt}}\ {\isacharat}{\kern0pt}\ {\isacharbrackleft}{\kern0pt}\isafree{a}{\isacharbrackright}{\kern0pt}}. 
Furthermore, we also know that \isa{\isaconst{rev}\ {\isacharparenleft}{\kern0pt}\isafree{X}\ {\isacharhash}{\kern0pt}\ \isactrlbsub {\isacharbar}{\kern0pt}\isafree{ps}{\isacharbar}{\kern0pt}\isactrlesub {\isacharbar}{\kern0pt}\isafree{{\isasymalpha}}{\isacharparenright}{\kern0pt}} is exactly 
\isa{\isaconst{rev}\ \isafree{{\isasymalpha}}\ {\isacharat}{\kern0pt}\ {\isacharbrackleft}{\kern0pt}\isaconst{Tm}\ \isafree{a}{\isacharbrackright}{\kern0pt}}. Therefore, the invariant is fulfilled.

If the step is a reducing transition, we know there is an item \isa{{\isacharbrackleft}{\kern0pt}\isafree{Y}\ {\isasymrightarrow}\ \isafree{{\isasymbeta}}\ {\isasymcdot}\ {\isacharbrackright}{\kern0pt}} in the topmost
state of \isa{\isafree{ps}\ {\isacharat}{\kern0pt}\ \isafree{q}\ {\isacharhash}{\kern0pt}\ \isafree{qs}} which signaled the reduction, i.e., \isa{{\isacharbar}{\kern0pt}\isafree{{\isasymbeta}}{\isacharbar}{\kern0pt}} stack states 
were removed. By Lemma~\ref{P0_nth_is_valids_of_nth_stack_word}, this means that \isa{{\isacharbrackleft}{\kern0pt}\isafree{Y}\ {\isasymrightarrow}\ \isafree{{\isasymbeta}}\ {\isasymcdot}\ {\isacharbrackright}{\kern0pt}}
is valid for \isa{rev\ {\isasymalpha}}. This implies that \isa{{\isasymalpha}} is of the form \isa{\isaconst{rev}\ \isafree{{\isasymbeta}}\ {\isacharat}{\kern0pt}\ \isafree{{\isasymalpha}{\isacharprime}{\kern0pt}}} for some
\isa{{\isasymalpha}{\isacharprime}{\kern0pt}\ {\isacharcolon}{\kern0pt}{\isacharcolon}{\kern0pt}\ syms}. Since we know that the length of \isa{{\isasymbeta}} equals the length of \isa{ps}, this means
that \isa{\isafree{{\isasymalpha}}\ {\isacharequal}{\kern0pt}\ \isaconst{rev}\ \isafree{{\isasymbeta}}\ {\isacharat}{\kern0pt}\ \isactrlbsub {\isacharbar}{\kern0pt}\isafree{ps}{\isacharbar}{\kern0pt}\isactrlesub {\isacharbar}{\kern0pt}\isafree{{\isasymalpha}}}. Since we know that after the reducing configuration,
\isa{\isaconst{nxt}\ \ensuremath{LR_0(G)}\ \isafree{q}\ {\isacharparenleft}{\kern0pt}\isaconst{Nt}\ \isafree{Y}{\isacharparenright}{\kern0pt}} is the topmost state, with Lemma~\ref{inj_nxt_dfa_LR0_if_nonempty} we know that
\isa{\isafree{X}\ {\isacharequal}{\kern0pt}\ \isaconst{Nt}\ \isafree{Y}} holds by our assumption that \isa{\isaconst{nxt}\ \ensuremath{LR_0(G)}\ \isafree{q}\ \isafree{X}} is nonempty. Moreover, from
the reducing transition assumptions we know that \isa{\isafree{u{\isacharprime}{\kern0pt}}\ {\isacharequal}{\kern0pt}\ \isafree{u}} must hold, since the remaining
input after the transition is \isa{v}, and we know that this transition does not consume symbols. From
all this, we can show the invariant is satisfied.

The last case, which is that of the finish transition, contradicts the structure of the stack word.
This is because the topmost state being \isa{\isaconst{nxt}\ \ensuremath{LR_0(G)}\ \isafree{q}\ \isafree{X}} means that it is either empty,
or a state of \isa{\ensuremath{LR_0(G)}}. This is a direct consequence of the definition of the transition
functions of \isa{\ensuremath{LR_0(G)}} and \isa{\ensuremath{\mathrm{char}(G)}}. Since \isa{{\isacharbraceleft}{\kern0pt}{\isacharbrackleft}{\kern0pt}\isaconst{S{\isacharprime}{\kern0pt}}\ {\isasymrightarrow}\ {\isasymcdot}\ {\isacharbrackright}{\kern0pt}{\isacharbraceright}{\kern0pt}} is not a state of
\isa{\ensuremath{LR_0(G)}}, this case cannot occur. The proof of the invariant is therefore complete.
\end{proof}
\end{lemma}

\begin{theorem}[Soundness of the \isa{LR{\isacharparenleft}{\kern0pt}{\isadigit{0}}{\isacharparenright}{\kern0pt}} parser]\label{P0_sound}
Every word accepted by \isa{\ensuremath{P_0(G)}} is a word in \isa{L{\isacharparenleft}{\kern0pt}\isafree{G}{\isacharparenright}{\kern0pt}}.
\begin{proof}
We begin by fixing \isa{w} and assuming that \isa{\isafree{w}\ {\isasymin}\ \ensuremath{L(P_0(G))}}. Therefore, the initial computation 
of \isa{\ensuremath{P_0(G)}} reaches the final configuration \isa{{\isacharparenleft}{\kern0pt}{\isacharbrackleft}{\kern0pt}\mbox{{\isacharbraceleft}{\kern0pt}{\isacharbrackleft}{\kern0pt}\isaconst{S{\isacharprime}{\kern0pt}}\ {\isasymrightarrow}\ {\isasymcdot}\ {\isacharbrackright}{\kern0pt}{\isacharbraceright}{\kern0pt}}{\isacharbrackright}{\kern0pt}{\isacharcomma}{\kern0pt}\ {\isacharbrackleft}{\kern0pt}{\isacharbrackright}{\kern0pt}{\isacharparenright}{\kern0pt}}. We know this 
computation has nonzero length because \isa{{\isacharbraceleft}{\kern0pt}{\isacharbrackleft}{\kern0pt}\isaconst{S{\isacharprime}{\kern0pt}}\ {\isasymrightarrow}\ {\isasymcdot}\ {\isacharbrackright}{\kern0pt}{\isacharbraceright}{\kern0pt}} is distinct from the initial state.
We also know that the last step of the computation must have been a finishing transition, meaning
there exists a state \isa{q} of \isa{\ensuremath{LR_0(G)}} such that \isa{{\isacharbrackleft}{\kern0pt}\isaconst{S{\isacharprime}{\kern0pt}}\ {\isasymrightarrow}\ {\isacharbrackleft}{\kern0pt}\isaconst{Nt}\ \isaconst{S}{\isacharbrackright}{\kern0pt}\ {\isasymcdot}\ {\isacharbrackright}{\kern0pt}\ {\isasymin}\ \isafree{q}} for which 
\isa{\ensuremath{P_0(G)}} computes \isa{\mbox{{\isacharparenleft}{\kern0pt}{\isacharbrackleft}{\kern0pt}\isaconst{init}\ \ensuremath{P_0(G)}{\isacharbrackright}{\kern0pt}{\isacharcomma}{\kern0pt}\ \isafree{w}{\isacharparenright}{\kern0pt}}\ {\isasymturnstile}P{\isacharasterisk}{\kern0pt}\ {\isacharparenleft}{\kern0pt}{\isacharbrackleft}{\kern0pt}\isafree{q}{\isacharcomma}{\kern0pt}\ \isaconst{init}\ \ensuremath{LR_0(G)}{\isacharbrackright}{\kern0pt}{\isacharcomma}{\kern0pt}\ {\isacharbrackleft}{\kern0pt}{\isacharbrackright}{\kern0pt}{\isacharparenright}{\kern0pt}}. With this, we know by 
Lemma~\ref{P0_steps_imp_stack_word} that there exists some \isa{{\isasymalpha}\ {\isacharcolon}{\kern0pt}{\isacharcolon}{\kern0pt}\ syms} with
\isa{\isafree{{\isasymalpha}}\ {\isasymTurnstile}\ {\isacharparenleft}{\kern0pt}{\isacharbrackleft}{\kern0pt}\isaconst{init}\ \ensuremath{P_0(G)}{\isacharbrackright}{\kern0pt}{\isacharcomma}{\kern0pt}\ \isafree{w}{\isacharparenright}{\kern0pt}\ {\isasymturnstile}P{\isacharasterisk}{\kern0pt}\ {\isacharparenleft}{\kern0pt}{\isacharbrackleft}{\kern0pt}\isafree{q}{\isacharcomma}{\kern0pt}\ \isaconst{init}\ \ensuremath{LR_0(G)}{\isacharbrackright}{\kern0pt}{\isacharcomma}{\kern0pt}\ {\isacharbrackleft}{\kern0pt}{\isacharbrackright}{\kern0pt}{\isacharparenright}{\kern0pt}}.

We also know by Lemma~\ref{P0_init_stack_word_length_inv} that this stack word is of the form 
\isa{\isafree{{\isasymalpha}}\ {\isacharequal}{\kern0pt}\ {\isacharbrackleft}{\kern0pt}\isafree{X}{\isacharbrackright}{\kern0pt}} for some symbol \isa{X}. By Lemma~\ref{P0_nth_is_valids_of_nth_stack_word}, we know 
that \isa{\isafree{q}\ {\isacharequal}{\kern0pt}\ \isaconst{valids}\ {\isacharbrackleft}{\kern0pt}\isafree{X}{\isacharbrackright}{\kern0pt}}, and since \isa{\mbox{{\isacharbrackleft}{\kern0pt}\isaconst{S{\isacharprime}{\kern0pt}}\ {\isasymrightarrow}\ {\isacharbrackleft}{\kern0pt}\isaconst{Nt}\ \isaconst{S}{\isacharbrackright}{\kern0pt}\ {\isasymcdot}\ {\isacharbrackright}{\kern0pt}\ {\isasymin}\ \isafree{q}}}, there exists a \isa{{\isasymgamma}} such 
that \isa{{\isacharbrackleft}{\kern0pt}\isafree{X}{\isacharbrackright}{\kern0pt}\ {\isacharequal}{\kern0pt}\ \isafree{{\isasymgamma}}\ {\isacharat}{\kern0pt}\ {\isacharbrackleft}{\kern0pt}\isaconst{Nt}\ \isaconst{S}{\isacharbrackright}{\kern0pt}} since \isa{{\isacharbrackleft}{\kern0pt}X{\isacharbrackright}{\kern0pt}} is a reliable prefix for this item. This forces 
\isa{\isafree{X}\ {\isacharequal}{\kern0pt}\ \isaconst{Nt}\ \isaconst{S}}, and by Lemma~\ref{P0_invariant} this stack word implies 
\isa{\isaconst{P{\isacharprime}{\kern0pt}}\ {\isasymturnstile}\ {\isacharbrackleft}{\kern0pt}\isaconst{Nt}\ \isaconst{S}{\isacharbrackright}{\kern0pt}\ {\isasymRightarrow}r{\isacharasterisk}{\kern0pt}\ \isaconst{map}\ \isaconst{Tm}\ \isafree{w}}. Since \isa{S{\isacharprime}{\kern0pt}} derives \isa{{\isacharbrackleft}{\kern0pt}Nt\ S{\isacharbrackright}{\kern0pt}} by definition, this means 
that \isa{\isafree{w}\ {\isasymin}\ L{\isacharparenleft}{\kern0pt}\isaconst{G{\isacharprime}{\kern0pt}}{\isacharparenright}{\kern0pt}}, and the proof is complete by Theorem~\ref{Lang_preserved}.
\end{proof}
\end{theorem}%
\end{isamarkuptext}\isamarkuptrue%
%
\isadelimdocument
%
\endisadelimdocument
%
\isatagdocument
%
\isamarkupsubsection{The Shift-Reduce Pushdown Automaton: Proving Completeness%
}
\isamarkuptrue%
%
\endisatagdocument
{\isafolddocument}%
%
\isadelimdocument
%
\endisadelimdocument
%
\isadelimproof
%
\endisadelimproof
%
\isatagproof
%
\endisatagproof
{\isafoldproof}%
%
\isadelimproof
%
\endisadelimproof
%
\begin{isamarkuptext}%
Now that we have proven the soundness of our parser, all that is left is proving its 
completeness, i.e., that every word in the grammar is accepted by the parser.

We will again need stack words to prove completeness, but constructing a valid parser computation 
from a rightmost derivation directly is quite challenging since the parser's construction, albeit
ideal for our original goal of achieving a deterministic parsing algorithm, is too far removed from 
the pure concept of rightmost derivations. We will overcome this by defining the
\concept{shift-reduce pushdown automaton} (SRPDA), proving the completeness of this machine w.r.t.\
the grammar, and then show the relation between SRPDA and \isa{\ensuremath{P_0(G)}} computations, which will allow
us to prove the parser complete. Our definition of the SRPDA is based on the lecture slides by
Petter~\cite{Petter}.

The SRPDA is an automaton that uses a CFG's symbols as its states. In order to distinguish an
initial or final state from symbols that are being parsed by the automaton, it is necessary for us
to extend the grammar's symbols by two more states. We achieve this with a new datatype
\isa{srpda{\isacharunderscore}{\kern0pt}state}:
\begin{quote}
\isa{\isakeywordONE{datatype}\ {\isacharparenleft}{\kern0pt}\isatfree{{\isacharprime}{\kern0pt}n}{\isacharcomma}{\kern0pt}\ \isatfree{{\isacharprime}{\kern0pt}t}{\isacharparenright}{\kern0pt}\ \isatconst{srpda{\isacharunderscore}{\kern0pt}state}\ {\isacharequal}{\kern0pt}\ \isaconst{Sym}\ {\isacharparenleft}{\kern0pt}{\isacharparenleft}{\kern0pt}\isatfree{{\isacharprime}{\kern0pt}n}{\isacharcomma}{\kern0pt}\ \isatfree{{\isacharprime}{\kern0pt}t}{\isacharparenright}{\kern0pt}\ \isatconst{sym}{\isacharparenright}{\kern0pt}\ {\isacharbar}{\kern0pt}\ \isaconst{Init}\ {\isacharbar}{\kern0pt}\ \isaconst{Final}}
\end{quote}

\begin{definition}[Shift-reduce pushdown automaton]
The shift-reduce pushdown automaton (SRPDA) to an arbitrary CFG \isa{\isafree{G}{\isacharcolon}{\kern0pt}{\isacharcolon}{\kern0pt}{\isacharparenleft}{\kern0pt}\isatfree{{\isacharprime}{\kern0pt}n}{\isacharcomma}{\kern0pt}\ \isatfree{{\isacharprime}{\kern0pt}t}{\isacharparenright}{\kern0pt}\ \isatconst{Cfg}}
with production set \isa{P} is the \isa{{\isacharparenleft}{\kern0pt}{\isacharparenleft}{\kern0pt}\isatfree{{\isacharprime}{\kern0pt}n}{\isacharcomma}{\kern0pt}\ \isatfree{{\isacharprime}{\kern0pt}t}{\isacharparenright}{\kern0pt}\ \isatconst{srpda{\isacharunderscore}{\kern0pt}state}{\isacharcomma}{\kern0pt}\ \isatfree{{\isacharprime}{\kern0pt}t}{\isacharparenright}{\kern0pt}\ \isatconst{gpda}}
\begin{multline*}
\isa{M\isactrlsub G\ {\isacharequal}{\kern0pt}\ {\isasymlparr}gpda{\isachardot}{\kern0pt}states\ {\isacharequal}{\kern0pt}\ UNIV{\isacharcomma}{\kern0pt}\ init\ {\isacharequal}{\kern0pt}\ Init{\isacharcomma}{\kern0pt}\ final\ {\isacharequal}{\kern0pt}\ {\isacharbraceleft}{\kern0pt}Final{\isacharbraceright}{\kern0pt}}\\
  \isa{nxt\ {\isacharequal}{\kern0pt}\ range\ {\isacharparenleft}{\kern0pt}{\isasymlambda}{\isacharparenleft}{\kern0pt}q{\isacharcomma}{\kern0pt}\ x{\isacharparenright}{\kern0pt}{\isachardot}{\kern0pt}\ {\isacharparenleft}{\kern0pt}{\isacharbrackleft}{\kern0pt}q{\isacharbrackright}{\kern0pt}{\isacharcomma}{\kern0pt}\ x{\isacharcomma}{\kern0pt}\ {\isacharbrackleft}{\kern0pt}Sym\ {\isacharparenleft}{\kern0pt}Tm\ x{\isacharparenright}{\kern0pt}{\isacharcomma}{\kern0pt}\ q{\isacharbrackright}{\kern0pt}{\isacharparenright}{\kern0pt}{\isacharparenright}{\kern0pt}{\isacharcomma}{\kern0pt}\ eps\ {\isacharequal}{\kern0pt}\ {\isasymE}\isactrlsub M{\isasymrparr}}
\end{multline*}
\end{definition}

Where \isa{UNIV} is the universe of the type \isa{{\isacharparenleft}{\kern0pt}\isatfree{{\isacharprime}{\kern0pt}n}{\isacharcomma}{\kern0pt}\ \isatfree{{\isacharprime}{\kern0pt}t}{\isacharparenright}{\kern0pt}\ \isatconst{srpda{\isacharunderscore}{\kern0pt}state}}, \isa{\isaconst{range}} denotes the
image of a function, and a term of the form (\isa{{\isasymlambda}x\isactrlsub {\isadigit{1}}\ x\isactrlsub {\isadigit{2}}\ {\isasymdots}\ x\isactrlsub n{\isachardot}{\kern0pt}\ y}) denotes a function that takes
parameters \isa{x\isactrlsub {\isadigit{1}}{\isacharcomma}{\kern0pt}\ x\isactrlsub {\isadigit{2}}{\isacharcomma}{\kern0pt}\ {\isasymdots}{\isacharcomma}{\kern0pt}\ x\isactrlsub n} and returns \isa{y}. 

We define three types of transition once again:
\begin{itemize}
\item In a \concept{reading} transition, \isa{M\isactrlsub G} reads a symbol from the input, and pushes it onto the 
stack. This corresponds, as is the case with the other automata we have defined until now, to the 
\isa{nxt} relation. Since we define the relation as the range of the mapping, this transition is
completely independent of the current state.
\item \concept{Reducing} transitions allow the SRPDA to reduce a topmost stack string \isa{{\isasymalpha}} to a 
nonterminal \isa{A} if the production \isa{{\isacharparenleft}{\kern0pt}A{\isacharcomma}{\kern0pt}\ {\isasymalpha}{\isacharparenright}{\kern0pt}} is in \isa{G}. Formally, this is the set
\[ \isa{{\isacharbraceleft}{\kern0pt}{\isacharparenleft}{\kern0pt}\isaconst{map}\ \isaconst{Sym}\ {\isacharparenleft}{\kern0pt}\isaconst{rev}\ \isabound{{\isasymalpha}}{\isacharparenright}{\kern0pt}\ {\isacharat}{\kern0pt}\ {\isacharbrackleft}{\kern0pt}\isabound{q}{\isacharbrackright}{\kern0pt}{\isacharcomma}{\kern0pt}\ {\isacharbrackleft}{\kern0pt}\isaconst{Sym}\ {\isacharparenleft}{\kern0pt}\isaconst{Nt}\ \isabound{A}{\isacharparenright}{\kern0pt}{\isacharcomma}{\kern0pt}\ \isabound{q}{\isacharbrackright}{\kern0pt}{\isacharparenright}{\kern0pt}\ {\isacharbar}{\kern0pt}\ {\isacharparenleft}{\kern0pt}\isabound{A}{\isacharcomma}{\kern0pt}\ \isabound{{\isasymalpha}}{\isacharparenright}{\kern0pt}\ {\isasymin}\ \isaconst{P}{\isacharbraceright}{\kern0pt}}. \]
\item Lastly, the \concept{finishing} transition, very similarly to the \isa{LR{\isacharparenleft}{\kern0pt}{\isadigit{0}}{\isacharparenright}{\kern0pt}} parser, 
signals that the consumed input has successfully been reduced to the start symbol. This corresponds
to the singleton set \isa{{\isacharbraceleft}{\kern0pt}{\isacharparenleft}{\kern0pt}{\isacharbrackleft}{\kern0pt}\isaconst{Sym}\ {\isacharparenleft}{\kern0pt}\isaconst{Nt}\ \isaconst{S}{\isacharparenright}{\kern0pt}{\isacharcomma}{\kern0pt}\ \isaconst{Init}{\isacharbrackright}{\kern0pt}{\isacharcomma}{\kern0pt}\ {\isacharbrackleft}{\kern0pt}\isaconst{Final}{\isacharbrackright}{\kern0pt}{\isacharparenright}{\kern0pt}{\isacharbraceright}{\kern0pt}}.
\end{itemize}
\isa{{\isasymE}\isactrlsub M} is once again the union of reducing transitions and the singleton finishing transition set.

The SRPDA, unlike our \isa{LR{\isacharparenleft}{\kern0pt}{\isadigit{0}}{\isacharparenright}{\kern0pt}} parser, computes reductions very similarly to how its grammar
derives words, essentially mimicking the grammar's derivation of a word in reverse with little overhead.
It is worth noting that this automaton is not a GPDA in the strict sense of our definition. This is due to the 
fact that we define the alphabet of a CFG through types rather than sets. Since our SRPDA states 
are based on the grammar's symbols, in order to satisfy the assumption that the set of states is finite, one 
would need to prove that the types for terminals and nonterminals are both finite. If we restricted
the SRPDA to finite types and wanted to use it on the same grammar as our parser, the parser's type 
variables would need to be finite as well. Applying such a constraint to a type is much more restrictive 
than to a set, and we purposefully avoid this so our parser can be used for grammars of arbitrary 
types. This is because since we work with grammars with finitely many productions, there are also
finitely many symbols that are actually produced by the grammar regardless of the finiteness of the
type itself. The fact that the SRPDA does not fulfill the finiteness condition as we specified it is
not problematic for our purposes, but it is a discrepancy with our original GPDA definition that is
worth justifying.

We will now prove that the SRPDA to our reduced grammar \isa{G} is complete, i.e., every word in 
\isa{L{\isacharparenleft}{\kern0pt}\isafree{G}{\isacharparenright}{\kern0pt}} is accepted by \isa{\isaconst{M\isactrlsub G}}. We denote \isa{M\isactrlsub G} steps and their reflexive transitive
closure as \isa{{\isasymturnstile}M} and \isa{{\isasymturnstile}M{\isacharasterisk}{\kern0pt}} respectively.

\begin{lemma}\label{prefix_consumable}
\isa{{\isacharparenleft}{\kern0pt}\isafree{X}\ {\isacharhash}{\kern0pt}\ \isafree{{\isasymalpha}}{\isacharcomma}{\kern0pt}\ \isafree{u}\ {\isacharat}{\kern0pt}\ \isafree{v}{\isacharparenright}{\kern0pt}\ {\isasymturnstile}M{\isacharasterisk}{\kern0pt}\ {\isacharparenleft}{\kern0pt}\isaconst{map}\ {\isacharparenleft}{\kern0pt}\isaconst{Sym}\ {\isasymcirc}\ \isaconst{Tm}{\isacharparenright}{\kern0pt}\ {\isacharparenleft}{\kern0pt}\isaconst{rev}\ \isafree{u}{\isacharparenright}{\kern0pt}\ {\isacharat}{\kern0pt}\ \isafree{X}\ {\isacharhash}{\kern0pt}\ \isafree{{\isasymalpha}}{\isacharcomma}{\kern0pt}\ \isafree{v}{\isacharparenright}{\kern0pt}} holds for any SRPDA state list 
\isa{X\ {\isacharhash}{\kern0pt}\ {\isasymalpha}} and input \isa{u\ {\isacharat}{\kern0pt}\ v\ {\isacharcolon}{\kern0pt}{\isacharcolon}{\kern0pt}\ {\isacharprime}{\kern0pt}t\ list}.
\begin{proof}
By induction on \isa{u} for arbitrary \isa{X} and \isa{{\isasymalpha}}.
\end{proof}
\end{lemma}

\begin{lemma}\label{Tms_on_stack_imp_consumed}
\isa{{\isacharparenleft}{\kern0pt}\isafree{{\isasymalpha}}{\isacharcomma}{\kern0pt}\ \isafree{u}\ {\isacharat}{\kern0pt}\ \isafree{v}\ {\isacharat}{\kern0pt}\ \isafree{w}{\isacharparenright}{\kern0pt}\ {\isasymturnstile}M{\isacharasterisk}{\kern0pt}\ {\isacharparenleft}{\kern0pt}\isaconst{map}\ {\isacharparenleft}{\kern0pt}\isaconst{Sym}\ {\isasymcirc}\ \isaconst{Tm}{\isacharparenright}{\kern0pt}\ {\isacharparenleft}{\kern0pt}\isaconst{rev}\ \isafree{v}{\isacharparenright}{\kern0pt}\ {\isacharat}{\kern0pt}\ \isafree{{\isasymbeta}}{\isacharcomma}{\kern0pt}\ \isafree{w}{\isacharparenright}{\kern0pt}} implies
\[ \isa{{\isacharparenleft}{\kern0pt}\isafree{{\isasymalpha}}{\isacharcomma}{\kern0pt}\ \isafree{u}\ {\isacharat}{\kern0pt}\ \isafree{v}\ {\isacharat}{\kern0pt}\ \isafree{w}{\isacharparenright}{\kern0pt}\ {\isasymturnstile}M{\isacharasterisk}{\kern0pt}\ {\isacharparenleft}{\kern0pt}\isafree{{\isasymbeta}}{\isacharcomma}{\kern0pt}\ \isafree{v}\ {\isacharat}{\kern0pt}\ \isafree{w}{\isacharparenright}{\kern0pt}}. \]
\begin{proof}
The proof is by induction on \isa{v} for arbitrary \isa{u} and \isa{w}.

The base case is trivial.

In the inductive case, \isa{v\ {\isacharequal}{\kern0pt}\ x\ {\isacharat}{\kern0pt}\ {\isacharbrackleft}{\kern0pt}a{\isacharbrackright}{\kern0pt}} for some \isa{x\ {\isacharcolon}{\kern0pt}{\isacharcolon}{\kern0pt}\ {\isacharprime}{\kern0pt}t\ list} and \isa{a\ {\isacharcolon}{\kern0pt}{\isacharcolon}{\kern0pt}\ {\isacharprime}{\kern0pt}t}. Therefore, we have
\isa{\mbox{{\isacharparenleft}{\kern0pt}\isafree{{\isasymalpha}}{\isacharcomma}{\kern0pt}\ \isafree{u}\ {\isacharat}{\kern0pt}\ {\isacharparenleft}{\kern0pt}\isafree{v}\ {\isacharat}{\kern0pt}\ {\isacharbrackleft}{\kern0pt}\isafree{a}{\isacharbrackright}{\kern0pt}{\isacharparenright}{\kern0pt}\ {\isacharat}{\kern0pt}\ \isafree{w}{\isacharparenright}{\kern0pt}\ {\isasymturnstile}M{\isacharasterisk}{\kern0pt}\ {\isacharparenleft}{\kern0pt}\isaconst{map}\ {\isacharparenleft}{\kern0pt}\isaconst{Sym}\ {\isasymcirc}\ \isaconst{Tm}{\isacharparenright}{\kern0pt}\ {\isacharparenleft}{\kern0pt}\isaconst{rev}\ {\isacharparenleft}{\kern0pt}\isafree{v}\ {\isacharat}{\kern0pt}\ {\isacharbrackleft}{\kern0pt}\isafree{a}{\isacharbrackright}{\kern0pt}{\isacharparenright}{\kern0pt}{\isacharparenright}{\kern0pt}\ {\isacharat}{\kern0pt}\ \isafree{{\isasymbeta}}{\isacharcomma}{\kern0pt}\ \isafree{w}{\isacharparenright}{\kern0pt}}}.

Since this computation is of nonzero length, the final step must have been a reading transition, 
meaning that
\isa{{\isacharparenleft}{\kern0pt}\isafree{{\isasymalpha}}{\isacharcomma}{\kern0pt}\ \isafree{u}\ {\isacharat}{\kern0pt}\ \isafree{v}\ {\isacharat}{\kern0pt}\ {\isacharbrackleft}{\kern0pt}\isafree{a}{\isacharbrackright}{\kern0pt}\ {\isacharat}{\kern0pt}\ \isafree{w}{\isacharparenright}{\kern0pt}\ {\isasymturnstile}M{\isacharasterisk}{\kern0pt}\ {\isacharparenleft}{\kern0pt}\isaconst{map}\ {\isacharparenleft}{\kern0pt}\isaconst{Sym}\ {\isasymcirc}\ \isaconst{Tm}{\isacharparenright}{\kern0pt}\ {\isacharparenleft}{\kern0pt}\isaconst{rev}\ \isafree{v}{\isacharparenright}{\kern0pt}\ {\isacharat}{\kern0pt}\ \isafree{{\isasymbeta}}{\isacharcomma}{\kern0pt}\ \isafree{a}\ {\isacharhash}{\kern0pt}\ \isafree{w}{\isacharparenright}{\kern0pt}} holds.

The claim then follows directly by the induction hypothesis.
\end{proof}
\end{lemma}

\begin{lemma}[SRPDA invariant]\label{srpda.invariant}
If \isa{\isaconst{P}\ {\isasymturnstile}\ \isafree{{\isasymalpha}}\ {\isasymRightarrow}r{\isacharasterisk}{\kern0pt}\ \isaconst{map}\ \isaconst{Tm}\ \isafree{w}}, then 
\[ \isa{{\isacharparenleft}{\kern0pt}{\isacharbrackleft}{\kern0pt}\isaconst{Init}{\isacharbrackright}{\kern0pt}{\isacharcomma}{\kern0pt}\ \isafree{w}{\isacharparenright}{\kern0pt}\ {\isasymturnstile}M{\isacharasterisk}{\kern0pt}\ {\isacharparenleft}{\kern0pt}\isaconst{map}\ \isaconst{Sym}\ {\isacharparenleft}{\kern0pt}\isaconst{rev}\ \isafree{{\isasymalpha}}{\isacharparenright}{\kern0pt}\ {\isacharat}{\kern0pt}\ {\isacharbrackleft}{\kern0pt}\isaconst{Init}{\isacharbrackright}{\kern0pt}{\isacharcomma}{\kern0pt}\ {\isacharbrackleft}{\kern0pt}{\isacharbrackright}{\kern0pt}{\isacharparenright}{\kern0pt}}. \]
\begin{proof}
We use reverse induction on the length of the derivation. 

The reflexive case is trivial by Lemma~\ref{prefix_consumable}.

For the transitive case, we consider the first step of the derivation.

We know that \isa{{\isasymalpha}\ {\isacharequal}{\kern0pt}\ {\isasymbeta}\ {\isacharat}{\kern0pt}\ Nt\ A\ {\isacharhash}{\kern0pt}\ map\ Tm\ v} for some \isa{A\ {\isacharcolon}{\kern0pt}{\isacharcolon}{\kern0pt}\ {\isacharprime}{\kern0pt}n}, \isa{{\isasymbeta}\ {\isacharcolon}{\kern0pt}{\isacharcolon}{\kern0pt}\ syms}, and \isa{v\ {\isacharcolon}{\kern0pt}{\isacharcolon}{\kern0pt}\ {\isacharprime}{\kern0pt}t\ list}, and the 
right sentential form produced is \isa{\isafree{{\isasymbeta}}\ {\isacharat}{\kern0pt}\ \isafree{{\isasymgamma}}\ {\isacharat}{\kern0pt}\ \isaconst{map}\ \isaconst{Tm}\ \isafree{v}} for some \isa{{\isasymgamma}\ {\isacharcolon}{\kern0pt}{\isacharcolon}{\kern0pt}\ syms} such that 
\isa{{\isacharparenleft}{\kern0pt}\isafree{A}{\isacharcomma}{\kern0pt}\ \isafree{{\isasymgamma}}{\isacharparenright}{\kern0pt}\ {\isasymin}\ \isaconst{P}}. By the induction hypothesis we know that 
\isa{{\isacharparenleft}{\kern0pt}{\isacharbrackleft}{\kern0pt}\isaconst{Init}{\isacharbrackright}{\kern0pt}{\isacharcomma}{\kern0pt}\ \isafree{w}{\isacharparenright}{\kern0pt}\ {\isasymturnstile}M{\isacharasterisk}{\kern0pt}\ {\isacharparenleft}{\kern0pt}\isaconst{map}\ \isaconst{Sym}\ {\isacharparenleft}{\kern0pt}\isaconst{rev}\ {\isacharparenleft}{\kern0pt}\isafree{{\isasymbeta}}\ {\isacharat}{\kern0pt}\ \isafree{{\isasymgamma}}\ {\isacharat}{\kern0pt}\ \isaconst{map}\ \isaconst{Tm}\ \isafree{v}{\isacharparenright}{\kern0pt}{\isacharparenright}{\kern0pt}\ {\isacharat}{\kern0pt}\ {\isacharbrackleft}{\kern0pt}\isaconst{Init}{\isacharbrackright}{\kern0pt}{\isacharcomma}{\kern0pt}\ {\isacharbrackleft}{\kern0pt}{\isacharbrackright}{\kern0pt}{\isacharparenright}{\kern0pt}} holds.
By Lemma~\ref{Tms_on_stack_imp_consumed}, this implies that the starting configuration for \isa{w} also 
reaches \isa{{\isacharparenleft}{\kern0pt}\isaconst{map}\ \isaconst{Sym}\ {\isacharparenleft}{\kern0pt}\isaconst{rev}\ \isafree{{\isasymdelta}}\ {\isacharat}{\kern0pt}\ \isaconst{rev}\ \isafree{{\isasymgamma}}{\isacharparenright}{\kern0pt}\ {\isacharat}{\kern0pt}\ {\isacharbrackleft}{\kern0pt}\isaconst{Init}{\isacharbrackright}{\kern0pt}{\isacharcomma}{\kern0pt}\ \isafree{v}{\isacharparenright}{\kern0pt}}. The SRPDA can then reduce
\isa{\mbox{\isaconst{rev}\ \isafree{{\isasymdelta}}}} to \isa{A}, and push \isa{\isaconst{rev}\ \isafree{v}} onto the stack by
Lemma~\ref{prefix_consumable}. Since this is equivalent to \isa{rev\ {\isasymalpha}}, the proof is now complete.
\end{proof}
\end{lemma}

We can now use this invariant to show completeness.

\begin{lemma}[Completeness of the shift-reduce PDA]\label{srpda_complete}
If a word is in \isa{L{\isacharparenleft}{\kern0pt}\isafree{G}{\isacharparenright}{\kern0pt}}, it is accepted by \isa{\isaconst{M\isactrlsub G}}.
\begin{proof}
We begin by fixing some \isa{w} in 
\isa{L{\isacharparenleft}{\kern0pt}\isafree{G}{\isacharparenright}{\kern0pt}}. This implies that \isa{\isaconst{P}\ {\isasymturnstile}\ {\isacharbrackleft}{\kern0pt}\isaconst{Nt}\ \isaconst{S}{\isacharbrackright}{\kern0pt}\ {\isasymRightarrow}r{\isacharasterisk}{\kern0pt}\ \isaconst{map}\ \isaconst{Tm}\ \isafree{w}}.
By Lemma~\ref{srpda.invariant}, this derivation implies the existence of the computation
\[ \isa{{\isacharparenleft}{\kern0pt}{\isacharbrackleft}{\kern0pt}\isaconst{Init}{\isacharbrackright}{\kern0pt}{\isacharcomma}{\kern0pt}\ \isafree{w}{\isacharparenright}{\kern0pt}\ {\isasymturnstile}M{\isacharasterisk}{\kern0pt}\ {\isacharparenleft}{\kern0pt}{\isacharbrackleft}{\kern0pt}\isaconst{Sym}\ {\isacharparenleft}{\kern0pt}\isaconst{Nt}\ \isaconst{S}{\isacharparenright}{\kern0pt}{\isacharcomma}{\kern0pt}\ \isaconst{Init}{\isacharbrackright}{\kern0pt}{\isacharcomma}{\kern0pt}\ {\isacharbrackleft}{\kern0pt}{\isacharbrackright}{\kern0pt}{\isacharparenright}{\kern0pt}}. \]
Trivially, the configuration \isa{{\isacharparenleft}{\kern0pt}{\isacharbrackleft}{\kern0pt}\isaconst{Sym}\ {\isacharparenleft}{\kern0pt}\isaconst{Nt}\ \isaconst{S}{\isacharparenright}{\kern0pt}{\isacharcomma}{\kern0pt}\ \isaconst{Init}{\isacharbrackright}{\kern0pt}{\isacharcomma}{\kern0pt}\ {\isacharbrackleft}{\kern0pt}{\isacharbrackright}{\kern0pt}{\isacharparenright}{\kern0pt}} reaches 
\isa{{\isacharparenleft}{\kern0pt}{\isacharbrackleft}{\kern0pt}\isaconst{Final}{\isacharbrackright}{\kern0pt}{\isacharcomma}{\kern0pt}\ {\isacharbrackleft}{\kern0pt}{\isacharbrackright}{\kern0pt}{\isacharparenright}{\kern0pt}} in a single step via a finishing transition, meaning \isa{\isafree{w}\ {\isasymin}\ \ensuremath{L(M_G)}}.
The proof is thus complete.
\end{proof}
\end{lemma}

We will now show the relation between \isa{\isaconst{M\isactrlsub G}} and \isa{\ensuremath{P_0(G)}} which will allow us to prove the 
parser's correctness. It is worth noting that we are using the SRPDA for \isa{G}, not for \isa{G{\isacharprime}{\kern0pt}}. 
We first show an auxiliary lemma.

\begin{lemma}\label{nonempty_valids_imp_nonempty_valids_prefix}
If \isa{\isaconst{valids}\ {\isacharparenleft}{\kern0pt}\isafree{{\isasymalpha}}\ {\isacharat}{\kern0pt}\ \isafree{{\isasymbeta}}{\isacharparenright}{\kern0pt}\ {\isasymnoteq}\ {\isasymemptyset}}, then \isa{\isaconst{valids}\ \isafree{{\isasymalpha}}\ {\isasymnoteq}\ {\isasymemptyset}}
\begin{proof}
Since \isa{\isaconst{valids}\ {\isacharparenleft}{\kern0pt}\isafree{{\isasymalpha}}\ {\isacharat}{\kern0pt}\ \isafree{{\isasymbeta}}{\isacharparenright}{\kern0pt}\ {\isasymnoteq}\ {\isasymemptyset}}, there exists an item \isa{{\isacharbrackleft}{\kern0pt}\isafree{X}\ {\isasymrightarrow}\ \isafree{{\isasymgamma}}\ {\isasymcdot}\ \isafree{{\isasymdelta}}{\isacharbrackright}{\kern0pt}} valid for \isa{{\isasymalpha}\ {\isacharat}{\kern0pt}\ {\isasymbeta}}. 
Therefore, there exists a \isa{{\isasymgamma}{\isacharprime}{\kern0pt}} such that \isa{{\isasymalpha}\ {\isacharat}{\kern0pt}\ {\isasymbeta}\ {\isacharequal}{\kern0pt}\ {\isasymgamma}{\isacharprime}{\kern0pt}\ {\isacharat}{\kern0pt}\ {\isasymgamma}}. We now distinguish cases on whether \isa{{\isasymalpha}} is 
a prefix of \isa{{\isasymgamma}{\isacharprime}{\kern0pt}}. If it is, we show there must exist a valid item for \isa{{\isasymalpha}} by 
Lemma~\ref{derivers_substring_reliable}. Otherwise, there exists an \isa{{\isasymalpha}{\isacharprime}{\kern0pt}} with \isa{{\isasymgamma}{\isacharprime}{\kern0pt}\ {\isacharat}{\kern0pt}\ {\isasymalpha}{\isacharprime}{\kern0pt}\ {\isacharequal}{\kern0pt}\ {\isasymalpha}} and 
\isa{{\isasymalpha}{\isacharprime}{\kern0pt}\ {\isacharat}{\kern0pt}\ {\isasymbeta}\ {\isacharequal}{\kern0pt}\ {\isasymgamma}}, and thus item \isa{{\isacharbrackleft}{\kern0pt}\isafree{X}\ {\isasymrightarrow}\ \isafree{{\isasymalpha}{\isacharprime}{\kern0pt}}\ {\isasymcdot}\ \isafree{{\isasymbeta}}\ {\isacharat}{\kern0pt}\ \isafree{{\isasymdelta}}{\isacharbrackright}{\kern0pt}} is valid for \isa{{\isasymalpha}}, completing the proof.
\end{proof}
\end{lemma}

\begin{lemma}\label{MG_steps_imp_stack_word}
If \isa{{\isacharparenleft}{\kern0pt}{\isacharbrackleft}{\kern0pt}\isaconst{Init}{\isacharbrackright}{\kern0pt}{\isacharcomma}{\kern0pt}\ \isafree{u}{\isacharparenright}{\kern0pt}\ {\isasymturnstile}M{\isacharasterisk}{\kern0pt}\ {\isacharparenleft}{\kern0pt}\isaconst{map}\ \isaconst{Sym}\ \isafree{{\isasymalpha}}\ {\isacharat}{\kern0pt}\ {\isacharbrackleft}{\kern0pt}\isaconst{Init}{\isacharbrackright}{\kern0pt}{\isacharcomma}{\kern0pt}\ \isafree{v}{\isacharparenright}{\kern0pt}} and 
\isa{\isaconst{valids}\ {\isacharparenleft}{\kern0pt}\isaconst{rev}\ \isafree{{\isasymalpha}}{\isacharparenright}{\kern0pt}\ {\isasymnoteq}\ {\isasymemptyset}}, there exists a stack \isa{qs} such that 
\isa{\isafree{{\isasymalpha}}\ {\isasymTurnstile}\ {\isacharparenleft}{\kern0pt}{\isacharbrackleft}{\kern0pt}\isaconst{init}\ \ensuremath{P_0(G)}{\isacharbrackright}{\kern0pt}{\isacharcomma}{\kern0pt}\ \isafree{u}{\isacharparenright}{\kern0pt}\ {\isasymturnstile}P{\isacharasterisk}{\kern0pt}\ {\isacharparenleft}{\kern0pt}\isaconst{valids}\ {\isacharparenleft}{\kern0pt}\isaconst{rev}\ \isafree{{\isasymalpha}}{\isacharparenright}{\kern0pt}\ {\isacharhash}{\kern0pt}\ \isafree{qs}{\isacharcomma}{\kern0pt}\ \isafree{v}{\isacharparenright}{\kern0pt}}.
\begin{proof}
We induct on the length of the computation for arbitrary \isa{{\isasymalpha}} and \isa{v}.

The reflexive case holds trivially for \isa{qs\ {\isacharequal}{\kern0pt}\ {\isacharbrackleft}{\kern0pt}{\isacharbrackright}{\kern0pt}}, since the initial state is \isa{\isaconst{valids}\ {\isacharbrackleft}{\kern0pt}{\isacharbrackright}{\kern0pt}}.

In the transitive case, we distinguish the type of transition for the final step of the computation.

If the final step is a reading transition, we know the second-to-last configuration is of the form 
\isa{c\ {\isacharequal}{\kern0pt}\ {\isacharparenleft}{\kern0pt}x\ {\isacharhash}{\kern0pt}\ xs{\isacharcomma}{\kern0pt}\ a\ {\isacharhash}{\kern0pt}\ v{\isacharparenright}{\kern0pt}}. We then distinguish cases based on whether \isa{xs} is empty. 

If it is, we know \isa{{\isasymalpha}} is \isa{{\isacharbrackleft}{\kern0pt}\isaconst{Sym}\ {\isacharparenleft}{\kern0pt}\isaconst{Tm}\ \isafree{a}{\isacharparenright}{\kern0pt}{\isacharbrackright}{\kern0pt}}. We can therefore show that 
\isa{\isafree{{\isasymalpha}}\ {\isasymTurnstile}\ {\isacharparenleft}{\kern0pt}\mbox{{\isacharbrackleft}{\kern0pt}\isaconst{init}\ \ensuremath{P_0(G)}{\isacharbrackright}{\kern0pt}}{\isacharcomma}{\kern0pt}\ \isafree{a}\ {\isacharhash}{\kern0pt}\ \isafree{v}{\isacharparenright}{\kern0pt}\ {\isasymturnstile}P{\isacharasterisk}{\kern0pt}\ {\isacharparenleft}{\kern0pt}{\isacharbrackleft}{\kern0pt}\isaconst{valids}\ {\isacharbrackleft}{\kern0pt}\isaconst{Tm}\ \isafree{a}{\isacharbrackright}{\kern0pt}{\isacharcomma}{\kern0pt}\ \isaconst{init}\ \ensuremath{P_0(G)}{\isacharbrackright}{\kern0pt}{\isacharcomma}{\kern0pt}\ \isafree{v}{\isacharparenright}{\kern0pt}}
by first showing that \isa{{\isacharparenleft}{\kern0pt}{\isacharbrackleft}{\kern0pt}\isaconst{init}\ \ensuremath{P_0(G)}{\isacharbrackright}{\kern0pt}{\isacharcomma}{\kern0pt}\ \isafree{a}\ {\isacharhash}{\kern0pt}\ \isafree{v}{\isacharparenright}{\kern0pt}\ {\isasymturnstile}P\ \mbox{{\isacharparenleft}{\kern0pt}{\isacharbrackleft}{\kern0pt}\isaconst{valids}\ {\isacharbrackleft}{\kern0pt}\isaconst{Tm}\ \isafree{a}{\isacharbrackright}{\kern0pt}{\isacharcomma}{\kern0pt}\ \isaconst{init}\ \ensuremath{P_0(G)}{\isacharbrackright}{\kern0pt}{\isacharcomma}{\kern0pt}\ \isafree{v}{\isacharparenright}{\kern0pt}}}
holds, using the fact that \isa{\isaconst{init}\ \ensuremath{P_0(G)}} is \isa{\isaconst{valids}\ {\isacharbrackleft}{\kern0pt}{\isacharbrackright}{\kern0pt}}, the fact that 
\isa{\isaconst{valids}\ \isafree{{\isasymalpha}}\ {\isasymnoteq}\ {\isasymemptyset}} and Lemma~\ref{nxt_dfa_LR0_shift_is_valids_app}. This completes the empty case.

If \isa{xs} on the other hand is nonempty, we know that \isa{c\ {\isacharequal}{\kern0pt}\ {\isacharparenleft}{\kern0pt}map\ Sym\ {\isacharparenleft}{\kern0pt}tl\ {\isasymalpha}{\isacharparenright}{\kern0pt}\ {\isacharat}{\kern0pt}\ {\isacharbrackleft}{\kern0pt}gpda{\isachardot}{\kern0pt}init\ M\isactrlsub G{\isacharbrackright}{\kern0pt}{\isacharcomma}{\kern0pt}\ a\ {\isacharhash}{\kern0pt}\ v{\isacharparenright}{\kern0pt}}.
From Lemma~\ref{nonempty_valids_imp_nonempty_valids_prefix} we know that 
\isa{\mbox{\isaconst{valids}\ {\isacharparenleft}{\kern0pt}\isaconst{rev}\ {\isacharparenleft}{\kern0pt}\isaconst{tl}\ \isafree{{\isasymalpha}}{\isacharparenright}{\kern0pt}{\isacharparenright}{\kern0pt}}\ {\isasymnoteq}\ {\isasymemptyset}} by our assumption that \isa{\mbox{\isaconst{valids}\ {\isacharparenleft}{\kern0pt}\isaconst{rev}\ \isafree{{\isasymalpha}}{\isacharparenright}{\kern0pt}}\ {\isasymnoteq}\ {\isasymemptyset}}. We can 
therefore use the induction hypothesis to obtain some \isa{qs} with
\isa{\isaconst{tl}\ \isafree{{\isasymalpha}}\ {\isasymTurnstile}\ {\isacharparenleft}{\kern0pt}{\isacharbrackleft}{\kern0pt}\isaconst{init}\ \ensuremath{P_0(G)}{\isacharbrackright}{\kern0pt}{\isacharcomma}{\kern0pt}\ \isafree{u}{\isacharparenright}{\kern0pt}\ {\isasymturnstile}P{\isacharasterisk}{\kern0pt}\ {\isacharparenleft}{\kern0pt}\isaconst{valids}\ {\isacharparenleft}{\kern0pt}\isaconst{rev}\ \isafree{{\isasymbeta}}{\isacharparenright}{\kern0pt}\ {\isacharhash}{\kern0pt}\ \isafree{qs}{\isacharcomma}{\kern0pt}\ \isafree{a}\ {\isacharhash}{\kern0pt}\ \isafree{v}{\isacharparenright}{\kern0pt}}.
The proof can then be finished for this case by showing
\[ \isa{\isaconst{Tm}\ \isafree{a}\ {\isacharhash}{\kern0pt}\ \isaconst{tl}\ \isafree{{\isasymalpha}}\ {\isasymTurnstile}\ {\isacharparenleft}{\kern0pt}{\isacharbrackleft}{\kern0pt}\isaconst{init}\ \ensuremath{P_0(G)}{\isacharbrackright}{\kern0pt}{\isacharcomma}{\kern0pt}\ \isafree{u}{\isacharparenright}{\kern0pt}\ {\isasymturnstile}P{\isacharasterisk}{\kern0pt}\ {\isacharparenleft}{\kern0pt}\isaconst{valids}\ {\isacharparenleft}{\kern0pt}\isaconst{rev}\ \isafree{{\isasymalpha}}{\isacharparenright}{\kern0pt}\ {\isacharhash}{\kern0pt}\ \isaconst{valids}\ {\isacharparenleft}{\kern0pt}\isaconst{rev}\ \isafree{{\isasymbeta}}{\isacharparenright}{\kern0pt}\ {\isacharhash}{\kern0pt}\ \isafree{qs}{\isacharcomma}{\kern0pt}\ \isafree{v}{\isacharparenright}{\kern0pt}} \]
analogously to the case where \isa{\isafree{xs}\ {\isacharequal}{\kern0pt}\ {\isacharbrackleft}{\kern0pt}{\isacharbrackright}{\kern0pt}}, completing the shifting transition case.

For the reducing transition case, we know that \isa{c\ {\isacharequal}{\kern0pt}\ {\isacharparenleft}{\kern0pt}map\ Sym\ {\isacharparenleft}{\kern0pt}rev\ {\isasymbeta}{\isacharparenright}{\kern0pt}\ {\isacharat}{\kern0pt}\ map\ Sym\ {\isasymgamma}\ {\isacharat}{\kern0pt}\ {\isacharbrackleft}{\kern0pt}gpda{\isachardot}{\kern0pt}init\ M\isactrlsub G{\isacharbrackright}{\kern0pt}{\isacharcomma}{\kern0pt}\ v{\isacharparenright}{\kern0pt}}.
and \isa{\isafree{{\isasymalpha}}\ {\isacharequal}{\kern0pt}\ \isaconst{Nt}\ \isafree{A}\ {\isacharhash}{\kern0pt}\ \isafree{{\isasymgamma}}}. Since \isa{A} produces \isa{{\isasymbeta}} and we originally assumed that \isa{\isaconst{valids}\ \isafree{{\isasymalpha}}\ {\isasymnoteq}\ {\isasymemptyset}},
we can use the definition of reliable prefixes to show that 
\isa{{\isacharbrackleft}{\kern0pt}\isafree{A}\ {\isasymrightarrow}\ \isafree{{\isasymbeta}}\ {\isasymcdot}\ {\isacharbrackright}{\kern0pt}\ {\isasymin}\ \isaconst{valids}\ {\isacharparenleft}{\kern0pt}\isaconst{rev}\ \isafree{{\isasymgamma}}\ {\isacharat}{\kern0pt}\ \isafree{{\isasymbeta}}{\isacharparenright}{\kern0pt}}. We can now use the IH once again to obtain some \isa{qs} 
such that 
\begin{equation}\label{MG_sw}
\isa{\isaconst{rev}\ \isafree{{\isasymbeta}}\ {\isacharat}{\kern0pt}\ \isafree{{\isasymgamma}}\ {\isasymTurnstile}\ {\isacharparenleft}{\kern0pt}{\isacharbrackleft}{\kern0pt}\isaconst{init}\ \ensuremath{P_0(G)}{\isacharbrackright}{\kern0pt}{\isacharcomma}{\kern0pt}\ \isafree{u}{\isacharparenright}{\kern0pt}\ {\isasymturnstile}P{\isacharasterisk}{\kern0pt}\ {\isacharparenleft}{\kern0pt}\isaconst{valids}\ {\isacharparenleft}{\kern0pt}\isaconst{rev}\ \isafree{{\isasymgamma}}\ {\isacharat}{\kern0pt}\ \isafree{{\isasymbeta}}{\isacharparenright}{\kern0pt}\ {\isacharhash}{\kern0pt}\ \isafree{qs}{\isacharcomma}{\kern0pt}\ \isafree{v}{\isacharparenright}{\kern0pt}}.
\end{equation}

In order to show the implication holds, it suffices to show now that 
\begin{multline*}
\isa{\isaconst{Nt}\ \isafree{A}\ {\isacharhash}{\kern0pt}\ \isafree{{\isasymgamma}}}\ \isa{{\isasymTurnstile}}\ \isa{{\isacharparenleft}{\kern0pt}{\isacharbrackleft}{\kern0pt}\isaconst{init}\ \ensuremath{P_0(G)}{\isacharbrackright}{\kern0pt}{\isacharcomma}{\kern0pt}\ \isafree{u}{\isacharparenright}{\kern0pt}}\\
      \isa{{\isasymturnstile}P{\isacharasterisk}{\kern0pt}}\ \isa{{\isacharparenleft}{\kern0pt}\isaconst{valids}\ {\isacharparenleft}{\kern0pt}\isaconst{rev}\ \isafree{{\isasymgamma}}\ {\isacharat}{\kern0pt}\ {\isacharbrackleft}{\kern0pt}\isaconst{Nt}\ \isafree{A}{\isacharbrackright}{\kern0pt}{\isacharparenright}{\kern0pt}\ {\isacharhash}{\kern0pt}\ \isactrlbsub {\isacharbar}{\kern0pt}\isafree{{\isasymbeta}}{\isacharbar}{\kern0pt}\isactrlesub {\isacharbar}{\kern0pt}{\isacharparenleft}{\kern0pt}\isaconst{valids}\ {\isacharparenleft}{\kern0pt}\isaconst{rev}\ \isafree{{\isasymgamma}}\ {\isacharat}{\kern0pt}\ \isafree{{\isasymbeta}}{\isacharparenright}{\kern0pt}\ {\isacharhash}{\kern0pt}\ \isafree{qs}{\isacharparenright}{\kern0pt}{\isacharcomma}{\kern0pt}\ \isafree{v}{\isacharparenright}{\kern0pt}}.
\end{multline*}

We first show the existence of the reduction step 
\begin{multline}\label{redstep}
  \isa{{\isacharparenleft}{\kern0pt}\isaconst{valids}\ {\isacharparenleft}{\kern0pt}\isaconst{rev}\ \isafree{{\isasymgamma}}\ {\isacharat}{\kern0pt}\ \isafree{{\isasymbeta}}{\isacharparenright}{\kern0pt}\ {\isacharhash}{\kern0pt}\ \isafree{qs}{\isacharcomma}{\kern0pt}\ \isafree{v}{\isacharparenright}{\kern0pt}}\\
  \isa{{\isasymturnstile}P}\ \isa{{\isacharparenleft}{\kern0pt}\isaconst{valids}\ {\isacharparenleft}{\kern0pt}\isaconst{rev}\ \isafree{{\isasymgamma}}\ {\isacharat}{\kern0pt}\ {\isacharbrackleft}{\kern0pt}\isaconst{Nt}\ \isafree{A}{\isacharbrackright}{\kern0pt}{\isacharparenright}{\kern0pt}\ {\isacharhash}{\kern0pt}\ \isactrlbsub {\isacharbar}{\kern0pt}\isafree{{\isasymbeta}}{\isacharbar}{\kern0pt}\isactrlesub {\isacharbar}{\kern0pt}{\isacharparenleft}{\kern0pt}\isaconst{valids}\ {\isacharparenleft}{\kern0pt}\isaconst{rev}\ \isafree{{\isasymgamma}}\ {\isacharat}{\kern0pt}\ \isafree{{\isasymbeta}}{\isacharparenright}{\kern0pt}\ {\isacharhash}{\kern0pt}\ \isafree{qs}{\isacharparenright}{\kern0pt}{\isacharcomma}{\kern0pt}\ \isafree{v}{\isacharparenright}{\kern0pt}}.
\end{multline}
This follows from the fact that \isa{{\isacharbrackleft}{\kern0pt}\isafree{A}\ {\isasymrightarrow}\ \isafree{{\isasymbeta}}\ {\isasymcdot}\ {\isacharbrackright}{\kern0pt}\ {\isasymin}\ \isaconst{valids}\ {\isacharparenleft}{\kern0pt}\isaconst{rev}\ \isafree{{\isasymgamma}}\ {\isacharat}{\kern0pt}\ \isafree{{\isasymbeta}}{\isacharparenright}{\kern0pt}}, and with 
Lemmas~\ref{P0_nth_is_valids_of_nth_stack_word} and \ref{P0_init_stack_word_length_inv} for 
\eqref{MG_sw}, which we use to prove that 
\isa{{\isacharparenleft}{\kern0pt}\isaconst{valids}\ {\isacharparenleft}{\kern0pt}\isaconst{rev}\ \isafree{{\isasymgamma}}\ {\isacharat}{\kern0pt}\ \isafree{{\isasymbeta}}{\isacharparenright}{\kern0pt}\ {\isacharhash}{\kern0pt}\ \isafree{qs}{\isacharparenright}{\kern0pt}\ensuremath{_{[\mathit{{\isacharbar}{\kern0pt}\isafree{{\isasymbeta}}{\isacharbar}{\kern0pt}}]}}\ {\isacharequal}{\kern0pt}\ \isaconst{valids}\ {\isacharparenleft}{\kern0pt}\isaconst{rev}\ \isafree{{\isasymgamma}}{\isacharparenright}{\kern0pt}}. The property that 
\isa{\isafree{A}\ {\isasymnoteq}\ \isaconst{S{\isacharprime}{\kern0pt}}} is guaranteed by the fact that \isa{A} is in a production of \isa{G}, since it originates
from the computation of \isa{\isaconst{M\isactrlsub G}}.

Next, we need to show that this reduction pops a topmost string off the stack, and pushes 
the successor of the topmost state after popping. We know this popped string is exactly 
\isa{{\isacharparenleft}{\kern0pt}\isaconst{valids}\ {\isacharparenleft}{\kern0pt}\isaconst{rev}\ \isafree{{\isasymgamma}}\ {\isacharat}{\kern0pt}\ \isafree{{\isasymbeta}}{\isacharparenright}{\kern0pt}\ {\isacharhash}{\kern0pt}\ \isafree{qs}{\isacharparenright}{\kern0pt}{\isacharbar}{\kern0pt}\isactrlbsub {\isacharbar}{\kern0pt}\isafree{{\isasymbeta}}{\isacharbar}{\kern0pt}\isactrlesub }. With 
Lemma~\ref{P0_init_stack_word_length_inv} again for \eqref{MG_sw}, and 
Lemma~\ref{nxt_dfa_LR0_shift_is_valids_app} we can rewrite the reduction step \eqref{redstep} as
\begin{multline*} 
\isa{{\isacharparenleft}{\kern0pt}{\isacharparenleft}{\kern0pt}\isaconst{valids}\ {\isacharparenleft}{\kern0pt}\isaconst{rev}\ \isafree{{\isasymgamma}}\ {\isacharat}{\kern0pt}\ \isafree{{\isasymbeta}}{\isacharparenright}{\kern0pt}\ {\isacharhash}{\kern0pt}\ \isafree{qs}{\isacharparenright}{\kern0pt}{\isacharbar}{\kern0pt}\isactrlbsub {\isacharbar}{\kern0pt}\isafree{{\isasymbeta}}{\isacharbar}{\kern0pt}\isactrlesub \ {\isacharat}{\kern0pt}\ \isaconst{valids}\ {\isacharparenleft}{\kern0pt}\isaconst{rev}\ \isafree{{\isasymgamma}}{\isacharparenright}{\kern0pt}\ {\isacharhash}{\kern0pt}\ \isactrlbsub {\isacharparenleft}{\kern0pt}{\isacharbar}{\kern0pt}\isafree{{\isasymbeta}}{\isacharbar}{\kern0pt}\ {\isacharplus}{\kern0pt}\ {\isadigit{1}}{\isacharparenright}{\kern0pt}\isactrlesub {\isacharbar}{\kern0pt}{\isacharparenleft}{\kern0pt}\isaconst{valids}\ {\isacharparenleft}{\kern0pt}\isaconst{rev}\ \isafree{{\isasymgamma}}\ {\isacharat}{\kern0pt}\ \isafree{{\isasymbeta}}{\isacharparenright}{\kern0pt}\ {\isacharhash}{\kern0pt}\ \isafree{qs}{\isacharparenright}{\kern0pt}{\isacharcomma}{\kern0pt}\ \isafree{v}{\isacharparenright}{\kern0pt}}\\
  \isa{{\isasymturnstile}P\ {\isacharparenleft}{\kern0pt}} \isa{\isaconst{nxt}\ \ensuremath{LR_0(G)}\ {\isacharparenleft}{\kern0pt}\isaconst{valids}\ {\isacharparenleft}{\kern0pt}\isaconst{rev}\ \isafree{{\isasymgamma}}{\isacharparenright}{\kern0pt}{\isacharparenright}{\kern0pt}\ {\isacharparenleft}{\kern0pt}\isaconst{Nt}\ \isafree{A}{\isacharparenright}{\kern0pt}}\ \isa{{\isacharhash}{\kern0pt}}\ \isa{\isaconst{valids}\ {\isacharparenleft}{\kern0pt}\isaconst{rev}\ \isafree{{\isasymgamma}}{\isacharparenright}{\kern0pt}}\\
        \isa{{\isacharhash}{\kern0pt}}\ \isa{\isactrlbsub {\isacharparenleft}{\kern0pt}{\isacharbar}{\kern0pt}\isafree{{\isasymbeta}}{\isacharbar}{\kern0pt}\ {\isacharplus}{\kern0pt}\ {\isadigit{1}}{\isacharparenright}{\kern0pt}\isactrlesub {\isacharbar}{\kern0pt}{\isacharparenleft}{\kern0pt}\isaconst{valids}\ {\isacharparenleft}{\kern0pt}\isaconst{rev}\ \isafree{{\isasymgamma}}\ {\isacharat}{\kern0pt}\ \isafree{{\isasymbeta}}{\isacharparenright}{\kern0pt}\ {\isacharhash}{\kern0pt}\ \isafree{qs}{\isacharparenright}{\kern0pt}} \isa{{\isacharcomma}{\kern0pt}\ v{\isacharparenright}{\kern0pt}}.
\end{multline*}
Finally, we can also use the same instance of the lemma to show that 
\isa{\isafree{{\isasymgamma}}\ {\isacharequal}{\kern0pt}\ \isactrlbsub {\isacharbar}{\kern0pt}\isafree{{\isasymbeta}}{\isacharbar}{\kern0pt}\isactrlesub {\isacharbar}{\kern0pt}{\isacharparenleft}{\kern0pt}\isaconst{rev}\ \isafree{{\isasymbeta}}\ {\isacharat}{\kern0pt}\ \isafree{{\isasymgamma}}{\isacharparenright}{\kern0pt}}, completing the 
proof for the computation with stack word \isa{\isaconst{Nt}\ \isafree{A}\ {\isacharhash}{\kern0pt}\ \isafree{{\isasymgamma}}\ {\isacharequal}{\kern0pt}\ \isafree{{\isasymalpha}}}. Since the topmost state of the 
configuration on the RHS of the stack word computation is equal to \isa{\isaconst{rev}\ \isafree{{\isasymalpha}}}, the proof for the
reduce case is complete.
\end{proof}
\end{lemma}

We can now prove our parser complete.

\begin{theorem}[Completeness of the \isa{LR{\isacharparenleft}{\kern0pt}{\isadigit{0}}{\isacharparenright}{\kern0pt}} parser]\label{P0_complete}
Every word in \isa{L{\isacharparenleft}{\kern0pt}\isafree{G}{\isacharparenright}{\kern0pt}} is accepted by \isa{\ensuremath{P_0(G)}}.
\begin{proof}
We begin by fixing some \isa{w} in \isa{L{\isacharparenleft}{\kern0pt}\isafree{G}{\isacharparenright}{\kern0pt}}. By Lemma~\ref{srpda_complete}, this implies
that \isa{\isafree{w}\ {\isasymin}\ \ensuremath{L(M_G)}}, i.e., \isa{{\isacharparenleft}{\kern0pt}{\isacharbrackleft}{\kern0pt}\isaconst{Init}{\isacharbrackright}{\kern0pt}{\isacharcomma}{\kern0pt}\ \isafree{w}{\isacharparenright}{\kern0pt}\ {\isasymturnstile}M{\isacharasterisk}{\kern0pt}\ {\isacharparenleft}{\kern0pt}{\isacharbrackleft}{\kern0pt}\isaconst{Final}{\isacharbrackright}{\kern0pt}{\isacharcomma}{\kern0pt}\ {\isacharbrackleft}{\kern0pt}{\isacharbrackright}{\kern0pt}{\isacharparenright}{\kern0pt}}. Since \isa{\isaconst{Init}}
and \isa{\isaconst{Final}} are distinct by definition, the computation has nonzero length, meaning 
\isa{{\isacharparenleft}{\kern0pt}{\isacharbrackleft}{\kern0pt}\isaconst{Init}{\isacharbrackright}{\kern0pt}{\isacharcomma}{\kern0pt}\ \isafree{w}{\isacharparenright}{\kern0pt}\ {\isasymturnstile}M{\isacharasterisk}{\kern0pt}\ {\isacharparenleft}{\kern0pt}{\isacharbrackleft}{\kern0pt}\isaconst{Sym}\ {\isacharparenleft}{\kern0pt}\isaconst{Nt}\ \isaconst{S}{\isacharparenright}{\kern0pt}{\isacharcomma}{\kern0pt}\ \isaconst{Init}{\isacharbrackright}{\kern0pt}{\isacharcomma}{\kern0pt}\ {\isacharbrackleft}{\kern0pt}{\isacharbrackright}{\kern0pt}{\isacharparenright}{\kern0pt}} holds. Furthermore, it is trivial to prove that
the set \isa{\isaconst{valids}\ {\isacharbrackleft}{\kern0pt}\isaconst{Nt}\ \isaconst{S}{\isacharbrackright}{\kern0pt}} is nonempty by showing that the item \isa{{\isacharbrackleft}{\kern0pt}\isaconst{S{\isacharprime}{\kern0pt}}\ {\isasymrightarrow}\ {\isacharbrackleft}{\kern0pt}\isaconst{Nt}\ \isaconst{S}{\isacharbrackright}{\kern0pt}\ {\isasymcdot}\ {\isacharbrackright}{\kern0pt}} is
in this set. We can now use Lemma~\ref{MG_steps_imp_stack_word} to obtain some \isa{qs} where
\[ \isa{{\isacharbrackleft}{\kern0pt}\isaconst{Nt}\ \isaconst{S}{\isacharbrackright}{\kern0pt}\ {\isasymTurnstile}\ {\isacharparenleft}{\kern0pt}{\isacharbrackleft}{\kern0pt}\isaconst{init}\ \ensuremath{P_0(G)}{\isacharbrackright}{\kern0pt}{\isacharcomma}{\kern0pt}\ \isafree{w}{\isacharparenright}{\kern0pt}\ {\isasymturnstile}P{\isacharasterisk}{\kern0pt}\ {\isacharparenleft}{\kern0pt}\isaconst{valids}\ {\isacharbrackleft}{\kern0pt}\isaconst{Nt}\ \isaconst{S}{\isacharbrackright}{\kern0pt}\ {\isacharhash}{\kern0pt}\ \isafree{qs}{\isacharcomma}{\kern0pt}\ {\isacharbrackleft}{\kern0pt}{\isacharbrackright}{\kern0pt}{\isacharparenright}{\kern0pt}}. \]
By Lemmas~\ref{P0_nth_is_valids_of_nth_stack_word} and \ref{P0_init_stack_word_length_inv}, we can 
then show that \isa{\isafree{qs}\ {\isacharequal}{\kern0pt}\ {\isacharbrackleft}{\kern0pt}\isaconst{init}\ \ensuremath{P_0(G)}{\isacharbrackright}{\kern0pt}}. With Lemma~\ref{stack_word_imp_P0_steps}, this implies
that \isa{{\isacharparenleft}{\kern0pt}{\isacharbrackleft}{\kern0pt}\isaconst{init}\ \ensuremath{P_0(G)}{\isacharbrackright}{\kern0pt}{\isacharcomma}{\kern0pt}\ \isafree{w}{\isacharparenright}{\kern0pt}\ {\isasymturnstile}P{\isacharasterisk}{\kern0pt}\ {\isacharparenleft}{\kern0pt}{\isacharbrackleft}{\kern0pt}\isaconst{valids}\ {\isacharbrackleft}{\kern0pt}\isaconst{Nt}\ \isaconst{S}{\isacharbrackright}{\kern0pt}{\isacharcomma}{\kern0pt}\ \isaconst{init}\ \ensuremath{P_0(G)}{\isacharbrackright}{\kern0pt}{\isacharcomma}{\kern0pt}\ {\isacharbrackleft}{\kern0pt}{\isacharbrackright}{\kern0pt}{\isacharparenright}{\kern0pt}} holds. By the fact that
\isa{{\isacharbrackleft}{\kern0pt}\isaconst{S{\isacharprime}{\kern0pt}}\ {\isasymrightarrow}\ {\isacharbrackleft}{\kern0pt}\isaconst{Nt}\ \isaconst{S}{\isacharbrackright}{\kern0pt}\ {\isasymcdot}\ {\isacharbrackright}{\kern0pt}} is in \isa{\isaconst{valids}\ {\isacharbrackleft}{\kern0pt}\isaconst{Nt}\ \isaconst{S}{\isacharbrackright}{\kern0pt}}, we know that the RHS configuration of this
computation can then perform a finishing transition, meaning that 
\[ \isa{{\isacharparenleft}{\kern0pt}{\isacharbrackleft}{\kern0pt}\isaconst{init}\ \ensuremath{P_0(G)}{\isacharbrackright}{\kern0pt}{\isacharcomma}{\kern0pt}\ \isafree{w}{\isacharparenright}{\kern0pt}\ {\isasymturnstile}P{\isacharasterisk}{\kern0pt}\ {\isacharparenleft}{\kern0pt}{\isacharbrackleft}{\kern0pt}{\isacharbraceleft}{\kern0pt}{\isacharbrackleft}{\kern0pt}\isaconst{S{\isacharprime}{\kern0pt}}\ {\isasymrightarrow}\ {\isasymcdot}\ {\isacharbrackright}{\kern0pt}{\isacharbraceright}{\kern0pt}{\isacharbrackright}{\kern0pt}{\isacharcomma}{\kern0pt}\ {\isacharbrackleft}{\kern0pt}{\isacharbrackright}{\kern0pt}{\isacharparenright}{\kern0pt}}. \]
This is equivalent to \isa{\isafree{w}\ {\isasymin}\ \ensuremath{L(P_0(G))}}, completing the proof.
\end{proof}
\end{theorem}

With the completeness of the parser, we can prove our final goal.

\begin{theorem}[Correctness of the \isa{LR{\isacharparenleft}{\kern0pt}{\isadigit{0}}{\isacharparenright}{\kern0pt}} parser]
The language accepted by the \isa{LR{\isacharparenleft}{\kern0pt}{\isadigit{0}}{\isacharparenright}{\kern0pt}} parser \isa{\ensuremath{P_0(G)}} is exactly the language of \isa{G}.
\begin{proof}
This is a consequence of Theorems~\ref{P0_sound} and \ref{P0_complete}.
\end{proof}
\end{theorem}

It is worth noting that the automaton's correctness is completely unrelated to its grammar's
satisfaction of the \isa{LR{\isacharparenleft}{\kern0pt}{\isadigit{0}}{\isacharparenright}{\kern0pt}} condition, meaning that \isa{\ensuremath{P_0(G)}} is a recognizer for any
CFG \isa{G}. With this result and Theorem~\ref{is_LR0_iff_no_LR0_inadequates}, we have proven
that \isa{\ensuremath{P_0(G)}} is a deterministic parser for \isa{LR{\isacharparenleft}{\kern0pt}{\isadigit{0}}{\isacharparenright}{\kern0pt}} grammars. 

This concludes our formalization of the \isa{LR{\isacharparenleft}{\kern0pt}{\isadigit{0}}{\isacharparenright}{\kern0pt}} parser and general \isa{LR{\isacharparenleft}{\kern0pt}{\isadigit{0}}{\isacharparenright}{\kern0pt}} parsing 
theory. We have successfully formalized the parsing algorithm presented by Wilhelm et al., and
verified the correctness of their theorems while additionally proving the parser itself correct.%
\end{isamarkuptext}\isamarkuptrue%
%
\isadelimdocument
%
\endisadelimdocument
%
\isatagdocument
%
\isamarkupchapter{Conclusion%
}
\isamarkuptrue%
%
\isamarkupsection{Results%
}
\isamarkuptrue%
%
\endisatagdocument
{\isafolddocument}%
%
\isadelimdocument
%
\endisadelimdocument
%
\begin{isamarkuptext}%
In this thesis, we formally verified the \isa{LR{\isacharparenleft}{\kern0pt}{\isadigit{0}}{\isacharparenright}{\kern0pt}} parsing theory presented by Wilhelm et al., 
in particular two main theorems: Theorem~\ref{char_derivers_ipda_iffs}, which shows how exactly the 
characteristic finite automaton, rightmost derivations, and the item pushdown automaton relate to
each other, and Theorem~\ref{is_LR0_iff_no_LR0_inadequates}, which states the equivalence between an
extended grammar satisfying the \isa{LR{\isacharparenleft}{\kern0pt}{\isadigit{0}}{\isacharparenright}{\kern0pt}} condition, and the canonical \isa{LR{\isacharparenleft}{\kern0pt}{\isadigit{0}}{\isacharparenright}{\kern0pt}} automaton for this grammar
having no \isa{LR{\isacharparenleft}{\kern0pt}{\isadigit{0}}{\isacharparenright}{\kern0pt}} inadequate states. Additionally, we provided a formal proof of correctness of the
\isa{LR{\isacharparenleft}{\kern0pt}{\isadigit{0}}{\isacharparenright}{\kern0pt}} parser, meaning that it accepts exactly the language of its grammar.

We began by formalizing the concept of a reduced grammar using the formalization of CFG theory by
Nipkow et al.~\cite{Nipkow-AFP, Nipkow-Paper}, and we formally introduced the concept of
extending such grammars by a new starting symbol, proving that the resulting grammars preserve the
language, and that the extension of a reduced grammar is itself reduced. We also formally defined generalized pushdown automata, which
operate similarly to ordinary PDAs, except for the fact that they are able to read topmost stack
strings of variable length instead of only the topmost stack symbol.

We then defined the item pushdown automaton \isa{I\isactrlsub G} and proved its correctness w.r.t.\ its grammar
\isa{G}. Our formalization of this property required us to strengthen a lemma that Wilhelm et al.\ used,
particularly in order to later prove one of the two main theorems we formalized. Our correctness
proof followed the original argument closely, with the main difference being that their weaker lemma 
was simply a special case of our stronger one, Lemma~\ref{derives_imp_completes}.

A major addition that was absent in the book was the introduction of rightmost chains. This allowed
us to formalize an intuitive manner of decomposing an arbitrary rightmost derivation, and we proved
the existence of such chains for any rightmost derivation. We also showed that an IPDA stack that
leads to an accepting state also corresponds to one of these rightmost chains. These properties were
essential to prove the first of the two theorems in the book. We proved this theorem, which
corresponds to our Theorem~\ref{char_derivers_ipda_iffs}, after defining the characteristic finite 
automaton \isa{\ensuremath{\mathrm{char}(G)}}. The proof for this theorem was divided into three separate implications
in a circular proof, the first of which contained multiple errors, most importantly in the
application of the induction hypothesis, and had to be corrected. The remaining two implications,
however, followed the original argumentation of Wilhelm et al., which was possible with our
aforementioned formalization of rightmost chains. With these results and the introduction of reliable
prefixes and valid items, we were able to establish an equivalence between
\isa{\ensuremath{\mathrm{char}(G)}} computations and reliable prefixes. We then defined the canonical \isa{LR{\isacharparenleft}{\kern0pt}{\isadigit{0}}{\isacharparenright}{\kern0pt}} automaton
\isa{\ensuremath{LR_0(G)}}, and proved two essential properties: every state of this automaton is the set of valid 
items for some string \isa{{\isasymgamma}}, \isa{\isaconst{valids}\ \isafree{{\isasymgamma}}}, and for its transition function \isa{nxt} holds 
\isa{\isaconst{nxt}\ \ensuremath{LR_0(G)}\ {\isacharparenleft}{\kern0pt}\isaconst{valids}\ \isafree{{\isasymgamma}}{\isacharparenright}{\kern0pt}\ \isafree{X}\ {\isacharequal}{\kern0pt}\ \isaconst{valids}\ {\isacharparenleft}{\kern0pt}\isafree{{\isasymgamma}}\ {\isacharat}{\kern0pt}\ {\isacharbrackleft}{\kern0pt}\isafree{X}{\isacharbrackright}{\kern0pt}{\isacharparenright}{\kern0pt}} for any string \isa{{\isasymgamma}} and symbol \isa{X}, provided
that \isa{valids\ {\isasymgamma}} is a state of \isa{\ensuremath{LR_0(G)}}.

Afterward, we defined the \isa{LR{\isacharparenleft}{\kern0pt}{\isadigit{0}}{\isacharparenright}{\kern0pt}} parser \isa{\ensuremath{P_0(G)}} with minor differences to the definition by
Wilhelm et al., in particular, restricting reducing transitions to items \isa{{\isacharbrackleft}{\kern0pt}\isafree{X}\ {\isasymrightarrow}\ \isafree{{\isasymalpha}}\ {\isasymcdot}\ \isafree{{\isasymbeta}}{\isacharbrackright}{\kern0pt}} with 
\isa{X\ {\isasymnoteq}\ S{\isacharprime}{\kern0pt}} to avoid reduce-reduce conflicts in finishing transitions. We also showed the
transition relations of the parser to be finite, which leaves the door open for the development of
an executable \isa{LR{\isacharparenleft}{\kern0pt}{\isadigit{0}}{\isacharparenright}{\kern0pt}} parser. Afterward, we defined the notion of \isa{LR{\isacharparenleft}{\kern0pt}{\isadigit{0}}{\isacharparenright}{\kern0pt}} inadequacy,
and presented a proof for Lemma~\ref{LR0_adequate_cases} which was absent from the book. With this
lemma and the subsequent definition of the \isa{LR{\isacharparenleft}{\kern0pt}k{\isacharparenright}{\kern0pt}} condition for the extended grammar \isa{G{\isacharprime}{\kern0pt}}, we showed
that \isa{G{\isacharprime}{\kern0pt}} is an \isa{LR{\isacharparenleft}{\kern0pt}{\isadigit{0}}{\isacharparenright}{\kern0pt}} grammar if and only if \isa{\ensuremath{LR_0(G)}} has no \isa{LR{\isacharparenleft}{\kern0pt}{\isadigit{0}}{\isacharparenright}{\kern0pt}}-inadequate states. Our
proof had major differences with that of Wilhelm et al., since we found certain errors in their
argumentation, and corrected them. 

Furthermore, we showed that the extension of an \isa{LR{\isacharparenleft}{\kern0pt}k{\isacharparenright}{\kern0pt}} grammar is not \isa{LR{\isacharparenleft}{\kern0pt}k{\isacharparenright}{\kern0pt}} in general. In
particular, the extension is \isa{LR{\isacharparenleft}{\kern0pt}{\isadigit{0}}{\isacharparenright}{\kern0pt}} if and only if the original grammar \isa{G} is \isa{LR{\isacharparenleft}{\kern0pt}{\isadigit{0}}{\isacharparenright}{\kern0pt}} and its
start symbol \isa{S} lacks a left-recursion of the form \isa{Prods\ G\ {\isasymturnstile}\ {\isacharbrackleft}{\kern0pt}Nt\ S{\isacharbrackright}{\kern0pt}\ {\isasymRightarrow}r{\isacharparenleft}{\kern0pt}Suc\ n{\isacharparenright}{\kern0pt}\ Nt\ S\ {\isacharhash}{\kern0pt}\ {\isasymalpha}}.
Similarly for \isa{k\ {\isachargreater}{\kern0pt}\ {\isadigit{0}}}, we found that the extension is \isa{LR{\isacharparenleft}{\kern0pt}k{\isacharparenright}{\kern0pt}} if and only if the original grammar 
is \isa{LR{\isacharparenleft}{\kern0pt}k{\isacharparenright}{\kern0pt}} and it has no cycles of the form \isa{Prods\ G\ {\isasymturnstile}\ {\isacharbrackleft}{\kern0pt}Nt\ S{\isacharbrackright}{\kern0pt}\ {\isasymRightarrow}r{\isacharparenleft}{\kern0pt}Suc\ n{\isacharparenright}{\kern0pt}\ mbox\ {\isacharbrackleft}{\kern0pt}Nt\ S{\isacharbrackright}{\kern0pt}}.
These results lead us to conclude that the act of arbitrarily extending a grammar as we did originally is
not problematic for \isa{k\ {\isachargreater}{\kern0pt}\ {\isadigit{0}}} in practice, since grammars with cycles that cause the extension to violate the
\isa{LR{\isacharparenleft}{\kern0pt}k{\isacharparenright}{\kern0pt}} condition should be avoided per Wilhelm et al. We can also conclude, however,
that indiscriminately extending an \isa{LR{\isacharparenleft}{\kern0pt}{\isadigit{0}}{\isacharparenright}{\kern0pt}} grammar can be troublesome, and since the \isa{LR{\isacharparenleft}{\kern0pt}{\isadigit{0}}{\isacharparenright}{\kern0pt}} parser
is only defined for extended grammars, it is possible that \isa{LR{\isacharparenleft}{\kern0pt}{\isadigit{0}}{\isacharparenright}{\kern0pt}} grammars exist which are
unsuitable for the parser after the extension.

Lastly, we formalized the concept of stack words, inspired by an intuitive explanation by Wilhelm et
al. of how the parser operates. We furthermore formalized the shift-reduce pushdown automaton
based on the definition by Petter~\cite{Petter}, and proved its completeness w.r.t.\ its underlying
grammar. With these two constructions, we finally proved the correctness of the canonical \isa{LR{\isacharparenleft}{\kern0pt}{\isadigit{0}}{\isacharparenright}{\kern0pt}}
parser.%
\end{isamarkuptext}\isamarkuptrue%
%
\isadelimdocument
%
\endisadelimdocument
%
\isatagdocument
%
\isamarkupsection{Future Work%
}
\isamarkuptrue%
%
\endisatagdocument
{\isafolddocument}%
%
\isadelimdocument
%
\endisadelimdocument
%
\begin{isamarkuptext}%
In terms of how our formalization could be extended, Wilhelm et al.\ briefly define determinism
in the PDA chapter of the book~\cite[p. 58]{Wilhelm}, but do not address the equivalence between
\isa{LR{\isacharparenleft}{\kern0pt}{\isadigit{0}}{\isacharparenright}{\kern0pt}}-adequate states of \isa{\ensuremath{LR_0(G)}} and the determinism of \isa{\ensuremath{P_0(G)}} anywhere
in the \isa{LR{\isacharparenleft}{\kern0pt}{\isadigit{0}}{\isacharparenright}{\kern0pt}} section of the book. Therefore, it may be interesting to verify this property, which
is of particular significance for Theorem~\ref{is_LR0_iff_no_LR0_inadequates}.

It would also be of interest to develop an executable \isa{LR{\isacharparenleft}{\kern0pt}{\isadigit{0}}{\isacharparenright}{\kern0pt}} parser.
Although we verified the parser's properties proven by Wilhelm et al., as well as its correctness, 
our definition is not executable. This gap could be closed in the future, as discussed in
Section~\ref{sec:finite}, which would allow the parser to be applied in practice for concrete
grammars. On a similar note, our variant of the parser acts only as a recognizer, meaning it
determines whether the input is in the grammar. However, parsers usually yield more information than
this, namely the derivation of the parsed word if it indeed is in the grammar. Wilhelm et al.
mention this in the IPDA section of the 
book~\cite[p. 63, \textit{Pushdown Automata with Output}]{Wilhelm}, but do not apply this concept to
the definition of the \isa{LR{\isacharparenleft}{\kern0pt}{\isadigit{0}}{\isacharparenright}{\kern0pt}} parser we formalized.

There is also a question out of the scope of parsing theory itself, but of relevance in automata
theory, particularly as an addition for the AFP. Standard pushdown automata have been formalized by
Taskin and Nipkow~\cite{Taskin} and submitted to the Archive, and one could fully formalize GPDAs
and prove their equivalence with the existing PDAs. As we mentioned in Section~\ref{sec:gpdas}, 
however, our definition must be slightly modified: one must explicitly assume that the transition
relations are finite. This is because for a PDA to simulate a GPDA, it must pop a topmost stack
string of length \isa{n}, which the GPDA pops in one transition step, in \isa{n} transition steps. In order
to perform this simulation, one must naturally define a nw state for each step, which is only
possible if the GPDA's transition relation is finite. Furthermore, PDAs as defined by Taskin and 
Nipkow treat the type of states as the set of states for the PDA, and since the set of states of a 
PDA must be finite, the type for states is finite, and the same applies for stack symbols, and even 
the input alphabet. Since our GPDA definition is not as restrictive in its types, this raises
numerous questions as to how one could show this equivalence without sacrificing generality.
Finite types also raise a separate question: how can the states of the GPDA be extended? Our answer
to this, which we used to define extended grammars, is to use Popescu's and Bauereiss's fresh
identifiers~\cite{Popescu}, but this solution is antithetical to finite types, since types of class
\isa{fresh} are infinite by definition.

Lastly, a broader topic that is also of great interest is the generalization of our formalization
for \isa{LR{\isacharparenleft}{\kern0pt}k{\isacharparenright}{\kern0pt}} parsing of arbitrary \isa{k}. This generalization raises questions as fundamental as how one
should model \isa{LR{\isacharparenleft}{\kern0pt}k{\isacharparenright}{\kern0pt}} items in Isabelle: since these items contain a \isa{k}-length lookahead, one must
be able to restrict the definition of \isa{LR{\isacharparenleft}{\kern0pt}k{\isacharparenright}{\kern0pt}} items to items with lookaheads of specific lengths. 
\isa{LR{\isacharparenleft}{\kern0pt}k{\isacharparenright}{\kern0pt}} items would also come with additional challenges: in order to guarantee that a grammar has
finitely many such items, it must have infinitely many symbols. Intuitively, this is not problematic,
since a grammar with finite productions trivially also has finitely many terminals and nonterminals.
However, as we could see in the case of the shift-reduce PDA, the convention of defining the set of
symbols through a type means one would need to restrict \isa{LR{\isacharparenleft}{\kern0pt}k{\isacharparenright}{\kern0pt}} items to work only with finite
types, which may be too restrictive. If one wishes to work with general types without imposing
such restrictions, it is perhaps necessary to define the grammar's symbols through sets instead of
types.

A general \isa{LR{\isacharparenleft}{\kern0pt}k{\isacharparenright}{\kern0pt}} formalization is not only of interest as the natural development of \isa{LR} parsing
theory in Isabelle. This generalization can also serve as the continuation of the formal verification
of the theory presented by Wilhelm et al., and therefore contribute in a potential future goal of
verifying larger portions of the book.%
\end{isamarkuptext}\isamarkuptrue%
%
\isadelimtheory
%
\endisadelimtheory
%
\isatagtheory
%
\endisatagtheory
{\isafoldtheory}%
%
\isadelimtheory
%
\endisadelimtheory
%
\end{isabellebody}%
\endinput
%:%file=~/Work/AIST/AIST/Escallon/LR_Parsing/thesis/Paper.thy%:%
%:%24=44%:%
%:%36=46%:%
%:%37=47%:%
%:%38=48%:%
%:%39=49%:%
%:%40=50%:%
%:%41=51%:%
%:%42=52%:%
%:%43=53%:%
%:%44=54%:%
%:%45=55%:%
%:%46=56%:%
%:%47=57%:%
%:%48=58%:%
%:%49=59%:%
%:%50=60%:%
%:%51=61%:%
%:%52=62%:%
%:%53=63%:%
%:%54=64%:%
%:%55=65%:%
%:%56=66%:%
%:%57=67%:%
%:%58=68%:%
%:%59=69%:%
%:%60=70%:%
%:%61=71%:%
%:%62=72%:%
%:%63=73%:%
%:%64=74%:%
%:%65=75%:%
%:%66=76%:%
%:%67=77%:%
%:%68=78%:%
%:%69=79%:%
%:%70=80%:%
%:%71=81%:%
%:%72=82%:%
%:%73=83%:%
%:%74=84%:%
%:%75=85%:%
%:%76=86%:%
%:%77=87%:%
%:%78=88%:%
%:%79=89%:%
%:%80=90%:%
%:%81=91%:%
%:%82=92%:%
%:%83=93%:%
%:%84=94%:%
%:%85=95%:%
%:%86=96%:%
%:%87=97%:%
%:%88=98%:%
%:%89=99%:%
%:%90=100%:%
%:%91=101%:%
%:%92=102%:%
%:%93=103%:%
%:%94=104%:%
%:%95=105%:%
%:%96=106%:%
%:%97=107%:%
%:%98=108%:%
%:%99=109%:%
%:%100=110%:%
%:%101=111%:%
%:%102=112%:%
%:%103=113%:%
%:%104=114%:%
%:%105=115%:%
%:%106=116%:%
%:%107=117%:%
%:%108=118%:%
%:%109=119%:%
%:%110=120%:%
%:%111=121%:%
%:%112=122%:%
%:%113=123%:%
%:%122=125%:%
%:%134=127%:%
%:%135=128%:%
%:%136=129%:%
%:%137=130%:%
%:%138=131%:%
%:%139=132%:%
%:%140=133%:%
%:%141=134%:%
%:%142=135%:%
%:%143=136%:%
%:%144=137%:%
%:%145=138%:%
%:%146=139%:%
%:%147=140%:%
%:%148=141%:%
%:%149=142%:%
%:%150=143%:%
%:%151=144%:%
%:%152=145%:%
%:%153=146%:%
%:%162=148%:%
%:%166=149%:%
%:%178=151%:%
%:%179=152%:%
%:%180=153%:%
%:%181=154%:%
%:%182=155%:%
%:%183=156%:%
%:%184=157%:%
%:%185=158%:%
%:%186=159%:%
%:%187=160%:%
%:%188=161%:%
%:%189=162%:%
%:%190=163%:%
%:%191=164%:%
%:%192=165%:%
%:%193=166%:%
%:%194=167%:%
%:%195=168%:%
%:%196=169%:%
%:%197=170%:%
%:%198=171%:%
%:%199=172%:%
%:%200=173%:%
%:%201=174%:%
%:%202=175%:%
%:%203=176%:%
%:%204=177%:%
%:%205=178%:%
%:%206=179%:%
%:%207=180%:%
%:%208=181%:%
%:%209=182%:%
%:%210=183%:%
%:%211=184%:%
%:%212=185%:%
%:%213=186%:%
%:%214=187%:%
%:%215=188%:%
%:%216=189%:%
%:%217=190%:%
%:%218=191%:%
%:%219=192%:%
%:%220=193%:%
%:%221=194%:%
%:%222=195%:%
%:%223=196%:%
%:%224=197%:%
%:%225=198%:%
%:%226=199%:%
%:%227=200%:%
%:%228=201%:%
%:%229=202%:%
%:%230=203%:%
%:%231=204%:%
%:%232=205%:%
%:%233=206%:%
%:%234=207%:%
%:%235=208%:%
%:%236=209%:%
%:%237=210%:%
%:%246=213%:%
%:%258=215%:%
%:%259=216%:%
%:%260=217%:%
%:%261=218%:%
%:%262=219%:%
%:%263=220%:%
%:%264=221%:%
%:%265=222%:%
%:%266=223%:%
%:%267=224%:%
%:%268=225%:%
%:%269=226%:%
%:%270=227%:%
%:%271=228%:%
%:%272=229%:%
%:%273=230%:%
%:%274=231%:%
%:%275=232%:%
%:%276=233%:%
%:%277=234%:%
%:%278=235%:%
%:%279=236%:%
%:%280=237%:%
%:%281=238%:%
%:%282=239%:%
%:%283=240%:%
%:%284=241%:%
%:%285=242%:%
%:%286=243%:%
%:%287=244%:%
%:%288=245%:%
%:%289=246%:%
%:%298=248%:%
%:%310=250%:%
%:%311=251%:%
%:%312=252%:%
%:%313=253%:%
%:%314=254%:%
%:%315=255%:%
%:%316=256%:%
%:%317=257%:%
%:%318=258%:%
%:%319=259%:%
%:%320=260%:%
%:%321=261%:%
%:%322=262%:%
%:%323=263%:%
%:%324=264%:%
%:%325=265%:%
%:%326=266%:%
%:%327=267%:%
%:%328=268%:%
%:%329=269%:%
%:%330=270%:%
%:%331=271%:%
%:%332=272%:%
%:%333=273%:%
%:%334=274%:%
%:%335=275%:%
%:%336=276%:%
%:%337=277%:%
%:%338=278%:%
%:%339=279%:%
%:%340=280%:%
%:%341=281%:%
%:%342=282%:%
%:%343=283%:%
%:%344=284%:%
%:%345=285%:%
%:%346=286%:%
%:%347=287%:%
%:%348=288%:%
%:%349=289%:%
%:%350=290%:%
%:%351=291%:%
%:%352=292%:%
%:%353=293%:%
%:%354=294%:%
%:%355=295%:%
%:%356=296%:%
%:%357=297%:%
%:%358=298%:%
%:%359=299%:%
%:%360=300%:%
%:%361=301%:%
%:%362=302%:%
%:%363=303%:%
%:%372=305%:%
%:%384=307%:%
%:%385=308%:%
%:%386=309%:%
%:%387=310%:%
%:%388=311%:%
%:%389=312%:%
%:%390=313%:%
%:%399=315%:%
%:%403=317%:%
%:%407=319%:%
%:%432=334%:%
%:%433=335%:%
%:%434=336%:%
%:%435=337%:%
%:%436=338%:%
%:%437=339%:%
%:%438=340%:%
%:%439=341%:%
%:%440=342%:%
%:%441=343%:%
%:%442=344%:%
%:%443=345%:%
%:%444=346%:%
%:%445=347%:%
%:%446=348%:%
%:%447=349%:%
%:%448=350%:%
%:%449=351%:%
%:%450=352%:%
%:%451=353%:%
%:%452=354%:%
%:%453=355%:%
%:%454=356%:%
%:%455=357%:%
%:%456=358%:%
%:%457=359%:%
%:%458=360%:%
%:%459=361%:%
%:%460=362%:%
%:%461=363%:%
%:%462=364%:%
%:%463=365%:%
%:%464=366%:%
%:%465=367%:%
%:%466=368%:%
%:%467=369%:%
%:%468=370%:%
%:%469=371%:%
%:%470=372%:%
%:%471=373%:%
%:%472=374%:%
%:%473=375%:%
%:%474=376%:%
%:%475=377%:%
%:%476=378%:%
%:%477=379%:%
%:%478=380%:%
%:%479=381%:%
%:%480=382%:%
%:%481=383%:%
%:%482=384%:%
%:%483=385%:%
%:%484=386%:%
%:%485=387%:%
%:%486=388%:%
%:%487=389%:%
%:%488=390%:%
%:%489=391%:%
%:%490=392%:%
%:%491=393%:%
%:%492=394%:%
%:%493=395%:%
%:%494=396%:%
%:%495=397%:%
%:%496=398%:%
%:%497=399%:%
%:%498=400%:%
%:%507=402%:%
%:%519=404%:%
%:%520=405%:%
%:%521=406%:%
%:%522=407%:%
%:%523=408%:%
%:%524=409%:%
%:%525=410%:%
%:%526=411%:%
%:%527=412%:%
%:%528=413%:%
%:%529=414%:%
%:%530=415%:%
%:%531=416%:%
%:%532=417%:%
%:%533=418%:%
%:%534=419%:%
%:%535=420%:%
%:%536=421%:%
%:%537=422%:%
%:%538=423%:%
%:%539=424%:%
%:%540=425%:%
%:%541=426%:%
%:%542=427%:%
%:%543=428%:%
%:%544=429%:%
%:%545=430%:%
%:%546=431%:%
%:%547=432%:%
%:%548=433%:%
%:%549=434%:%
%:%550=435%:%
%:%551=436%:%
%:%552=437%:%
%:%553=438%:%
%:%554=439%:%
%:%555=440%:%
%:%556=441%:%
%:%557=442%:%
%:%558=443%:%
%:%559=444%:%
%:%560=445%:%
%:%561=446%:%
%:%562=447%:%
%:%563=448%:%
%:%564=449%:%
%:%565=450%:%
%:%566=451%:%
%:%567=452%:%
%:%568=453%:%
%:%569=454%:%
%:%570=455%:%
%:%571=456%:%
%:%572=457%:%
%:%573=458%:%
%:%574=459%:%
%:%575=460%:%
%:%576=461%:%
%:%577=462%:%
%:%578=463%:%
%:%579=464%:%
%:%580=465%:%
%:%581=466%:%
%:%582=467%:%
%:%583=468%:%
%:%584=469%:%
%:%585=470%:%
%:%586=471%:%
%:%587=472%:%
%:%588=473%:%
%:%589=474%:%
%:%590=475%:%
%:%591=476%:%
%:%592=477%:%
%:%593=478%:%
%:%594=479%:%
%:%595=480%:%
%:%596=481%:%
%:%597=482%:%
%:%598=483%:%
%:%599=484%:%
%:%600=485%:%
%:%601=486%:%
%:%602=487%:%
%:%603=488%:%
%:%604=489%:%
%:%605=490%:%
%:%606=491%:%
%:%607=492%:%
%:%608=493%:%
%:%609=494%:%
%:%610=495%:%
%:%611=496%:%
%:%612=497%:%
%:%613=498%:%
%:%614=499%:%
%:%615=500%:%
%:%616=501%:%
%:%617=502%:%
%:%626=504%:%
%:%638=506%:%
%:%639=507%:%
%:%640=508%:%
%:%641=509%:%
%:%642=510%:%
%:%643=511%:%
%:%644=512%:%
%:%645=513%:%
%:%646=514%:%
%:%647=515%:%
%:%648=516%:%
%:%649=517%:%
%:%650=518%:%
%:%651=519%:%
%:%652=520%:%
%:%653=521%:%
%:%654=522%:%
%:%655=523%:%
%:%656=524%:%
%:%657=525%:%
%:%658=526%:%
%:%659=527%:%
%:%660=528%:%
%:%661=529%:%
%:%662=530%:%
%:%663=531%:%
%:%664=532%:%
%:%665=533%:%
%:%666=534%:%
%:%667=535%:%
%:%668=536%:%
%:%669=537%:%
%:%670=538%:%
%:%671=539%:%
%:%672=540%:%
%:%673=541%:%
%:%674=542%:%
%:%675=543%:%
%:%676=544%:%
%:%677=545%:%
%:%678=546%:%
%:%679=547%:%
%:%680=548%:%
%:%681=549%:%
%:%682=550%:%
%:%683=551%:%
%:%684=552%:%
%:%685=553%:%
%:%686=554%:%
%:%687=555%:%
%:%688=556%:%
%:%689=557%:%
%:%690=558%:%
%:%691=559%:%
%:%692=560%:%
%:%693=561%:%
%:%694=562%:%
%:%703=564%:%
%:%715=575%:%
%:%716=576%:%
%:%717=577%:%
%:%718=578%:%
%:%719=579%:%
%:%720=580%:%
%:%721=581%:%
%:%722=582%:%
%:%723=583%:%
%:%724=584%:%
%:%725=585%:%
%:%726=586%:%
%:%727=587%:%
%:%728=588%:%
%:%729=589%:%
%:%730=590%:%
%:%731=591%:%
%:%732=592%:%
%:%733=593%:%
%:%734=594%:%
%:%735=595%:%
%:%736=596%:%
%:%737=597%:%
%:%738=598%:%
%:%739=599%:%
%:%740=600%:%
%:%741=601%:%
%:%742=602%:%
%:%743=603%:%
%:%744=604%:%
%:%745=605%:%
%:%746=606%:%
%:%747=607%:%
%:%748=608%:%
%:%749=609%:%
%:%750=610%:%
%:%751=611%:%
%:%752=612%:%
%:%753=613%:%
%:%754=614%:%
%:%755=615%:%
%:%756=616%:%
%:%757=617%:%
%:%758=618%:%
%:%759=619%:%
%:%760=620%:%
%:%761=621%:%
%:%762=622%:%
%:%763=623%:%
%:%764=624%:%
%:%765=625%:%
%:%766=626%:%
%:%767=627%:%
%:%768=628%:%
%:%769=629%:%
%:%770=630%:%
%:%771=631%:%
%:%772=632%:%
%:%773=633%:%
%:%782=636%:%
%:%812=656%:%
%:%824=658%:%
%:%825=659%:%
%:%826=660%:%
%:%827=661%:%
%:%828=662%:%
%:%829=663%:%
%:%830=664%:%
%:%831=665%:%
%:%832=666%:%
%:%833=667%:%
%:%834=668%:%
%:%835=669%:%
%:%836=670%:%
%:%837=671%:%
%:%838=672%:%
%:%839=673%:%
%:%840=674%:%
%:%841=675%:%
%:%842=676%:%
%:%843=677%:%
%:%844=678%:%
%:%845=679%:%
%:%846=680%:%
%:%847=681%:%
%:%848=682%:%
%:%849=683%:%
%:%850=684%:%
%:%851=685%:%
%:%852=686%:%
%:%853=687%:%
%:%854=688%:%
%:%855=689%:%
%:%856=690%:%
%:%857=691%:%
%:%858=692%:%
%:%859=693%:%
%:%860=694%:%
%:%861=695%:%
%:%862=696%:%
%:%863=697%:%
%:%864=698%:%
%:%865=699%:%
%:%866=700%:%
%:%867=701%:%
%:%868=702%:%
%:%869=703%:%
%:%870=704%:%
%:%871=705%:%
%:%872=706%:%
%:%873=707%:%
%:%874=708%:%
%:%875=709%:%
%:%876=710%:%
%:%877=711%:%
%:%878=712%:%
%:%879=713%:%
%:%880=714%:%
%:%881=715%:%
%:%882=716%:%
%:%883=717%:%
%:%884=718%:%
%:%885=719%:%
%:%886=720%:%
%:%887=721%:%
%:%888=722%:%
%:%889=723%:%
%:%890=724%:%
%:%891=725%:%
%:%892=726%:%
%:%893=727%:%
%:%894=728%:%
%:%895=729%:%
%:%896=730%:%
%:%897=731%:%
%:%898=732%:%
%:%899=733%:%
%:%900=734%:%
%:%909=736%:%
%:%921=738%:%
%:%922=739%:%
%:%923=740%:%
%:%924=741%:%
%:%925=742%:%
%:%926=743%:%
%:%927=744%:%
%:%928=745%:%
%:%929=746%:%
%:%930=747%:%
%:%931=748%:%
%:%932=749%:%
%:%933=750%:%
%:%934=751%:%
%:%935=752%:%
%:%936=753%:%
%:%937=754%:%
%:%938=755%:%
%:%939=756%:%
%:%940=757%:%
%:%941=758%:%
%:%942=759%:%
%:%943=760%:%
%:%944=761%:%
%:%945=762%:%
%:%946=763%:%
%:%947=764%:%
%:%948=765%:%
%:%949=766%:%
%:%950=767%:%
%:%951=768%:%
%:%952=769%:%
%:%953=770%:%
%:%954=771%:%
%:%955=772%:%
%:%956=773%:%
%:%957=774%:%
%:%958=775%:%
%:%959=776%:%
%:%960=777%:%
%:%961=778%:%
%:%962=779%:%
%:%963=780%:%
%:%964=781%:%
%:%965=782%:%
%:%966=783%:%
%:%967=784%:%
%:%968=785%:%
%:%969=786%:%
%:%970=787%:%
%:%971=788%:%
%:%972=789%:%
%:%973=790%:%
%:%974=791%:%
%:%975=792%:%
%:%976=793%:%
%:%977=794%:%
%:%978=795%:%
%:%979=796%:%
%:%980=797%:%
%:%981=798%:%
%:%982=799%:%
%:%983=800%:%
%:%984=801%:%
%:%985=802%:%
%:%986=803%:%
%:%987=804%:%
%:%988=805%:%
%:%989=806%:%
%:%990=807%:%
%:%991=808%:%
%:%992=809%:%
%:%993=810%:%
%:%994=811%:%
%:%995=812%:%
%:%996=813%:%
%:%997=814%:%
%:%998=815%:%
%:%999=816%:%
%:%1000=817%:%
%:%1001=818%:%
%:%1002=819%:%
%:%1003=820%:%
%:%1004=821%:%
%:%1005=822%:%
%:%1006=823%:%
%:%1007=824%:%
%:%1008=825%:%
%:%1009=826%:%
%:%1010=827%:%
%:%1011=828%:%
%:%1012=829%:%
%:%1013=830%:%
%:%1014=831%:%
%:%1015=832%:%
%:%1016=833%:%
%:%1017=834%:%
%:%1018=835%:%
%:%1019=836%:%
%:%1020=837%:%
%:%1021=838%:%
%:%1022=839%:%
%:%1023=840%:%
%:%1024=841%:%
%:%1025=842%:%
%:%1026=843%:%
%:%1027=844%:%
%:%1028=845%:%
%:%1029=846%:%
%:%1030=847%:%
%:%1031=848%:%
%:%1032=849%:%
%:%1033=850%:%
%:%1034=851%:%
%:%1035=852%:%
%:%1036=853%:%
%:%1037=854%:%
%:%1038=855%:%
%:%1039=856%:%
%:%1040=857%:%
%:%1041=858%:%
%:%1042=859%:%
%:%1043=860%:%
%:%1044=861%:%
%:%1045=862%:%
%:%1046=863%:%
%:%1047=864%:%
%:%1048=865%:%
%:%1049=866%:%
%:%1050=867%:%
%:%1051=868%:%
%:%1052=869%:%
%:%1053=870%:%
%:%1054=871%:%
%:%1055=872%:%
%:%1056=873%:%
%:%1057=874%:%
%:%1058=875%:%
%:%1059=876%:%
%:%1060=877%:%
%:%1061=878%:%
%:%1062=879%:%
%:%1063=880%:%
%:%1064=881%:%
%:%1065=882%:%
%:%1066=883%:%
%:%1067=884%:%
%:%1068=885%:%
%:%1069=886%:%
%:%1070=887%:%
%:%1071=888%:%
%:%1072=889%:%
%:%1073=890%:%
%:%1074=891%:%
%:%1075=892%:%
%:%1076=893%:%
%:%1077=894%:%
%:%1078=895%:%
%:%1079=896%:%
%:%1080=897%:%
%:%1081=898%:%
%:%1082=899%:%
%:%1083=900%:%
%:%1084=901%:%
%:%1085=902%:%
%:%1086=903%:%
%:%1087=904%:%
%:%1088=905%:%
%:%1089=906%:%
%:%1090=907%:%
%:%1091=908%:%
%:%1092=909%:%
%:%1093=910%:%
%:%1094=911%:%
%:%1095=912%:%
%:%1104=914%:%
%:%1116=922%:%
%:%1117=923%:%
%:%1118=924%:%
%:%1119=925%:%
%:%1120=926%:%
%:%1121=927%:%
%:%1122=928%:%
%:%1123=929%:%
%:%1124=930%:%
%:%1125=931%:%
%:%1126=932%:%
%:%1127=933%:%
%:%1128=934%:%
%:%1129=935%:%
%:%1130=936%:%
%:%1131=937%:%
%:%1132=938%:%
%:%1133=939%:%
%:%1134=940%:%
%:%1135=941%:%
%:%1136=942%:%
%:%1137=943%:%
%:%1138=944%:%
%:%1139=945%:%
%:%1140=946%:%
%:%1141=947%:%
%:%1142=948%:%
%:%1143=949%:%
%:%1144=950%:%
%:%1145=951%:%
%:%1146=952%:%
%:%1147=953%:%
%:%1148=954%:%
%:%1149=955%:%
%:%1150=956%:%
%:%1151=957%:%
%:%1152=958%:%
%:%1153=959%:%
%:%1154=960%:%
%:%1155=961%:%
%:%1156=962%:%
%:%1157=963%:%
%:%1158=964%:%
%:%1159=965%:%
%:%1160=966%:%
%:%1161=967%:%
%:%1162=968%:%
%:%1163=969%:%
%:%1164=970%:%
%:%1165=971%:%
%:%1166=972%:%
%:%1167=973%:%
%:%1168=974%:%
%:%1169=975%:%
%:%1170=976%:%
%:%1171=977%:%
%:%1172=978%:%
%:%1173=979%:%
%:%1174=980%:%
%:%1175=981%:%
%:%1176=982%:%
%:%1177=983%:%
%:%1178=984%:%
%:%1179=985%:%
%:%1180=986%:%
%:%1181=987%:%
%:%1182=988%:%
%:%1183=989%:%
%:%1184=990%:%
%:%1185=991%:%
%:%1186=992%:%
%:%1187=993%:%
%:%1188=994%:%
%:%1189=995%:%
%:%1190=996%:%
%:%1191=997%:%
%:%1192=998%:%
%:%1193=999%:%
%:%1194=1000%:%
%:%1195=1001%:%
%:%1196=1002%:%
%:%1197=1003%:%
%:%1198=1004%:%
%:%1199=1005%:%
%:%1200=1006%:%
%:%1201=1007%:%
%:%1202=1008%:%
%:%1203=1009%:%
%:%1204=1010%:%
%:%1205=1011%:%
%:%1206=1012%:%
%:%1207=1013%:%
%:%1208=1014%:%
%:%1209=1015%:%
%:%1210=1016%:%
%:%1211=1017%:%
%:%1212=1018%:%
%:%1213=1019%:%
%:%1214=1020%:%
%:%1215=1021%:%
%:%1216=1022%:%
%:%1217=1023%:%
%:%1218=1024%:%
%:%1219=1025%:%
%:%1220=1026%:%
%:%1221=1027%:%
%:%1222=1028%:%
%:%1223=1029%:%
%:%1224=1030%:%
%:%1225=1031%:%
%:%1226=1032%:%
%:%1227=1033%:%
%:%1228=1034%:%
%:%1229=1035%:%
%:%1230=1036%:%
%:%1231=1037%:%
%:%1232=1038%:%
%:%1233=1039%:%
%:%1234=1040%:%
%:%1235=1041%:%
%:%1236=1042%:%
%:%1237=1043%:%
%:%1238=1044%:%
%:%1239=1045%:%
%:%1240=1046%:%
%:%1241=1047%:%
%:%1242=1048%:%
%:%1243=1049%:%
%:%1244=1050%:%
%:%1245=1051%:%
%:%1246=1052%:%
%:%1247=1053%:%
%:%1248=1054%:%
%:%1249=1055%:%
%:%1250=1056%:%
%:%1251=1057%:%
%:%1252=1058%:%
%:%1253=1059%:%
%:%1254=1060%:%
%:%1255=1061%:%
%:%1256=1062%:%
%:%1257=1063%:%
%:%1258=1064%:%
%:%1259=1065%:%
%:%1260=1066%:%
%:%1261=1067%:%
%:%1262=1068%:%
%:%1263=1069%:%
%:%1264=1070%:%
%:%1265=1071%:%
%:%1266=1072%:%
%:%1267=1073%:%
%:%1268=1074%:%
%:%1269=1075%:%
%:%1270=1076%:%
%:%1271=1077%:%
%:%1272=1078%:%
%:%1273=1079%:%
%:%1274=1080%:%
%:%1275=1081%:%
%:%1276=1082%:%
%:%1277=1083%:%
%:%1278=1084%:%
%:%1279=1085%:%
%:%1280=1086%:%
%:%1281=1087%:%
%:%1282=1088%:%
%:%1283=1089%:%
%:%1284=1090%:%
%:%1285=1091%:%
%:%1286=1092%:%
%:%1287=1093%:%
%:%1288=1094%:%
%:%1289=1095%:%
%:%1290=1096%:%
%:%1291=1097%:%
%:%1292=1098%:%
%:%1293=1099%:%
%:%1310=1115%:%
%:%1311=1116%:%
%:%1312=1117%:%
%:%1313=1118%:%
%:%1314=1119%:%
%:%1315=1120%:%
%:%1316=1121%:%
%:%1317=1122%:%
%:%1318=1123%:%
%:%1319=1124%:%
%:%1320=1125%:%
%:%1321=1126%:%
%:%1322=1127%:%
%:%1323=1128%:%
%:%1324=1129%:%
%:%1325=1130%:%
%:%1326=1131%:%
%:%1327=1132%:%
%:%1328=1133%:%
%:%1329=1134%:%
%:%1330=1135%:%
%:%1331=1136%:%
%:%1332=1137%:%
%:%1333=1138%:%
%:%1334=1139%:%
%:%1335=1140%:%
%:%1336=1141%:%
%:%1337=1142%:%
%:%1338=1143%:%
%:%1339=1144%:%
%:%1340=1145%:%
%:%1341=1146%:%
%:%1342=1147%:%
%:%1343=1148%:%
%:%1344=1149%:%
%:%1345=1150%:%
%:%1346=1151%:%
%:%1347=1152%:%
%:%1348=1153%:%
%:%1349=1154%:%
%:%1350=1155%:%
%:%1351=1156%:%
%:%1352=1157%:%
%:%1353=1158%:%
%:%1354=1159%:%
%:%1355=1160%:%
%:%1356=1161%:%
%:%1357=1162%:%
%:%1358=1163%:%
%:%1359=1164%:%
%:%1360=1165%:%
%:%1361=1166%:%
%:%1362=1167%:%
%:%1363=1168%:%
%:%1364=1169%:%
%:%1365=1170%:%
%:%1366=1171%:%
%:%1367=1172%:%
%:%1368=1173%:%
%:%1369=1174%:%
%:%1370=1175%:%
%:%1371=1176%:%
%:%1372=1177%:%
%:%1373=1178%:%
%:%1374=1179%:%
%:%1375=1180%:%
%:%1376=1181%:%
%:%1377=1182%:%
%:%1378=1183%:%
%:%1379=1184%:%
%:%1380=1185%:%
%:%1381=1186%:%
%:%1382=1187%:%
%:%1383=1188%:%
%:%1384=1189%:%
%:%1385=1190%:%
%:%1386=1191%:%
%:%1387=1192%:%
%:%1388=1193%:%
%:%1389=1194%:%
%:%1390=1195%:%
%:%1391=1196%:%
%:%1392=1197%:%
%:%1393=1198%:%
%:%1394=1199%:%
%:%1395=1200%:%
%:%1396=1201%:%
%:%1397=1202%:%
%:%1398=1203%:%
%:%1399=1204%:%
%:%1400=1205%:%
%:%1401=1206%:%
%:%1402=1207%:%
%:%1403=1208%:%
%:%1404=1209%:%
%:%1405=1210%:%
%:%1406=1211%:%
%:%1407=1212%:%
%:%1408=1213%:%
%:%1409=1214%:%
%:%1410=1215%:%
%:%1411=1216%:%
%:%1412=1217%:%
%:%1413=1218%:%
%:%1414=1219%:%
%:%1415=1220%:%
%:%1416=1221%:%
%:%1417=1222%:%
%:%1418=1223%:%
%:%1419=1224%:%
%:%1420=1225%:%
%:%1421=1226%:%
%:%1422=1227%:%
%:%1423=1228%:%
%:%1424=1229%:%
%:%1425=1230%:%
%:%1426=1231%:%
%:%1427=1232%:%
%:%1428=1233%:%
%:%1429=1234%:%
%:%1430=1235%:%
%:%1439=1237%:%
%:%1451=1251%:%
%:%1452=1252%:%
%:%1461=1254%:%
%:%1473=1256%:%
%:%1474=1257%:%
%:%1475=1258%:%
%:%1476=1259%:%
%:%1477=1260%:%
%:%1478=1261%:%
%:%1479=1262%:%
%:%1480=1263%:%
%:%1481=1264%:%
%:%1482=1265%:%
%:%1483=1266%:%
%:%1484=1267%:%
%:%1485=1268%:%
%:%1486=1269%:%
%:%1487=1270%:%
%:%1488=1271%:%
%:%1489=1272%:%
%:%1496=1275%:%
%:%1497=1276%:%
%:%1498=1277%:%
%:%1499=1278%:%
%:%1500=1279%:%
%:%1501=1280%:%
%:%1502=1281%:%
%:%1503=1282%:%
%:%1504=1283%:%
%:%1505=1284%:%
%:%1506=1285%:%
%:%1507=1286%:%
%:%1508=1287%:%
%:%1509=1288%:%
%:%1510=1289%:%
%:%1511=1290%:%
%:%1512=1291%:%
%:%1513=1292%:%
%:%1514=1293%:%
%:%1515=1294%:%
%:%1516=1295%:%
%:%1517=1296%:%
%:%1518=1297%:%
%:%1519=1298%:%
%:%1520=1299%:%
%:%1521=1300%:%
%:%1522=1301%:%
%:%1523=1302%:%
%:%1524=1303%:%
%:%1525=1304%:%
%:%1526=1305%:%
%:%1527=1306%:%
%:%1528=1307%:%
%:%1529=1308%:%
%:%1530=1309%:%
%:%1531=1310%:%
%:%1532=1311%:%
%:%1533=1312%:%
%:%1534=1313%:%
%:%1535=1314%:%
%:%1536=1315%:%
%:%1537=1316%:%
%:%1538=1317%:%
%:%1539=1318%:%
%:%1540=1319%:%
%:%1541=1320%:%
%:%1542=1321%:%
%:%1543=1322%:%
%:%1544=1323%:%
%:%1545=1324%:%
%:%1546=1325%:%
%:%1547=1326%:%
%:%1548=1327%:%
%:%1549=1328%:%
%:%1550=1329%:%
%:%1551=1330%:%
%:%1552=1331%:%
%:%1553=1332%:%
%:%1554=1333%:%
%:%1555=1334%:%
%:%1556=1335%:%
%:%1557=1336%:%
%:%1558=1337%:%
%:%1559=1338%:%
%:%1560=1339%:%
%:%1561=1340%:%
%:%1562=1341%:%
%:%1563=1342%:%
%:%1564=1343%:%
%:%1565=1344%:%
%:%1566=1345%:%
%:%1567=1346%:%
%:%1568=1347%:%
%:%1569=1348%:%
%:%1570=1349%:%
%:%1571=1350%:%
%:%1572=1351%:%
%:%1573=1352%:%
%:%1574=1353%:%
%:%1575=1354%:%
%:%1576=1355%:%
%:%1577=1356%:%
%:%1578=1357%:%
%:%1579=1358%:%
%:%1580=1359%:%
%:%1581=1360%:%
%:%1582=1361%:%
%:%1583=1362%:%
%:%1584=1363%:%
%:%1585=1364%:%
%:%1586=1365%:%
%:%1587=1366%:%
%:%1588=1367%:%
%:%1589=1368%:%
%:%1590=1369%:%
%:%1591=1370%:%
%:%1592=1371%:%
%:%1593=1372%:%
%:%1594=1373%:%
%:%1595=1374%:%
%:%1596=1375%:%
%:%1597=1376%:%
%:%1598=1377%:%
%:%1599=1378%:%
%:%1600=1379%:%
%:%1601=1380%:%
%:%1602=1381%:%
%:%1603=1382%:%
%:%1604=1383%:%
%:%1605=1384%:%
%:%1606=1385%:%
%:%1607=1386%:%
%:%1608=1387%:%
%:%1609=1388%:%
%:%1610=1389%:%
%:%1611=1390%:%
%:%1612=1391%:%
%:%1613=1392%:%
%:%1614=1393%:%
%:%1615=1394%:%
%:%1616=1395%:%
%:%1617=1396%:%
%:%1618=1397%:%
%:%1619=1398%:%
%:%1620=1399%:%
%:%1621=1400%:%
%:%1622=1401%:%
%:%1623=1402%:%
%:%1624=1403%:%
%:%1625=1404%:%
%:%1626=1405%:%
%:%1627=1406%:%
%:%1628=1407%:%
%:%1629=1408%:%
%:%1630=1409%:%
%:%1631=1410%:%
%:%1632=1411%:%
%:%1633=1412%:%
%:%1634=1413%:%
%:%1635=1414%:%
%:%1636=1415%:%
%:%1637=1416%:%
%:%1638=1417%:%
%:%1639=1418%:%
%:%1640=1419%:%
%:%1641=1420%:%
%:%1642=1421%:%
%:%1643=1422%:%
%:%1644=1423%:%
%:%1645=1424%:%
%:%1646=1425%:%
%:%1647=1426%:%
%:%1648=1427%:%
%:%1649=1428%:%
%:%1650=1429%:%
%:%1651=1430%:%
%:%1652=1431%:%
%:%1653=1432%:%
%:%1654=1433%:%
%:%1655=1434%:%
%:%1656=1435%:%
%:%1657=1436%:%
%:%1658=1437%:%
%:%1659=1438%:%
%:%1660=1439%:%
%:%1661=1440%:%
%:%1662=1441%:%
%:%1663=1442%:%
%:%1664=1443%:%
%:%1665=1444%:%
%:%1666=1445%:%
%:%1667=1446%:%
%:%1668=1447%:%
%:%1669=1448%:%
%:%1670=1449%:%
%:%1671=1450%:%
%:%1672=1451%:%
%:%1673=1452%:%
%:%1674=1453%:%
%:%1675=1454%:%
%:%1676=1455%:%
%:%1677=1456%:%
%:%1678=1457%:%
%:%1679=1458%:%
%:%1680=1459%:%
%:%1681=1460%:%
%:%1682=1461%:%
%:%1683=1462%:%
%:%1684=1463%:%
%:%1685=1464%:%
%:%1686=1465%:%
%:%1687=1466%:%
%:%1688=1467%:%
%:%1689=1468%:%
%:%1690=1469%:%
%:%1691=1470%:%
%:%1692=1471%:%
%:%1693=1472%:%
%:%1694=1473%:%
%:%1695=1474%:%
%:%1696=1475%:%
%:%1697=1476%:%
%:%1698=1477%:%
%:%1699=1478%:%
%:%1700=1479%:%
%:%1701=1480%:%
%:%1702=1481%:%
%:%1703=1482%:%
%:%1704=1483%:%
%:%1705=1484%:%
%:%1706=1485%:%
%:%1707=1486%:%
%:%1708=1487%:%
%:%1709=1488%:%
%:%1710=1489%:%
%:%1711=1490%:%
%:%1712=1491%:%
%:%1713=1492%:%
%:%1714=1493%:%
%:%1715=1494%:%
%:%1716=1495%:%
%:%1717=1496%:%
%:%1718=1497%:%
%:%1719=1498%:%
%:%1720=1499%:%
%:%1721=1500%:%
%:%1722=1501%:%
%:%1723=1502%:%
%:%1724=1503%:%
%:%1725=1504%:%
%:%1726=1505%:%
%:%1727=1506%:%
%:%1728=1507%:%
%:%1729=1508%:%
%:%1730=1509%:%
%:%1731=1510%:%
%:%1732=1511%:%
%:%1733=1512%:%
%:%1734=1513%:%
%:%1735=1514%:%
%:%1736=1515%:%
%:%1737=1516%:%
%:%1738=1517%:%
%:%1739=1518%:%
%:%1740=1519%:%
%:%1741=1520%:%
%:%1742=1521%:%
%:%1743=1522%:%
%:%1744=1523%:%
%:%1745=1524%:%
%:%1746=1525%:%
%:%1747=1526%:%
%:%1748=1527%:%
%:%1749=1528%:%
%:%1750=1529%:%
%:%1751=1530%:%
%:%1752=1531%:%
%:%1753=1532%:%
%:%1754=1533%:%
%:%1755=1534%:%
%:%1756=1535%:%
%:%1757=1536%:%
%:%1758=1537%:%
%:%1759=1538%:%
%:%1760=1539%:%
%:%1761=1540%:%
%:%1762=1541%:%
%:%1763=1542%:%
%:%1764=1543%:%
%:%1765=1544%:%
%:%1766=1545%:%
%:%1767=1546%:%
%:%1768=1547%:%
%:%1769=1548%:%
%:%1770=1549%:%
%:%1771=1550%:%
%:%1772=1551%:%
%:%1773=1552%:%
%:%1774=1553%:%
%:%1775=1554%:%
%:%1776=1555%:%
%:%1777=1556%:%
%:%1778=1557%:%
%:%1779=1558%:%
%:%1780=1559%:%
%:%1781=1560%:%
%:%1782=1561%:%
%:%1783=1562%:%
%:%1784=1563%:%
%:%1785=1564%:%
%:%1786=1565%:%
%:%1787=1566%:%
%:%1788=1567%:%
%:%1789=1568%:%
%:%1790=1569%:%
%:%1791=1570%:%
%:%1792=1571%:%
%:%1793=1572%:%
%:%1794=1573%:%
%:%1795=1574%:%
%:%1796=1575%:%
%:%1797=1576%:%
%:%1798=1577%:%
%:%1799=1578%:%
%:%1800=1579%:%
%:%1801=1580%:%
%:%1802=1581%:%
%:%1803=1582%:%
%:%1804=1583%:%
%:%1805=1584%:%
%:%1806=1585%:%
%:%1807=1586%:%
%:%1808=1587%:%
%:%1809=1588%:%
%:%1810=1589%:%
%:%1811=1590%:%
%:%1812=1591%:%
%:%1813=1592%:%
%:%1814=1593%:%
%:%1815=1594%:%
%:%1816=1595%:%
%:%1817=1596%:%
%:%1818=1597%:%
%:%1819=1598%:%
%:%1820=1599%:%
%:%1821=1600%:%
%:%1822=1601%:%
%:%1823=1602%:%
%:%1824=1603%:%
%:%1825=1604%:%
%:%1826=1605%:%
%:%1827=1606%:%
%:%1828=1607%:%
%:%1829=1608%:%
%:%1830=1609%:%
%:%1831=1610%:%
%:%1832=1611%:%
%:%1833=1612%:%
%:%1834=1613%:%
%:%1835=1614%:%
%:%1836=1615%:%
%:%1837=1616%:%
%:%1838=1617%:%
%:%1839=1618%:%
%:%1840=1619%:%
%:%1841=1620%:%
%:%1842=1621%:%
%:%1843=1622%:%
%:%1844=1623%:%
%:%1845=1624%:%
%:%1846=1625%:%
%:%1847=1626%:%
%:%1848=1627%:%
%:%1849=1628%:%
%:%1850=1629%:%
%:%1851=1630%:%
%:%1852=1631%:%
%:%1861=1633%:%
%:%1873=1635%:%
%:%1874=1636%:%
%:%1875=1637%:%
%:%1876=1638%:%
%:%1877=1639%:%
%:%1878=1640%:%
%:%1879=1641%:%
%:%1880=1642%:%
%:%1881=1643%:%
%:%1882=1644%:%
%:%1883=1645%:%
%:%1884=1646%:%
%:%1885=1647%:%
%:%1886=1648%:%
%:%1887=1649%:%
%:%1888=1650%:%
%:%1889=1651%:%
%:%1890=1652%:%
%:%1891=1653%:%
%:%1892=1654%:%
%:%1893=1655%:%
%:%1894=1656%:%
%:%1895=1657%:%
%:%1896=1658%:%
%:%1897=1659%:%
%:%1898=1660%:%
%:%1899=1661%:%
%:%1900=1662%:%
%:%1901=1663%:%
%:%1902=1664%:%
%:%1903=1665%:%
%:%1904=1666%:%
%:%1905=1667%:%
%:%1906=1668%:%
%:%1907=1669%:%
%:%1908=1670%:%
%:%1909=1671%:%
%:%1910=1672%:%
%:%1911=1673%:%
%:%1912=1674%:%
%:%1913=1675%:%
%:%1914=1676%:%
%:%1915=1677%:%
%:%1916=1678%:%
%:%1917=1679%:%
%:%1918=1680%:%
%:%1919=1681%:%
%:%1920=1682%:%
%:%1921=1683%:%
%:%1922=1684%:%
%:%1923=1685%:%
%:%1924=1686%:%
%:%1925=1687%:%
%:%1926=1688%:%
%:%1927=1689%:%
%:%1928=1690%:%
%:%1929=1691%:%
%:%1930=1692%:%
%:%1931=1693%:%
%:%1932=1694%:%
%:%1933=1695%:%
%:%1934=1696%:%
%:%1935=1697%:%
%:%1936=1698%:%
%:%1937=1699%:%
%:%1938=1700%:%
%:%1939=1701%:%
%:%1940=1702%:%
%:%1941=1703%:%
%:%1942=1704%:%
%:%1943=1705%:%
%:%1944=1706%:%
%:%1945=1707%:%
%:%1946=1708%:%
%:%1947=1709%:%
%:%1948=1710%:%
%:%1949=1711%:%
%:%1950=1712%:%
%:%1951=1713%:%
%:%1952=1714%:%
%:%1953=1715%:%
%:%1954=1716%:%
%:%1955=1717%:%
%:%1956=1718%:%
%:%1957=1719%:%
%:%1958=1720%:%
%:%1959=1721%:%
%:%1960=1722%:%
%:%1961=1723%:%
%:%1962=1724%:%
%:%1963=1725%:%
%:%1964=1726%:%
%:%1965=1727%:%
%:%1966=1728%:%
%:%1967=1729%:%
%:%1968=1730%:%
%:%1969=1731%:%
%:%1970=1732%:%
%:%1971=1733%:%
%:%1972=1734%:%
%:%1973=1735%:%
%:%1974=1736%:%
%:%1975=1737%:%
%:%1976=1738%:%
%:%1977=1739%:%
%:%1978=1740%:%
%:%1979=1741%:%
%:%1980=1742%:%
%:%1989=1744%:%
%:%2019=1754%:%
%:%2031=1756%:%
%:%2032=1757%:%
%:%2033=1758%:%
%:%2034=1759%:%
%:%2035=1760%:%
%:%2036=1761%:%
%:%2037=1762%:%
%:%2038=1763%:%
%:%2039=1764%:%
%:%2040=1765%:%
%:%2041=1766%:%
%:%2042=1767%:%
%:%2043=1768%:%
%:%2044=1769%:%
%:%2045=1770%:%
%:%2046=1771%:%
%:%2047=1772%:%
%:%2048=1773%:%
%:%2049=1774%:%
%:%2050=1775%:%
%:%2051=1776%:%
%:%2052=1777%:%
%:%2053=1778%:%
%:%2054=1779%:%
%:%2055=1780%:%
%:%2056=1781%:%
%:%2057=1782%:%
%:%2058=1783%:%
%:%2059=1784%:%
%:%2060=1785%:%
%:%2061=1786%:%
%:%2062=1787%:%
%:%2063=1788%:%
%:%2064=1789%:%
%:%2065=1790%:%
%:%2066=1791%:%
%:%2067=1792%:%
%:%2068=1793%:%
%:%2069=1794%:%
%:%2070=1795%:%
%:%2071=1796%:%
%:%2072=1797%:%
%:%2073=1798%:%
%:%2074=1799%:%
%:%2075=1800%:%
%:%2076=1801%:%
%:%2077=1802%:%
%:%2078=1803%:%
%:%2079=1804%:%
%:%2080=1805%:%
%:%2081=1806%:%
%:%2082=1807%:%
%:%2083=1808%:%
%:%2084=1809%:%
%:%2093=1811%:%
%:%2105=1813%:%
%:%2106=1814%:%
%:%2107=1815%:%
%:%2108=1816%:%
%:%2109=1817%:%
%:%2110=1818%:%
%:%2111=1819%:%
%:%2112=1820%:%
%:%2113=1821%:%
%:%2114=1822%:%
%:%2115=1823%:%
%:%2116=1824%:%
%:%2117=1825%:%
%:%2118=1826%:%
%:%2119=1827%:%
%:%2120=1828%:%
%:%2121=1829%:%
%:%2122=1830%:%
%:%2123=1831%:%
%:%2124=1832%:%
%:%2125=1833%:%
%:%2126=1834%:%
%:%2127=1835%:%
%:%2128=1836%:%
%:%2129=1837%:%
%:%2130=1838%:%
%:%2131=1839%:%
%:%2132=1840%:%
%:%2133=1841%:%
%:%2134=1842%:%
%:%2135=1843%:%
%:%2136=1844%:%
%:%2137=1845%:%
%:%2138=1846%:%
%:%2139=1847%:%
%:%2140=1848%:%
%:%2141=1849%:%
%:%2142=1850%:%
%:%2143=1851%:%
%:%2144=1852%:%
%:%2145=1853%:%
%:%2146=1854%:%
%:%2147=1855%:%
%:%2148=1856%:%
%:%2149=1857%:%
%:%2150=1858%:%
%:%2151=1859%:%
%:%2152=1860%:%
%:%2153=1861%:%
%:%2154=1862%:%
%:%2155=1863%:%
%:%2156=1864%:%
%:%2157=1865%:%
%:%2158=1866%:%
%:%2159=1867%:%
%:%2160=1868%:%
%:%2161=1869%:%
%:%2162=1870%:%
%:%2163=1871%:%
%:%2164=1872%:%
%:%2165=1873%:%
%:%2166=1874%:%
%:%2167=1875%:%
%:%2168=1876%:%
%:%2169=1877%:%
%:%2170=1878%:%
%:%2171=1879%:%
%:%2172=1880%:%
%:%2173=1881%:%
%:%2174=1882%:%
%:%2175=1883%:%
%:%2176=1884%:%
%:%2177=1885%:%
%:%2178=1886%:%
%:%2179=1887%:%
%:%2180=1888%:%
%:%2181=1889%:%
%:%2182=1890%:%
%:%2183=1891%:%
%:%2184=1892%:%
%:%2185=1893%:%
%:%2186=1894%:%
%:%2187=1895%:%
%:%2188=1896%:%
%:%2189=1897%:%
%:%2190=1898%:%
%:%2199=1900%:%
%:%2211=1902%:%
%:%2212=1903%:%
%:%2213=1904%:%
%:%2214=1905%:%
%:%2215=1906%:%
%:%2216=1907%:%
%:%2217=1908%:%
%:%2218=1909%:%
%:%2219=1910%:%
%:%2220=1911%:%
%:%2221=1912%:%
%:%2222=1913%:%
%:%2223=1914%:%
%:%2224=1915%:%
%:%2225=1916%:%
%:%2226=1917%:%
%:%2227=1918%:%
%:%2228=1919%:%
%:%2229=1920%:%
%:%2230=1921%:%
%:%2231=1922%:%
%:%2232=1923%:%
%:%2233=1924%:%
%:%2234=1925%:%
%:%2235=1926%:%
%:%2236=1927%:%
%:%2237=1928%:%
%:%2238=1929%:%
%:%2239=1930%:%
%:%2240=1931%:%
%:%2241=1932%:%
%:%2242=1933%:%
%:%2243=1934%:%
%:%2244=1935%:%
%:%2245=1936%:%
%:%2246=1937%:%
%:%2247=1938%:%
%:%2248=1939%:%
%:%2249=1940%:%
%:%2250=1941%:%
%:%2251=1942%:%
%:%2252=1943%:%
%:%2253=1944%:%
%:%2254=1945%:%
%:%2255=1946%:%
%:%2256=1947%:%
%:%2257=1948%:%
%:%2258=1949%:%
%:%2259=1950%:%
%:%2260=1951%:%
%:%2261=1952%:%
%:%2262=1953%:%
%:%2263=1954%:%
%:%2264=1955%:%
%:%2265=1956%:%
%:%2266=1957%:%
%:%2267=1958%:%
%:%2268=1959%:%
%:%2269=1960%:%
%:%2270=1961%:%
%:%2271=1962%:%
%:%2272=1963%:%
%:%2273=1964%:%
%:%2274=1965%:%
%:%2275=1966%:%
%:%2276=1967%:%
%:%2277=1968%:%
%:%2278=1969%:%
%:%2279=1970%:%
%:%2280=1971%:%
%:%2281=1972%:%
%:%2282=1973%:%
%:%2283=1974%:%
%:%2284=1975%:%
%:%2285=1976%:%
%:%2286=1977%:%
%:%2287=1978%:%
%:%2288=1979%:%
%:%2289=1980%:%
%:%2290=1981%:%
%:%2291=1982%:%
%:%2292=1983%:%
%:%2293=1984%:%
%:%2294=1985%:%
%:%2295=1986%:%
%:%2296=1987%:%
%:%2297=1988%:%
%:%2298=1989%:%
%:%2299=1990%:%
%:%2300=1991%:%
%:%2301=1992%:%
%:%2302=1993%:%
%:%2303=1994%:%
%:%2304=1995%:%
%:%2305=1996%:%
%:%2306=1997%:%
%:%2307=1998%:%
%:%2308=1999%:%
%:%2309=2000%:%
%:%2310=2001%:%
%:%2311=2002%:%
%:%2312=2003%:%
%:%2313=2004%:%
%:%2314=2005%:%
%:%2315=2006%:%
%:%2316=2007%:%
%:%2317=2008%:%
%:%2318=2009%:%
%:%2319=2010%:%
%:%2320=2011%:%
%:%2321=2012%:%
%:%2322=2013%:%
%:%2323=2014%:%
%:%2324=2015%:%
%:%2325=2016%:%
%:%2326=2017%:%
%:%2327=2018%:%
%:%2328=2019%:%
%:%2329=2020%:%
%:%2330=2021%:%
%:%2331=2022%:%
%:%2332=2023%:%
%:%2333=2024%:%
%:%2334=2025%:%
%:%2335=2026%:%
%:%2336=2027%:%
%:%2337=2028%:%
%:%2338=2029%:%
%:%2339=2030%:%
%:%2340=2031%:%
%:%2341=2032%:%
%:%2342=2033%:%
%:%2343=2034%:%
%:%2344=2035%:%
%:%2345=2036%:%
%:%2346=2037%:%
%:%2347=2038%:%
%:%2348=2039%:%
%:%2349=2040%:%
%:%2350=2041%:%
%:%2351=2042%:%
%:%2352=2043%:%
%:%2353=2044%:%
%:%2354=2045%:%
%:%2355=2046%:%
%:%2356=2047%:%
%:%2365=2049%:%
%:%2377=2051%:%
%:%2378=2052%:%
%:%2379=2053%:%
%:%2380=2054%:%
%:%2381=2055%:%
%:%2382=2056%:%
%:%2383=2057%:%
%:%2384=2058%:%
%:%2385=2059%:%
%:%2386=2060%:%
%:%2387=2061%:%
%:%2388=2062%:%
%:%2389=2063%:%
%:%2390=2064%:%
%:%2391=2065%:%
%:%2392=2066%:%
%:%2393=2067%:%
%:%2394=2068%:%
%:%2395=2069%:%
%:%2396=2070%:%
%:%2397=2071%:%
%:%2398=2072%:%
%:%2399=2073%:%
%:%2400=2074%:%
%:%2401=2075%:%
%:%2402=2076%:%
%:%2403=2077%:%
%:%2404=2078%:%
%:%2405=2079%:%
%:%2406=2080%:%
%:%2407=2081%:%
%:%2408=2082%:%
%:%2409=2083%:%
%:%2410=2084%:%
%:%2411=2085%:%
%:%2412=2086%:%
%:%2413=2087%:%
%:%2414=2088%:%
%:%2415=2089%:%
%:%2416=2090%:%
%:%2417=2091%:%
%:%2418=2092%:%
%:%2419=2093%:%
%:%2420=2094%:%
%:%2421=2095%:%
%:%2422=2096%:%
%:%2423=2097%:%
%:%2424=2098%:%
%:%2425=2099%:%
%:%2426=2100%:%
%:%2427=2101%:%
%:%2428=2102%:%
%:%2429=2103%:%
%:%2430=2104%:%
%:%2431=2105%:%
%:%2432=2106%:%
%:%2433=2107%:%
%:%2434=2108%:%
%:%2435=2109%:%
%:%2436=2110%:%
%:%2437=2111%:%
%:%2438=2112%:%
%:%2439=2113%:%
%:%2440=2114%:%
%:%2441=2115%:%
%:%2442=2116%:%
%:%2443=2117%:%
%:%2444=2118%:%
%:%2445=2119%:%
%:%2446=2120%:%
%:%2447=2121%:%
%:%2448=2122%:%
%:%2449=2123%:%
%:%2450=2124%:%
%:%2451=2125%:%
%:%2452=2126%:%
%:%2453=2127%:%
%:%2454=2128%:%
%:%2455=2129%:%
%:%2456=2130%:%
%:%2457=2131%:%
%:%2458=2132%:%
%:%2459=2133%:%
%:%2460=2134%:%
%:%2461=2135%:%
%:%2462=2136%:%
%:%2463=2137%:%
%:%2464=2138%:%
%:%2465=2139%:%
%:%2466=2140%:%
%:%2467=2141%:%
%:%2468=2142%:%
%:%2469=2143%:%
%:%2470=2144%:%
%:%2471=2145%:%
%:%2472=2146%:%
%:%2473=2147%:%
%:%2474=2148%:%
%:%2475=2149%:%
%:%2476=2150%:%
%:%2477=2151%:%
%:%2478=2152%:%
%:%2479=2153%:%
%:%2480=2154%:%
%:%2481=2155%:%
%:%2482=2156%:%
%:%2483=2157%:%
%:%2484=2158%:%
%:%2485=2159%:%
%:%2486=2160%:%
%:%2487=2161%:%
%:%2488=2162%:%
%:%2489=2163%:%
%:%2490=2164%:%
%:%2491=2165%:%
%:%2492=2166%:%
%:%2493=2167%:%
%:%2494=2168%:%
%:%2495=2169%:%
%:%2496=2170%:%
%:%2497=2171%:%
%:%2498=2172%:%
%:%2499=2173%:%
%:%2500=2174%:%
%:%2501=2175%:%
%:%2502=2176%:%
%:%2503=2177%:%
%:%2504=2178%:%
%:%2505=2179%:%
%:%2506=2180%:%
%:%2507=2181%:%
%:%2508=2182%:%
%:%2509=2183%:%
%:%2510=2184%:%
%:%2511=2185%:%
%:%2512=2186%:%
%:%2513=2187%:%
%:%2514=2188%:%
%:%2515=2189%:%
%:%2516=2190%:%
%:%2517=2191%:%
%:%2518=2192%:%
%:%2519=2193%:%
%:%2520=2194%:%
%:%2521=2195%:%
%:%2522=2196%:%
%:%2523=2197%:%
%:%2524=2198%:%
%:%2525=2199%:%
%:%2526=2200%:%
%:%2527=2201%:%
%:%2528=2202%:%
%:%2529=2203%:%
%:%2530=2204%:%
%:%2531=2205%:%
%:%2532=2206%:%
%:%2533=2207%:%
%:%2534=2208%:%
%:%2535=2209%:%
%:%2536=2210%:%
%:%2537=2211%:%
%:%2538=2212%:%
%:%2539=2213%:%
%:%2540=2214%:%
%:%2541=2215%:%
%:%2542=2216%:%
%:%2543=2217%:%
%:%2544=2218%:%
%:%2545=2219%:%
%:%2546=2220%:%
%:%2547=2221%:%
%:%2548=2222%:%
%:%2549=2223%:%
%:%2550=2224%:%
%:%2551=2225%:%
%:%2552=2226%:%
%:%2553=2227%:%
%:%2554=2228%:%
%:%2555=2229%:%
%:%2556=2230%:%
%:%2557=2231%:%
%:%2558=2232%:%
%:%2559=2233%:%
%:%2560=2234%:%
%:%2561=2235%:%
%:%2562=2236%:%
%:%2563=2237%:%
%:%2564=2238%:%
%:%2565=2239%:%
%:%2566=2240%:%
%:%2567=2241%:%
%:%2568=2242%:%
%:%2569=2243%:%
%:%2570=2244%:%
%:%2571=2245%:%
%:%2572=2246%:%
%:%2573=2247%:%
%:%2574=2248%:%
%:%2575=2249%:%
%:%2576=2250%:%
%:%2577=2251%:%
%:%2578=2252%:%
%:%2579=2253%:%
%:%2580=2254%:%
%:%2581=2255%:%
%:%2582=2256%:%
%:%2583=2257%:%
%:%2584=2258%:%
%:%2585=2259%:%
%:%2586=2260%:%
%:%2587=2261%:%
%:%2588=2262%:%
%:%2589=2263%:%
%:%2590=2264%:%
%:%2591=2265%:%
%:%2592=2266%:%
%:%2593=2267%:%
%:%2594=2268%:%
%:%2595=2269%:%
%:%2596=2270%:%
%:%2597=2271%:%
%:%2598=2272%:%
%:%2599=2273%:%
%:%2600=2274%:%
%:%2601=2275%:%
%:%2602=2276%:%
%:%2603=2277%:%
%:%2604=2278%:%
%:%2605=2279%:%
%:%2606=2280%:%
%:%2607=2281%:%
%:%2608=2282%:%
%:%2609=2283%:%
%:%2610=2284%:%
%:%2611=2285%:%
%:%2612=2286%:%
%:%2613=2287%:%
%:%2614=2288%:%
%:%2615=2289%:%
%:%2616=2290%:%
%:%2617=2291%:%
%:%2618=2292%:%
%:%2619=2293%:%
%:%2620=2294%:%
%:%2621=2295%:%
%:%2622=2296%:%
%:%2623=2297%:%
%:%2624=2298%:%
%:%2625=2299%:%
%:%2626=2300%:%
%:%2627=2301%:%
%:%2628=2302%:%
%:%2629=2303%:%
%:%2630=2304%:%
%:%2631=2305%:%
%:%2632=2306%:%
%:%2633=2307%:%
%:%2634=2308%:%
%:%2635=2309%:%
%:%2636=2310%:%
%:%2637=2311%:%
%:%2638=2312%:%
%:%2639=2313%:%
%:%2640=2314%:%
%:%2641=2315%:%
%:%2642=2316%:%
%:%2643=2317%:%
%:%2644=2318%:%
%:%2645=2319%:%
%:%2646=2320%:%
%:%2647=2321%:%
%:%2648=2322%:%
%:%2649=2323%:%
%:%2650=2324%:%
%:%2651=2325%:%
%:%2652=2326%:%
%:%2653=2327%:%
%:%2654=2328%:%
%:%2655=2329%:%
%:%2656=2330%:%
%:%2657=2331%:%
%:%2658=2332%:%
%:%2659=2333%:%
%:%2660=2334%:%
%:%2661=2335%:%
%:%2662=2336%:%
%:%2663=2337%:%
%:%2664=2338%:%
%:%2665=2339%:%
%:%2666=2340%:%
%:%2667=2341%:%
%:%2668=2342%:%
%:%2669=2343%:%
%:%2670=2344%:%
%:%2671=2345%:%
%:%2672=2346%:%
%:%2673=2347%:%
%:%2674=2348%:%
%:%2675=2349%:%
%:%2676=2350%:%
%:%2677=2351%:%
%:%2678=2352%:%
%:%2679=2353%:%
%:%2680=2354%:%
%:%2681=2355%:%
%:%2682=2356%:%
%:%2683=2357%:%
%:%2684=2358%:%
%:%2685=2359%:%
%:%2686=2360%:%
%:%2687=2361%:%
%:%2688=2362%:%
%:%2689=2363%:%
%:%2690=2364%:%
%:%2691=2365%:%
%:%2692=2366%:%
%:%2693=2367%:%
%:%2694=2368%:%
%:%2695=2369%:%
%:%2696=2370%:%
%:%2697=2371%:%
%:%2698=2372%:%
%:%2699=2373%:%
%:%2700=2374%:%
%:%2701=2375%:%
%:%2702=2376%:%
%:%2703=2377%:%
%:%2704=2378%:%
%:%2705=2379%:%
%:%2706=2380%:%
%:%2707=2381%:%
%:%2708=2382%:%
%:%2709=2383%:%
%:%2710=2384%:%
%:%2719=2386%:%
%:%2731=2388%:%
%:%2732=2389%:%
%:%2733=2390%:%
%:%2734=2391%:%
%:%2735=2392%:%
%:%2736=2393%:%
%:%2737=2394%:%
%:%2738=2395%:%
%:%2739=2396%:%
%:%2740=2397%:%
%:%2750=2398%:%
%:%2751=2399%:%
%:%2752=2400%:%
%:%2753=2401%:%
%:%2763=2402%:%
%:%2764=2403%:%
%:%2765=2404%:%
%:%2766=2405%:%
%:%2767=2406%:%
%:%2768=2407%:%
%:%2769=2408%:%
%:%2770=2409%:%
%:%2771=2410%:%
%:%2772=2411%:%
%:%2773=2412%:%
%:%2774=2413%:%
%:%2775=2414%:%
%:%2776=2415%:%
%:%2777=2416%:%
%:%2778=2417%:%
%:%2779=2418%:%
%:%2780=2419%:%
%:%2781=2420%:%
%:%2782=2421%:%
%:%2783=2422%:%
%:%2784=2423%:%
%:%2785=2424%:%
%:%2786=2425%:%
%:%2787=2426%:%
%:%2788=2427%:%
%:%2789=2428%:%
%:%2790=2429%:%
%:%2791=2430%:%
%:%2792=2431%:%
%:%2793=2432%:%
%:%2794=2433%:%
%:%2795=2434%:%
%:%2796=2435%:%
%:%2797=2436%:%
%:%2798=2437%:%
%:%2799=2438%:%
%:%2800=2439%:%
%:%2801=2440%:%
%:%2802=2441%:%
%:%2803=2442%:%
%:%2804=2443%:%
%:%2805=2444%:%
%:%2806=2445%:%
%:%2807=2446%:%
%:%2808=2447%:%
%:%2809=2448%:%
%:%2810=2449%:%
%:%2811=2450%:%
%:%2812=2451%:%
%:%2813=2452%:%
%:%2814=2453%:%
%:%2815=2454%:%
%:%2816=2455%:%
%:%2817=2456%:%
%:%2818=2457%:%
%:%2819=2458%:%
%:%2820=2459%:%
%:%2821=2460%:%
%:%2822=2461%:%
%:%2823=2462%:%
%:%2824=2463%:%
%:%2825=2464%:%
%:%2826=2465%:%
%:%2827=2466%:%
%:%2828=2467%:%
%:%2829=2468%:%
%:%2830=2469%:%
%:%2831=2470%:%
%:%2832=2471%:%
%:%2833=2472%:%
%:%2834=2473%:%
%:%2835=2474%:%
%:%2836=2475%:%
%:%2837=2476%:%
%:%2838=2477%:%
%:%2839=2478%:%
%:%2840=2479%:%
%:%2841=2480%:%
%:%2842=2481%:%
%:%2843=2482%:%
%:%2844=2483%:%
%:%2845=2484%:%
%:%2846=2485%:%
%:%2847=2486%:%
%:%2848=2487%:%
%:%2849=2488%:%
%:%2850=2489%:%
%:%2851=2490%:%
%:%2852=2491%:%
%:%2853=2492%:%
%:%2854=2493%:%
%:%2855=2494%:%
%:%2856=2495%:%
%:%2857=2496%:%
%:%2858=2497%:%
%:%2859=2498%:%
%:%2860=2499%:%
%:%2861=2500%:%
%:%2862=2501%:%
%:%2863=2502%:%
%:%2864=2503%:%
%:%2865=2504%:%
%:%2866=2505%:%
%:%2867=2506%:%
%:%2868=2507%:%
%:%2869=2508%:%
%:%2870=2509%:%
%:%2871=2510%:%
%:%2872=2511%:%
%:%2873=2512%:%
%:%2874=2513%:%
%:%2875=2514%:%
%:%2876=2515%:%
%:%2877=2516%:%
%:%2878=2517%:%
%:%2879=2518%:%
%:%2880=2519%:%
%:%2881=2520%:%
%:%2882=2521%:%
%:%2883=2522%:%
%:%2884=2523%:%
%:%2885=2524%:%
%:%2886=2525%:%
%:%2887=2526%:%
%:%2888=2527%:%
%:%2889=2528%:%
%:%2890=2529%:%
%:%2891=2530%:%
%:%2892=2531%:%
%:%2893=2532%:%
%:%2894=2533%:%
%:%2895=2534%:%
%:%2896=2535%:%
%:%2897=2536%:%
%:%2898=2537%:%
%:%2899=2538%:%
%:%2900=2539%:%
%:%2901=2540%:%
%:%2902=2541%:%
%:%2903=2542%:%
%:%2904=2543%:%
%:%2905=2544%:%
%:%2906=2545%:%
%:%2907=2546%:%
%:%2908=2547%:%
%:%2909=2548%:%
%:%2910=2549%:%
%:%2911=2550%:%
%:%2912=2551%:%
%:%2913=2552%:%
%:%2914=2553%:%
%:%2915=2554%:%
%:%2916=2555%:%
%:%2917=2556%:%
%:%2918=2557%:%
%:%2919=2558%:%
%:%2920=2559%:%
%:%2921=2560%:%
%:%2922=2561%:%
%:%2923=2562%:%
%:%2924=2563%:%
%:%2925=2564%:%
%:%2926=2565%:%
%:%2927=2566%:%
%:%2928=2567%:%
%:%2929=2568%:%
%:%2930=2569%:%
%:%2931=2570%:%
%:%2932=2571%:%
%:%2933=2572%:%
%:%2934=2573%:%
%:%2935=2574%:%
%:%2936=2575%:%
%:%2937=2576%:%
%:%2938=2577%:%
%:%2939=2578%:%
%:%2940=2579%:%
%:%2949=2581%:%
%:%2961=2586%:%
%:%2962=2587%:%
%:%2963=2588%:%
%:%2964=2589%:%
%:%2965=2590%:%
%:%2974=2592%:%
%:%2986=2594%:%
%:%2987=2595%:%
%:%2988=2596%:%
%:%2989=2597%:%
%:%2990=2598%:%
%:%2991=2599%:%
%:%2992=2600%:%
%:%2993=2601%:%
%:%2994=2602%:%
%:%2995=2603%:%
%:%2996=2604%:%
%:%2997=2605%:%
%:%2998=2606%:%
%:%2999=2607%:%
%:%3000=2608%:%
%:%3001=2609%:%
%:%3002=2610%:%
%:%3003=2611%:%
%:%3004=2612%:%
%:%3005=2613%:%
%:%3006=2614%:%
%:%3007=2615%:%
%:%3008=2616%:%
%:%3009=2617%:%
%:%3010=2618%:%
%:%3011=2619%:%
%:%3012=2620%:%
%:%3013=2621%:%
%:%3014=2622%:%
%:%3015=2623%:%
%:%3016=2624%:%
%:%3017=2625%:%
%:%3018=2626%:%
%:%3019=2627%:%
%:%3020=2628%:%
%:%3021=2629%:%
%:%3022=2630%:%
%:%3023=2631%:%
%:%3024=2632%:%
%:%3025=2633%:%
%:%3026=2634%:%
%:%3027=2635%:%
%:%3028=2636%:%
%:%3029=2637%:%
%:%3030=2638%:%
%:%3031=2639%:%
%:%3032=2640%:%
%:%3033=2641%:%
%:%3034=2642%:%
%:%3035=2643%:%
%:%3036=2644%:%
%:%3037=2645%:%
%:%3038=2646%:%
%:%3039=2647%:%
%:%3040=2648%:%
%:%3041=2649%:%
%:%3042=2650%:%
%:%3043=2651%:%
%:%3044=2652%:%
%:%3045=2653%:%
%:%3046=2654%:%
%:%3047=2655%:%
%:%3048=2656%:%
%:%3049=2657%:%
%:%3050=2658%:%
%:%3051=2659%:%
%:%3052=2660%:%
%:%3053=2661%:%
%:%3054=2662%:%
%:%3055=2663%:%
%:%3056=2664%:%
%:%3057=2665%:%
%:%3058=2666%:%
%:%3059=2667%:%
%:%3060=2668%:%
%:%3061=2669%:%
%:%3062=2670%:%
%:%3063=2671%:%
%:%3064=2672%:%
%:%3065=2673%:%
%:%3066=2674%:%
%:%3067=2675%:%
%:%3068=2676%:%
%:%3069=2677%:%
%:%3070=2678%:%
%:%3071=2679%:%
%:%3072=2680%:%
%:%3073=2681%:%
%:%3074=2682%:%
%:%3075=2683%:%
%:%3076=2684%:%
%:%3077=2685%:%
%:%3078=2686%:%
%:%3079=2687%:%
%:%3080=2688%:%
%:%3081=2689%:%
%:%3082=2690%:%
%:%3083=2691%:%
%:%3084=2692%:%
%:%3085=2693%:%
%:%3086=2694%:%
%:%3087=2695%:%
%:%3088=2696%:%
%:%3089=2697%:%
%:%3090=2698%:%
%:%3091=2699%:%
%:%3092=2700%:%
%:%3093=2701%:%
%:%3094=2702%:%
%:%3095=2703%:%
%:%3096=2704%:%
%:%3097=2705%:%
%:%3098=2706%:%
%:%3099=2707%:%
%:%3100=2708%:%
%:%3101=2709%:%
%:%3102=2710%:%
%:%3103=2711%:%
%:%3104=2712%:%
%:%3105=2713%:%
%:%3106=2714%:%
%:%3107=2715%:%
%:%3108=2716%:%
%:%3109=2717%:%
%:%3110=2718%:%
%:%3111=2719%:%
%:%3112=2720%:%
%:%3113=2721%:%
%:%3114=2722%:%
%:%3115=2723%:%
%:%3116=2724%:%
%:%3117=2725%:%
%:%3118=2726%:%
%:%3119=2727%:%
%:%3120=2728%:%
%:%3121=2729%:%
%:%3122=2730%:%
%:%3123=2731%:%
%:%3124=2732%:%
%:%3125=2733%:%
%:%3126=2734%:%
%:%3127=2735%:%
%:%3128=2736%:%
%:%3129=2737%:%
%:%3130=2738%:%
%:%3131=2739%:%
%:%3132=2740%:%
%:%3133=2741%:%
%:%3134=2742%:%
%:%3135=2743%:%
%:%3136=2744%:%
%:%3137=2745%:%
%:%3138=2746%:%
%:%3139=2747%:%
%:%3140=2748%:%
%:%3141=2749%:%
%:%3142=2750%:%
%:%3143=2751%:%
%:%3144=2752%:%
%:%3145=2753%:%
%:%3146=2754%:%
%:%3147=2755%:%
%:%3148=2756%:%
%:%3149=2757%:%
%:%3150=2758%:%
%:%3151=2759%:%
%:%3152=2760%:%
%:%3153=2761%:%
%:%3154=2762%:%
%:%3155=2763%:%
%:%3156=2764%:%
%:%3157=2765%:%
%:%3158=2766%:%
%:%3159=2767%:%
%:%3160=2768%:%
%:%3161=2769%:%
%:%3162=2770%:%
%:%3163=2771%:%
%:%3164=2772%:%
%:%3165=2773%:%
%:%3166=2774%:%
%:%3167=2775%:%
%:%3168=2776%:%
%:%3169=2777%:%
%:%3170=2778%:%
%:%3171=2779%:%
%:%3172=2780%:%
%:%3173=2781%:%
%:%3174=2782%:%
%:%3183=2784%:%
%:%3208=2798%:%
%:%3209=2799%:%
%:%3210=2800%:%
%:%3211=2801%:%
%:%3212=2802%:%
%:%3213=2803%:%
%:%3214=2804%:%
%:%3215=2805%:%
%:%3216=2806%:%
%:%3217=2807%:%
%:%3218=2808%:%
%:%3219=2809%:%
%:%3220=2810%:%
%:%3221=2811%:%
%:%3222=2812%:%
%:%3223=2813%:%
%:%3224=2814%:%
%:%3225=2815%:%
%:%3226=2816%:%
%:%3227=2817%:%
%:%3228=2818%:%
%:%3229=2819%:%
%:%3230=2820%:%
%:%3231=2821%:%
%:%3232=2822%:%
%:%3233=2823%:%
%:%3234=2824%:%
%:%3235=2825%:%
%:%3236=2826%:%
%:%3237=2827%:%
%:%3238=2828%:%
%:%3239=2829%:%
%:%3240=2830%:%
%:%3241=2831%:%
%:%3242=2832%:%
%:%3243=2833%:%
%:%3244=2834%:%
%:%3245=2835%:%
%:%3246=2836%:%
%:%3247=2837%:%
%:%3248=2838%:%
%:%3249=2839%:%
%:%3250=2840%:%
%:%3251=2841%:%
%:%3252=2842%:%
%:%3253=2843%:%
%:%3254=2844%:%
%:%3255=2845%:%
%:%3256=2846%:%
%:%3257=2847%:%
%:%3258=2848%:%
%:%3259=2849%:%
%:%3260=2850%:%
%:%3261=2851%:%
%:%3262=2852%:%
%:%3263=2853%:%
%:%3264=2854%:%
%:%3265=2855%:%
%:%3266=2856%:%
%:%3267=2857%:%
%:%3268=2858%:%
%:%3269=2859%:%
%:%3270=2860%:%
%:%3271=2861%:%
%:%3272=2862%:%
%:%3273=2863%:%
%:%3274=2864%:%
%:%3275=2865%:%
%:%3276=2866%:%
%:%3277=2867%:%
%:%3278=2868%:%
%:%3279=2869%:%
%:%3280=2870%:%
%:%3281=2871%:%
%:%3282=2872%:%
%:%3283=2873%:%
%:%3284=2874%:%
%:%3285=2875%:%
%:%3286=2876%:%
%:%3287=2877%:%
%:%3288=2878%:%
%:%3289=2879%:%
%:%3290=2880%:%
%:%3291=2881%:%
%:%3292=2882%:%
%:%3293=2883%:%
%:%3294=2884%:%
%:%3295=2885%:%
%:%3296=2886%:%
%:%3297=2887%:%
%:%3298=2888%:%
%:%3299=2889%:%
%:%3300=2890%:%
%:%3301=2891%:%
%:%3302=2892%:%
%:%3303=2893%:%
%:%3304=2894%:%
%:%3305=2895%:%
%:%3306=2896%:%
%:%3307=2897%:%
%:%3308=2898%:%
%:%3309=2899%:%
%:%3310=2900%:%
%:%3311=2901%:%
%:%3312=2902%:%
%:%3313=2903%:%
%:%3314=2904%:%
%:%3315=2905%:%
%:%3316=2906%:%
%:%3317=2907%:%
%:%3318=2908%:%
%:%3319=2909%:%
%:%3320=2910%:%
%:%3321=2911%:%
%:%3322=2912%:%
%:%3323=2913%:%
%:%3324=2914%:%
%:%3325=2915%:%
%:%3326=2916%:%
%:%3327=2917%:%
%:%3328=2918%:%
%:%3329=2919%:%
%:%3330=2920%:%
%:%3331=2921%:%
%:%3332=2922%:%
%:%3333=2923%:%
%:%3334=2924%:%
%:%3335=2925%:%
%:%3336=2926%:%
%:%3337=2927%:%
%:%3338=2928%:%
%:%3339=2929%:%
%:%3340=2930%:%
%:%3341=2931%:%
%:%3342=2932%:%
%:%3343=2933%:%
%:%3344=2934%:%
%:%3345=2935%:%
%:%3346=2936%:%
%:%3347=2937%:%
%:%3348=2938%:%
%:%3349=2939%:%
%:%3350=2940%:%
%:%3351=2941%:%
%:%3352=2942%:%
%:%3353=2943%:%
%:%3354=2944%:%
%:%3355=2945%:%
%:%3356=2946%:%
%:%3357=2947%:%
%:%3358=2948%:%
%:%3359=2949%:%
%:%3360=2950%:%
%:%3361=2951%:%
%:%3362=2952%:%
%:%3363=2953%:%
%:%3364=2954%:%
%:%3365=2955%:%
%:%3366=2956%:%
%:%3367=2957%:%
%:%3368=2958%:%
%:%3369=2959%:%
%:%3370=2960%:%
%:%3371=2961%:%
%:%3372=2962%:%
%:%3373=2963%:%
%:%3374=2964%:%
%:%3375=2965%:%
%:%3376=2966%:%
%:%3377=2967%:%
%:%3378=2968%:%
%:%3379=2969%:%
%:%3380=2970%:%
%:%3381=2971%:%
%:%3382=2972%:%
%:%3383=2973%:%
%:%3384=2974%:%
%:%3385=2975%:%
%:%3386=2976%:%
%:%3387=2977%:%
%:%3388=2978%:%
%:%3389=2979%:%
%:%3390=2980%:%
%:%3391=2981%:%
%:%3392=2982%:%
%:%3393=2983%:%
%:%3394=2984%:%
%:%3395=2985%:%
%:%3396=2986%:%
%:%3397=2987%:%
%:%3398=2988%:%
%:%3399=2989%:%
%:%3400=2990%:%
%:%3401=2991%:%
%:%3402=2992%:%
%:%3403=2993%:%
%:%3404=2994%:%
%:%3405=2995%:%
%:%3406=2996%:%
%:%3407=2997%:%
%:%3408=2998%:%
%:%3409=2999%:%
%:%3410=3000%:%
%:%3411=3001%:%
%:%3412=3002%:%
%:%3413=3003%:%
%:%3414=3004%:%
%:%3415=3006%:%
%:%3416=3007%:%
%:%3417=3008%:%
%:%3418=3009%:%
%:%3419=3010%:%
%:%3420=3011%:%
%:%3421=3012%:%
%:%3422=3013%:%
%:%3423=3014%:%
%:%3424=3015%:%
%:%3425=3016%:%
%:%3426=3017%:%
%:%3427=3018%:%
%:%3428=3019%:%
%:%3429=3020%:%
%:%3430=3021%:%
%:%3431=3022%:%
%:%3432=3023%:%
%:%3433=3024%:%
%:%3434=3025%:%
%:%3435=3026%:%
%:%3436=3027%:%
%:%3437=3028%:%
%:%3438=3029%:%
%:%3439=3030%:%
%:%3440=3031%:%
%:%3441=3032%:%
%:%3442=3033%:%
%:%3443=3034%:%
%:%3444=3035%:%
%:%3445=3036%:%
%:%3446=3037%:%
%:%3447=3038%:%
%:%3448=3039%:%
%:%3449=3040%:%
%:%3450=3041%:%
%:%3451=3042%:%
%:%3452=3043%:%
%:%3453=3044%:%
%:%3454=3045%:%
%:%3455=3046%:%
%:%3456=3047%:%
%:%3457=3048%:%
%:%3458=3049%:%
%:%3459=3050%:%
%:%3460=3051%:%
%:%3461=3052%:%
%:%3462=3053%:%
%:%3463=3054%:%
%:%3464=3055%:%
%:%3465=3056%:%
%:%3474=3058%:%
%:%3478=3059%:%
%:%3490=3061%:%
%:%3491=3062%:%
%:%3492=3063%:%
%:%3493=3064%:%
%:%3494=3065%:%
%:%3495=3066%:%
%:%3496=3067%:%
%:%3497=3068%:%
%:%3498=3069%:%
%:%3499=3070%:%
%:%3500=3071%:%
%:%3501=3072%:%
%:%3502=3073%:%
%:%3503=3074%:%
%:%3504=3075%:%
%:%3505=3076%:%
%:%3506=3077%:%
%:%3507=3078%:%
%:%3508=3079%:%
%:%3509=3080%:%
%:%3510=3081%:%
%:%3511=3082%:%
%:%3512=3083%:%
%:%3513=3084%:%
%:%3514=3085%:%
%:%3515=3086%:%
%:%3516=3087%:%
%:%3517=3088%:%
%:%3518=3089%:%
%:%3519=3090%:%
%:%3520=3091%:%
%:%3521=3092%:%
%:%3522=3093%:%
%:%3523=3094%:%
%:%3524=3095%:%
%:%3525=3096%:%
%:%3526=3097%:%
%:%3527=3098%:%
%:%3528=3099%:%
%:%3529=3100%:%
%:%3530=3101%:%
%:%3531=3102%:%
%:%3532=3103%:%
%:%3533=3104%:%
%:%3534=3105%:%
%:%3535=3106%:%
%:%3536=3107%:%
%:%3537=3108%:%
%:%3538=3109%:%
%:%3539=3110%:%
%:%3540=3111%:%
%:%3541=3112%:%
%:%3542=3113%:%
%:%3543=3114%:%
%:%3544=3115%:%
%:%3545=3116%:%
%:%3546=3117%:%
%:%3547=3118%:%
%:%3548=3119%:%
%:%3549=3120%:%
%:%3550=3121%:%
%:%3551=3122%:%
%:%3552=3123%:%
%:%3553=3124%:%
%:%3554=3125%:%
%:%3555=3126%:%
%:%3556=3127%:%
%:%3565=3129%:%
%:%3577=3131%:%
%:%3578=3132%:%
%:%3579=3133%:%
%:%3580=3134%:%
%:%3581=3135%:%
%:%3582=3136%:%
%:%3583=3137%:%
%:%3584=3138%:%
%:%3585=3139%:%
%:%3586=3140%:%
%:%3587=3141%:%
%:%3588=3142%:%
%:%3589=3143%:%
%:%3590=3144%:%
%:%3591=3145%:%
%:%3592=3146%:%
%:%3593=3147%:%
%:%3594=3148%:%
%:%3595=3149%:%
%:%3596=3150%:%
%:%3597=3151%:%
%:%3598=3152%:%
%:%3599=3153%:%
%:%3600=3154%:%
%:%3601=3155%:%
%:%3602=3156%:%
%:%3603=3157%:%
%:%3604=3158%:%
%:%3605=3159%:%
%:%3606=3160%:%
%:%3607=3161%:%
%:%3608=3162%:%
%:%3609=3163%:%
%:%3610=3164%:%
%:%3611=3165%:%
%:%3612=3166%:%
%:%3613=3167%:%
%:%3614=3168%:%
%:%3615=3169%:%
%:%3616=3170%:%
%:%3617=3171%:%
%:%3618=3172%:%
%:%3619=3173%:%
%:%3620=3174%:%
%:%3621=3175%:%
%:%3622=3176%:%
%:%3623=3177%:%
%:%3624=3178%:%
%:%3625=3179%:%
%:%3626=3180%:%
%:%3627=3181%:%
%:%3628=3182%:%
